\documentclass[11pt,a4paper]{amsart}
\usepackage[margin=1in]{geometry}
\usepackage[T1]{fontenc}
\usepackage[utf8]{inputenc}
\usepackage{lmodern}
\usepackage{microtype}
\emergencystretch=3em
\usepackage{amsmath,amssymb,amsthm,mathtools}
\usepackage{enumitem}
\usepackage{booktabs}
\usepackage{array}
\usepackage{float}
\usepackage{aliascnt}
\usepackage[colorlinks=true,linkcolor=blue,citecolor=blue,urlcolor=blue]{hyperref}
\usepackage[capitalize,nameinlink]{cleveref}
\numberwithin{equation}{section}

\newtheorem{theorem}{Theorem}[section]
\newaliascnt{proposition}{theorem}\newtheorem{proposition}[proposition]{Proposition}\aliascntresetthe{proposition}
\newaliascnt{lemma}{theorem}\newtheorem{lemma}[lemma]{Lemma}\aliascntresetthe{lemma}
\newaliascnt{corollary}{theorem}\newtheorem{corollary}[corollary]{Corollary}\aliascntresetthe{corollary}
\newaliascnt{definition}{theorem}\newtheorem{definition}[definition]{Definition}\aliascntresetthe{definition}
\newtheorem{hypothesis}[theorem]{Hypothesis}
\newtheorem{problem}[theorem]{Problem}
\theoremstyle{remark}
\newaliascnt{remark}{theorem}\newtheorem{remark}[remark]{Remark}\aliascntresetthe{remark}
\crefname{theorem}{Theorem}{Theorems}\Crefname{theorem}{Theorem}{Theorems}
\crefname{proposition}{Proposition}{Propositions}\Crefname{proposition}{Proposition}{Propositions}
\crefname{lemma}{Lemma}{Lemmas}\Crefname{lemma}{Lemma}{Lemmas}
\crefname{corollary}{Corollary}{Corollaries}\Crefname{corollary}{Corollary}{Corollaries}
\crefname{definition}{Definition}{Definitions}\Crefname{definition}{Definition}{Definitions}
\crefname{remark}{Remark}{Remarks}\Crefname{remark}{Remark}{Remarks}
\crefname{hypothesis}{Hypothesis}{Hypotheses}\Crefname{hypothesis}{Hypothesis}{Hypotheses}
\crefname{problem}{Problem}{Problems}\Crefname{problem}{Problem}{Problems}

\DeclareMathOperator{\RS}{RS}\DeclareMathOperator{\MCA}{MCA}\DeclareMathOperator{\CA}{CA}
\DeclareMathOperator{\rank}{rank}\DeclareMathOperator{\Span}{Span}\DeclareMathOperator{\Spec}{Spec}
\DeclareMathOperator{\Hom}{Hom}\DeclareMathOperator{\im}{im}\DeclareMathOperator{\codim}{codim}
\DeclareMathOperator{\supp}{supp}\DeclareMathOperator{\Fib}{Fib}\DeclareMathOperator{\Sym}{Sym}
\DeclareMathOperator{\Gr}{Gr}\DeclareMathOperator{\Ent}{H}\DeclareMathOperator{\ord}{ord}
\DeclareMathOperator{\Tr}{Tr}\DeclareMathOperator{\Gal}{Gal}\DeclareMathOperator{\Res}{Res}
\DeclareMathOperator{\Aut}{Aut}\DeclareMathOperator{\diag}{diag}\DeclareMathOperator{\wt}{wt}
\DeclareMathOperator{\Q}{Q}\DeclareMathOperator{\poly}{poly}\DeclareMathOperator{\wdeg}{wdeg}
\DeclareMathOperator{\Cen}{Cen}\DeclareMathOperator{\LineRay}{LineRay}\DeclareMathOperator{\dist}{dist}

\newcommand{\F}{\mathbb F}\newcommand{\B}{\mathbb B}\newcommand{\K}{\mathbb K}\newcommand{\Z}{\mathbb Z}
\newcommand{\A}{\mathbb A}\newcommand{\Pj}{\mathbb P}\newcommand{\C}{\mathbb C}\newcommand{\Prob}{\mathbb P}
\newcommand{\E}{\mathbb E}\newcommand{\eps}{\varepsilon}\newcommand{\emca}{\varepsilon_{\mathrm{mca}}}
\newcommand{\eca}{\varepsilon_{\mathrm{ca}}}\newcommand{\Htwo}{H_2}\newcommand{\Omc}{\Omega^{\circ}}
\newcommand{\Ccal}{\mathcal C}\newcommand{\Dcal}{\mathcal D}\newcommand{\Ecal}{\mathcal E}
\newcommand{\Fcal}{\mathcal F}\newcommand{\Gcal}{\mathcal G}\newcommand{\Hcal}{\mathcal H}
\newcommand{\Ical}{\mathcal I}\newcommand{\Jcal}{\mathcal J}\newcommand{\Lcal}{\mathcal L}
\newcommand{\Mcal}{\mathcal M}\newcommand{\Ncal}{\mathcal N}\newcommand{\Pcal}{\mathcal P}
\newcommand{\Qcal}{\mathcal Q}\newcommand{\Rcal}{\mathcal R}\newcommand{\Scal}{\mathcal S}
\newcommand{\Tcal}{\mathcal T}\newcommand{\Ucal}{\mathcal U}\newcommand{\Vcal}{\mathcal V}
\newcommand{\Wcal}{\mathcal W}\newcommand{\Xcal}{\mathcal X}\newcommand{\Ycal}{\mathcal Y}
\newcommand{\Zcal}{\mathcal Z}\newcommand{\Gord}{\Gamma^{\rm ord}}\newcommand{\Gsid}{\Gamma^{\rm sid}}
\newcommand{\Gprim}{\Gamma^{\rm prim}}\newcommand{\barN}{\overline N}\newcommand{\Renyi}{R\'enyi}
\newcommand{\Hb}{H_2}\newcommand{\UU}{\mathbb U}\newcommand{\abs}[1]{\lvert #1\rvert}
\newcommand{\norm}[1]{\lVert #1\rVert}\newcommand{\floor}[1]{\left\lfloor #1\right\rfloor}
\newcommand{\ceil}[1]{\left\lceil #1\right\rceil}\newcommand{\defeq}{\coloneqq}
\newcommand{\one}{\mathbf 1}\newcommand{\Fq}{\mathbb F_q}
\newcommand{\cA}{\mathcal A}\newcommand{\cB}{\mathcal B}\newcommand{\cC}{\mathcal C}
\newcommand{\cE}{\mathcal E}\newcommand{\cH}{\mathcal H}\newcommand{\cK}{\mathcal K}
\newcommand{\cL}{\mathcal L}\newcommand{\cM}{\mathcal M}\newcommand{\cP}{\mathcal P}
\newcommand{\cS}{\mathcal S}\newcommand{\cZ}{\mathcal Z}
\newcommand{\vander}[2]{(1,#1,\ldots,#1^{#2-1})}
\newcommand{\fall}[2]{#1^{\underline{#2}}}\newcommand{\rise}[2]{#1^{\overline{#2}}}
\newcommand{\Bbad}{B_{\mathrm{MCA}}}




\title[Grande Finale v4]{Grande Finale v4:\\
Exact Threshold Compilers, Primitive Q, and Order-$32$ Reed--Solomon MCA}
\author{Przemek Chojecki}
\hypersetup{
pdftitle={Grande Finale v4: Exact Threshold Compilers, Primitive Q, and Order-32 Reed--Solomon MCA},
pdfauthor={Przemek Chojecki},
pdfsubject={A complete companion to CAP25 v13.2 integrating exact threshold refinements, primitive Q, support compilers, order-32 rank, rational atoms, owner localization, and spread-core incidence},
pdfkeywords={Reed--Solomon codes, mutual correlated agreement, primitive Q, Sidon moments, generalized Hamming weights, order-32 rank, rational atoms, spread cores}
}
\subjclass[2020]{94B35, 11T71, 05B40}
\keywords{Reed--Solomon codes, mutual correlated agreement, primitive Q, Sidon moments, generalized Hamming weights, order-$32$ rank, rational atoms, spread cores}
\date{July 2026}

\begin{document}
\begin{abstract}
This paper is a complete companion to CAP25 v13.2 \cite{CAP25v132}.  It imports by citation, and does not reprove, the support-wise MCA conventions, universal cap, four deployed unsafe endpoints, identity-prefix and simple-pole constructions, aperiodic and subfield floors, prefix rigidity and exact second moment, quotient and extension reductions, balanced-core reduction, and normalized finite ledger established there.

The first half develops the full safe-side compiler.  It proves exact support sparsification, syndrome--secant and fixed-union compilers, recursive affine-span and generalized-weight bounds, common-zero and directional Johnson estimates, codimension-one and arbitrary-codimension support-flag recursions, a minimum-distance sharpening using the actual direction-code support, fixed-mismatch puncturing, a quadratic mean-overlap staircase, exact quotient--remainder prefix normal forms, effective-image Fourier normalization, primitive Q after an explicit Sidon payment, saturation, a moving-root theorem for one-parameter split pencils, and the implication that a row-sharp Q theorem pays the primitive shift-pair ledger.

The second half develops a different MCA completion mechanism.  A rank-two reduced intersection matrix gives an exact kernel dictionary and shows that the deployed critical order is $32$.  Exact dyadic descent identifies the quotient branch.  Maximal-support interpolation excludes all primitive slope degrees below $18$, while support-collapsed interpolation extracts every low-complexity explanation into a scalar-locator rational atom.  Rational atoms satisfy collision rigidity and local-to-global coherence.  A finite critical example shows that primitive spread cores can exist, so the remaining theorem is abundance rather than nonexistence.  Canonical owner localization pays every small-owner atom, component-free owner pencils contribute at most seven rich points, core-compatible slopes contribute at most $58$, a complete correction ray contributes at most $2(n-m+1)+963$, and proper or clone-tolerant correction spaces are paid in the explicit deployed ranges.

The manuscript does not assert the adjacent safe rows unconditionally.  For MCA, the remaining inputs are the classification and payment of high-dimensional positive-dimensional correction components, the row-sharp image bound for one coherent large-owner atom---most sharply the Mersenne-31 factor-$9.57219783037$ prefix target---and complete routing of denominator-root, extension, clone, imperfect-quotient, and near-sunflower exceptions.  For the list rows, the remaining inputs are the row-sharp high-nullity prefix and arbitrary-word interior bounds, extension payment, and one exact summed integer certificate.
\end{abstract}
\maketitle
\tableofcontents

\section{Scope and relation to CAP25 v13.2}\label{sec:relation-v132}
Throughout, CAP25 v13.2 \cite{CAP25v132} is the foundational source for the deployed unsafe frontier and the normalized first-match ledger.  Its theorems are invoked by citation rather than re-numbered here.  This paper contains all of the additional threshold, compiler, rank, atom, and spread-incidence results needed for the current completion programme.

\begin{center}
\begin{tabular}{@{}>{\raggedright\arraybackslash}p{0.39\textwidth}>{\raggedright\arraybackslash}p{0.53\textwidth}@{}}
\toprule
Imported from CAP25 v13.2 & Proved in this paper\\
\midrule
MCA conventions, universal cap, active unsafe endpoints & Exact sparsification, syndrome--secant, fixed-union, affine-span, generalized-weight, support-flag, and minimum-distance compilers\\
Identity-prefix, pole-line, aperiodic, subfield, quotient, and extension packages & Quadratic threshold staircases, exact profile normal forms, effective Fourier normalization, and primitive Q after the stated Sidon payment\\
Prefix rigidity, exact second moment, balanced-core reduction, normalized ledger & Saturation, moving-root BC, Q $\Rightarrow$ SP, order-$32$ rank, rational-atom extraction and coherence, owner localization, and spread-core correction incidence\\
\bottomrule
\end{tabular}
\end{center}

The paper uses the closed integer-ball convention of CAP25 v13.2.  An adjacent certificate consists of the imported unsafe inequality at agreement $a_0$ and a proved upper inequality at $a_0+1$.  Structural reductions, asymptotic estimates, and route-cut boundaries are not called adjacent certificates unless they fit the literal row budget.

\section{Main theorem package}\label{sec:new-contributions}
\begin{theorem}[Integrated theorem package]\label{thm:v4-package}
The following statements hold.
\begin{enumerate}[label=\textup{(\roman*)},leftmargin=*]
\item The exact support, syndrome--secant, fixed-union, affine-span, generalized-weight, common-zero, codimension-one, support-flag, minimum-distance, puncturing, directional-Johnson, paving, mean-overlap, packing, and explicit prime-row results of Part~I are unconditional.
\item The quotient--remainder normal form, partial-occupancy disintegration, effective-image Fourier identities, and profile add-back of Part~II are unconditional.  Primitive Q follows from the Sidon-moment payment of Part~III, and the shallow Fourier theorem proves that payment in its displayed range.
\item Saturation, the one-parameter moving-root theorem, and Q $\Rightarrow$ SP are unconditional.  The finite compilers pay the displayed low-nullity, low-affine-span, common-zero, support-flag, and direction-defect ranges.
\item The rank-two kernel dictionary, critical order $32$, global affine-block bound, dyadic pair descent, rational-profile converse, slope-degree barrier, low-complexity atom extraction, collision rigidity, atom coherence, owner localization, owner-pencil incidence bound, core-compatible abundance bound, complete correction-ray bound, and the stated proper and clone-tolerant correction-space bounds are unconditional.
\item The two deployed MCA adjacent rows follow from the three explicit inputs in \cref{thm:conditional-final}; the two list adjacent rows follow from the exact summed certificate in \cref{thm:exact-completion-certificate}.  Those remaining inputs are not asserted in this paper.
\end{enumerate}
\end{theorem}

\part{Exact threshold refinements beyond CAP25 v13.2}

\section{Exact decomposition of the mutual layer}\label{sec:mutual-decomposition}
Let $C\subseteq\F^D$ be linear, let $a$ be an agreement threshold, and put $t=|D|-a$.  We use the foundation's definitions of support-wise MCA-badness and of the challenge-restricted numerator $B^{\rm MCA}_{C,\Gamma}(a)$ \cite[Secs.~1--2]{CAP25v132}.

A received pair is \emph{column-far at agreement $a$} when no support of size at least $a$ simultaneously explains its two coordinates.  Let $B^{\rm CA}_{C,\Gamma}(a)$ be the maximum number of challenge slopes whose affine combination is explained on some support of size at least $a$, maximized over column-far pairs.  Define the sparse mutual numerator
\[
 \sigma_{C,\Gamma}(a)=\max\#\{\gamma\in\Gamma:\gamma\text{ is MCA-bad for }(e_0,e_1)\},
\]
where the maximum is over pairs with $|\supp(e_0)\cup\supp(e_1)|\le t$.

\begin{theorem}[Exact sparsification of the mutual layer]\label{thm:exact-sparsification}
For every linear code, every nonempty challenge set $\Gamma$, and every $a$,
\begin{equation}
 B^{\rm MCA}_{C,\Gamma}(a)=
 \max\{B^{\rm CA}_{C,\Gamma}(a),\sigma_{C,\Gamma}(a)\}.\label{eq:exact-sparsification}
\end{equation}
\end{theorem}
\begin{proof}
Both terms on the right maximize over subclasses of received pairs, so the left side dominates them.  Conversely fix a pair.  If it is column-far, its MCA-bad and CA-bad challenge slopes coincide.  Otherwise translate it by a codeword pair that explains it on one common support of size at least $a$.  The translated pair is supported on at most $t$ coordinates, and translation preserves, slope by slope and support by support, both affine explanation and simultaneous explanation.  Its bad set is therefore bounded by $\sigma_{C,\Gamma}(a)$.
\end{proof}

The foundation supplies the exact-agreement reduction, support uniqueness, tangent floor, and deep-regime theorem used below \cite[Secs.~13--14]{CAP25v132}.  We invoke those interfaces only by citation.

\section{Syndrome--secant geometry and finite ray compilers}
\label{sec:syndrome-exact}

The \emph{syndrome} of a word is the result of applying a parity-check
matrix; it vanishes exactly on codewords.  Thus syndrome space forgets the
part of a received word already lying in the code and retains only its error
class.  This description separates possible error supports from distinct
line parameters.  Put \(R=n-k\), the redundancy (and syndrome-space
dimension), and choose the generalized Vandermonde parity-check matrix
\(H\) with columns
\begin{equation}
 h_x=\lambda_x(1,x,\ldots,x^{R-1})^{\mathsf T},\qquad
 \lambda_x=\left(\prod_{y\in D\setminus\{x\}}(x-y)\right)^{-1}.\label{eq:parity-columns}
\end{equation}
Lagrange interpolation gives
\(\sum_{x\in D}\lambda_xx^j=0\) for \(0\le j\le n-2\), so this matrix
annihilates every degree-less-than-\(k\) evaluation.  Every set of at most
\(R\) columns is independent.  For \(E\subseteq D\), write
\(V_E=\Span\{h_x:x\in E\}\); it is exactly the set of syndromes of words
supported on \(E\).  We write \(s(u)=Hu^{\mathsf T}\) for the syndrome of
\(u\).

\begin{proposition}[Syndrome-line normal form]
\label{prop:syndrome-line-normal-form}
Let \(S=D\setminus E\).  A word \(u\) is explained on \(S\) if and only if
\(s(u)\in V_E\).  Consequently \(\gamma\) is MCA-bad on \(S\) for a pair
\((u_0,u_1)\) if and only if
\begin{equation}
 s(u_0)+\gamma s(u_1)\in V_E,
 \qquad
 \{s(u_0),s(u_1)\}\not\subseteq V_E.\label{eq:syndrome-line}
\end{equation}
For fixed \(E\), there is at most one such finite slope.
\end{proposition}

\begin{proof}
If \(u=c+e\), where \(c\in C\) and \(e\) is supported on \(E\), then
\(s(u)=s(e)\in V_E\).  Conversely, if
\(s(u)=\sum_{x\in E}e_xh_x\), subtract the word supported on \(E\) with
values \(e_x\); the result has zero syndrome and belongs to \(C\).  Apply
this equivalence to the line point and to both coordinates.  Passing
\eqref{eq:syndrome-line} to \(\F^R/V_E\) shows that two distinct solutions
would put both projected syndromes at zero, contradicting badness.
\end{proof}

Since \(\lambda_x\ne0\), one has
\([h_x]=[1:x:\cdots:x^{R-1}]\); these are the points of the rational
normal curve in projective syndrome space, while the weights
\(\lambda_x\) merely choose vector representatives.  A \emph{secant span} is the linear span
\(V_E\) of a selected set of those points.  We call the meeting of a
syndrome line with \(V_E\) \emph{transverse} when the entire line is not
contained in \(V_E\), equivalently when its two generators are not both in
\(V_E\).  The next result is a \emph{compiler} in the sense that it
translates the bad-slope problem, exactly and in both directions, into this
line--secant incidence problem.  A \emph{ray} means a one-dimensional
direction in syndrome or coefficient space; different raw witnesses may
represent the same ray and hence the same slope.

\begin{theorem}[Exact syndrome--secant compiler]
\label{thm:syndrome-secant-exact}
Put \(t=n-a\), and for a syndrome line
\(\ell_{y_0,y_1}=\{y_0+\gamma y_1:\gamma\in\F\}\) define
\begin{equation}
 \Theta_t(y_0,y_1)=
 \abs{\{\gamma\in\F:\exists E\subseteq D,\ \abs E\le t,
 y_0+\gamma y_1\in V_E,
 \ \{y_0,y_1\}\not\subseteq V_E\}}.\label{eq:transverse-secant-count}
\end{equation}
Then
\begin{equation}
 B_C^{\rm MCA}(a)=\max_{y_0,y_1\in\F^{n-k}}
                    \Theta_t(y_0,y_1).\label{eq:exact-secant-numerator}
\end{equation}
For each fixed \(E\), the affine line has at most one transverse
intersection parameter with \(V_E\).  For a Reed--Solomon code the
spaces \(V_E\) are the spans of at most \(t\) points of the weighted
rational normal curve in \eqref{eq:parity-columns}.  Thus the
higher-dimensional ray problem is exactly a transverse line--secant
incidence problem, with repeated intersection parameters deduplicated.
\end{theorem}

\begin{proof}
Apply \cref{prop:syndrome-line-normal-form} with \(S=D\setminus E\), and
take the union over \(\abs E\le t\).  The syndrome map is surjective, so
maximizing over received pairs is the same as maximizing over
\((y_0,y_1)\).  In the quotient \(\F^{n-k}/V_E\), a transverse
intersection solves
\(\overline y_0+\gamma\overline y_1=0\) with the two projected vectors
not both zero, and therefore has at most one finite solution.  The last
assertion is the Vandermonde description of the parity columns.
\end{proof}

A \emph{parity-column circuit} is a minimally linearly dependent set of
columns of \(H\).  The next lemma says that many distinct transverse rays
cannot all be supported inside an independent union of columns.

\begin{lemma}[Independent-union ray bound]
\label{lem:independent-union-rays}
Let every set of at most \(R\) parity columns be independent, put
\(t<R\), and fix a syndrome line.  For each slope in a transverse set
\(Z\), choose an error set \(E_\gamma\), \(\abs{E_\gamma}\le t\),
witnessing \eqref{eq:transverse-secant-count}.  If
\(U=\bigcup_{\gamma\in Z}E_\gamma\) has size at most \(R\), then
\(\abs Z\le t+1\).  Consequently, more than \(t+1\) transverse slopes
force the witness union to contain a parity-column circuit.
\end{lemma}

\begin{proof}
If \(\abs Z\le1\), there is nothing to prove; assume that \(Z\) contains
two distinct slopes.  Two distinct line points lie in \(V_U\), hence
\(y_0,y_1\in V_U\).
Column independence gives unique coefficient lifts
\(e_0,e_1\in\F^U\), and every chosen lift is
\(e_\gamma=e_0+\gamma e_1\).  Let \(U'\) be the set of \(s\) coordinates
where the pair \((e_0,e_1)\) is nonzero.  If \(s\le t\) and
\(E_\gamma\) contained all of \(U'\), then
\(y_0,y_1\in V_{U'}\subseteq V_{E_\gamma}\), contradicting the
transversality of this particular witness.  Hence some
\(x\in U'\setminus E_\gamma\) exists, and uniqueness of the lift gives
\(e_0(x)+\gamma e_1(x)=0\).  Each nonzero affine coordinate function has
at most one root, so \(\abs Z\le s\le t\).  If \(s>t\), every chosen lift of
weight at most \(t\) has at least \(s-t\) zero coordinates.  Double
counting the pairs \((\gamma,x)\) with
\(e_0(x)+\gamma e_1(x)=0\) gives
\(\abs Z(s-t)\le s\), whence
\(\abs Z\le\floor{s/(s-t)}\le t+1\).  If \(\abs U\le R\), independence
holds; the contrapositive proves the circuit assertion.
\end{proof}

\begin{theorem}[Bounded residual-kernel ray compiler]
\label{thm:bounded-residual-kernel-ray}
Let \(H\) be a parity-check matrix over \(\F\), let \(U\) be a set of
columns, and put
\[
 \kappa(U)=\dim_\F\ker(H_U:\F^U\longrightarrow\F^{n-k}).
\]
Fix a syndrome line \(y_0+\gamma y_1\), an integer \(t\ge0\), and a set
\(Z\) of finite transverse slopes.  Suppose that for every \(\gamma\in Z\)
there is \(E_\gamma\subseteq U\), \(\abs{E_\gamma}\le t\), such that
\[
 y_0+\gamma y_1\in V_{E_\gamma},\qquad
 \{y_0,y_1\}\not\subseteq V_{E_\gamma}.
\]
Then
\[
                  \abs Z\le(t+1)\abs{\F}^{\kappa(U)}. \tag{RC$_{\rm ker}$}
\]
For a Reed--Solomon parity check of redundancy \(R=n-k\),
\(\kappa(U)=\max\{0,\abs U-R\}\).  Hence a certified residual chart with
\((\abs U-R)_+\log\abs\F=o(n)\) has a direct subexponential ray payment
with the displayed exact finite loss.
\end{theorem}

\begin{proof}
The assertion is immediate when \(\abs Z\le1\).  Otherwise two line
points belong to \(V_U\), so \(y_0,y_1\in V_U\).  Choose lifts
\(b_0,b_1\in\F^U\) with \(H_Ub_i=y_i\).  For every \(\gamma\), choose a
lift \(c_\gamma\) supported in \(E_\gamma\) and put
\(z_\gamma=c_\gamma-b_0-\gamma b_1\in\ker H_U\).

Fix \(z\in\ker H_U\), and consider the slopes for which
\(z_\gamma=z\).  Put
\[
 U_z=\{x\in U:(b_0(x)+z(x),b_1(x))\ne(0,0)\},
 \qquad s=\abs{U_z}.
\]
Every active coordinate is a nonzero affine function of \(\gamma\), and
therefore vanishes for at most one slope.  Hence the total number of active
zero incidences in this kernel class is at most \(s\).  If \(s>t\), the
weight bound on \(c_\gamma\) forces at least \(s-t\) active zeros, so
\[
 \abs{Z_z}(s-t)\le s,\qquad
 \abs{Z_z}\le\floor{s/(s-t)}\le t+1.
\]
If \(s\le t\), transversality forces
at least one active zero: without one,
\(U_z=\supp c_\gamma\subseteq E_\gamma\), while both \(b_0+z\) and
\(b_1\) are supported in \(U_z\), putting \(y_0,y_1\) in
\(V_{E_\gamma}\).  Thus this kernel class contains at most \(s\le t\)
slopes.  Summing over the \(\abs{\F}^{\kappa(U)}\) kernel elements proves
the bound.  The MDS column property gives the Reed--Solomon formula for
\(\kappa(U)\).
\end{proof}


\begin{theorem}[Multiplicity-amplified fixed-union ray compiler]
\label{thm:fixed-union-ray}
Let \(H\) be a Reed--Solomon parity check of redundancy \(R\), let
\(U\subseteq D\) have size \(N=R+\nu\) with \(\nu\ge1\), and put
\(C_U=\ker H_U\).  Fix \(t<R\), a syndrome line
\(y_0+\gamma y_1\), and a set \(Z\) of distinct finite transverse slopes.
Suppose that for every \(\gamma\in Z\) there is
\(e_\gamma\in\F^U\) such that
\[
 H_Ue_\gamma=y_0+\gamma y_1,
 \qquad \wt(e_\gamma)\le t,
 \qquad
 \{y_0,y_1\}\not\subseteq V_{\supp e_\gamma}.
\]
Then
\begin{equation}
 |Z|\le
 \left\lfloor
 \frac{\binom{R+\nu}{\nu+1}}{R-t}
 \right\rfloor .
 \label{eq:fixed-union-ray}
\end{equation}
The estimate is independent of \(|\F|\); it replaces the factor
\(|\F|^\nu\) in \cref{thm:bounded-residual-kernel-ray} by an exact
incidence multiplicity.
\end{theorem}

\begin{proof}
There is nothing to prove when \(|Z|\le1\).  Otherwise two line points lie
in \(V_U\), hence \(y_0,y_1\in V_U\).  Choose lifts
\(b_0,b_1\in\F^U\), and choose a basis of the \(\nu\)-dimensional kernel
\(C_U\).  Write its generator matrix as \(K\), with row \(K_x\) at
\(x\in U\).  Every chosen error has a unique expression
\[
 e_\gamma=b_0+\gamma b_1+Kc_\gamma,
 \qquad c_\gamma\in\F^\nu.
\]
Put \(p_\gamma=(\gamma,c_\gamma)\in\F^{\nu+1}\).  A zero coordinate of
\(e_\gamma\) is the affine hyperplane incidence
\begin{equation}
 b_0(x)+\gamma b_1(x)+K_xc_\gamma=0,
 \qquad x\in U,                                     \label{eq:fixed-union-hyperplane}
\end{equation}
whose normal is \(n_x=(b_1(x),K_x)\in\F^{\nu+1}\).  Since
\(\wt(e_\gamma)\le t\), the point \(p_\gamma\) lies on at least
\(h=N-t\) of these hyperplanes.

We claim that the normals incident with \(p_\gamma\) span
\(\F^{\nu+1}\).  Otherwise a nonzero pair \((\delta,v)\) is orthogonal
to all of them.  Then
\[
 q=\delta b_1+Kv
\]
vanishes at every zero coordinate of \(e_\gamma\), so
\(\supp q\subseteq\supp e_\gamma\) and \(\wt(q)\le t\).  If
\(\delta=0\), then \(q=Kv\ne0\) is a kernel word of weight at most
\(t<R+1\), contradicting the MDS distance of
\(C_U=[R+\nu,\nu,R+1]\).  If \(\delta\ne0\), divide by \(\delta\).
The word \(q/\delta\) represents \(y_1\), and
\(e_\gamma-\gamma q/\delta\) represents \(y_0\); both are supported
inside \(\supp e_\gamma\), contradicting transversality.

Every normal \(n_x\) is nonzero.  Indeed, if the coordinate functional of
\(C_U\) vanished at \(x\), then puncturing at \(x\) would leave kernel
dimension \(\nu\), whereas the MDS rank of \(H_{U\setminus\{x\}}\) makes
that dimension \(\nu-1\).  Thus the incident normals form a loopless
rank-\((\nu+1)\) matroid on at least \(h\) elements.  Fix one basis.
Each of the other incident elements enters a basis by a fundamental-circuit
exchange, producing a distinct basis.  Hence \(p_\gamma\) is incident with
at least
\[
 h-(\nu+1)+1=R-t
\]
independent \((\nu+1)\)-subsets of hyperplanes.  An independent subset
has at most one affine intersection point.  Counting pairs consisting of a
slope point and one of its incident bases gives
\[
 |Z|(R-t)\le\binom{R+\nu}{\nu+1},
\]
which is \eqref{eq:fixed-union-ray}.
\end{proof}

\begin{corollary}[Single MDS-circuit ray compiler, multiplicity form]
\label{thm:single-mds-circuit-ray}
Under the hypotheses of the former single-circuit chart, namely
\(|U|=R+1\),
\begin{equation}
 |Z|\le
 \left\lfloor\frac{R(R+1)}{2(R-t)}\right\rfloor .
 \label{eq:single-circuit-multiplicity}
\end{equation}
Thus the pair-charge \(\binom{R+1}{2}\) is divided by the exact number
\(R-t\) of independent zero-coordinate pairs forced by every transverse
slope.
\end{corollary}

\begin{proof}
Apply \cref{thm:fixed-union-ray} with \(\nu=1\).
\end{proof}

\begin{theorem}[Fixed-syndrome GRS--Johnson compiler]
\label{thm:fixed-union-list-johnson}
Let \(H\) be a Reed--Solomon parity check of redundancy \(R\), let
\(U\subseteq D\) have size \(N=R+\nu\), and fix a syndrome \(y\).  Put
\[
 \mathcal L_U(y,t)=
 \{e\in\F^U:H_Ue=y,\ \wt(e)\le t\},
 \qquad h=N-t,
\]
where \(t<R\).  If \(\nu=0\), then \(|\mathcal L_U(y,t)|\le1\).  If
\(\nu\ge1\) and
\begin{equation}
 h^2>N(\nu-1),                                      \label{eq:fixed-union-johnson-condition}
\end{equation}
then
\begin{equation}
 |\mathcal L_U(y,t)|
 \le
 \left\lfloor
 \frac{N(h-\nu+1)}{h^2-N(\nu-1)}
 \right\rfloor .                                  \label{eq:fixed-union-list-johnson}
\end{equation}
Equivalently, the same bound holds for the number of codewords in one
Reed--Solomon list whose error supports are all contained in the fixed
union \(U\).
\end{theorem}

\begin{proof}
For \(\nu=0\), the columns in \(U\) are independent, so a syndrome has at
most one lift.  Assume \(\nu\ge1\), choose one lift \(b\) of \(y\), and
put \(C_U=\ker H_U\).  This is an
\([N,\nu,R+1]\) generalized Reed--Solomon code.  The set
\(\mathcal L_U(y,t)\) is the low-weight part of the affine coset
\(b+C_U\).  Thus its members are in bijection with codewords of \(C_U\)
that agree with the fixed word \(-b\) on at least \(h=N-t\) coordinates.

For each such codeword choose an \(h\)-subset \(A_i\) of its agreement
set.  Distinct codewords of the \([N,\nu]\) GRS code agree on at most
\(\nu-1\) coordinates, so \(|A_i\cap A_j|\le\nu-1\).  If there are \(L\)
sets, and \(d_x\) is the number containing \(x\), then
\[
 \sum_xd_x=Lh,
 \qquad
 \sum_x\binom{d_x}{2}\le\binom L2(\nu-1).
\]
Cauchy--Schwarz gives
\[
 \sum_x\binom{d_x}{2}
 \ge\frac12\left(\frac{L^2h^2}{N}-Lh\right).
\]
Combining the two inequalities and using
\eqref{eq:fixed-union-johnson-condition} yields
\[
 L\bigl(h^2-N(\nu-1)\bigr)
 \le N(h-\nu+1),
\]
which proves \eqref{eq:fixed-union-list-johnson}.
\end{proof}

\begin{theorem}[Recursive affine-span Reed--Solomon list compiler]
\label{thm:affine-span-list}
Let \(C=\RS[\F,D,K]\) have length \(n\), let \(u\in\F^D\), and put
\(m=K+w\le n\).  Let \(\mathcal A=c_0+C'\subseteq C\) be an affine
subspace of codewords whose direction space has dimension \(s\ge0\).  Then
\begin{equation}
 |\mathcal L_{\mathcal A}(u,m)|
 \le
 \left\lfloor
 \frac{(n-K+s)^{\underline s}}
      {(w+1)^{\overline s}}
 \right\rfloor
 =
 \left\lfloor
 \frac{\binom{n-K+s}{s}}
      {\binom{w+s}{s}}
 \right\rfloor .
 \label{eq:affine-span-list}
\end{equation}
The empty products at \(s=0\) are one.  In particular, a one-dimensional
codeword pencil contains at most
\(\lfloor(n-K+1)/(w+1)\rfloor\) members of the list.
\end{theorem}

\begin{proof}
We induct on \(s\).  The case \(s=0\) is immediate.  Assume \(s\ge1\),
and let
\[
 G=\{x\in D:c(x)=0\text{ for every }c\in C'\},
 \qquad z=|G|.
\]
The polynomial space representing \(C'\) has dimension \(s\).  Division
by the locator of \(G\) embeds it into the polynomials of degree less than
\(K-z\), so
\begin{equation}
                         z\le K-s.                 \label{eq:affine-common-zero}
\end{equation}
Let \(g\le z\) be the number of coordinates in \(G\) at which the common
value of the affine flat agrees with \(u\).  Every member of
\(\mathcal L_{\mathcal A}(u,m)\) therefore has at least \(m-g\)
agreements outside \(G\).

For \(x\in D\setminus G\), the evaluation functional on \(C'\) is
nonzero.  Hence
\[
 \mathcal A_x=\{c\in\mathcal A:c(x)=u(x)\}
\]
is either empty or an affine subspace of dimension \(s-1\).  By the
induction hypothesis, the number of listed codewords in \(\mathcal A_x\)
is at most
\[
 B_{s-1}:=
 \frac{\binom{n-K+s-1}{s-1}}
      {\binom{w+s-1}{s-1}}.
\]
Double-counting pairs \((c,x)\) with
\(c\in\mathcal L_{\mathcal A}(u,m)\), \(x\notin G\), and \(c(x)=u(x)\)
gives
\begin{equation}
 |\mathcal L_{\mathcal A}(u,m)|(m-g)
 \le (n-z)B_{s-1}.                                \label{eq:affine-recursion}
\end{equation}

Put \(b=z-g\) and \(q=K-s-z\).  Both are nonnegative by
\eqref{eq:affine-common-zero}, and the exact identity
\begin{equation}
 \begin{split}
 &(m-g)(n-K+s)-(n-z)(w+s)\\
 &\qquad=b(n-K+s)+q(n-m)\ge0
 \end{split}                                      \label{eq:affine-ratio-identity}
\end{equation}
shows that
\[
 \frac{n-z}{m-g}\le\frac{n-K+s}{w+s}.
\]
Combining this with \eqref{eq:affine-recursion} yields
\[
 |\mathcal L_{\mathcal A}(u,m)|
 \le B_{s-1}\frac{n-K+s}{w+s}
 =\frac{\binom{n-K+s}{s}}{\binom{w+s}{s}}.
\]
Taking the integer floor proves \eqref{eq:affine-span-list}.
\end{proof}


\begin{theorem}[Generalized-weight rank-flat list compiler]
\label{thm:rank-flat-list}
Retain the hypotheses of \cref{thm:affine-span-list}, assume \(s\ge1\),
and let
\[
 d_j=d_j(C'):=
 \min_{\substack{V\le C'\\ \dim V=j}}|\supp(V)|,
 \qquad 1\le j\le s,
\]
be the generalized Hamming weights of the direction code.  Put
\[
 R=n-K,\qquad t=n-m,
\]
let \(G\) be the common-zero set of \(C'\), let \(z=|G|=n-d_s\),
let \(g\) be the number of common agreements with \(u\) on \(G\), and
put \(b=z-g\).  Then
\begin{equation}
 |\mathcal L_{\mathcal A}(u,m)|
 \le
 \left\lfloor
 \frac{d_s^{\underline s}}
      {\prod_{j=1}^{s}(d_j-t+b)}
 \right\rfloor .
 \label{eq:rank-flat-list}
\end{equation}
Every factor in the denominator is positive.  More precisely,
\begin{equation}
                    d_j\ge R+j,
 \qquad
                    d_j-t+b\ge w+j+b.
 \label{eq:rank-flat-singleton}
\end{equation}
\end{theorem}

\begin{proof}
Choose a basis of \(C'\) and identify the affine flat with \(\F^s\).
For each active coordinate \(x\in D\setminus G\), agreement with \(u\)
is an affine hyperplane whose nonzero normal is the evaluation functional
\(\ell_x\in(C')^*\).  Every listed point is incident with at least
\[
 h=m-g=d_s-t+b
\]
of these hyperplanes.

Suppose that \(r\) independent incident normals have already been chosen
and span \(W\le(C')^*\).  The coordinates whose normals lie in \(W\) are
common zeros of the \((s-r)\)-dimensional subcode \(W^\perp\).  This
subcode has support at least \(d_{s-r}\); after the common-zero set \(G\)
is removed, at most
\[
 (n-d_{s-r})-z=d_s-d_{s-r}
\]
active normals lie in \(W\).  Hence there are at least
\[
 h-(d_s-d_{s-r})=d_{s-r}-t+b
\]
choices for the next independent incident normal.  Each listed point is
therefore incident with at least
\[
 \prod_{j=1}^{s}(d_j-t+b)
\]
ordered bases of coordinate hyperplanes.  There are at most
\(d_s^{\underline s}\) ordered coordinate bases globally, and an
independent affine-hyperplane basis has at most one common point.  Double
counting proves \eqref{eq:rank-flat-list}.

Finally, if \(V\le C'\) has dimension \(j\) and common-zero set of size
\(z_V\), division by its locator embeds the representing polynomial space
into the polynomials of degree less than \(K-z_V\).  Thus
\(j\le K-z_V\), so \(|\supp(V)|=n-z_V\ge n-K+j=R+j\).  This proves
\eqref{eq:rank-flat-singleton} and positivity.
\end{proof}

\begin{corollary}[Over-budget rank-flat concentration]
\label{cor:rank-flat-concentration}
In \cref{thm:rank-flat-list}, let \(s\ge2\) and let \(B\) be a nonnegative
integer.  Define
\begin{equation}
 \Theta_{s,B}(d_s,b)=
 \max\left\{q\in\mathbb Z_{\ge0}:
 q^{s-1}\le
 \left\lfloor
 \frac{d_s^{\underline s}}
 {(B+1)(d_s-t+b)}
 \right\rfloor
 \right\}.
 \label{eq:rank-flat-theta}
\end{equation}
If \(|\mathcal L_{\mathcal A}(u,m)|>B\), then some proper direction
subcode has dimension \(1\le j<s\) and
\begin{equation}
                     d_j\le t-b+\Theta_{s,B}(d_s,b).
 \label{eq:rank-flat-concentration}
\end{equation}
A row-independent coarse version is obtained by replacing
\(d_s^{\underline s}\) by \(n^{\underline s}\) and
\(d_s-t+b\) by \(w+s\).
\end{corollary}

\begin{proof}
If the list contains at least \(B+1\) points, \eqref{eq:rank-flat-list}
implies
\[
 \prod_{j=1}^{s-1}(d_j-t+b)
 \le
 \left\lfloor
 \frac{d_s^{\underline s}}
 {(B+1)(d_s-t+b)}
 \right\rfloor .
\]
At least one of the \(s-1\) integer factors is at most the integer root in
\eqref{eq:rank-flat-theta}, proving \eqref{eq:rank-flat-concentration}.
The coarse form follows from \(d_s\le n\) and
\eqref{eq:rank-flat-singleton}.
\end{proof}

\begin{theorem}[Common-zero Johnson list compiler]
\label{thm:common-zero-johnson}
In the notation of \cref{thm:rank-flat-list}, put
\[
 N=d_s,
 \qquad h=m-g=d_s-t+b,
 \qquad \alpha=d_s-d_1.
\]
If
\begin{equation}
                         h^2>N\alpha,
 \label{eq:common-zero-johnson-condition}
\end{equation}
then
\begin{equation}
 |\mathcal L_{\mathcal A}(u,m)|
 \le
 \left\lfloor
 \frac{N(h-\alpha)}{h^2-N\alpha}
 \right\rfloor .
 \label{eq:common-zero-johnson}
\end{equation}
Equivalently, write \(z=K-\nu\).  Then \(N=R+\nu\),
\(h=w+\nu+b\), and the generalized Singleton bound gives the universal
specialization
\begin{equation}
 |\mathcal L_{\mathcal A}(u,m)|
 \le
 \left\lfloor
 \frac{(R+\nu)(w+b+1)}
 {(w+\nu+b)^2-(R+\nu)(\nu-1)}
 \right\rfloor
 \label{eq:common-zero-nu}
\end{equation}
whenever the displayed denominator is positive.
\end{theorem}

\begin{proof}
All codewords of the affine flat have the same values on \(G\).  Delete
those coordinates and, for each listed codeword, choose an \(h\)-subset of
its active agreement set.  If two distinct codewords are chosen, their
difference is a nonzero word of \(C'\), has weight at least \(d_1\), and
therefore vanishes on at most \(N-d_1=\alpha\) active coordinates.  The
standard Johnson incidence count now gives
\[
 L(h^2-N\alpha)\le N(h-\alpha),
\]
which proves \eqref{eq:common-zero-johnson}.

If \(z=K-\nu\), then \(N=n-z=R+\nu\) and
\(h=m-(z-b)=w+\nu+b\).  By \eqref{eq:rank-flat-singleton},
\(d_1\ge R+1\), so \(\alpha\le\nu-1\).  The Johnson fraction is
nondecreasing in \(\alpha\) while its denominator is positive.  Replacing
\(\alpha\) by \(\nu-1\) gives \eqref{eq:common-zero-nu}.
\end{proof}

\section{Coset-free codimension-one rank-flat recursion}
\label{sec:codim-one-recursion}
The rank-flat concentration theorem produces a low-support direction
subcode, but a naive partition into its affine cosets introduces the field
size.  The next theorem removes that loss when the support closure has
codimension one.  It counts the whole affine chart at once by two disjoint
resources: bases using an agreement in the new support layer, and bases
whose final coordinate lies in the old support.

\begin{theorem}[Two-resource codimension-one recursion]
\label{thm:codim-one-recursion}
Let
\[
 C=\RS[\F,D,K],\qquad n=|D|,\qquad m=K+w,\qquad t=n-m,
\]
and let \(\mathcal A=c_0+C'\subseteq C\) be an affine codeword flat with
\(\dim C'=j+1\), where \(j\ge1\).  Let \(V<C'\) be a hyperplane, put
\[
 U=\supp(V),\qquad S=\supp(C'),\qquad
 T=S\setminus U,\qquad H=D\setminus S,
\]
and assume \(e:=|T|\ge1\).  Write \(d=|U|\).  Then
\begin{equation}
 V=\{v\in C':\supp(v)\subseteq U\};
 \label{eq:codim-one-saturation}
\end{equation}
in particular every word of \(C'\setminus V\) is nonzero at every
coordinate of \(T\).

Fix a received word \(u\in\F^D\).  Let \(b_0\) be the number of
coordinates in \(H\) on which the common value of \(\mathcal A\)
disagrees with \(u\).  For \(0\le b\le e\), let \(L_b\) count the
members \(c\in\mathcal A\) satisfying
\[
 \operatorname{agr}(u,c)\ge m,
 \qquad
 |\{x\in T:c(x)\ne u(x)\}|=b.
\]
Let \(f_i=d_i(V)\), \(1\le i\le j\), be the generalized Hamming weights
of \(V\), and put
\begin{equation}
 P_b=\prod_{i=1}^{j}(f_i-t+b_0+b),
 \qquad Q=w+1+b_0.
 \label{eq:codim-one-profile}
\end{equation}
Then every factor in \(P_b\) is positive and
\begin{align}
 \sum_{b=0}^{e}L_b(e-b)P_b
 &\le e\,d^{\underline j},
 \label{eq:codim-one-T-resource}\\
 \sum_{b=0}^{e}L_b(Q-e+b)_+P_b
 &\le d^{\underline{j+1}}.
 \label{eq:codim-one-U-resource}
\end{align}
Moreover, with \(L=\sum_bL_b\),
\begin{equation}
 L\le
 \left\lfloor
 \frac{(n-j)d^{\underline j}}{QP_0}
 \right\rfloor .
 \label{eq:codim-one-all-extenders}
\end{equation}
More generally, if \(x,y\ge0\) satisfy
\begin{equation}
 x(e-b)P_b+y(Q-e+b)_+P_b\ge1
 \qquad(0\le b\le e),
 \label{eq:codim-one-dual-feasible}
\end{equation}
then
\begin{equation}
 L\le
 \left\lfloor
 xe\,d^{\underline j}+y\,d^{\underline{j+1}}
 \right\rfloor .
 \label{eq:codim-one-dual-bound}
\end{equation}
Thus the affine cosets of \(V\) are never summed separately.
\end{theorem}

\begin{proof}
For \(x\in T\), every word of \(V\) vanishes at \(x\), whereas some word
of \(C'\) is nonzero there.  Evaluation at \(x\) therefore induces a
nonzero functional on the one-dimensional quotient \(C'/V\).  It is
nonzero on every nonzero quotient class, proving
\eqref{eq:codim-one-saturation}.

Parameterize \(\mathcal A\) by the direction space \(C'\).  Agreement at
a coordinate \(x\) is an affine hyperplane with normal
\(\ell_x\in(C')^*\).  A listed point counted by \(L_b\) has at least
\[
 m-(|H|-b_0)-(e-b)=d-t+b_0+b
\]
agreements in \(U\).  Restrict the corresponding normals to \(V\).
After \(r\) independent restricted normals have been selected, their span
annihilates a \((j-r)\)-dimensional subcode of \(V\), whose support has
size at least \(f_{j-r}\).  Hence at most \(d-f_{j-r}\) normals of
\(U\) lie in the old span, and there are at least
\[
 f_{j-r}-t+b_0+b
\]
choices for the next independent normal.  Every listed point counted by
\(L_b\) is therefore incident with at least \(P_b\) ordered
\(V\)-normal bases.  The generalized Singleton inequalities
\(f_i\ge n-K+i\) show at the same time that all factors in \(P_b\) are
at least \(w+i+b_0+b\).

The span in \((C')^*\) of a lifted \(V\)-normal basis meets the
one-dimensional annihilator \(V^\perp\) trivially.  Every agreeing
coordinate of \(T\) has a nonzero normal in \(V^\perp\), and therefore
extends the basis.  A point in class \(b\) supplies at least
\((e-b)P_b\) ordered bases of this type.  Globally there are at most
\(e\,d^{\underline j}\) ordered choices, and an independent affine
hyperplane basis cuts out at most one point.  This proves
\eqref{eq:codim-one-T-resource}.

Fix now one lifted \(V\)-normal basis and let \(h\in C'\setminus V\)
generate its one-dimensional annihilator.  A coordinate normal extends
the basis exactly when \(h\) is nonzero at that coordinate.  The
nonzero Reed--Solomon word \(h\) has at most \(K-1\) zeros.  All
coordinates of \(H\) are zeros of \(h\), and exactly \(|H|-b_0\) of
them are agreements.  Thus among the at least \(m\) agreement coordinates
of the listed point, at most \(K-1-b_0\) fail to extend the basis.  There
are consequently at least
\[
 m-(K-1-b_0)=Q
\]
extending agreement coordinates.  All \(e-b\) agreeing coordinates of
\(T\) already extend, so at least \((Q-e+b)_+\) extending agreements lie
in \(U\).  Double-counting the resulting ordered \((j+1)\)-tuples inside
\(U\) proves \eqref{eq:codim-one-U-resource}.

The same argument without separating \(U\) from \(T\) gives at least
\(QP_b\ge QP_0\) full bases through every listed point.  There are at
most \((n-j)d^{\underline j}\) ordered choices of a \(V\)-basis and one
additional coordinate, proving \eqref{eq:codim-one-all-extenders}.
Finally multiply \eqref{eq:codim-one-T-resource} by \(x\), multiply
\eqref{eq:codim-one-U-resource} by \(y\), and use
\eqref{eq:codim-one-dual-feasible}; this gives
\eqref{eq:codim-one-dual-bound}.
\end{proof}

\begin{lemma}[Profile interpolation]
\label{lem:codim-one-profile-interpolation}
Let
\[
 P(z)=\prod_{i=1}^{j}(a_i+z),\qquad a_i>0,
\]
and let \(e,Q>0\).  For every integer \(0\le b\le e\),
\begin{equation}
 \frac{(e-b)P(b)}{eP(0)}+
 \frac{(Q-e+b)_+P(b)}{QP(e)}\ge1
 \label{eq:profile-interpolation}
\end{equation}
under either of the following conditions:
\begin{enumerate}[label=\textup{(\alph*)},leftmargin=*]
\item \(e\le Q\);
\item \(e>Q\) and
\begin{equation}
 \left(\frac{P(e)}{P(0)}\right)^{1-Q/e}\ge\frac eQ.
 \label{eq:profile-ratio-condition}
\end{equation}
\end{enumerate}
If \(j\ge2\), \(A=\max_i a_i\), and \((j-1)Q\ge A\), then
\eqref{eq:profile-ratio-condition} holds for every \(e\ge Q\).
\end{lemma}

\begin{proof}
Put \(r=P(e)/P(0)\ge1\).  Since \(\log P\) is concave, with
\(\theta=b/e\) one has
\[
 \frac{P(b)}{P(0)}\ge r^\theta,
 \qquad
 \frac{P(b)}{P(e)}\ge r^{-(1-\theta)}.
\]
If \(e\le Q\), then \((Q-e+b)/Q\ge\theta\), and weighted AM--GM gives
\[
 (1-\theta)r^\theta+\theta r^{-(1-\theta)}\ge1.
\]
This proves \eqref{eq:profile-interpolation} in case \textup{(a)}.

Assume \(e>Q\), put \(c=Q/e\), and first take
\(0\le\theta\le1-c\).  The second summand in
\eqref{eq:profile-interpolation} vanishes, while concavity gives the lower
bound \((1-\theta)r^\theta\).  The logarithm of this expression is
concave in \(\theta\); it is zero at \(0\), and at \(1-c\) it is
nonnegative exactly when \eqref{eq:profile-ratio-condition} holds.

For \(1-c\le\theta\le1\), write \(\alpha=1-\theta\in[0,c]\).  After
multiplication by \(r^\alpha\), the lower bound supplied by log-concavity
is
\[
 \alpha r+1-\frac\alpha c.
\]
Convexity of \(r^\alpha\) on \([0,c]\) gives
\[
 r^\alpha\le1+\frac\alpha c(r^c-1).
\]
The former expression dominates the latter precisely when
\(cr\ge r^c\), which is again
\eqref{eq:profile-ratio-condition}.  This proves case \textup{(b)}.

Finally, if \(A=\max_i a_i\), then
\[
 r\ge(1+e/A)^j.
\]
For \(e\ge Q\), set
\[
 \phi(e)=j(1-Q/e)\log(1+e/A)-\log(e/Q).
\]
One has \(\phi(Q)=0\), and the elementary inequality
\(\log(1+x)\ge x/(1+x)\) gives
\[
 \phi'(e)\ge \frac{j}{A+e}-\frac1e
 =\frac{(j-1)e-A}{e(A+e)}\ge0.
\]
Thus \(\phi(e)\ge0\), which is
\eqref{eq:profile-ratio-condition}.
\end{proof}

\begin{corollary}[MDS-soft codimension-one compiler]
\label{cor:codim-one-mds-soft}
Retain the hypotheses of \cref{thm:codim-one-recursion}, and put
\begin{equation}
 \Pi_b=(d-t+b_0+b)
 \prod_{i=1}^{j-1}(w+i+b_0+b).
 \label{eq:codim-one-mds-profile}
\end{equation}
Suppose the profile interpolation inequality
\eqref{eq:profile-interpolation} holds for \(P=\Pi\), with
\(Q=w+1+b_0\).  Then
\begin{equation}
 L\le
 \left\lfloor
 \frac{d^{\underline j}}{\Pi_0}+
 \frac{d^{\underline{j+1}}}{Q\Pi_e}
 \right\rfloor .
 \label{eq:codim-one-closed}
\end{equation}
Independently of profile interpolation,
\begin{equation}
 L\le
 \left\lfloor
 \frac{(n-j)d^{\underline j}}{Q\Pi_0}
 \right\rfloor .
 \label{eq:codim-one-all-extenders-mds}
\end{equation}
In particular, \eqref{eq:codim-one-closed} is a bound for the whole affine
chart, not a per-coset estimate.
\end{corollary}

\begin{proof}
The generalized Singleton inequalities give \(f_i\ge n-K+i\) for
\(i<j\), while \(f_j=d\).  Hence \(P_b\ge\Pi_b\).  The two resource
inequalities of \cref{thm:codim-one-recursion} remain true after replacing
\(P_b\) by \(\Pi_b\).  Under profile interpolation, the choices
\[
 x=\frac1{e\Pi_0},\qquad y=\frac1{Q\Pi_e}
\]
are feasible in \eqref{eq:codim-one-dual-feasible}, and
\eqref{eq:codim-one-dual-bound} gives
\eqref{eq:codim-one-closed}.  The all-extender estimate follows directly
from \eqref{eq:codim-one-all-extenders}.
\end{proof}



\section{Higher-codimension support-flag recursion}
\label{sec:support-flag-recursion}
The codimension-one theorem removes the field-size loss for a single
support-saturated hyperplane.  The same principle extends to arbitrary
codimension, but one must keep all possible choices of ``new-layer'' and
``old-support'' extenders.  The resulting object is a finite exact linear
program with one resource inequality for each binary flag pattern.  In
particular, no affine coset is ever counted separately.

\begin{definition}[Support saturation]
\label{def:support-saturation}
If $V\le W\le C$, the \emph{support closure of $V$ in $W$} is
\[
 \overline V^{\,W}
 =\{c\in W:\supp(c)\subseteq\supp(V)\}.
\]
We call $V$ \emph{support-saturated in $W$} when
$\overline V^{\,W}=V$.
\end{definition}

\begin{lemma}[Generalized-weight minimizers are support-saturated]
\label{lem:weight-minimizer-saturated}
Let $W$ be a finite-length linear code, and let $V\le W$ have dimension
$j$ and support size $d_j(W)$.  Then $V$ is support-saturated in $W$.
More generally, the generalized Hamming weights satisfy
\[
 d_1(W)<d_2(W)<\cdots<d_{\dim W}(W).
\]
\end{lemma}

\begin{proof}
Let $Y\le W$ have dimension $i+1$ and support size $d_{i+1}(W)$.
Choose a coordinate on which some word of $Y$ is nonzero.  The kernel of
evaluation at that coordinate is an $i$-dimensional subcode whose support
omits that coordinate, so
\[
 d_i(W)\le d_{i+1}(W)-1.
\]
This proves strict increase.  If the support closure of a
$j$-dimensional minimizer $V$ had dimension $k>j$, then
\[
 d_k(W)\le|\supp(V)|=d_j(W),
\]
contradicting strict increase.
\end{proof}

\begin{lemma}[Existence of a support-saturated flag]
\label{lem:support-saturated-flag}
Let $V$ be support-saturated in $W$ and put
$r=\dim W-\dim V$.  There is a flag
\[
 V=V_0<V_1<\cdots<V_r=W,
 \qquad \dim V_\ell=\dim V+\ell,
\]
such that $V_{\ell-1}$ is support-saturated in $V_\ell$ for every
$1\le\ell\le r$.  Consequently, with
\[
 U_\ell=\supp(V_\ell),\qquad
 T_\ell=U_\ell\setminus U_{\ell-1},
\]
every layer $T_\ell$ is nonempty.
\end{lemma}

\begin{proof}
Construct the flag downward.  Suppose $V<X\le W$.  Since $V$ is
support-saturated in $W$, some word of $X\setminus V$ is nonzero at a
coordinate $x\notin\supp(V)$.  Put
\[
 X'=\{c\in X:c(x)=0\}.
\]
Then $V\le X'$, $X'$ is a hyperplane of $X$, and $x\notin\supp(X')$.
If $c\in X$ has support contained in $\supp(X')$, then $c(x)=0$, hence
$c\in X'$.  Thus $X'$ is support-saturated in $X$.  Repeating until the
dimension reaches $\dim V$ gives the claimed flag; the terminal subspace
contains $V$ and has the same dimension, so it is $V$.  A proper
support-saturated inclusion cannot have equal supports, proving
$T_\ell\ne\varnothing$.
\end{proof}

\begin{theorem}[Exact support-flag resource recursion]
\label{thm:support-flag-recursion}
Let
\[
 C=\RS[\F,D,K],\qquad n=|D|,\qquad m=K+w,\qquad t=n-m,
\]
and let $\mathcal A=c_0+C'\subseteq C$ be an affine codeword flat.
Assume that
\[
 V_0<V_1<\cdots<V_r=C'
\]
is a support-saturated flag with
\[
 \dim V_0=j\ge1,
 \qquad \dim V_\ell=j+\ell.
\]
Put
\[
 U_\ell=\supp(V_\ell),\quad d_\ell=|U_\ell|,
 \quad T_\ell=U_\ell\setminus U_{\ell-1},
 \quad e_\ell=|T_\ell|,
\]
and $H=D\setminus U_r$.  Fix a received word $u\in\F^D$, and let
$b_0$ be the number of coordinates of $H$ on which the common value of
$\mathcal A$ disagrees with $u$.

For a profile
$\mathbf b=(b_1,\ldots,b_r)$ with $0\le b_\ell\le e_\ell$, let
$L_{\mathbf b}$ count the words $c\in\mathcal A$ satisfying
\[
 \operatorname{agr}(u,c)\ge m,
 \qquad
 |\{x\in T_\ell:c(x)\ne u(x)\}|=b_\ell
 \quad(1\le\ell\le r).
\]
Let $f_i=d_i(V_0)$, $1\le i\le j$, and define
\begin{align}
 P_{\mathbf b}
 &=\prod_{i=1}^{j}
   \left(f_i-t+b_0+\sum_{q=1}^{r}b_q\right),
 \label{eq:flag-profile-P}\\
 Q_\ell(\mathbf b)
 &=w+1+b_0+\sum_{q>\ell}b_q,
 \label{eq:flag-profile-Q}\\
 A_\ell(\mathbf b)&=e_\ell-b_\ell,
 \qquad
 B_\ell(\mathbf b)
 =\bigl(Q_\ell(\mathbf b)-e_\ell+b_\ell\bigr)_+.
 \label{eq:flag-profile-AB}
\end{align}
For a binary pattern $\sigma\in\{T,U\}^r$, put
\begin{align}
 C_\sigma(\mathbf b)
 &=P_{\mathbf b}
   \prod_{\ell:\sigma_\ell=T}A_\ell(\mathbf b)
   \prod_{\ell:\sigma_\ell=U}B_\ell(\mathbf b),
 \label{eq:flag-profile-resource}\\
 G_\sigma
 &=d_0^{\underline j}
   \prod_{\ell:\sigma_\ell=T}e_\ell
   \prod_{\ell:\sigma_\ell=U}
   \bigl(d_{\ell-1}-(j+\ell-1)\bigr).
 \label{eq:flag-global-resource}
\end{align}
Then every factor in $P_{\mathbf b}$ is positive and, for every
$\sigma\in\{T,U\}^r$,
\begin{equation}
 \boxed{
 \sum_{\mathbf b}L_{\mathbf b}C_\sigma(\mathbf b)
 \le G_\sigma.}
 \label{eq:flag-resource-inequality}
\end{equation}
Consequently, if nonnegative weights $x_\sigma$ satisfy
\begin{equation}
 \sum_{\sigma\in\{T,U\}^r}
 x_\sigma C_\sigma(\mathbf b)\ge1
 \qquad
 \text{for every profile }\mathbf b,
 \label{eq:flag-dual-feasible}
\end{equation}
then, with $L=\sum_{\mathbf b}L_{\mathbf b}$,
\begin{equation}
 \boxed{
 L\le
 \left\lfloor
 \sum_{\sigma\in\{T,U\}^r}x_\sigma G_\sigma
 \right\rfloor.}
 \label{eq:flag-dual-bound}
\end{equation}
The bound is for the whole affine chart and contains no factor depending
on $|\F|$ or on the number of affine cosets of $V_0$.
\end{theorem}

\begin{proof}
Parameterize $\mathcal A$ by $C'$.  Agreement at a coordinate $x$ is an
affine hyperplane with normal $\ell_x\in(C')^*$.  A point of profile
$\mathbf b$ has at least
\[
 d_0-t+b_0+\sum_{q=1}^{r}b_q
\]
agreements in $U_0$.  Applying the ordered-basis argument of
\cref{thm:rank-flat-list} inside $V_0$ shows that it is incident with at
least $P_{\mathbf b}$ ordered coordinate bases whose restrictions form a
basis of $V_0^*$.

Suppose an adapted basis has already been chosen through level
$V_{\ell-1}$.  Its annihilator in $V_\ell$ is one-dimensional; let
$h_\ell\in V_\ell\setminus V_{\ell-1}$ generate it.  At a coordinate
of $T_\ell$, evaluation vanishes on $V_{\ell-1}$ but not on $V_\ell$;
on the one-dimensional quotient $V_\ell/V_{\ell-1}$ it is therefore a
nonzero functional.  Hence $h_\ell$ is nonzero at every coordinate of
$T_\ell$.
Moreover, $h_\ell$ vanishes outside $U_\ell$.  The number of mismatches
outside $U_\ell$ is
\[
 b_0+\sum_{q>\ell}b_q.
\]
Since a nonzero degree-$<K$ word has at most $K-1$ zeros, at most
$K-1-b_0-\sum_{q>\ell}b_q$ agreement coordinates fail to extend the
current basis.  Hence there are at least
$Q_\ell(\mathbf b)$ extending agreement coordinates in total.  The
$A_\ell(\mathbf b)$ agreeing coordinates of $T_\ell$ all extend, so at
least $B_\ell(\mathbf b)$ extending agreement coordinates lie in
$U_{\ell-1}$.

Fix a pattern $\sigma$.  Multiplying the preceding lower bounds over the
flag levels gives at least $C_\sigma(\mathbf b)$ adapted ordered bases
of pattern $\sigma$ through every point of profile $\mathbf b$.  Globally,
the initial ordered basis has at most $d_0^{\underline j}$ choices.  At a
$T$-level there are at most $e_\ell$ coordinate choices.  At a $U$-level,
all $j+\ell-1$ previously selected coordinates lie in $U_{\ell-1}$, and
an extending normal cannot repeat one of them; there are therefore at
most $d_{\ell-1}-(j+\ell-1)$ choices.  An independent affine-hyperplane
basis cuts out at most one parameter point.  Double counting proves
\eqref{eq:flag-resource-inequality}.  Multiplying those inequalities by
$x_\sigma$, summing over $\sigma$, and using
\eqref{eq:flag-dual-feasible} proves \eqref{eq:flag-dual-bound}.
Positivity follows from $f_i\ge n-K+i$.
\end{proof}

\begin{corollary}[All-extender flag bound]
\label{cor:flag-all-extender}
In the notation of \cref{thm:support-flag-recursion}, put
\[
 P_0=\prod_{i=1}^{j}(f_i-t+b_0),
 \qquad Q=w+1+b_0.
\]
Then
\begin{equation}
 \boxed{
 L\le
 \left\lfloor
 \frac{d_0^{\underline j}
       \prod_{\ell=1}^{r}
       \bigl(d_\ell-(j+\ell-1)\bigr)}
      {P_0Q^r}
 \right\rfloor.}
 \label{eq:flag-all-extender}
\end{equation}
In particular, the numerator may be replaced by
$d_0^{\underline j}(n-j)^{\underline r}$.
\end{corollary}

\begin{proof}
Every point has at least $P_0$ initial bases.  At level $\ell$ it has at
least $Q_\ell(\mathbf b)\ge Q$ extending agreement coordinates, without
separating the two resources.  The global number of adapted extensions at
that level is at most $d_\ell-(j+\ell-1)$.  Multiplication and double
counting give the first bound.  Since $d_\ell\le n$, each factor is at most
$n-j-\ell+1$, giving the second statement.
\end{proof}

\begin{lemma}[One-level interpolation defect]
\label{lem:flag-interpolation-defect}
For $e,Q>0$ and $\rho\ge1$, define
\begin{equation}
 \eta(e,Q,\rho)=
 \begin{cases}
 1,&e\le Q,\\
 \displaystyle
 \min\left\{1,\frac Qe\rho^{1-Q/e}\right\},&e>Q.
 \end{cases}
 \label{eq:flag-eta}
\end{equation}
Then, for every real $0\le b\le e$,
\begin{equation}
 \rho^{b/e}
 \left(
  \frac{e-b}{e}
  +\frac{(Q-e+b)_+}{Q\rho}
 \right)
 \ge\eta(e,Q,\rho).
 \label{eq:flag-one-level-interpolation}
\end{equation}
Thus the right side is $1$ precisely in the profile-interpolation range
used in \cref{lem:codim-one-profile-interpolation}.
\end{lemma}

\begin{proof}
If $e\le Q$, the proof is the weighted AM--GM argument of
\cref{lem:codim-one-profile-interpolation}.  Suppose $e>Q$, put
$c=Q/e$ and $\theta=b/e$.  On $0\le\theta\le1-c$, the logarithm of the
left side is
$\log(1-\theta)+\theta\log\rho$, a concave function, so its minimum is
at an endpoint and equals at least
$\min\{1,c\rho^{1-c}\}$.  On $1-c\le\theta\le1$, put
$\alpha=1-\theta\in[0,c]$.  The left side becomes
\[
 \rho^{-\alpha}\left(\alpha\rho+1-\frac\alpha c\right).
\]
Its logarithm is a linear function plus the logarithm of a positive
affine function, hence is concave.  Its endpoint values are again
$1$ and $c\rho^{1-c}$.  This proves the claim.
\end{proof}

\begin{corollary}[Tensorized support-flag certificate]
\label{cor:flag-tensor-certificate}
Retain the notation of \cref{thm:support-flag-recursion}.  Put
\[
 E=\sum_{\ell=1}^{r}e_\ell,
 \qquad
 P(z)=\prod_{i=1}^{j}(f_i-t+b_0+z),
 \qquad
 \rho=\frac{P(E)}{P(0)},
\]
and define
\[
 \rho_\ell=\rho^{e_\ell/E},
 \qquad
 \eta_\ell=\eta(e_\ell,Q,\rho_\ell),
 \qquad Q=w+1+b_0.
\]
Then
\begin{equation}
 \boxed{
 L\le
 \left\lfloor
 \frac{d_0^{\underline j}}
      {P(0)\prod_{\ell=1}^{r}\eta_\ell}
 \prod_{\ell=1}^{r}
 \left(
 1+\frac{d_{\ell-1}-(j+\ell-1)}{Q\rho_\ell}
 \right)
 \right\rfloor.}
 \label{eq:flag-tensor-bound}
\end{equation}
If every $\eta_\ell=1$, this is a closed higher-codimension analogue of
\eqref{eq:codim-one-closed}.  For $r=1$ it specializes to the
codimension-one two-resource certificate, with the explicit defect factor
when profile interpolation fails.
\end{corollary}

\begin{proof}
Let $B=\sum b_\ell$.  Concavity of $\log P$ on $[0,E]$ gives
\[
 \frac{P(B)}{P(0)}
 \ge\rho^{B/E}
 =\prod_{\ell=1}^{r}\rho_\ell^{b_\ell/e_\ell}.
\]
Also $Q_\ell(\mathbf b)\ge Q$, so
$B_\ell(\mathbf b)\ge(Q-e_\ell+b_\ell)_+$.  By
\cref{lem:flag-interpolation-defect},
\[
 \frac{P_{\mathbf b}}{P(0)}
 \prod_{\ell=1}^{r}
 \left(
  \frac{A_\ell(\mathbf b)}{e_\ell}
  +\frac{B_\ell(\mathbf b)}{Q\rho_\ell}
 \right)
 \ge\prod_{\ell=1}^{r}\eta_\ell.
\]
Use the tensor-product dual weights
\[
 x_\sigma=
 \frac1{P(0)\prod_\ell\eta_\ell}
 \prod_{\ell:\sigma_\ell=T}\frac1{e_\ell}
 \prod_{\ell:\sigma_\ell=U}\frac1{Q\rho_\ell}.
\]
They satisfy \eqref{eq:flag-dual-feasible}; substituting them into
\eqref{eq:flag-dual-bound} and factoring the sum over binary patterns
gives \eqref{eq:flag-tensor-bound}.
\end{proof}

\begin{remark}[What the flag theorem closes]
\label{rem:flag-theorem-scope}
The uncontrolled field-size coset sum is now removed in every codimension.
For a concrete finite chart, \eqref{eq:flag-resource-inequality} is an
exact rational linear-program certificate, while
\eqref{eq:flag-tensor-bound} is a closed analytic certificate.  What can
still fail is numerical payment: a low-mismatch flag with several short
layers may have a certified bound above the row budget.  Such a flag is a
specific finite residual atom, not an unproved multiplicity principle.
\end{remark}



\begin{proposition}[Counterexample to the affine-span transverse MCA compiler]
\label{prop:affine-span-mca-counterexample}
The affine-span MCA bound formerly printed here is false even under the
direction-separation hypothesis.  There is a received line for
\[
 C=\RS[\F_{1009},\{0,\ldots,99\},1],
 \qquad (n,K,m,w,s)=(100,1,21,20,1),
\]
with \(31\) distinct selected slopes and explanations in the
one-dimensional constant code, while
\[
 \max_{c\in C}\operatorname{agr}(r_1,c)=20<21=m
\]
and the claimed bound was
\[
 \left\lfloor
 \max\left\{
 \frac{100^{\underline2}}{21\cdot20},
 \frac{100^{\underline2}}{20^{\overline2}}
 \right\}
 \right\rfloor=23<31.
\]
Every selected maximal agreement support has size exactly \(21\) and is
same-support pair-noncontained.
\end{proposition}

\begin{proof}
Partition the domain into sets \(C_0,E,T,W\) of sizes \(20,30,21,29\).
Put \((r_0,r_1)=(0,0)\) on \(C_0\).  Index \(E\) by
\(1\le i\le30\) and put \((r_0,r_1)=(-i^2,i)\) at its \(i\)-th point.
For slope \(i\), the zero codeword then agrees exactly on \(C_0\) and
that point of \(E\).

On \(T\), put \(r_0=1\) and choose \(21\) distinct fresh nonzero direction
values avoiding \(1+i r_1=0\) for every \(1\le i\le30\).  At slope zero,
the constant-one codeword agrees exactly on \(T\).  On \(W\), choose fresh
direction values and choose each base value outside
\(\{0,1\}\cup\{-i r_1:1\le i\le30\}\).  All choices exist in
\(\F_{1009}\), and they remove every unintended selected agreement.

The direction is zero on \(C_0\) and takes distinct nonzero values elsewhere,
so its maximum agreement with a constant is exactly \(20\).  The selected
explanations are zero and one, hence have affine span one.  On a
zero-explanation support, the core forces a simultaneous constant pair to be
\((0,0)\), but the remaining direction value is nonzero.  On \(T\), the
direction values are distinct.  Thus every support is pair-noncontained.

The proof failure is the asserted proper-subspace occupancy bound for the
incident normals.  On each zero-explanation support, \(20\) normals lie on
one line and one is transverse.  There are only \(40\) ordered bases, not the
claimed lower count \(m w=420\).  The complete deterministic reconstruction
is in
\texttt{experimental/verify\_mca\_affine\_span\_incidence\_counterexample\_v1.py}.
\end{proof}

\begin{remark}[Scope of the route cut]
The counterexample retracts only the MCA incidence compiler and statements
that use its ordered-basis denominator.  It does not refute the ordinary
affine-span list theorem, common-core cancellation, the directional Johnson
compiler, codeword-direction gauge equivalence, or the selector-free
all-LineRay error-affine-core set-pair theorem, whose proof uses zero-mask
restriction and the Bollobas set-pair inequality instead.
\end{remark}

\begin{theorem}[Pair-noncontained proper-subspace MCA compiler]
\label{thm:proper-subspace-mca}
Let \(C=\RS[\F,D,K]\) have length \(n\), let
\(r_\gamma=r_0+\gamma r_1\), put \(w=m-K>0\), and let
\(\mathcal A=c_0+C'\subseteq C\) have \(\dim C'=s\), where
\(1\le s\le K\).  Let \(Z\subseteq\F\) be a set of distinct slopes such
that for every \(\gamma\in Z\), some \(c_\gamma\in\mathcal A\) has
agreement at least \(m\), and its maximal agreement support is
same-support pair-noncontained.  Define
\[
 e=\min_{b\in C}\wt(r_1-b),\qquad
 L=\max\{1,e-(n-m)\}.
\]
Then
\begin{equation}
 |Z|\le\left\lfloor\frac1L\max\left\{
 \frac{n^{\underline{s+1}}}
      {m(w+1)^{\overline{s-1}}},
 \frac{(n-K+s)^{\underline{s+1}}}
      {(w+1)^{\overline s}}
 \right\}\right\rfloor .
 \label{eq:proper-subspace-mca}
\end{equation}
For \(s=1\), the empty rising product in the first denominator is one.
\end{theorem}

\begin{proof}
Choose a basis \(c_1,\ldots,c_s\) of \(C'\).  As in the rejected proof,
the agreement hyperplane at coordinate \(x\) has normal
\[
 v_x=(r_1(x),-c_1(x),\ldots,-c_s(x))\in\F^{s+1}.
\]
The normals incident with every selected parameter point span
\(\F^{s+1}\).  Indeed, an annihilator with zero slope coordinate would
give a nonzero codeword with at least \(m>K-1\) roots.  An annihilator with
nonzero slope coordinate would give simultaneous codeword explanations for
\(r_0\) and \(r_1\) on the same maximal support, contrary to
pair noncontainment.

Let \(W<\F^{s+1}\) have dimension \(j<s\), and put
\(X_W=\{x:v_x\in W\}\).  If the slope coordinate vanishes on
\(W^\perp\), then \(s+1-j\) independent codewords vanish on \(X_W\).
Otherwise the kernel of that coordinate supplies \(s-j\) independent
codewords.  The MDS generalized-weight formula gives in either case
\[
 |X_W|\le K-s+j.
\]
Consequently, after \(j=1,\ldots,s-1\) independent incident normals have
been selected, there are at least \(w+s-j\) choices outside their span.
These factors multiply to \((w+1)^{\overline{s-1}}\).

After \(s\) independent normals have been selected, local full rank leaves
at least one final choice.  If their hyperplane has an annihilator with
nonzero slope coordinate, it contains at most \(n-e\) coordinate normals
and leaves at least \(e-(n-m)\) incident choices.  If that slope coordinate
is zero, the ordinary root bound leaves at least \(w+1\).  Interpolation on
\(K\) coordinates gives \(e\le n-K\), so every case leaves at least \(L\)
final choices.

Let \(z\) be the number of zero normals and let \(g\le z\) be those whose
agreement hyperplane contains every parameter point.  Common-zero MDS gives
\(z\le K-s\).  Every selected point is therefore incident with at least
\[
 (m-g)(w+1)^{\overline{s-1}}L
\]
ordered normal bases, while there are at most
\((n-z)^{\underline{s+1}}\) ordered independent coordinate tuples.
For fixed \(z\), the resulting ratio is largest at \(g=z\).  Its successive
ratio in \(z\) has the sign of
\(n-(s+1)m+sz\), so its maximum on \(0\le z\le K-s\) is at an endpoint.
The endpoints \(z=0\) and \(z=K-s\) are precisely the two terms in
\eqref{eq:proper-subspace-mca}.
\end{proof}

\begin{theorem}[Support-transverse affine-span MCA compiler]
\label{thm:support-transverse-affine-span-mca}
Retain the code, received line, affine explanation space, and parameters of
\cref{thm:proper-subspace-mca}, with \(n\ge m=K+w\), \(w>0\), and
\(1\le s\le K\).  Let \(\varnothing\ne Z\subseteq\F\) be a set of
distinct slopes.  Suppose that for every \(\gamma\in Z\) there are a
codeword \(c_\gamma\in\mathcal A\) and an exact \(m\)-set
\(S_\gamma\subseteq D\) such that \(r_\gamma=c_\gamma\) on
\(S_\gamma\), while no pair in \(C^2\) explains \((r_0,r_1)\)
simultaneously on that same set.  Define
\begin{equation}
 \theta_0=
 \min_{\substack{\gamma\in Z\\ b\in C'}}
 |\{x\in S_\gamma:r_1(x)\ne b(x)\}|,
 \qquad
 \theta=\min\{\theta_0,w+1\}.
 \label{eq:support-local-mca-theta}
\end{equation}
Then \(\theta\ge1\), and
\begin{equation}
 |Z|\le
 \left\lfloor
 \max\left\{
 \frac{n^{\underline{s+1}}}
      {m\,\theta\,(w+1)^{\overline{s-1}}},
 \frac{(n-K+s)^{\underline{s+1}}}
      {\theta\,(w+1)^{\overline s}}
 \right\}
 \right\rfloor .
 \label{eq:support-local-proper-subspace-mca}
\end{equation}
The first rising product is empty when \(s=1\).  If \(e,L\) are the
global direction-coset quantities of \cref{thm:proper-subspace-mca}, then
\(\theta\ge L\).  Thus \eqref{eq:support-local-proper-subspace-mca}
refines that theorem whenever the selected supports have more local
transversality than is forced by the global direction distance.
\end{theorem}

\begin{proof}
Choose a basis \(c_1,\ldots,c_s\) of \(C'\) and use the same incident
normal count as in \cref{thm:proper-subspace-mca}.  The lower-dimensional
proper-normal-subspace argument is unchanged and supplies the factors
\((w+1)^{\overline{s-1}}\).  At the final step, let
\((\delta,\mu)\) span the annihilator of the selected \(s\)-space.  If
\(\delta=0\), the ordinary root bound leaves at least
\(w+1\ge\theta\) incident choices outside it.  If \(\delta\ne0\), put
\(b=\delta^{-1}\sum_i\mu_i c_i\in C'\).  Every point of the actual
support \(S_\gamma\) on which \(r_1\ne b\) has normal outside the
selected space, so \eqref{eq:support-local-mca-theta} supplies at least
\(\theta\) final choices.  Pair noncontainment gives \(\theta_0\ge1\):
if \(r_1=b\) on \(S_\gamma\), then
\((c_\gamma-\gamma b,b)\) explains \((r_0,r_1)\) there.

Keeping the zero-normal parameters $z,g$ from the preceding proof, every
selected point therefore owns at least
\((m-g)\theta(w+1)^{\overline{s-1}}\) ordered bases, while there are at
most $(n-z)^{\underline{s+1}}$ ordered independent coordinate tuples.
For fixed $z$, the ratio is largest at $g=z$; its successive ratio in
$z$ again has the sign of $n-(s+1)m+sz$.  The two endpoints
$z=0,K-s$ give \eqref{eq:support-local-proper-subspace-mca}.
Finally, for every \(b\in C'\subseteq C\),
\[
 |\{x\in S_\gamma:r_1(x)\ne b(x)\}|
\ge \wt(r_1-b)-(n-m)\ge e-(n-m).
\]
Moreover \(e\le n-K\), so \(e-(n-m)\le w\).  If this difference is at
most one, then $L=1\le\theta$; otherwise the preceding inequality and
the cap $w+1$ give \(\theta\ge e-(n-m)=L\).  Hence \(\theta\ge L\).
\end{proof}

\begin{theorem}[Support-wise MCA error-rank gauge]
\label{thm:mca-error-rank-gauge}
Let \(C\le\F^D\) be a linear code, let
\(r_\gamma=r_0+\gamma r_1\), and let \(Z\) contain at least two distinct
slopes.  For every \(\gamma\in Z\), choose \(h_\gamma\in C\) and a
support \(S_\gamma\) on which \(r_\gamma=h_\gamma\), with no
simultaneous \(C^2\)-explanation of \((r_0,r_1)\) on that same support.
Put
\[
 e_\gamma=r_\gamma-h_\gamma,
 \qquad
 a=\dim\Span\{e_\gamma-e_{\gamma_0}:\gamma\in Z\}.
\]
Then \(a\ge1\), and there is a codeword \(b\in C\) such that the
two-way gauge
\begin{equation}
 r'_0=r_0,\qquad r'_1=r_1-b,
 \qquad h'_\gamma=h_\gamma-\gamma b
 \label{eq:mca-error-rank-gauge}
\end{equation}
preserves every finite slope, error word, agreement support, and
same-support containment predicate, while
\[
 \dim\operatorname{Aff}\{h'_\gamma:\gamma\in Z\}=a-1.
\]
This is a reversible codeword translation, not identity of the received-line
representative.
\end{theorem}

\begin{proof}
First \(r_1\notin C\).  Otherwise, on every selected support the pair
\((h_\gamma-\gamma r_1,r_1)\) would explain \((r_0,r_1)\), contrary
to the hypothesis.  Fix \(\gamma_0\), and consider
\[
 E=\Span\{(\gamma-\gamma_0,h_\gamma-h_{\gamma_0}):\gamma\in Z\}
 \le\F\oplus C.
\]
The map \((\delta,c)\mapsto\delta r_1-c\) is injective on \(E\): a
kernel element with \(\delta\ne0\) would put \(r_1\) in \(C\), while
\(\delta=0\) forces \(c=0\).  Its image is the affine error-difference
space, so \(\dim E=a\).  The slope projection on \(E\) is nonzero;
choose \((1,b)\in E\).  Under \eqref{eq:mca-error-rank-gauge}, the
explanation differences are the image of \(E\) under
\((\delta,c)\mapsto c-\delta b\), namely the kernel of the slope
projection.  Their dimension is \(a-1\).

The identity
\[
 r_0+\gamma(r_1-b)-(h_\gamma-\gamma b)
 =r_0+\gamma r_1-h_\gamma
\]
proves error and support preservation.  Subtracting \(b\) from the second
codeword of a simultaneous explanation, and adding it back for the inverse,
proves preservation of same-support containment.
\end{proof}

\begin{remark}[Relation to the full-explanation lifted-rank split]
When the selected explanations have full affine rank (K), the lifted pair
space (E) above has dimension (K) or (K+1).  The gauge theorem gives
explanation rank (K-1) in the first branch and (K) in the second.  Thus it
is the arbitrary-rank existence adapter underlying the two numerical cases
of \cref{thm:full-explanation-lifted-rank-dichotomy} below.  That theorem is
sharper at full rank: it classifies every gauge, rather than selecting one.
\end{remark}

\begin{remark}[Exact first shortened-row calibration]
For a shortened row \((n,K,m)=(R+K,K,d+K)\), the last factor is
\(L=\max\{1,e-(R-d)\}\le d\).  At the first residual dimensions from the
common-core staircase, exact integer evaluation gives
\[
\begin{array}{c|c|c|c}
\text{row}&K&\text{paid for every }e&\text{additional paid walls}\\ \hline
\text{KoalaBear}&14&s\le9&
s=10,11,12,13:\ e\ge981108,981153,981861,992852\\
\text{Mersenne-31}&6&s\le1&
s=2,3,4,5:\ e\ge981144,981363,984779,1037876
\end{array}
\]
At top rank, the required factors are \(182530>d=67472\) and
\(882143>d=67448\), respectively, so this compiler does not pay those
cells.  The exact replay is in
\texttt{experimental/verify\_mca\_proper\_subspace\_occupancy\_compiler\_v1.py}.
\end{remark}

\begin{theorem}[Sparse-direction punctured Johnson MCA profile]
\label{thm:sparse-direction-punctured-johnson-mca}
Let \(C=\RS[\F,D,K]\) have
\[
  N=R+K,\qquad m=d+K,
\]
and let \(r_\gamma=r_0+\gamma r_1\).  Suppose
\[
  r_1=b+q,\qquad b\in C,\qquad
  E=\operatorname{supp}(q),\qquad |E|=e<d.
\]
For every selected slope, assume an exact size-\(m\), same-support
pair-noncontained agreement witness with explanation \(c_\gamma\), and put
\(a_\gamma=c_\gamma-\gamma b\).  For each distinct selected transformed
explanation define
\[
 h_a=m-|\{x\in D\setminus E:a(x)=r_0(x)\}|.
\]
Then \(1\le h_a\le e\), and one deficit-\(h\) explanation owns at most
\(\lfloor e/h\rfloor\) selected slopes.

Assume
\begin{equation}
  \Delta_e=(m-e)^2-(N-e)(K-1)>0.
  \label{eq:punctured-johnson-positive}
\end{equation}
Put \(J_0=0\) and, for \(1\le h\le e\),
\begin{equation}
 J_h=\left\lfloor
 \frac{(N-e)(m-h-K+1)}{(m-h)^2-(N-e)(K-1)}
 \right\rfloor .
 \label{eq:punctured-johnson-cap}
\end{equation}
The selected slope set \(Z\) satisfies
\begin{equation}
 |Z|\le \sum_{h=1}^e(J_h-J_{h-1})\left\lfloor\frac eh\right\rfloor
 \le (e-1)J_{\lfloor e/2\rfloor}+J_e.
 \label{eq:punctured-johnson-profile}
\end{equation}
For \(e=1\), the term indexed by zero in the last expression is zero.
\end{theorem}

\begin{proof}
Outside \(E\), the direction equals \(b\).  If \(a\) agreed with \(r_0\)
on at least \(m\) outside coordinates, an \(m\)-subset there would be
simultaneously explained by \((a,b)\), contrary to same-support
pair noncontainment.  Hence \(h_a\ge1\).  A size-\(m\) witness can use at
most \(e\) coordinates of \(E\), so \(h_a\le e\).

For \(x\in E\), agreement at slope \(\gamma\) is equivalent to
\[
  \frac{a(x)-r_0(x)}{q(x)}=\gamma.
\]
These ratio fibers are disjoint.  A deficit-\(h\) explanation needs at
least \(h\) inside agreements, and therefore owns at most
\(\lfloor e/h\rfloor\) slopes.

Puncture \(E\) and fix \(h\).  Each explanation of deficit at most \(h\)
has an agreement set of size at least \(A=m-h\) in length \(n'=N-e\).
Agreement sets of two distinct degree-\(<K\) polynomials intersect in at
most \(K-1\) coordinates.  If there are \(L\) explanations and \(s_x\)
is the incidence multiplicity at coordinate \(x\), then
\[
 \sum_xs_x\ge LA,
 \qquad
 \sum_x\binom{s_x}{2}\le(K-1)\binom L2.
\]
Cauchy--Schwarz gives
\[
 \sum_x\binom{s_x}{2}\ge
 \frac12\left(\frac{L^2A^2}{n'}-LA\right).
\]
Since \(A^2>n'(K-1)\), comparison yields \(L\le J_h\).  Condition
\eqref{eq:punctured-johnson-positive} is the weakest, \(h=e\), denominator and hence makes
every \(J_h\) legal.  The Johnson fractions are nondecreasing in \(h\).
Applying their cumulative caps to the nonincreasing owner weights proves
the first inequality in \eqref{eq:punctured-johnson-profile}.

Finally put \(u=\lfloor e/2\rfloor\).  Bound all weights with \(h\le u\)
by \(e\); every weight with \(h>u\) equals one.  The two telescoping blocks
give
\[
 eJ_u+(J_e-J_u)=(e-1)J_u+J_e.
\]
\end{proof}

\begin{remark}[Exact sparse-direction support walls]
At the first shortened KoalaBear and Mersenne-31 rows, exact integer
evaluation gives
\[
\begin{array}{c|c|c|c|c}
\text{row}&K&\text{last paid }e&\Delta_e&
 (e-1)J_{\lfloor e/2\rfloor}+J_e\\ \hline
\text{KoalaBear}&14&63908&1218&4607583\\
\text{Mersenne-31}&6&65236&2794&2605443
\end{array}
\]
against budgets \(274980728111395087\) and \(16777215\).  At the adjacent
supports the denominators are \(-5924\) and \(-1636\), so this theorem
makes no claim there.  The exact scan is in
\texttt{experimental/verify\_mca\_sparse\_direction\_punctured\_johnson\_profile\_v1.py}.
\end{remark}

\begin{theorem}[Near-Johnson centered-Gram list compiler]
\label{thm:near-johnson-centered-gram}
Let \(S_1,\ldots,S_L\) be distinct \(A\)-subsets of an \(n\)-set such
that \(|S_i\cap S_j|\le c\) for \(i\ne j\), where \(A>c\).  Put
\[
  g=nc-A^2\ge0,
  \qquad
  G=(A-c)^2-cg.
\]
If \(G>0\), then
\begin{equation}
  L\le\left\lfloor\frac{nA(A-c)}{G}\right\rfloor.
  \label{eq:near-johnson-centered-gram}
\end{equation}
\end{theorem}

\begin{proof}
Let \(B\) be the \(L\)-by-\(n\) incidence matrix and put
\[
  H=BB^{\mathsf T}-c\mathbf1\mathbf1^{\mathsf T}.
\]
Every row of \(B\) has sum \(A\), so
\(\mathbf1=A^{-1}B\mathbf1_n\).  The images of both summands defining
\(H\) therefore lie in \(\operatorname{col}(B)\), and
\begin{equation}
  \operatorname{rank}(H)\le\operatorname{rank}(B)\le n.
  \label{eq:centered-gram-rank}
\end{equation}

Write \(\delta_{ij}=c-|S_i\cap S_j|\).  Then
\(0\le\delta_{ij}\le c\) is integral and
\[
 \operatorname{tr}(H)=L(A-c),\qquad
 \operatorname{tr}(H^2)=L(A-c)^2+2\sum_{i<j}\delta_{ij}^2.
\]
If \(s_x\) is the incidence multiplicity at coordinate \(x\), then
Cauchy--Schwarz gives
\[
 2\sum_{i<j}|S_i\cap S_j|
 =\sum_xs_x^2-LA
 \ge \frac{L^2A^2}{n}-LA.
\]
Thus, for \(\Delta=\sum_{i<j}\delta_{ij}\),
\[
 2\Delta\le L\left(\frac{Lg}{n}+A-c\right).
\]
Since \(\delta_{ij}^2\le c\delta_{ij}\),
\[
 \operatorname{tr}(H^2)
 \le L\left((A-c)^2+c(A-c)+\frac{cLg}{n}\right).
\]
The symmetric trace-rank inequality
\(\operatorname{rank}(H)\ge\operatorname{tr}(H)^2/
\operatorname{tr}(H^2)\), together with
\eqref{eq:centered-gram-rank}, yields
\[
  L\bigl((A-c)^2-cg\bigr)\le nA(A-c),
\]
which is \eqref{eq:near-johnson-centered-gram}.
\end{proof}

\begin{corollary}[Sparse-direction near-Johnson MCA payment]
\label{cor:sparse-direction-near-johnson-mca}
Retain the hypotheses of
\cref{thm:sparse-direction-punctured-johnson-mca}, put
\[
 u=\lfloor e/2\rfloor,
 \qquad n=N-e,
 \qquad c=K-1,
 \qquad A=m-e,
\]
and assume
\[
 (m-u)^2-nc>0,qquad
 g=nc-A^2\ge0,qquad
 G=(A-c)^2-cg>0.
\]
Define
\[
 J_u=\left\lfloor\frac{n(m-u-K+1)}{(m-u)^2-nc}\right\rfloor,
 \qquad
 Q_e=\left\lfloor\frac{nA(A-c)}G\right\rfloor.
\]
Then
\begin{equation}
  |Z|\le(e-1)J_u+Q_e.                              \label{eq:near-johnson-mca}
\end{equation}
\end{corollary}

\begin{proof}
Let \(N_u\) be the number of transformed explanations with outside deficit
at most \(u\), and let \(N_e\) be their total number.  The deficit-owner
bound gives
\[
 |Z|\le eN_u+(N_e-N_u)=(e-1)N_u+N_e.
\]
The ordinary Johnson argument on the punctured row gives \(N_u\le J_u\).
For each explanation counted by \(N_e\), choose exactly \(A=m-e\) outside
agreement coordinates.  Distinct degree-\(<K\) explanations give distinct
blocks with intersections at most \(c=K-1\), so
\cref{thm:near-johnson-centered-gram} gives \(N_e\le Q_e\).
\end{proof}

\begin{remark}[Exact centered-Gram support walls]
Together with the preceding Johnson prefix, exact integer evaluation of
\eqref{eq:near-johnson-mca} expands the low-support walls to
\[
\begin{array}{c|c|c|c}
\text{row}&K&\text{last paid }e&\text{endpoint bound}\\ \hline
\text{KoalaBear}&14&64037&198047217\\
\text{Mersenne-31}&6&65418&16759641
\end{array}
\]
against budgets \(274980728111395087\) and \(16777215\).  At KoalaBear
\(e=64038\), the centered-Gram denominator is \(-36911\).  At Mersenne
\(e=65419\), it remains positive, but the valid bound \(18212004\) exceeds
budget.  The exact replay is in
\path{experimental/verify_mca_sparse_direction_near_johnson_gram_rank_v1.py}.
\end{remark}

\begin{theorem}[Mean-centered Gram list compiler]
\label{thm:mean-centered-gram-list}
Let \(S_1,\ldots,S_L\) be distinct \(A\)-subsets of an \(n\)-set with
\(|S_i\cap S_j|\le c\) for \(i\ne j\), where \(A>c\).  Put
\[
 g=nc-A^2,
 \qquad
 T=(n-A)^2-(n-1)g.
\]
If
\[
 g\ge0,
 \qquad
 2A^2\ge nc,
 \qquad
 T>0,
\]
then
\begin{equation}
 L\le\left\lfloor
 \frac{(n-1)n^2(A-c)}{AT}
 \right\rfloor.
 \label{eq:mean-centered-gram-list}
\end{equation}
\end{theorem}

\begin{proof}
Let \(B\) be the incidence matrix and put
\[
 P=I_n-\frac1nJ_n,
 \qquad
 H=BPB^{\mathsf T}.
\]
Then \(H\) is positive semidefinite and has rank at most \(n-1\).  Set
\(p=A^2/n\).  The diagonal entries are \(A-p\), while an off-diagonal
entry is
\[
 x_{ij}=|S_i\cap S_j|-p\in[-p,c-p].
\]
Convexity gives the endpoint-chord inequality
\[
 x_{ij}^2\le(c-2p)x_{ij}+p(c-p).
\]
The assumption \(2A^2\ge nc\) makes the chord slope nonpositive.  Since
\(H\) is positive semidefinite,
\[
 2\sum_{i<j}x_{ij}\ge-L(A-p).
\]
With \(C=p(c-p)=A^2g/n^2\), the preceding two displays imply
\[
 \operatorname{tr}(H^2)\le CL^2+A(A-c)L.
\]
Now \(\operatorname{tr}(H)=L(A-p)\).  The PSD trace-rank inequality gives
\[
 L\bigl((A-p)^2-(n-1)C\bigr)\le(n-1)A(A-c).
\]
Finally,
\[
 (A-p)^2-(n-1)C=\frac{A^2T}{n^2},
\]
which proves \eqref{eq:mean-centered-gram-list}.
\end{proof}

\begin{corollary}[Exact mean-centered sparse-direction profile]
\label{cor:mean-centered-sparse-direction-profile}
Retain the sparse-direction notation of
\cref{thm:sparse-direction-punctured-johnson-mca}.  For each deficit
threshold \(1\le h\le e\), put \(A_h=m-h\) and let \(C_h\) be the
ordinary Johnson cap when \(A_h^2>(N-e)(K-1)\), and otherwise the cap in
\cref{thm:mean-centered-gram-list} whenever its three hypotheses hold.
Assume every \(C_h\) is defined, and put
\[
 B_0=0,
 \qquad
 B_h=\min_{h\le v\le e}C_v.
\]
Then
\begin{equation}
 |Z|\le\sum_{h=1}^e(B_h-B_{h-1})
                    \left\lfloor\frac eh\right\rfloor.
 \label{eq:mean-centered-sparse-direction-profile}
\end{equation}
\end{corollary}

\begin{proof}
Let \(N_h\) count distinct transformed explanations with outside deficit at
most \(h\).  Puncturing the direction support and applying the appropriate
ordinary-list theorem gives \(N_h\le C_h\).  For every \(v\ge h\), also
\(N_h\le N_v\le C_v\), hence \(N_h\le B_h\).  The suffix minima are
nondecreasing.  Apply the deficit owner weight \(\lfloor e/h\rfloor\) and
sum by parts.
\end{proof}

\begin{remark}[Exact mean-centered support walls]
The complete exact profile in
\eqref{eq:mean-centered-sparse-direction-profile} expands the walls to
\[
\begin{array}{c|c|c|c}
\text{row}&K&\text{last paid }e&\text{endpoint profile}\\ \hline
\text{KoalaBear}&14&64047&181731868\\
\text{Mersenne-31}&6&65454&16101127
\end{array}
\]
At KoalaBear \(e=64048\), the endpoint denominator is
\(T=-1499457466\).  At Mersenne \(e=65455\), every cap remains defined,
but the exact profile is \(17120123>16777215\).  The replay is in
\path{experimental/verify_mca_sparse_direction_mean_centered_gram_profile_v1.py}.
\end{remark}

\begin{theorem}[Terminal-deficit affine-line cap]
\label{thm:terminal-deficit-affine-line}
Retain the sparse-direction notation of
\cref{thm:sparse-direction-punctured-johnson-mca}, assume \(e\ge K\), and
put
\[
 n=N-e,
 \qquad A=m-e,
 \qquad c=K-1.
\]
If \(A>c\), then the number \(L_e\) of transformed explanations that own a
selected slope and have exact outside deficit \(e\) satisfies
\begin{equation}
 L_e\le\left\lfloor\frac{n-c}{A-c}\right\rfloor.
 \label{eq:terminal-deficit-affine-line}
\end{equation}
\end{theorem}

\begin{proof}
An exact-deficit-\(e\) explanation has exactly \(A=m-e\) outside
agreements.  To own a selected slope it must therefore agree at all \(e\)
coordinates of the direction support \(E\).  If two such explanations have
slopes \(\gamma\ne\delta\), their normalized difference is a nonzero
degree-\(<K\) codeword \(p\) whose restriction to \(E\) is the gauged
direction \(q|_E\).  Every further terminal explanation gives the same
normalized difference: the two codewords agree on \(E\), and evaluation on
\(E\) is injective because \(e\ge K\).  Thus all terminal explanations lie
on one affine codeword line.

Outside \(E\), let \(G\) be the common agreement core on the zeros of
\(p\), and write \(g=|G|\le K-1=c\).  Away from \(G\), a coordinate can be
an agreement for at most one affine-line parameter.  Since each terminal
explanation has \(A-g\) agreements away from the core,
\[
 L_e(A-g)\le n-g.
\]
The ratio \((n-g)/(A-g)\) increases with \(g\), so \(g\le c\) gives
\eqref{eq:terminal-deficit-affine-line}.  The case \(L_e\le1\) is
immediate.
\end{proof}

\begin{corollary}[Prefix-plus-terminal sparse-direction profile]
\label{cor:prefix-terminal-sparse-direction-profile}
For every \(1\le h<e\), let \(C_h\) be a proved cumulative cap on the
number \(N_h\) of transformed explanations with outside deficit at most
\(h\), and put
\[
 B_0=0,
 \qquad B_h=\min_{h\le v<e}C_v.
\]
Then
\begin{equation}
 |Z|\le
 \sum_{h=1}^{e-1}(B_h-B_{h-1})\left\lfloor\frac eh\right\rfloor
 +\left\lfloor\frac{N-e-(K-1)}{m-e-(K-1)}\right\rfloor.
 \label{eq:prefix-terminal-sparse-direction-profile}
\end{equation}
\end{corollary}

\begin{proof}
For \(h<e\), monotonicity gives
\(N_h\le\min_{h\le v<e}C_v=B_h\).  Sum the nonincreasing owner weights
\(\lfloor e/h\rfloor\) by parts over the prefix layers, then apply
\cref{thm:terminal-deficit-affine-line} to the exact terminal layer, whose
owner weight is one.
\end{proof}

\begin{remark}[Exact terminal-line support walls]
Using the Johnson/mean-centered cumulative caps for \(h<e\),
\eqref{eq:prefix-terminal-sparse-direction-profile} pays one further
support in each deployed row:
\[
\begin{array}{c|c|c|r|r}
\text{row}&K&\text{last paid }e&\text{terminal cap}&\text{total profile}\\ \hline
\text{KoalaBear}&14&64048&287&181326343\\
\text{Mersenne-31}&6&65455&493&16100647
\end{array}
\]
At KoalaBear \(e=64049\), the cumulative cap at \(h=e-1\) is unavailable.
At Mersenne \(e=65456\), the valid profile is
\(17119507>16777215\).  The remaining full-lift intervals are therefore
\(64049\le e\le1044238\) and \(65456\le e\le1044241\).  This is a
support-wall refinement, not a proof of either deployed row.  The exact
replay is in
\path{experimental/verify_mca_sparse_direction_terminal_deficit_line_payment_v1.py}.
\end{remark}

\begin{theorem}[Top-third exact-layer affine-line cap]
\label{thm:top-third-affine-line}
Retain the sparse-direction notation of
\cref{thm:sparse-direction-punctured-johnson-mca}.  Assume \(e\ge K\), put
\[
 s=\left\lfloor\frac{e-K}{3}\right\rfloor,
 \qquad H=e-s-1,
\]
and assume \(N-m>s\).  For \(0\le r\le s\), let \(L_r\) be the number of
selected explanations with exact outside deficit \(h=e-r\).  Then
\begin{equation}
 L_r\le
 \left\lfloor\frac{N-e-(K-1)}{m-e+r-(K-1)}\right\rfloor.
 \label{eq:top-third-affine-line}
\end{equation}
Every explanation counted by \(L_r\) owns at most one selected slope.
\end{theorem}

\begin{proof}
For an explanation \(a_i\) in this layer, let \(S_i\subseteq E\) be its
inside agreement set with its assigned slope.  Then
\[
 |E\setminus S_i|\le r.
\]
For explanations at distinct slopes \(\gamma_i,\gamma_j\), the normalized
difference
\[
 p_{ij}=\frac{a_i-a_j}{\gamma_i-\gamma_j}
\]
is a degree-\(<K\) codeword and equals the gauged direction \(q\) on
\(S_i\cap S_j\).  Any three inside agreement sets have intersection at
least \(e-3r\ge K\).  Hence \(p_{ij}\) and \(p_{ik}\) agree on at least
\(K\) evaluation coordinates and are equal.  Fixing two anchors puts the
entire exact layer on one affine codeword line.

Each member has \(A_r=m-e+r\) outside agreements.  The line has a common
agreement core of size \(g\le K-1\), and outside that core each coordinate
agrees for at most one line parameter.  Thus
\[
 L_r(A_r-g)\le N-e-g.
\]
The hypothesis \(N-m>s\) gives \(N-e>A_r\), so the ratio increases with
\(g\); substituting \(g\le K-1\) proves
\eqref{eq:top-third-affine-line}.  Finally \(r<e/2\), so the deficit-owner
bound is \(\lfloor e/(e-r)\rfloor=1\).
\end{proof}

\begin{corollary}[Top-third sparse-direction profile]
\label{cor:top-third-sparse-direction-profile}
For \(1\le h\le H\), let \(C_h\) be a proved cumulative explanation cap
and put
\[
 B_0=0,
 \qquad B_h=\min_{h\le v\le H}C_v.
\]
Then
\begin{align}
 |Z|
 &\le\sum_{h=1}^{H}(B_h-B_{h-1})
       \left\lfloor\frac eh\right\rfloor
   +\sum_{r=0}^{s}
       \left\lfloor\frac{N-e-(K-1)}{m-e+r-(K-1)}\right\rfloor.
 \label{eq:top-third-sparse-direction-profile}
\end{align}
If the positive punctured-Johnson caps \(J_h\) are defined through \(H\)
and \(u=\lfloor e/2\rfloor\), the first sum is at most
\[
 (e-1)J_u+J_H.
\]
\end{corollary}

\begin{proof}
Apply cumulative summation by parts through \(H\).  The remaining exact
layers are \(h=e-r\), \(0\le r\le s\); apply
\cref{thm:top-third-affine-line}.  For the coarse prefix estimate, give
layers through \(u\) weight at most \(e\) and all later prefix layers
weight one.
\end{proof}

\begin{remark}[Complete sparse-direction branch payment]
For the Johnson rational cap
\[
 F(n,A)=\frac{n(A-(K-1))}{A^2-n(K-1)},
\]
the function increases with \(n\) and decreases with \(A\).  Uniformly over
\(K\le e<d\), conservative endpoint substitution gives
\[
 J_{\lfloor e/2\rfloor}\le31,
 \qquad J_H\le47
\]
on both deployed rows.  Each affine-line summand is nondecreasing in \(e\),
so its sum is maximal at \(e=d-1\).  The exact endpoint calculations are
\[
\begin{array}{c|r|r|r}
\text{row}&\text{line sum}&\text{uniform total}&\text{budget}\\ \hline
\text{KoalaBear}&9405342&11496959&274980728111395087\\
\text{Mersenne-31}&9405365&11496238&16777215
\end{array}
\]
Consequently every sparse-direction support \(1\le e<d\) is paid.  The
remaining full-lift intervals are
\[
 67472\le e\le1044238
 \quad\text{and}\quad
 67448\le e\le1044241.
\]
This closes the low-support branch, not either deployed row.  The exact
replay is in
\path{experimental/verify_mca_sparse_direction_top_third_affine_line_payment_v1.py}.
\end{remark}

\begin{theorem}[Pair-noncontained total-core line cap]
\label{thm:full-lift-total-core-line}
Retain the notation of \cref{thm:top-third-affine-line}, but allow
\(e\ge d\).  Put
\[
 r_{\min}=\max\{0,e-m\},
 \qquad t=N-m.
\]
For \(r_{\min}\le r\le s\), set \(A_r=m-e+r\).  Then the exact-layer count
satisfies
\begin{equation}
 L_r\le Q_r,\qquad
 Q_r=
 \begin{cases}
 t+1,&A_r\le K-1,\\[2mm]
 \left\lfloor\dfrac{N-e-(K-1)}{A_r-(K-1)}\right\rfloor,
     &A_r>K-1.
 \end{cases}
 \label{eq:full-lift-total-core-line}
\end{equation}
\end{theorem}

\begin{proof}
The triple-overlap proof of \cref{thm:top-third-affine-line} still puts the
exact layer on an affine explanation line
\(a_\gamma=A+\gamma p\).  A coordinate agrees for every line parameter
exactly when \(A=r_0\) and \(p=q\) there.  If this total common agreement
core had at least \(m\) coordinates, the codeword pair \((A,b+p)\) would
explain the received pair on an \(m\)-set, contradicting pair
noncontainment.  Thus its size \(g\) is at most \(m-1\).  Away from the
core, each coordinate agrees for at most one parameter, and therefore
\[
 L_r(m-g)\le N-g,
 \qquad L_r\le N-m+1=t+1.
\]
When \(A_r>K-1\), outside-coordinate packing gives the second branch as in
\cref{thm:top-third-affine-line}.  That branch is no larger than \(t+1\):
after clearing its positive denominator, the difference is
\(t(A_r-K)+r\ge0\).
\end{proof}

\begin{corollary}[Full-lift common-core profile]
\label{cor:full-lift-common-core-profile}
If the positive punctured-Johnson caps exist through \(H=e-s-1\), then
\begin{equation}
 |Z|\le
 (e-1)J_{\lfloor e/2\rfloor}+J_H
 +\sum_{r=r_{\min}}^{s}Q_r.
 \label{eq:full-lift-common-core-profile}
\end{equation}
\end{corollary}

\begin{proof}
Deficits through \(H\) use the two-threshold Johnson prefix.  Deficits
above \(H\) and at most \(m\) are exactly the layers
\(r_{\min}\le r\le s\); apply \cref{thm:full-lift-total-core-line}.
\end{proof}

\begin{remark}[Exact full-lift common-core walls]
Exact integer evaluation of
\eqref{eq:full-lift-common-core-profile} extends the complete walls to
\[
\begin{array}{c|r|r|r}
\text{row}&\text{last paid }e&\text{endpoint bound}&\text{next bound}\\ \hline
\text{KoalaBear}&95943&27414298&\text{undefined}\\
\text{Mersenne-31}&67452&16266965&17248067
\end{array}
\]
At KoalaBear \(e=95944\), the prefix Johnson denominator at \(H\) is
\(-1037\).  The adjacent Mersenne bound exceeds budget by \(470852\).
The remaining full-lift intervals are
\[
 95944\le e\le1044238
 \quad\text{and}\quad
 67453\le e\le1044241.
\]
Neither adjacent failure is an unsafe certificate.  The replay is in
\path{experimental/verify_mca_full_lift_top_third_common_core_payment_v1.py}.
\end{remark}

\begin{theorem}[Cross-layer top-third global-line cap]
\label{thm:full-lift-top-third-global-line}
Retain the full-lift notation above and put
\[
 s=\left\lfloor\frac{e-K}{3}\right\rfloor,
 \qquad r_{\min}=\max\{0,e-m\}.
\]
All selected explanations in the union of exact layers
\[
 r_{\min}\le r\le s,
 \qquad h=e-r,
\]
lie on one common affine codeword line.  The number $L_{\mathrm{high}}$
of explanations in this entire union satisfies
\begin{equation}
 L_{\mathrm{high}}\le N-m+1.
 \label{eq:full-lift-top-third-global-line}
\end{equation}
Each explanation in the union owns at most one selected slope.
\end{theorem}

\begin{proof}
For an explanation $a_i$, let $r_i\le s$ be its missed-coordinate
allowance on the direction support $E$, and let $S_i\subseteq E$ be
its inside agreement set.  For any three explanations,
\[
 |S_i\cap S_j\cap S_k|
 \ge e-(r_i+r_j+r_k)
 \ge e-3s
 \ge K.
\]
For distinct slopes, the normalized pair difference
\[
 p_{ij}=\frac{a_i-a_j}{\gamma_i-\gamma_j}
\]
is a degree-$<K$ codeword equal to the gauged direction $q$ on
$S_i\cap S_j$.  Hence $p_{ij}=p_{ik}$ by restriction injectivity on
the displayed triple intersection.  Fixing two anchors synchronizes all
pair directions and puts the entire cross-layer union on one affine
codeword line.  Families of size at most two are immediate.

Let $g$ be the line's total common agreement core.  As in
\cref{thm:full-lift-total-core-line}, pair noncontainment gives
$g\le m-1$.  Away from this core, each coordinate agrees for at most one
line parameter.  Every selected explanation has at least $m-g$
agreements off the common core when all $N$ coordinates are counted, so
\[
 L_{\mathrm{high}}(m-g)\le N-g.
\]
The right-hand ratio increases with $g$, and $g\le m-1$ yields
$L_{\mathrm{high}}\le N-m+1$.  Finally $r_i\le s<e/2$, so the
deficit-owner bound is one.
\end{proof}

\begin{corollary}[Full-lift global-line profile]
\label{cor:full-lift-global-line-profile}
Put $H=e-s-1$ and $u=\lfloor e/2\rfloor$.  If the positive
punctured-Johnson caps $J_u,J_H$ exist, then
\begin{equation}
 |Z|\le (e-1)J_u+J_H+(N-m+1).
 \label{eq:full-lift-global-line-profile}
\end{equation}
\end{corollary}

\begin{proof}
Deficits at most $H$ form the Johnson prefix and contribute at most
$(e-1)J_u+J_H$.  Every remaining possible deficit belongs to the
cross-layer union in
\cref{thm:full-lift-top-third-global-line}; charge its line cap once.
\end{proof}

\begin{remark}[Exact cross-layer global-line walls]
Exact integer evaluation of
\eqref{eq:full-lift-global-line-profile} gives
\[
\begin{array}{c|r|r|r}
\text{row}&\text{last paid }e&\text{endpoint bound}&\text{largest bound}\\ \hline
\text{KoalaBear}&95943&6336049&6336049\text{ at }e=95943\\
\text{Mersenne-31}&97908&6682339&6683188\text{ at }e=97907
\end{array}
\]
At the adjacent supports $e=95944$ and $e=97909$, the respective
$H$-prefix Johnson denominators are $-1037$ and $-965$.  The
remaining full-lift intervals are therefore
\[
 95944\le e\le1044238
 \quad\text{and}\quad
 97909\le e\le1044241.
\]
This replaces the per-layer top-third charge by one global charge.  It does
not repair the failed prefix cap, certify either adjacent support unsafe,
or close a deployed row.  The exact replay is in
\path{experimental/verify_mca_full_lift_top_third_global_line_payment_v1.py}.
\end{remark}

\begin{corollary}[Full-lift mean-centered global-line profile]
\label{cor:full-lift-mean-centered-global-line-profile}
Retain the notation of
\cref{thm:full-lift-top-third-global-line}, put
\[
 H=e-s-1,
 \qquad c=K-1,
\]
and assume $m-H>c$.  For every $1\le h\le H$, put
$A_h=m-h$, and let $C_h$ be the punctured ordinary-Johnson cap when
$A_h^2>(N-e)c$, and otherwise the cap in
\cref{thm:mean-centered-gram-list} whenever its three hypotheses hold.
Assume every $C_h$ is defined, and put
\[
 B_0=0,
 \qquad B_h=\min_{h\le v\le H}C_v.
\]
Then
\begin{equation}
 |Z|\le
 \sum_{h=1}^{H}(B_h-B_{h-1})
       \left\lfloor\frac eh\right\rfloor
 +(N-m+1).
 \label{eq:full-lift-mean-centered-global-line-profile}
\end{equation}
\end{corollary}

\begin{proof}
Let $N_h$ count transformed explanations with outside deficit at most
$h$.  The earlier sparse-direction corollary was calibrated for $e<d$,
but its ordinary set-system proof has no such hypothesis.  After puncturing
the direction support, choose exactly $A_h=m-h$ outside agreement
coordinates for every explanation.  The guard $m-H>c$ gives $A_h>c$
throughout the prefix, and distinct degree-$<K$ explanations intersect in
at most $c$ coordinates.  Thus $N_h\le C_h$.  For $v\ge h$, also
$N_h\le N_v\le C_v$, so $N_h\le B_h$.  Summation by parts against the
owner weights $\lfloor e/h\rfloor$ gives the displayed prefix charge.
Every remaining deficit belongs to the cross-layer high union; apply
\cref{thm:full-lift-top-third-global-line} once.
\end{proof}

\begin{remark}[Exact full-lift mean-centered walls]
Exact evaluation of
\eqref{eq:full-lift-mean-centered-global-line-profile} gives
\[
\begin{array}{c|r|r|r}
\text{row}&\text{last paid }e&\text{endpoint bound}&\text{largest bound}\\ \hline
\text{KoalaBear}&96150&479693401&479693401\text{ at }e=96150\\
\text{Mersenne-31}&98229&16488216&16489118\text{ at }e=98228
\end{array}
\]
At KoalaBear $e=96151$, the first undefined cap is at $h=H=64105$,
where the mean-centered denominator is $-4625043784$.  At Mersenne-31
$e=98230$, every cap remains defined but the profile is
$17415873$, exceeding budget by $638658$.  The remaining full-lift
intervals are
\[
 96151\le e\le1044238
 \quad\text{and}\quad
 98230\le e\le1044241.
\]
Neither adjacent failure is an unsafe certificate, and no deployed v4 atom
moves.  The exact replay is in
\path{experimental/verify_mca_full_lift_mean_centered_global_line_profile_v1.py}.
\end{remark}

\begin{corollary}[One-layer boundary-anchor continuation]
\label{cor:full-lift-boundary-anchor-continuation}
Retain the hypotheses and notation of
\cref{cor:full-lift-mean-centered-global-line-profile}, and put
\[
 q=e-K-3s.
\]
For $0\le J\le H$, define
\[
 B^{(J)}_0=0,
 \qquad B^{(J)}_h=\min_{h\le v\le J}C_v,
 \qquad
 P_J=\sum_{h=1}^{J}
       (B^{(J)}_h-B^{(J)}_{h-1})
       \left\lfloor\frac eh\right\rfloor,
 \qquad P_0=0.
\]
If $H\ge2$, $q\ge1$, and $2(s+1)<e$, then
\begin{equation}
 |Z|\le\max\{P_H+1,\ P_{H-1}+N-m+1\}.
 \label{eq:full-lift-boundary-anchor-continuation}
\end{equation}
\end{corollary}

\begin{proof}
Let $\mathcal A$ be the selected explanations with missed-coordinate
allowance at most $s$, and let $\mathcal D$ be the exact boundary layer
with allowance $s+1$.  If $|\mathcal A|\le1$, the complete prefix through
$H$ and the one-slope ownership guard give $|Z|\le P_H+1$.

Suppose $|\mathcal A|\ge2$, and choose distinct anchors $a,b\in\mathcal A$.
For each $c\in\mathcal D$,
\[
 |S_a\cap S_b\cap S_c|
 \ge e-s-s-(s+1)=K+q-1\ge K.
\]
The normalized pair differences for $(c,a)$ and $(c,b)$ agree with the
gauged direction on this intersection.  Restriction injectivity for
degree-$<K$ codewords makes them identical, so $c$ lies on the affine
codeword line already containing $\mathcal A$.  Thus
$\mathcal A\cup\mathcal D$ is one pair-noncontained line.  The guard
$2(s+1)<e$ gives one-slope ownership throughout this union, and
\cref{thm:full-lift-total-core-line} charges it by $N-m+1$.  The remaining
prefix is bounded independently by $P_{H-1}$.
\end{proof}

\begin{remark}[First Mersenne boundary support]
For Mersenne-31 at $e=98230$,
\[
 (s,q,H)=(32741,1,65488),\qquad
 P_H=16434744,\qquad P_{H-1}=15506184,
\]
and $N-m+1=981129$.  Hence
\eqref{eq:full-lift-boundary-anchor-continuation} gives
\[
 |Z|\le\max\{16434745,16487313\}=16487313<16777215,
\]
with slack $289902$.  At $e=98231$ the same legal bound is $17492173$,
which exceeds budget by $714958$; the residue-two refinement below repairs
that proof-method failure.
The exact replay is in
\path{experimental/verify_mca_full_lift_boundary_anchor_continuation_v1.py}.
\end{remark}

\begin{corollary}[Residue-two boundary-layer continuation]
\label{cor:full-lift-two-boundary-layer-continuation}
Retain the hypotheses and notation of
\cref{cor:full-lift-boundary-anchor-continuation}.  Suppose in addition
that
\[
 q=2,\qquad H\ge3,\qquad 2(s+2)<e,\qquad m-H>c,
\]
and put
\[
 Q_H=\left\lfloor\frac{N-e-c}{m-H-c}\right\rfloor,
 \qquad
 \Delta_H=\left\lfloor\frac{e}{s+1}\right\rfloor.
\]
Then
\begin{equation}
 |Z|\le\max\left\{
 \begin{array}{l}
 P_{H-2}+N-m+1,\\
 P_{H-1}+Q_H+1,\\
 P_{H-1}+2,\\
 P_{H-1}+Q_H,\\
 P_{H-1}+\Delta_H
 \end{array}\right\}.
 \label{eq:full-lift-two-boundary-layer-continuation}
\end{equation}
\end{corollary}

\begin{proof}
Let $\mathcal A$ be the top-third union, and let
$\mathcal D_1,\mathcal D_2$ be the exact layers of deficits $H,H-1$.
Choose inside agreement sets at their layer sizes, so their missed sets
have sizes $s+1,s+2$.  The guard $2(s+2)<e$ gives one-slope ownership.

If $|\mathcal A|\ge2$, two top anchors meet every member of
$\mathcal D_1\cup\mathcal D_2$ on at least
\[
 e-s-s-(s+2)=K+q-2=K
\]
common coordinates.  Restriction injectivity puts all three layers on one
affine codeword line.  Total-core packing charges it by $N-m+1$, leaving
$P_{H-2}$.

Suppose $\mathcal A$ has one member.  If $|\mathcal D_1|\ge2$, that member
and any two boundary members have at least
$e-s-2(s+1)=K$ common coordinates, so $\mathcal D_1$ is one affine line.
Charge the top member separately.  Every boundary member has exactly
$m-H$ outside agreements, and outside-core line packing bounds the layer by
$Q_H$.  This gives $P_{H-1}+Q_H+1$.  If
$|\mathcal D_1|\le1$, direct charging gives $P_{H-1}+2$.

Finally suppose $\mathcal A$ is empty, and write $R_i$ for the size-$(s+1)$
missed set of a member of $\mathcal D_1$.  If some $R_a,R_b$ intersect,
then every third member has
\[
 |S_a\cap S_b\cap S_j|
 \ge e-(3(s+1)-1)=K+q-2=K.
\]
The fixed pair synchronizes the whole layer, giving $P_{H-1}+Q_H$.
Otherwise the $R_i$ are pairwise disjoint, whence
$|\mathcal D_1|(s+1)\le e$ and the charge is
$P_{H-1}+\Delta_H$.  These cases are exhaustive.
\end{proof}

\begin{remark}[Second Mersenne boundary support]
At Mersenne-31 $e=98231$,
\[
 (s,q,H)=(32741,2,65489),\quad
 P_{H-2}=15505282,\quad P_{H-1}=16433719,
\]
while $N-m+1=981129$, $Q_H=484$, and $\Delta_H=3$.  The five cases in
\eqref{eq:full-lift-two-boundary-layer-continuation} are
\[
 16486411,\quad16434204,\quad16433721,\quad16434203,\quad16433722.
\]
Thus $|Z|\le16486411<16777215$, with slack $290804$.  At $e=98232$ the
residue resets to $q=0$, so this corollary is inapplicable; that is not an
unsafe certificate.  The Mersenne full-lift residual interval is now
$98232\le e\le1044241$.  The exact replay is in
\path{experimental/verify_mca_full_lift_two_boundary_layer_continuation_v1.py}.
\end{remark}

\begin{corollary}[Boundary direction-class router]
\label{cor:full-lift-boundary-direction-router}
Retain the pair-noncontained full-lift hypotheses and notation of
\cref{thm:full-lift-top-third-global-line}, put
\[
 H=e-s-1,\qquad c=K-1,
\]
For $1\le h\le H-1$, define $C_h$ by the punctured Johnson or
mean-centered rule in
\cref{cor:full-lift-mean-centered-global-line-profile}, and suppose every
such cap is defined.  Define the independently truncated prefix
\[
 B_0=0,\qquad B_h=\min_{h\le v\le H-1}C_v,
 \qquad
 P_{H-1}=\sum_{h=1}^{H-1}(B_h-B_{h-1})
                 \left\lfloor\frac eh\right\rfloor.
\]
Let $\mathcal D$ be the exact deficit-$H$ layer, and write $q_{\rm dir}$
for the gauged direction on its support $E$.  Put
\[
 A=2H-e,
 \qquad
 J=\left\lfloor\frac{e(A-c)}{A^2-ec}\right\rfloor,
 \qquad
 Q=\left\lfloor\frac{N-e-c}{m-H-c}\right\rfloor.
\]
If
\[
 A>0,\qquad A^2>ec,\qquad N-e>m-H>c,
\]
then
\begin{equation}
 |\mathcal D|\le 1+J(Q-1).
 \label{eq:full-lift-boundary-direction-cap}
\end{equation}
If $\mathcal T$ is the synchronized union of deficits greater than $H$,
then consequently
\begin{equation}
 |Z|\le P_{H-1}+1+J(Q-1)+|\mathcal T|.
 \label{eq:full-lift-boundary-direction-router}
\end{equation}
\end{corollary}

\begin{proof}
The claim is immediate when $\mathcal D$ is empty.  Otherwise fix an anchor
$(\gamma_0,c_0)\in\mathcal D$.  For every other member put
\[
 p_\gamma=\frac{c_\gamma-c_0}{\gamma-\gamma_0}.
\]
The chosen size-$H$ inside agreement sets meet in at least $A=2H-e$
coordinates, and on this intersection $p_\gamma=q_{\rm dir}$.  Thus the
intrinsic set
\[
 U_p=\{x\in E:q_{\rm dir}(x)=p(x)\}
\]
has size at least $A$ for every represented direction $p$.  Such a $p$ is
nonzero: otherwise the two explanations would be the same codeword, while
on their nonempty inside intersection that codeword would equal received
words whose difference is $(\gamma-\gamma_0)q_{\rm dir}\ne0$.

For distinct represented directions $p,p'$, the intersection
$U_p\cap U_{p'}$ lies in the zero set of the nonzero degree-$<K$ codeword
$p-p'$, and hence has size at most $c$.  Choosing $A$ points from every
$U_p$ and applying the constant-block Johnson double count gives at most
$J$ represented directions.

For a fixed $p$, its members and the anchor lie on the affine codeword line
$c_0+(\gamma-\gamma_0)p$.  Outside $E$ the received direction vanishes, so
the common agreement core of this line lies in the zero set of the nonzero
codeword $p$ and has size at most $c$.  Every member has $m-H$ outside
agreements.  Since $N-e>m-H>c$, outside-core packing bounds the line by
$Q$.  The anchor is repeated in each line closure, so each of the at most
$J$ classes contributes at most $Q-1$ nonanchors.  This proves
\eqref{eq:full-lift-boundary-direction-cap}.  The boundary and top layers
have one-slope ownership; the independently truncated prefix charges all
lower deficits, proving \eqref{eq:full-lift-boundary-direction-router}.
\end{proof}

\begin{remark}[First residue-zero Mersenne router]
At Mersenne-31 $e=98232$,
\[
 (s,H,A)=(32742,65489,32746),\qquad J=3,\qquad Q=484,
\]
so $|\mathcal D|\le1450$.  The exact prefix is
$P_{H-1}=16432695$, and therefore
\[
 |Z|\le16434145+|\mathcal T|.
\]
An unsafe family would need $|\mathcal T|\ge343071$.  If $g$ is the total
common core of this affine line, line packing gives
\[
 g\ge
 \left\lceil\frac{|\mathcal T|m-N}{|\mathcal T|-1}\right\rceil
 \ge67452=m-2.
\]
Thus the first residue-zero obstruction is reduced to a near-maximal-core
line.  The router alone is not a safety or unsafety certificate at
$e=98232$; the following corollary absorbs the remaining layers into this
line.  The exact router replay is in
\path{experimental/verify_mca_full_lift_residue_zero_direction_router_v1.py},
with an independent endpoint audit in
\path{experimental/audit_mca_full_lift_residue_zero_direction_router_v1.py}.
\end{remark}

\begin{corollary}[First residue-zero Mersenne payment]
\label{cor:full-lift-residue-zero-core-absorption}
Under the pair-noncontained full-lift hypotheses, the Mersenne-31 support
$e=98232$ is safe:
\[
 |Z|\le16777215.
\]
Consequently its full-lift residual interval starts at $e=98233$.
\end{corollary}

\begin{proof}
Suppose instead that $|Z|>16777215$.  By the preceding router, the
synchronized top line $\mathcal T$ has at least $343071$ members and total
common agreement core $G$ of size at least $67452=m-2$.  Write the line as
$c_\gamma=A_0+\gamma p$.  Its direction $p$ is nonzero.  Outside the
gauged direction support $E$, a coordinate common to two line members must
be a zero of $p$, so
\[
 |G\cap E|\ge67452-(K-1)=67447=:u.                 \tag{1}
\]

Take any selected slope $\delta$ whose assigned explanation has outside
deficit $h\ge e-u+K=30791$.  Ownership supplies an inside agreement set
$S_\delta$ of size $h$.  Choose two distinct top anchors.  By (1),
\[
 |G\cap E\cap S_\delta|\ge u+h-e\ge K.
\]
On these coordinates, both normalized differences from $c_\delta$ to the
two anchors equal the gauged direction.  Restriction injectivity makes the
two degree-$<K$ differences equal.  Subtracting their defining equations
then identifies this common difference with $p$, and hence
$c_\delta=A_0+\delta p$.  Thus every assigned pair with $h\ge30791$ joins
the same affine line.  The nonzero direction makes its slope parameter
injective, and total-core line packing charges all these slopes once by
\[
 N-m+1=981129.                                      \tag{2}
\]

Every remaining assigned explanation has $h\le30790$, so after puncturing
$E$ it agrees with the outside received word on at least
$m-30790=36664$ coordinates.  The ordinary Johnson count gives at most
\[
 \left\lfloor
 \frac{950350(36664-5)}{36664^2-950350\cdot5}
 \right\rfloor=26                                  \tag{3}
\]
distinct explanations.  Each owns at most $e=98232$ slopes.  Combining
(2) and (3) yields
\[
 |Z|\le98232\cdot26+981129
     =3535161<16777215,
\]
the desired contradiction.
\end{proof}

\begin{remark}[Exact absorption margin]
The contradiction bound has margin $13242054$.  This payment does not
claim safety at $e=98233$ or provide an unsafe certificate.  Its exact
replay and independent audit are
\path{experimental/verify_mca_full_lift_residue_zero_core_absorption_v1.py}
and
\path{experimental/audit_mca_full_lift_residue_zero_core_absorption_v1.py}.
\end{remark}

\begin{corollary}[Fixed-cutoff boundary-stack continuation]
\label{cor:full-lift-fixed-cutoff-boundary-stack}
Retain the pair-noncontained full-lift hypotheses and notation above.  Fix
a target slope budget $B$ and $h_0<H$, and let $P_{h_0}(e)$ be the independently truncated
Johnson/mean-centered prefix through $h_0$.  For every $h_0<h\le H$, put
\[
 A_h=2h-e,
 \qquad
 J_h=\left\lfloor\frac{e(A_h-c)}{A_h^2-ec}\right\rfloor,
 \qquad
 Q_h=\left\lfloor\frac{N-e-c}{m-h-c}\right\rfloor,
\]
and suppose
\[
 2h>e,\qquad A_h^2>ec,\qquad N-e>m-h>c
\]
throughout this range.  Define
\[
 F_e=P_{h_0}(e)+\sum_{h=h_0+1}^{H}\{1+J_h(Q_h-1)\}.
 \tag{BS1}
\]
Then
\[
 |Z|\le F_e+|\mathcal T|.                           \tag{BS2}
\]
In particular, $F_e+N-m+1\le B$ implies safety.  Otherwise suppose
$F_e<B$, and put
\[
 L_e=B-F_e+1,
 \qquad
 g_e=\left\lceil\frac{L_em-N}{L_e-1}\right\rceil,
 \qquad
 u_e=g_e-c,
 \qquad
 a_e=e-u_e+K.
\]
If the punctured ordinary-Johnson cap
\[
 M_e=\left\lfloor
 \frac{(N-e)(m-a_e+1-c)}{(m-a_e+1)^2-(N-e)c}
 \right\rfloor                                      \tag{BS3}
\]
is defined and
\[
 eM_e+N-m+1\le B,                                  \tag{BS4}
\]
then the support is safe.
\end{corollary}

\begin{proof}
The prefix handles deficits through $h_0$.  Fix an exact layer
$h_0<h\le H$.  Every explanation owns one slope.  Relative to an anchor,
its nonzero normalized codeword direction agrees with the gauged direction
on at least $A_h=2h-e$ coordinates.  Distinct directions have intrinsic
agreement sets meeting in at most $c$ coordinates, so the constant-block
Johnson count gives at most $J_h$ classes.  Each class and the anchor form
one nonzero affine line.  Its outside common core has size at most $c$, and
outside-core packing gives $Q_h$ members.  Subtracting the repeated anchor
prices the layer by $1+J_h(Q_h-1)$.  The unpriced top union is the one
synchronized line $\mathcal T$, proving (BS2).

If (BS2) does not pay directly and the family is unsafe, then
$|\mathcal T|\ge L_e$ and total-core line packing forces core at least
$g_e$.  At most $c$ core coordinates lie outside the gauged direction
support, so at least $u_e$ lie inside it.  The two-anchor argument from
\cref{cor:full-lift-residue-zero-core-absorption} absorbs every assigned
pair of deficit at least $a_e$ into the same line.  The remaining
explanations have outside agreement at least $m-a_e+1$; (BS3) counts them,
the conservative owner multiplier $e$ prices their slopes, and (BS4) pays.
\end{proof}

\begin{remark}[Exact fixed-cutoff Mersenne interval]
For Mersenne-31, take $h_0=65200$.  Exact integer replay pays every support
\[
 98232\le e\le101155.
\]
The direct branch applies through $e=101149$; the final six supports use
core absorption.  At $e=101155$,
\[
 \begin{aligned}
 F_e&=16667033,& L_e&=110183,& g_e&=67446,\\
 M_e&=28,& eM_e+N-m+1&=3813469.
 \end{aligned}
\]
The final bound has slack $12963746$.  At adjacent $e=101156$, the same
fixed-cutoff charge is $16951223$, above budget by $174008$.  This is a
fixed-cutoff method wall, not an unsafe certificate.  The Mersenne
full-lift residual interval is now $101156\le e\le1044241$.

The endpoint verifier and full constant-memory C replay are
\path{experimental/verify_mca_full_lift_fixed_cutoff_boundary_stack_v1.py}
and
\path{experimental/verify_mca_full_lift_fixed_cutoff_boundary_stack_v1.c}.
\end{remark}

\begin{corollary}[Residue-two repair at the fixed-cutoff wall]
\label{cor:full-lift-fixed-cutoff-q2-anchor-repair}
Retain the pair-noncontained Mersenne-31 full-lift hypotheses.  At
\[
 e=101156,\qquad (s,q,H)=(33716,2,67439),
\]
use the fixed-cutoff charge \((\mathrm{BS1})\) with \(h_0=65258\).
Let \(\mathcal T\) be the synchronized top union, and let
\(\mathcal D_1,\mathcal D_2\) be the exact layers of deficits
\(H,H-1\).  The resulting charges are
\[
 F=16895280,\qquad
 D_1=284224,\qquad D_2=258385.
\]
Then the complete slope family satisfies
\[
 |Z|\le16705799<16777215.
\]
Consequently this support is safe, and the residual Mersenne full-lift
interval is \(101157\le e\le1044241\).
\end{corollary}

\begin{proof}
Every member of \(\mathcal T\cup\mathcal D_1\cup\mathcal D_2\)
owns at most one slope.  If \(|\mathcal T|\ge2\), two top anchors and
any boundary member have at least
\[
 e-s-s-(s+2)=K+q-2=K
\]
common inside agreement coordinates.  Thus the top union and both
boundary layers lie on one affine codeword line.  Removing both coarse
boundary charges leaves
\[
 G=F-D_1-D_2=16352671.
\]
Unsafety would force at least \(16777215-G+1=424545\) members on that
line.  Total-core packing then forces common core \(67452=m-2\), with
at least \(67447\) core coordinates inside the gauged direction
support.  The two-anchor absorption argument therefore synchronizes every
assigned explanation of deficit at least \(33715\) onto the line.  The
remaining punctured ordinary list has agreement at least \(33740\) and
Johnson cap \(28\), whence
\[
 |Z|\le101156\cdot28+981129=3813497.
\]

Suppose next that \(|\mathcal T|=1\).  If
\(|\mathcal D_1|\ge2\), the top anchor and any two first-boundary
members share at least
\[
 e-s-2(s+1)=K
\]
inside coordinates, so \(\mathcal D_1\) is one affine codeword line.
Its outside agreement is \(m-H=15\), and outside-core packing gives
\(|\mathcal D_1|\le94742\).  The two subcases give
\[
 F-D_1+94742+1=16705799,\qquad
 F-D_1+2=16611058.
\]

Finally suppose \(\mathcal T\) is empty.  If two first-boundary
missed sets intersect, those two members and any third share at least
\[
 e-\{3(s+1)-1\}=K
\]
inside coordinates.  The layer is one affine line and contributes at
most \(94742\), giving \(F-D_1+94742=16705798\).  Otherwise its
missed sets are pairwise disjoint, so
\(|\mathcal D_1|\le\lfloor e/(s+1)\rfloor=3\), giving
\(F-D_1+3=16611059\).  These five exhaustive cases have maximum
\(16705799\), as claimed.
\end{proof}

\begin{remark}[Scope and replay]
At adjacent \(e=101157\), the residue resets to zero, so the preceding
repair is inapplicable; no unsafe certificate is claimed.  The source-bound
replay and independent constant audit are
\path{experimental/verify_mca_full_lift_fixed_cutoff_q2_anchor_repair_v1.py}
and
\path{experimental/audit_mca_full_lift_fixed_cutoff_q2_anchor_repair_v1.py}.
\end{remark}

\begin{corollary}[Boundary direction-class affine-line bank]
\label{cor:full-lift-boundary-line-bank}
Retain the pair-noncontained full-lift hypotheses and put
\[
 s=\left\lfloor\frac{e-K}{3}\right\rfloor,\qquad
 H=e-s-1,\qquad U=\min(H,m),\qquad c=K-1.
\]
Fix \(h_0<U\), and let \(P_{h_0}(e)\) be the independently truncated
Johnson/mean-centered prefix through \(h_0\).  For
\(h_0<h\le U\), put
\[
 A_h=2h-e,\qquad
 J_h=\left\lfloor\frac{e(A_h-c)}{A_h^2-ec}\right\rfloor,
\]
and suppose
\[
 2h>e,\qquad A_h>c,\qquad A_h^2>ec
\]
throughout this range.  Define
\[
 G_e=\mathbf 1_{\{H<m\}}+\sum_{h=h_0+1}^{U}J_h,\qquad
 C_e=P_{h_0}(e)+\sum_{h=h_0+1}^{U}(1-J_h).
\]
Then there are \(G_e\) affine explanation-line slots
\(\mathcal L_i\) such that
\[
 |Z|\le C_e+\sum_{i=1}^{G_e}|\mathcal L_i|.        \tag{LB1}
\]
In particular, for target budget \(B\), unsafety forces one slot to
have size at least
\[
 \lambda_e=
 \left\lceil\frac{B-C_e+1}{G_e}\right\rceil.       \tag{LB2}
\]
If \(\lambda_e\ge2\), put
\[
 g_e=\left\lceil\frac{\lambda_em-N}{\lambda_e-1}\right\rceil,
 \quad u_e=g_e-c,\quad a_e=e-u_e+K.
\]
When the punctured ordinary-Johnson cap
\[
 M_e=\left\lfloor
 \frac{(N-e)(m-a_e+1-c)}
 {(m-a_e+1)^2-(N-e)c}
 \right\rfloor
\]
is defined and \(eM_e+N-m+1\le B\), the support is safe.
\end{corollary}

\begin{proof}
Fix one anchor in an exact layer \(h_0<h\le U\).  Every selected
explanation owns at most one slope because \(2h>e\).  Relative to the
anchor, each nonzero normalized codeword direction agrees with the gauged
direction on at least \(A_h\) inside coordinates.  Distinct directions
have agreement blocks meeting in at most \(c\), so the constant-block
Johnson count permits at most \(J_h\) direction classes.

Each class together with the anchor is one affine explanation line.  Pad
to exactly \(J_h\) classes with anchor-only slots.  The exact layer then
satisfies
\[
 |\mathcal D_h|
 =1-J_h+\sum_{j=1}^{J_h}|\mathcal L_{h,j}|.
\]
If \(H<m\), the cross-layer top union contributes one further affine line.
Summing these identities and the prefix proves (LB1).  Notice that this
argument uses no outside-core denominator.

If the family is unsafe, pigeonhole gives (LB2).  Total-core line packing
forces a common core of size at least \(g_e\), of which at least \(u_e\)
coordinates lie inside the gauged direction support.  Two anchors from
the large slot synchronize every assigned explanation of deficit at least
\(a_e\) onto that same line.  The remaining explanations have outside
agreement at least \(m-a_e+1\); the punctured ordinary-Johnson cap counts
them, and the conservative owner multiplier \(e\) completes the payment.
\end{proof}

\begin{remark}[Exact Mersenne line-bank interval]
For Mersenne-31, take \(h_0=65272\).  Exact replay pays every support
\[
 101157\le e\le124805.
\]
At the endpoint,
\[
 \begin{aligned}
 P_{h_0}&=1636955,& C_e&=1604577,&G_e&=34560,\\
 \lambda_e&=440,&g_e&=65220,&M_e&=126,\\
 eM_e+N-m+1&=16706559.
 \end{aligned}
\]
The final margin is \(70656\).  At adjacent \(e=124806\), the same legal
compiler gives \(16831491\), above budget by \(54276\).  This is a method
wall, not an unsafe certificate.  The Mersenne residual interval is now
\(124806\le e\le1044241\).

The source-bound endpoint replay, independent audit, and full
constant-memory interval replay are
\path{experimental/verify_mca_full_lift_boundary_line_bank_absorption_v1.py},
\path{experimental/audit_mca_full_lift_boundary_line_bank_absorption_v1.py},
and
\path{experimental/verify_mca_full_lift_boundary_line_bank_absorption_v1.c}.
\end{remark}

\begin{corollary}[Recursive affine-line peeling and inside-core packing]
\label{cor:full-lift-recursive-line-peeling}
Assume the hypotheses of
Corollary~\ref{cor:full-lift-boundary-line-bank}, suppose \(H\ge m\), and
put \(Q=N-m+1\).  Fix a cutoff \(b<m\) for which the prefix and every
line-bank guard are valid.  Let \(\mathcal Z_r\) be the residual assigned
slopes after \(r\) affine lines have been removed, with
\(\mathcal Z_0=Z\) and deficit ceiling \(U_0=m\), and charge the removed
lines by \(rQ\).
Put
\[
 T_r=B-rQ,\qquad
 G_r=\sum_{h=b+1}^{U_r}J_h,\qquad
 C_r=P_b(e)+(U_r-b)-G_r.
\]
If the suffix-minimum prefix through \(U_r\) is defined and at most \(T_r\),
then the support is safe.  Otherwise, whenever \(G_r>0\) and
\(T_r-C_r+1>0\), an unsafe residual forces an affine line of size at least
\[
 \lambda_r=\left\lceil\frac{T_r-C_r+1}{G_r}\right\rceil.       \tag{RP1}
\]
For \(\lambda_r\ge2\), define
\[
 g_r=\max\left(0,
       \left\lceil\frac{\lambda_rm-N}{\lambda_r-1}\right\rceil\right),
 \qquad u_r=\max(g_r-c,0),
 \qquad U_{r+1}=\min(U_r,e-u_r+K-1).             \tag{RP2}
\]
Remove the whole selected line and repeat.  The selected lines are distinct,
and their inside-core sets \(I_i\) satisfy
\[
 |I_i|\ge u_i,\qquad |I_i\cap I_j|\le c.
\]
Consequently, if at any stage
\[
 \sum_{i=0}^{r}u_i-\binom{r+1}{2}c>e,             \tag{RP3}
\]
the assumed unsafe family is impossible.  Thus termination either at the
weighted prefix or at (RP3) certifies safety.
\end{corollary}

\begin{proof}
Apply the exact-layer line-bank identity to the residual assigned family.
Empty layers only improve the resulting inequality, so
\[
 |\mathcal Z_r|\le C_r+\sum_{i=1}^{G_r}|\mathcal L_i|.
\]
If the original family is unsafe after \(r\) line charges, then
\(|\mathcal Z_r|>T_r\), and pigeonhole gives (RP1).  The guard
\(\lambda_r\ge2\) makes the selected slot an actual affine line rather than
a padded singleton.  Total-core line packing gives \(g_r\).  At most \(c\)
core coordinates lie outside the gauged direction support, and the
common-core absorption argument places every residual explanation of
deficit at least \(e-u_r+K\) on the same line.  Removing the whole line costs
at most \(Q\) and gives (RP2).  Even when (RP2) does not lower the deficit
ceiling, the selected line is absent from the next residual.

Write a selected line as \(c_\gamma=a_i+\gamma b_i\).  Its total common
core consists of coordinates where
\((r_0,r_1)=(a_i,b_i)\).  Distinct selected lines have distinct codeword
pairs.  At least one of \(a_i-a_j\) and \(b_i-b_j\) is therefore a nonzero
degree-\(<K\) codeword, so two total cores, and hence two inside cores, meet
in at most \(c=K-1\) coordinates.  Since every inside core lies in the same
\(e\)-coordinate support, first-order inclusion--exclusion gives the reverse
of (RP3), proving the contradiction.
\end{proof}

\begin{remark}[Exact Mersenne recursive-peeling interval]
For Mersenne-31, start from \(b=65304\).  When its first line-bank layer
fails a guard, take the least legal cutoff and add two guard layers.  Exact
replay checks every prefix and line-bank denominator under this printed
rule.  It pays all
\[
 124806\le e\le130198.
\]
Of these \(5393\) supports, \(3837\) terminate at the weighted prefix and
\(1556\) at (RP3), with at most five peeled lines.  The first packing
termination is \(e=128340\), where two inside cores give
\(62822+66574-5=129391>128340\).  At the last support, the five inside-core
lower bounds give
\[
 37718+33617+28204+20729+12942-10\cdot5
 =133160>130198.
\]
At adjacent \(e=130199\), nine legal peels give packing lower bound
\(126052\le130199\); the next residual target is \(7947054\), while the
certified base charge is \(8154082\), so (RP1) no longer forces a line.  This
is a method wall, not an unsafe certificate.  The Mersenne residual interval
is now \(130199\le e\le1044241\).

The endpoint replay, independent audit, and full constant-memory interval
replay are
\path{experimental/verify_mca_full_lift_recursive_line_peeling_core_packing_v1.py},
\path{experimental/audit_mca_full_lift_recursive_line_peeling_core_packing_v1.py},
and
\path{experimental/verify_mca_full_lift_recursive_line_peeling_core_packing_v1.c}.
\end{remark}

\begin{corollary}[Joint-core charge for peeled explanation lines]
\label{cor:full-lift-joint-core-charge}
In the setting of
Corollary~\ref{cor:full-lift-recursive-line-peeling}, suppose \(r\)
distinct parameterized affine explanation lines have been removed, and let
\(g_i\) be their actual total common-core sizes.  Put
\[
 S_r=\min\left\{r(m-1),e+\binom{r+1}{2}c\right\},
 \quad q_r=\left\lfloor\frac{S_r}{m-1}\right\rfloor,
 \quad z_r=S_r-q_r(m-1),
\]
and \(f(g)=(N-g)/(m-g)\), \(Q=N-m+1\).  Their total number of assigned
slopes is at most
\[
 L_r=
 \begin{cases}
  rQ,&q_r=r,\\
  \left\lfloor
   q_rQ+f(z_r)+(r-q_r-1)f(0)
  \right\rfloor,&q_r<r.
 \end{cases}                                      \tag{JC1}
\]
Consequently the residual peeling target may be strengthened from
\(B-rQ\) to \(B-L_r\).
\end{corollary}

\begin{proof}
Let \(I_i\) be the part of the \(i\)th total core inside the fixed
\(e\)-coordinate gauged direction support.  Distinct parameterized lines
have distinct codeword pairs, so
\(|I_i\cap I_j|\le c\).  Each line direction is a nonzero degree-\(<K\)
codeword, so at most \(c\) further core coordinates lie outside that
support.  Therefore
\[
 \sum_i g_i\le e+\binom r2c+rc
              =e+\binom{r+1}{2}c,
\]
and \(g_i\le m-1\) by pair noncontainment.

Off-core agreement sets on one line are disjoint, giving the rational cap
\(f(g_i)\).  The function
\[
 f(g)=1+\frac{N-m}{m-g}
\]
is increasing and convex.  Moving core mass from the smaller of two
interior arguments to the larger cannot decrease their sum.  Repeating
this exchange concentrates all available mass at \(m-1\), at most one
remainder \(z_r\), and zero.  The total number of slopes is integral, so
flooring the resulting rational sum gives (JC1).
\end{proof}

\begin{remark}[Exact Mersenne joint-charge interval]
Using \(B-L_r\) in the recursive line bank pays every support
\[
 130199\le e\le130219.
\]
All 21 supports terminate by inside-core packing, using between four and
thirteen lines.  At the endpoint, the thresholds are \(21\) and then twelve
copies of \(18\), giving
\[
 18393+12\cdot9736-\binom{13}{2}\cdot5
 =134835>130219.
\]
At adjacent \(e=130220\), the first threshold is \(20\) and the next 42 are
\(16\).  Their packing lower bound is only
\[
 15811+42\cdot2041-\binom{43}{2}\cdot5
 =97018.
\]
For \(r=43\), the joint core allowance is
\(134950=2(m-1)+44\), the exact joint charge is \(1962895\), and the next
threshold drops to \(13\), whose forced-core lower bound is zero.  Later
thresholds cannot increase.  This is a method wall, not an unsafe
certificate; the Mersenne residual is now
\(130220\le e\le1044241\).

The endpoint replay, independent rational audit, and full interval replay
are
\path{experimental/verify_mca_full_lift_joint_core_charge_peeling_v1.py},
\path{experimental/audit_mca_full_lift_joint_core_charge_peeling_v1.py},
and
\path{experimental/verify_mca_full_lift_joint_core_charge_peeling_v1.c}.
\end{remark}

\begin{corollary}[Lower-aware joint-core charge]
\label{cor:full-lift-lower-aware-joint-core-charge}
In the setting of
Corollary~\ref{cor:full-lift-joint-core-charge}, suppose the actual total
cores also have certified lower bounds
\(0\le\ell_i\le g_i\).  Sort the \(\ell_i\) decreasingly, start with
\(x_i=\ell_i\), and distribute
\[
 S_r-\sum_i\ell_i
\]
units by filling \(x_1\) to \(m-1\), then \(x_2\), and so on.  Whenever
the lower bounds are feasible, the removed lines contain at most
\[
 L_r^{\rm low}=\left\lfloor\sum_{i=1}^r
          \frac{N-x_i}{m-x_i}\right\rfloor              \tag{LA1}
\]
assigned slopes.  Consequently the residual peeling target may be
strengthened to \(B-L_r^{\rm low}\).
\end{corollary}

\begin{proof}
Since line labels do not occur in the objective, sort both the lower bounds
and a feasible core vector decreasingly.  This preserves the constraints:
if \(\ell_i\ge\ell_j\) but \(g_i<g_j\), then swapping the two core values
preserves both lower bounds.

For two nonterminal coordinates \(x_i\ge x_j\), transfer mass from
\(x_j-\ell_j\) to \(x_i\), stopping when \(x_i=m-1\) or
\(x_j=\ell_j\).  Convexity of
\(f(g)=(N-g)/(m-g)\) gives
\[
 f(x_i+d)+f(x_j-d)\ge f(x_i)+f(x_j).
\]
Iterating the exchange, and using monotonicity to spend all of the allowed
mass, gives exactly the greedy vector in the statement.  Equation (JC1)
supplies the common sum budget \(S_r\), so flooring the rational sum proves
(LA1).
\end{proof}

\begin{remark}[Exact Mersenne lower-aware payment]
At \(e=130220\) and \(e=130221\), the first 37 removed lines have forced
total-core lower-bound runs \(15816\cdot4,2046\cdot33\).  The maximizing
vectors are respectively
\[
 18769,15816\cdot3,2046\cdot33
 \quad\hbox{and}\quad
 18770,15816\cdot3,2046\cdot33,
\]
and both charges are \(609\).  One final threshold \(20\) is forced.  The
38 inside-core lower bounds then give
\[
 5\cdot15811+33\cdot2041-\binom{38}{2}\cdot5
 =142893>e.
\]
Thus both supports are safe.

At adjacent \(e=130222\), the compiler reaches 288 removed lines.  The
maximizing vector is \(67453\cdot5,1037,0\cdot282\), its charge is
\(4910044\), and the residual target \(11867171\) lies below the certified
base \(12148280\).  This is a method wall, not an unsafe certificate; the
Mersenne residual is now \(130222\le e\le1044241\).

The endpoint replay and independent rational audit are
\path{experimental/verify_mca_full_lift_lower_aware_joint_core_charge_v1.py}
and
\path{experimental/audit_mca_full_lift_lower_aware_joint_core_charge_v1.py}.
\end{remark}

\begin{corollary}[Core-dichotomy capped charge]
\label{cor:full-lift-core-dichotomy-capped-charge}
Fix an integer absorption cutoff $b$.  At every line selected by the
recursive bank, split on its actual total core $g$.

If $g\ge e+10-b$, every explanation of deficit at least $b+1$ lies on
that line, and the original family has size at most
\[
 P_e(b)+(N-m+1),                                  \tag{CD1}
\]
where $P_e(b)$ is the exact weighted prefix through $b$.

In the complementary branch, every peeled core has the individual cap
\[
 g_i\le G_e:=e+9-b.                               \tag{CD2}
\]
The lower-aware convex envelope remains valid after replacing its fill
ceiling by $G_e$ and its common budget by
\[
 S_r=\min\left\{rG_e,e+\binom{r+1}{2}c\right\}.   \tag{CD3}
\]
\end{corollary}

\begin{proof}
A selected line has at least $g-c$ core coordinates inside the gauged
direction support.  An explanation with inside agreement size $h$ and
two line anchors shares at least $(g-c)+h-e$ of those coordinates.
Restriction injectivity therefore puts it on the line whenever
\[
 h\ge e-g+c+K=e-g+11.
\]
The high-core threshold in the statement makes this at most $b+1$,
proving (CD1).

The integer complement of the high-core branch is (CD2).  Under that cap,
sort the forced core lower bounds and greedily spend (CD3), filling each
coordinate only to $G_e$.  The convex exchange in
Corollary~\ref{cor:full-lift-lower-aware-joint-core-charge} applies
unchanged and proves the capped joint charge.
\end{proof}

\begin{remark}[Exact Mersenne core-dichotomy payment]
Take $b=65450$.  On $130222\le e\le130226$, the high-core bound (CD1)
is at most $5161307<16777215$.  In the complementary branch,
$e=130222,130223$ force fourteen threshold-18 lines and
\[
 14\cdot9736-\binom{14}{2}\cdot5=135849>e.
\]
The next two supports force seventy threshold-16 lines and
\[
 70\cdot2041-\binom{70}{2}\cdot5=130795>e.
\]
Thus every support $130222\le e\le130225$ is safe.

At adjacent $e=130226$, the first threshold is 14 and has zero forced
core.  After 14763 zero-lower-bound peels, the capped budget is
$545032556=8412\cdot64785+61136$, the charge is $3199542$, and the
next pigeonhole threshold is one.  This is a method wall, not an unsafe
certificate; the Mersenne residual is now
$130226\le e\le1044241$.

The endpoint replay and independent rational audit are
\path{experimental/verify_mca_full_lift_core_dichotomy_capped_charge_v1.py}
and
\path{experimental/audit_mca_full_lift_core_dichotomy_capped_charge_v1.py}.
\end{remark}

\begin{corollary}[Exact-layer slot-core incidence]
\label{cor:full-lift-exact-layer-slot-core}
Suppose a line slot selected by the recursive bank belongs to exact layer
$h$ and contains $\lambda\ge2$ explanations.  If $u$ is their common core
inside the gauged $e$-coordinate direction support, then
\[
 u\ge\left\lceil\frac{\lambda h-e}{\lambda-1}\right\rceil. \tag{EL1}
\]
Consequently, in the low-core branch of
Corollary~\ref{cor:full-lift-core-dichotomy-capped-charge}, (EL1) may be used
as the forced lower bound in the capped convex charge.
\end{corollary}

\begin{proof}
Write the parameterized affine explanation line as
$c_\gamma=a+\gamma b$.  At an inside coordinate outside the common core,
the equation
\[
 r_0+\gamma r_1=a+\gamma b
\]
holds for at most one line parameter $\gamma$.  Every one of the $\lambda$
members has at least $h$ inside agreements, whereas a core coordinate is
counted $\lambda$ times.  Incidence counting in the fixed inside support
therefore gives
\[
 \lambda h\le e+(\lambda-1)u,
\]
which is equivalent to (EL1).  This is already an inside-core bound and
does not spend the outside zero allowance.  In the calibrated support
interval the parent ceiling has $H\ge m$, so the optional cross-layer top
slot is absent and every recursive-bank slot has one exact-layer owner.
\end{proof}

\begin{remark}[Exact Mersenne layer-aware payment]
Retain absorption cutoff $65450$.  In the complementary capped-core branch,
legal cutoffs $65516,\ldots,65521$ and (EL1) force three line cores of the
following sizes:
\[
\begin{array}{c|c|c|c}
e&b&\lambda&u\\ \hline
130226&65516&14&60540\\
130227&65516&5&49340\\
130228&65517&13&60126\\
130229&65517&4&43948\\
130230&65518&11&59048\\
130231&65518&4&43949\\
130232&65519&9&57431\\
130233,130234&65520&7&54736\\
130235,130236&65521&4&43951
\end{array}
\]
The displayed $\lambda$ is the forced minimum slot size and $h=b+1$ is the
minimum possible exact layer.  In this range the right side of (EL1) is
nondecreasing in both variables, so these substitutions are conservative.
Distinct selected lines have pairwise inside-core intersection at most
$c=5$.  Exact capped-charge replay therefore makes three lines violate
inside-core packing on every support $130226\le e\le130236$; the smallest
printed lower bound is
\[
 3\cdot43948-\binom32\cdot5=131829>130229.
\]

At adjacent $e=130237$, cutoff $65521$ forces only size-two affine
explanation slots.  Equation (EL1) gives $u=807$, and the first-order
packing expression has maximum
\[
 \max_s\left\{807s-5\binom{s}{2}\right\}=65529<e.
\]
After 7583 such lower bounds the capped charge is $882245$, the residual
target is $15894970$, and the next threshold is one.  This is a method wall,
not an unsafe certificate.  The Mersenne full-lift residual is now
$130237\le e\le1044241$; primitive shift-pair control is the next structural
obstruction.

The endpoint replay and independent rational audit are
\path{experimental/verify_mca_full_lift_exact_layer_slot_core_packing_v1.py}
and
\path{experimental/audit_mca_full_lift_exact_layer_slot_core_packing_v1.py}.
\end{remark}

\begin{corollary}[First-wall interpolation common-factor router]
\label{cor:full-lift-interpolation-common-factor-router}
At Mersenne full-lift support $e=130237$, give $(X,Y,Z)$ weights
$(1,5,5)$ and let $I_{264}$ be the vector space of weighted-degree at most
$264$ polynomials vanishing at all gauged inside received points
$(x,r_0(x),r_1(x))$.  Then
\[
 \dim I_{264}\ge 938.                              \tag{IF1}
\]
If the selected family is unsafe, the nonzero members of $I_{264}$ have a
common factor of positive $(Y,Z)$-degree over the algebraic closure of
$\F(X)$.
\end{corollary}

\begin{proof}
The weight-$(1,5,5)$ monomial count is
\[
 \sum_{s=0}^{52}(s+1)(264-5s+1)=131175.
\]
Evaluation at the $130237$ received points has rank at most $130237$, which
proves (IF1).

In the complementary capped-core branch, every selected line at cutoff
$65521$ has exact layer at least $65522$ and at least two members.  By
Corollary~\ref{cor:full-lift-exact-layer-slot-core}, its inside core has
size at least $807$.  Write the line as
$c_\gamma=a+\gamma b$.  Its inside core is precisely the set on which
$(r_0,r_1)=(a,b)$.

For $Q\in I_{264}$, substitution gives a polynomial
$Q(X,a(X),b(X))$ of degree at most $264$.  It has at least $807$ distinct
roots, hence vanishes identically.  Thus every selected pair $(a,b)$ is a
common $\F(X)$-rational zero of the interpolation kernel.

After $2704$ removed lines, the capped convex charge is $132203$ and the
residual target is $16645012$.  With base $13961576$ and group count
$1933560$, the next strict pigeonhole threshold remains two.  Therefore an
unsafe family forces $2705$ distinct polynomial-pair cores.

Suppose instead that the kernel has no common factor of positive
$(Y,Z)$-degree.  Over the algebraic closure of $\F(X)$, two generic kernel
members are then coprime.  Their total $(Y,Z)$-degrees are at most
$\lfloor264/5\rfloor=52$.  Affine Bezout bounds their common zero set by
$52^2=2704$, contradicting the $2705$ distinct selected pairs.
\end{proof}

\begin{remark}[Common-factor residual]
The corollary pays the coprime interpolation branch but does not identify
the residual factor with a split pencil.  A further theorem must classify
its ruled components or charge their line/core mass.  The Mersenne support
residual remains $130237\le e\le1044241$, with its first support reduced to
this common-factor branch.

The exact calibration and independent audit are
\path{experimental/verify_mca_full_lift_interpolation_common_factor_router_v1.py}
and
\path{experimental/audit_mca_full_lift_interpolation_common_factor_router_v1.py}.
\end{remark}

\begin{corollary}[Common-factor mass concentration]
\label{cor:full-lift-common-factor-mass}
In the unsafe branch of
Corollary~\ref{cor:full-lift-interpolation-common-factor-router}, let $P$ be
the full common interpolation factor and put
$d=\deg_{Y,Z}P$.  At least
\[
 7583-(52-d)^2\ge4982                              \tag{FM1}
\]
distinct selected polynomial pairs lie on $P$.  Moreover,
\[
 \#\{x\in E:P(x,r_0(x),r_1(x))=0\}\ge126188.       \tag{FM2}
\]
Thus at most $4049$ of the $130237$ inside coordinates are exceptional to
the factor relation.
\end{corollary}

\begin{proof}
After $7582$ low-core lines have been removed, the capped charge is
$881897$, the target is $15895318$, and the strict line-bank threshold is
still two.  Hence unsafety forces $7583$ distinct polynomial-pair cores.

Divide every interpolation-kernel member by the full gcd $P$.  The cofactor
family has gcd one and total $(Y,Z)$-degree at most $52-d$.  Two generic
cofactors are coprime, so affine Bezout permits at most $(52-d)^2$ selected
pairs outside $P=0$.  This proves (FM1), whose right side is minimized at
$d=1$.

For each of $t$ captured pairs, its inside core $U_i$ has size at least
$u=807$ and lies in the set on the left of (FM2).  Distinct pair cores meet
in at most five coordinates.  If $m_x$ is the core-incidence multiplicity,
then
\[
 \sum_xm_x\ge tu,
 \qquad
 \sum_x\binom{m_x}{2}\le5\binom{t}{2}.
\]
Cauchy's inequality therefore gives
\[
 \left|\bigcup_iU_i\right|
 \ge\left\lceil\frac{tu^2}{u+5(t-1)}\right\rceil.
\]
This is increasing in $t$.  At $t=4982$ it equals $126188$, proving (FM2).
\end{proof}

\begin{remark}
The set in (FM2) is not asserted to be a common core shared by all captured
sections.  Nor does the corollary assert that $P$ is irreducible, rational,
or a split pencil.  Its output is a degree-at-most-$52$ relation holding on
at least $96.89\%$ of the inside support and carrying at least $4982$
degree-five polynomial sections.

The exact calibration and all-degree audit are
\path{experimental/verify_mca_full_lift_common_factor_mass_router_v1.py}
and
\path{experimental/audit_mca_full_lift_common_factor_mass_router_v1.py}.
\end{remark}

\begin{corollary}[Common-factor weighted-degree bound]
\label{cor:full-lift-common-factor-weighted-degree}
In the setting of Corollary~\ref{cor:full-lift-common-factor-mass}, the
primitive full gcd satisfies
\[
 \operatorname{wdeg}_{(1,5,5)}P\le217,
 \qquad \deg_{Y,Z}P\le43.                       \tag{WD1}
\]
Consequently, if $\deg_{Y,Z}P\ge2$, then at least $5083$ selected
degree-five polynomial pairs lie on $P$, and
\[
 \#\{x\in E:P(x,r_0(x),r_1(x))=0\}\ge126266. \tag{WD2}
\]
Thus this higher-degree branch has at most $3971$ exceptional inside
coordinates.
\end{corollary}

\begin{proof}
Work over the algebraic closure.  Gauss's lemma makes the primitive
representative $P$ a polynomial divisor of every member of $I_{264}$.  If
$w=\operatorname{wdeg}_{(1,5,5)}P$, division by $P$ embeds the
at-least-$938$-dimensional kernel into weighted degree at most $264-w$.
The exact weighted-monomial count
\[
 M(D)=\sum_{s=0}^{\lfloor D/5\rfloor}(s+1)(D-5s+1)
\]
satisfies
\[
 M(46)=935<938\le990=M(47).
\]
Hence $264-w\ge47$, proving $w\le217$ and
$\deg_{Y,Z}P\le\lfloor217/5\rfloor=43$.

In the higher-degree branch, the factor-mass bound is minimized at $d=2$:
\[
 7583-(52-d)^2\ge7583-50^2=5083.
\]
Applying the same core-incidence inequality as in
Corollary~\ref{cor:full-lift-common-factor-mass} gives
\[
 \left\lceil
 \frac{5083\cdot807^2}{807+5(5083-1)}
 \right\rceil=126266,
\]
which proves (WD2).
\end{proof}

\begin{remark}
The full gcd in this corollary may be reducible.  No irreducible-component
classification or safety claim is implicit.  The exact threshold replay and
independent enumeration are
\path{experimental/verify_mca_full_lift_common_factor_weighted_degree_bound_v1.py}
and
\path{experimental/audit_mca_full_lift_common_factor_weighted_degree_bound_v1.py}.
\end{remark}

\begin{corollary}[Base-field component descent]
\label{cor:full-lift-common-factor-base-field-descent}
Retain the higher-degree branch of
Corollary~\ref{cor:full-lift-common-factor-weighted-degree}.  Over the
deployed Mersenne field
\[
 \F=\F_{p^4},\qquad p=2^{31}-1,
\]
let $P_{\rm bf}$ be the product of the geometric irreducible components of
$P$ that are individually defined over $\F(X)$.  At least $5079$ selected
polynomial pairs lie on $P_{\rm bf}$, and one absolutely irreducible
$\F(X)$-component contains at least $132$ selected pairs.  Moreover,
\[
 \#\{x\in E:P_{\rm bf}(x,r_0(x),r_1(x))=0\}\ge126263. \tag{FD1}
\]
Thus at most $3974$ inside coordinates are exceptional to the base-field
component union.
\end{corollary}

\begin{proof}
Factor the radical of $P$ geometrically as $\prod_iR_i$, with
$\delta_i=\deg_{Y,Z}R_i$ and $\sum_i\delta_i\le d\le43$.  Since
$p>d$, a component field of definition cannot have nontrivial purely
inseparable degree.  If $R_i$ is not defined over $\F(X)$, it therefore
has a distinct conjugate $R_i^\sigma$.  Every $\F(X)$-rational point of
$R_i$ also lies on $R_i^\sigma$, so Bezout bounds their number by
$\delta_i^2$.

Assign each selected pair outside the base-field component union to one
non-base-field component containing it.  At most
\[
 \sum_{R_i\ {
m not\ over}\ \F(X)}\delta_i^2\le d^2
\]
pairs are assigned.  Hence the retained number is at least
\[
 7583-(52-d)^2-d^2=4879+104d-2d^2.
\]
This concave expression has endpoint values $5079$ and $5653$ on
$2\le d\le43$, proving the uniform $5079$ bound.  There are at most $d$
components, and exact integer minimization gives
\[
 \min_{2\le d\le43}
 \left\lceil\frac{7583-(52-d)^2-d^2}{d}\right\rceil=132.
\]
Finally, the core-incidence bound at $t=5079$ gives
\[
 \left\lceil
 \frac{5079\cdot807^2}{807+5(5079-1)}
 \right\rceil=126263,
\]
which proves (FD1).
\end{proof}

\begin{remark}
This is a base-field normalization, not an irreducibility or split-pencil
theorem.  The exact degree scan and independent degree-partition audit are
\path{experimental/verify_mca_full_lift_common_factor_base_field_descent_v1.py}
and
\path{experimental/audit_mca_full_lift_common_factor_base_field_descent_v1.py}.
\end{remark}

\begin{corollary}[Linear factor gives a projective star]
\label{cor:full-lift-linear-factor-projective-star}
In the setting of Corollary~\ref{cor:full-lift-common-factor-mass}, suppose
the full common interpolation factor has total $(Y,Z)$-degree one.  Then,
after projective rescaling, it has a primitive equation
\[
 P(X,Y,Z)=AY+BZ+C(X),
 \qquad A,B\in\F,\quad (A,B)\ne(0,0),\quad \deg C\le5. \tag{LS1}
\]
All captured affine explanation lines share one $\F$-rational projective
slope--codeword center.  If $A\ne0$, the center is finite:
\[
 \gamma_*=B/A,qquad c_*=-C/A.
\]
If $A=0$, every captured line has the same direction codeword $-C/B$, the
center at projective slope infinity.
\end{corollary}

\begin{proof}
Initially work over the algebraic closure and write
$P=A(X)Y+B(X)Z+C(X)$ primitively.  The factor-mass corollary supplies at
least $4982$ distinct sections $(a_i,b_i)\in\F[X]_{<6}^2$.  Subtracting two
section equations and removing $\gcd(A,B)$ gives
\[
 (a_i,b_i)=(a_0+Bt_i,b_0-At_i),
 \qquad
 \deg t_i\le s:=5-\max(\deg A,\deg B).             \tag{LS2}
\]
Indeed, one section shows that $\gcd(A,B)$ also divides $C$, so primitivity
makes it a unit.  In particular $\deg A,\deg B\le5$.

Since $A$ and $B$ are coprime, the received pair induces a scalar word
$\tau$ on every inside coordinate: use
$(r_0-a_0)/B$ where $B\ne0$, and $-(r_1-b_0)/A$ otherwise.  Each $t_i$
agrees with $\tau$ on its core of at least $807$ coordinates.  The ordinary
constant-block Johnson cap for degree $s$ is
\[
 J_s=\left\lfloor
 \frac{130237(807-s)}{807^2-130237s}\right\rfloor.
\]
For $s=0,1,2,3,4$ its exact values are
$161,201,268,401,802$.  If either $A$ or $B$ were nonconstant, then
$s\le4$, contradicting $4982>802$.  Hence $A,B$ are constants and
$\deg C\le5$.

If $A\ne0$, two distinct captured $\F$-rational pairs give
$B/A=-(a_i-a_j)/(b_i-b_j)\in\F$, and then
$-C/A=a_i+(B/A)b_i\in\F[X]_{<6}$.  Every captured line passes through that
finite center.  If $A=0$, the common direction $-C/B$ lies in
$\F[X]_{<6}$, proving the projective-infinity case.
\end{proof}

\begin{remark}
This identifies the degree-one branch with the primitive projective-star
shape; it does not prove the star population bound.  The complementary
common-factor branch has $(Y,Z)$-degree at least two.

The exact Johnson replay and independent audit are
\path{experimental/verify_mca_full_lift_linear_factor_projective_star_router_v1.py}
and
\path{experimental/audit_mca_full_lift_linear_factor_projective_star_router_v1.py}.
\end{remark}

\begin{theorem}[Full-explanation lifted-rank gauge dichotomy]
\label{thm:full-explanation-lifted-rank-dichotomy}
Let \(C\leq\F^D\) have dimension \(K\), let
\(r_\gamma=r_0+\gamma r_1\), and let \(Z\subseteq\F\) contain distinct
slopes.  For every \(\gamma\in Z\), choose \(c_\gamma\in C\), and assume
that the \(c_\gamma\) have full affine rank \(K\).  Fix
\(\gamma_0\in Z\) and put
\[
 V=\operatorname{span}\{(\gamma-\gamma_0,
                 c_\gamma-c_{\gamma_0}):\gamma\in Z\}
       \leq \F\oplus C.
\]
Then \(\dim V\in\{K,K+1\}\), and exactly one of the following holds.
\begin{enumerate}
\item If \(\dim V=K\), there is a unique nonzero linear functional
  \(\ell:C\to\F\) such that
  \(V=\{(\ell(u),u):u\in C\}\).  The gauge
  \[
     r_1\longmapsto r_1-b,
     \qquad c_\gamma\longmapsto c_\gamma-\gamma b
  \]
  drops the explanation affine rank to exactly \(K-1\) precisely for
  \(b\in C\) satisfying \(\ell(b)=1\).  Thus the rank-dropping gauges form
  an affine hyperplane in \(C\).
\item If \(\dim V=K+1\), then \(V=\F\oplus C\), and every \(b\in C\)
  leaves the explanation affine rank equal to \(K\).
\end{enumerate}
If, in addition, every chosen maximal agreement support is same-support
pair-noncontained, then \(r_1\notin C\), and the selected error vectors
\(r_\gamma-c_\gamma\) have affine rank exactly \(\dim V\).
\end{theorem}

\begin{proof}
Projection of \(V\) to \(C\) is onto because the explanations have full
affine rank.  Hence \(K\leq\dim V\leq K+1\).  If \(\dim V=K\), projection
is an isomorphism and \(V\) is the graph of a unique functional \(\ell\).
It is nonzero because the slopes are not all equal.  On anchored
differences, the gauge is the map
\[
  T_b:V\longrightarrow C,\qquad T_b(a,u)=u-ab.
\]
For \(V=\{(\ell(u),u):u\in C\}\), this map is
\(u\mapsto u-\ell(u)b\).  It has rank \(K-1\) exactly when
\(\ell(b)=1\), in which case its image is \(\ker\ell\); otherwise it is
invertible.  If \(\dim V=K+1\), then \(V=\F\oplus C\), and \(T_b\) is
surjective for every \(b\).

Suppose now that \(r_1\in C\).  On any selected agreement support,
\(r_1\) explains the direction and \(c_\gamma-\gamma r_1\in C\) explains
the base, contradicting pair noncontainment.  Thus \(r_1\notin C\).  The
linear map
\[
  (a,u)\longmapsto ar_1-u
\]
is consequently injective on \(\F\oplus C\), and it sends the anchored
lifted differences to the anchored error-vector differences.  Their affine
rank is therefore \(\dim V\).
\end{proof}

\begin{theorem}[Full-lift near-MDS extension reduction]
\label{thm:full-lift-near-mds-extension}
Retain the hypotheses of
\cref{thm:full-explanation-lifted-rank-dichotomy}, assume
\(\dim V=K+1\), write \(N=|D|\), and put
\[
  W=C+\langle r_1\rangle,\qquad
  e=\min_{b\in C}\wt(r_1-b).
\]
Then \(\dim W=K+1\), and the generalized Hamming weights of \(W\) are
\[
  d_1(W)=e,\qquad
  d_j(W)=N-K+j-1\quad(2\le j\le K+1).
  \tag{NM}
\]
The selected errors \(r_\gamma-c_\gamma\) lie in one affine coset of
\(W\), have full affine rank \(K+1\), and determine their slopes uniquely.
Moreover, if every selected maximal agreement support has size greater than
\(K\), restriction of \(W\) to each such support is injective.
\end{theorem}

\begin{proof}
Pair noncontainment gives \(r_1\notin C\), so \(W\) has dimension \(K+1\).
Interpolation on \(K\) coordinates gives \(e\le N-K\).  Every nonzero word
of \(W\) is either a nonzero RS word, of weight at least \(N-K+1\), or a
nonzero scalar multiple of \(r_1-b\), of weight at least \(e\); a nearest
\(b\) attains \(e\).  Thus \(d_1(W)=e\).

Let \(A\le W\) have dimension \(j\ge2\).  If \(A\le C\), the RS MDS
hierarchy gives \(|\operatorname{supp}A|\ge N-K+j\).  Otherwise
\(\dim(A\cap C)=j-1\), and that RS subcode has at most \(K-j+1\) common
zero coordinates.  Hence
\[
  |\operatorname{supp}A|\ge N-K+j-1.
\]
The generalized Singleton bound for the \([N,K+1]\) code \(W\) is the
reverse inequality, proving \((\mathrm{NM})\).

The map \((a,u)\mapsto ar_1-u\) is an isomorphism
\(\F\oplus C\to W\), so the full lifted rank is exactly the error affine
rank and equal errors have equal slopes.  Finally, a nonzero
\(ar_1-u\in W\) vanishing on a selected support either contradicts the RS
root bound when \(a=0\), or supplies simultaneous direction and base
explanations when \(a\ne0\).  Pair noncontainment excludes both cases.
\end{proof}

\begin{remark}[Exact top-rank split at the first shortened rows]
\label{rem:full-explanation-top-rank-split}
The gauge preserves agreement supports, pair noncontainment, and
\(e=\min_{b\in C}\wt(r_1-b)\).  Therefore the \(\dim V=K\) branch may use
\cref{thm:proper-subspace-mca} at explanation rank \(K-1\).  Combining
this with the already separate low-support and recursive-direction gates
gives the following exact routing table; \(q\) is explanation affine rank
and \(h=\dim V\), equivalently the error affine rank.
\[
\begin{array}{c|c|c|c}
\text{row}&q&h&\text{currently paid support range}\\ \hline
\text{KoalaBear}&14&14&e\le64047\ \text{or}\ e\ge992852\\
\text{KoalaBear}&14&15&e\le96150\ \text{or}\ e\ge1044239\\
\text{Mersenne-31}&6&6&e\le65454\ \text{or}\ e\ge1037876\\
\text{Mersenne-31}&6&7&e\le130236\ \text{or}\ e\ge1044242
\end{array}
\]
The proper-subspace/gauge improvement is only in the \(h=K\) rows.  The
separate full-lift continuations extend the printed low walls in the
\(h=K+1\) rows, but their remaining middle intervals stay unpaid.
After the sparse-direction, centered-Gram, exact-layer, and interpolation
continuations, the full-lift intervals are
\(96151\le e\le1044238\) for KoalaBear and
\(130237\le e\le1044241\) for Mersenne-31.  The latter first support is
not paid: \cref{cor:full-lift-interpolation-common-factor-router} routes it
to the separately declared common-factor branch.
\Cref{thm:full-lift-near-mds-extension} shows that every generalized
weight from rank two onward is already MDS-sharp.  Even at the best endpoint
\(e=N-K\), the generic corrected compiler gives
\[
  743896698428332665>274980728111395087
  \quad\text{and}\quad
  219426634>16777215
\]
on KoalaBear and Mersenne-31, respectively.  Closing these cells therefore
requires structure of the codimension-one RS extension beyond its weight
hierarchy, or a sharper row-specific affine-list theorem.  A
finite-field audit and the exact adjacent arithmetic are in
\texttt{experimental/verify\_mca\_full\_explanation\_lifted\_rank\_gauge\_dichotomy\_v1.py}.
\end{remark}

\begin{theorem}[Directional Johnson ray compiler]
\label{thm:directional-johnson-ray}
Retain the hypotheses and notation of \cref{thm:fixed-union-ray}, but allow
\(\nu\ge0\).  Choose a lift \(b_1\in\F^U\) of \(y_1\), put
\(C_U=\ker H_U\), and define the lift-invariant directional agreement
\begin{equation}
 \alpha_U(y_1)=
 \max_{c\in C_U}|\{x\in U:b_1(x)+c(x)=0\}|.
 \label{eq:directional-agreement}
\end{equation}
If \(h=N-t\) and
\begin{equation}
 h^2>N\alpha_U(y_1),                                \label{eq:directional-johnson-condition}
\end{equation}
then
\begin{equation}
 |Z|\le
 \left\lfloor
 \frac{N(h-\alpha_U(y_1))}
 {h^2-N\alpha_U(y_1)}
 \right\rfloor .                                  \label{eq:directional-johnson-ray}
\end{equation}
Moreover, with
\[
 d_U(y_1)=\min\{\wt(e):e\in\F^U,\ H_Ue=y_1\},
\]
one has
\begin{equation}
 \alpha_U(y_1)=N-d_U(y_1)\ge\nu.                  \label{eq:direction-distance}
\end{equation}
Equality holds exactly when the direction syndrome has no representation
on fewer than \(R\) columns of \(U\).
\end{theorem}

\begin{proof}
Changing the lift \(b_1\) by a kernel word merely permutes the maximization
in \eqref{eq:directional-agreement}, so \(\alpha_U(y_1)\) is intrinsic.
Choose a lift \(b_0\) of \(y_0\).  For every slope write
\[
 e_\gamma=b_0+\gamma b_1+k_\gamma,
 \qquad k_\gamma\in C_U,
\]
and choose an \(h\)-set \(A_\gamma\) of zero coordinates of
\(e_\gamma\).  On \(A_\gamma\), the codeword \(-k_\gamma\) agrees with
\(b_0+\gamma b_1\).  If \(\gamma\ne\delta\), then on
\(A_\gamma\cap A_\delta\),
\[
 b_1=\frac{-k_\gamma+k_\delta}{\gamma-\delta},
\]
which is a codeword of \(C_U\).  Hence
\(|A_\gamma\cap A_\delta|\le\alpha_U(y_1)\).  The same Johnson
incidence count as in \cref{thm:fixed-union-list-johnson}, with
\(\alpha_U(y_1)\) in place of \(\nu-1\), gives
\eqref{eq:directional-johnson-ray}.

Finally, the words \(b_1+c\), \(c\in C_U\), are exactly all lifts of
\(y_1\) supported in \(U\).  Maximizing their number of zero coordinates
is therefore equivalent to minimizing their weight, proving the equality
in \eqref{eq:direction-distance}.  Any \(R\) columns of \(H_U\) span
syndrome space, so \(d_U(y_1)\le R\) and \(\alpha_U(y_1)\ge N-R=\nu\).
The equality criterion is immediate.
\end{proof}

\begin{corollary}[Rank-regular fixed-union ray bound]
\label{cor:rank-regular-fixed-union-ray}
In \cref{thm:directional-johnson-ray}, suppose the direction syndrome is
rank-regular on \(U\), meaning \(d_U(y_1)=R\).  If
\((R+\nu-t)^2>(R+\nu)\nu\), then
\begin{equation}
 |Z|\le
 \left\lfloor
 \frac{(R+\nu)(R-t)}
 {(R+\nu-t)^2-(R+\nu)\nu}
 \right\rfloor .                                  \label{eq:rank-regular-fixed-union-ray}
\end{equation}
Thus every failure of this stronger low-nullity payment is either a
high-nullity chart or an explicit direction-rank-defect cell
\(d_U(y_1)<R\).
\end{corollary}

\begin{proof}
By \eqref{eq:direction-distance}, rank regularity is exactly
\(\alpha_U(y_1)=\nu\).  Substitute this value into
\eqref{eq:directional-johnson-ray}.
\end{proof}

\begin{corollary}[Quantitative direction-rank-defect payment]
\label{cor:quantitative-direction-defect}
Retain the notation of \cref{thm:directional-johnson-ray}, and write
\[
 d=R-t,\qquad
 j=R-d_U(y_1),\qquad
 N=R+\nu,\qquad
 D_0=(d+\nu)^2-N\nu.
\]
Thus \(j\ge0\) measures the direction-rank defect and
\(\alpha_U(y_1)=\nu+j\).  Whenever \(D_0-Nj>0\),
\begin{equation}
 |Z|\le
 \left\lfloor\frac{N(d-j)}{D_0-Nj}\right\rfloor .
 \label{eq:direction-defect-ray}
\end{equation}
For an integer budget \(B>1\), define
\begin{equation}
 J_B(\nu)=\min\left\{
 d-1,
 \left\lfloor\frac{D_0-1}{N}\right\rfloor,
 \left\lfloor\frac{BD_0-Nd}{N(B-1)}\right\rfloor
 \right\}.
 \label{eq:direction-defect-threshold}
\end{equation}
If \(0\le j\le J_B(\nu)\), then the whole fixed-union ray chart costs at
most \(B\) slopes.  Hence the residual direction cell at a fixed nullity
can be stated quantitatively as \(j>J_B(\nu)\), rather than merely as
\(j>0\).
\end{corollary}

\begin{proof}
By \eqref{eq:direction-distance},
\(\alpha_U(y_1)=N-d_U(y_1)=\nu+j\).  Substitution into
\eqref{eq:directional-johnson-ray} gives
\eqref{eq:direction-defect-ray}.  The second term in
\eqref{eq:direction-defect-threshold} makes its denominator positive,
and the first makes its numerator positive.  The final term is exactly the
rearrangement
\[
 N(d-j)\le B(D_0-Nj).
\]
Thus the real fraction in \eqref{eq:direction-defect-ray} is at most
\(B\), and so is its floor.
\end{proof}


\begin{theorem}[Paving-basis direction-defect ray compiler]
\label{thm:paving-ray}
Retain the fixed-union MCA notation of
\cref{thm:fixed-union-ray,thm:directional-johnson-ray}.  Write
\[
 N=R+\nu,
 \qquad h=N-t,
 \qquad d=R-t,
 \qquad j=R-d_U(y_1).
\]
Assume \(0\le j<d\), equivalently that the direction word is outside the
agreement ball on this chart.  Then
\begin{equation}
 |Z|\le
 \left\lfloor
 \frac{(\nu+1)\binom{R+\nu}{\nu+1}}
 {(R-t-j)\binom{R+\nu-t}{\nu}}
 \right\rfloor .
 \label{eq:paving-ray}
\end{equation}
This estimate is independent of the Johnson denominator and is strongest
for small nullity, including large direction defect.
\end{theorem}

\begin{proof}
Choose a minimum-weight lift \(q\in\F^U\) of the direction syndrome, so
\(\wt(q)=R-j\), and choose a basis \(k_1,\ldots,k_\nu\) of the MDS kernel
\(C_U=\ker H_U\).  In the parameter space with coordinates
\((\gamma,\lambda_1,\ldots,\lambda_\nu)\), agreement at coordinate \(x\)
has normal
\[
                       v_x=(q_x,k_1(x),\ldots,k_\nu(x)).
\]
Every set of at most \(\nu\) such normals is independent, because its
projection to the last \(\nu\) coordinates is a set of columns of a
generator matrix of the \([R+\nu,\nu,R+1]\) MDS code \(C_U\).

Every rank-\(\nu\) flat of these normals contains at most \(\nu+j\)
coordinates.  Indeed, such a flat is annihilated by a nonzero functional
\((\delta,a)\).  If \(\delta\ne0\), the vanishing coordinates are zeros of
a lift \(\delta q+c\) of \(\delta y_1\), which has weight at least
\(R-j\).  If \(\delta=0\), they are zeros of a nonzero MDS kernel word and
number at most \(\nu-1\).

Fix a slope point and choose \(h\) incident coordinate hyperplanes.  Put
\(r=\nu+1\) and \(\alpha=\nu+j\).  Every \((r-1)\)-subset of their normals
is independent.  A dependent \(r\)-subset is counted by each of its \(r\)
subsets of size \(r-1\), while the closure of any such subset contains at
most \(\alpha\) coordinate normals.  Hence the number of independent
\(r\)-subsets among the \(h\) incident normals is at least
\[
 \binom hr-\frac{\alpha-r+1}{r}\binom h{r-1}
 =\frac{h-\alpha}{r}\binom h{r-1}
 =\frac{d-j}{\nu+1}\binom{R+\nu-t}{\nu}.
\]
Globally there are at most \(\binom{R+\nu}{\nu+1}\) independent coordinate
bases, and each basis cuts out at most one parameter point.  Double
counting proves \eqref{eq:paving-ray}.
\end{proof}

\begin{corollary}[Exact paving defect threshold]
\label{cor:paving-threshold}
For an integer budget \(B\ge0\), define
\begin{equation}
 J_B^{\rm pav}(\nu)=
 \max\left\{-1,
 d-1-
 \left\lfloor
 \frac{(\nu+1)\binom{R+\nu}{\nu+1}}
 {(B+1)\binom{R+\nu-t}{\nu}}
 \right\rfloor
 \right\}.
 \label{eq:paving-threshold}
\end{equation}
Every separated fixed-union MCA chart with
\(0\le j\le J_B^{\rm pav}(\nu)\) costs at most \(B\) slopes.  Combining
this with \cref{cor:quantitative-direction-defect}, the exact certified
defect range is
\[
                0\le j\le
                \max\{J_B(\nu),J_B^{\rm pav}(\nu)\}.
\]
\end{corollary}

\begin{proof}
The floor in \eqref{eq:paving-ray} is at most \(B\) exactly when its
numerator is strictly smaller than \(B+1\) times its denominator.  Solving
this integer inequality for \(j\) gives \eqref{eq:paving-threshold}.
\end{proof}




\section{Minimum-distance sharpening and the exact primitive wedge}
\label{sec:min-distance-sharpening}
Retain the notation of \cref{thm:support-flag-recursion}.  Let
\[
 \mu=d_1(C')=\min_{0\ne c\in C'}|\supp(c)|.
\]
The generalized Singleton bound gives
\[
 \mu\ge n-K+1=R+1,
 \qquad R=n-K.
\]
The general flag theorem uses only the ambient MDS lower bound for each flag-annihilator direction.  Every such direction actually belongs to $C'$, so the true minimum support gives a sharper resource count.

\begin{theorem}[Minimum-distance enhanced flag resources]
\label{thm:min-distance-flag-resources}
Retain the hypotheses and notation of \cref{thm:support-flag-recursion}.  For a profile $\mathbf b=(b_1,\ldots,b_r)$ define
\begin{align}
 P_{\mathbf b}^{\sharp}
 &=\prod_{i=1}^{j}
 \left(f_i-t+b_0+\sum_{q=1}^{r}b_q\right),\label{eq:min-flag-P}\\
 Q_{\ell}^{\sharp}(\mathbf b)
 &=\mu-t+b_0+\sum_{q>\ell}b_q,\label{eq:min-flag-Q}\\
 A_\ell(\mathbf b)&=e_\ell-b_\ell,
 \qquad
 B_{\ell}^{\sharp}(\mathbf b)
 =\bigl(Q_{\ell}^{\sharp}(\mathbf b)-e_\ell+b_\ell\bigr)_+ .\label{eq:min-flag-AB}
\end{align}
For $\sigma\in\{T,U\}^r$, put
\begin{align}
 C_{\sigma}^{\sharp}(\mathbf b)
 &=P_{\mathbf b}^{\sharp}
 \prod_{\ell:\sigma_\ell=T}A_\ell(\mathbf b)
 \prod_{\ell:\sigma_\ell=U}B_{\ell}^{\sharp}(\mathbf b),\label{eq:min-flag-C}\\
 G_\sigma
 &=d_0^{\underline j}
 \prod_{\ell:\sigma_\ell=T}e_\ell
 \prod_{\ell:\sigma_\ell=U}
 \bigl(d_{\ell-1}-(j+\ell-1)\bigr).\label{eq:min-flag-G}
\end{align}
Then
\begin{equation}
 \boxed{
 \sum_{\mathbf b}L_{\mathbf b}C_{\sigma}^{\sharp}(\mathbf b)
 \le G_\sigma
 }
 \label{eq:min-flag-resource}
\end{equation}
for every binary resource pattern.  Consequently every nonnegative family $x_\sigma$ satisfying
\[
 \sum_\sigma x_\sigma C_{\sigma}^{\sharp}(\mathbf b)\ge1
 \quad\text{for every profile }\mathbf b
\]
gives the whole-chart bound
\begin{equation}
 \boxed{
 L\le\left\lfloor\sum_\sigma x_\sigma G_\sigma\right\rfloor .
 }
 \label{eq:min-flag-dual}
\end{equation}
The tensor and all-extender corollaries remain valid after replacing $w+1+b_0$ by $\mu-t+b_0$ and replacing every old-support factor by \eqref{eq:min-flag-AB}.
\end{theorem}

\begin{proof}
Fix a flag level $\ell$ and an adapted coordinate basis through $V_{\ell-1}$.  Its annihilator in $V_\ell$ is generated by a nonzero word
\[
 h_\ell\in V_\ell\setminus V_{\ell-1}\subseteq C'.
\]
Therefore $\wt(h_\ell)\ge\mu$.  The word vanishes outside $U_\ell$.  The number of mismatches outside $U_\ell$ is
\[
 b_{\rm out}=b_0+\sum_{q>\ell}b_q,
\]
and all of those coordinates are zeros of $h_\ell$.  Among the zero coordinates of $h_\ell$, at most
\[
 (n-\mu)-b_{\rm out}
\]
are agreements.  Since every listed point has at least $m$ agreements, at least
\[
 m-\bigl((n-\mu)-b_{\rm out}\bigr)
 =\mu-t+b_{\rm out}
 =Q_\ell^{\sharp}(\mathbf b)
\]
agreement coordinates are nonzeros of $h_\ell$ and hence extend the adapted basis.  The $A_\ell=e_\ell-b_\ell$ agreeing coordinates in the new layer all extend, so at least $B_\ell^{\sharp}$ extending agreements lie in the old support.  The global basis capacities are unchanged.  The double count and finite dual are therefore identical to the proof of \cref{thm:support-flag-recursion} with the sharpened factors.
\end{proof}

\begin{corollary}[Minimum-word all-extender bound]
\label{cor:min-word-all-extender}
Choose $V_0=\langle h\rangle$, where $h$ is a minimum-weight word of $C'$.  If $s=\dim C'$, then
\begin{equation}
 \boxed{
 L\le
 \left\lfloor
 \frac{\mu\prod_{\ell=1}^{s-1}(d_\ell-\ell)}
      {(\mu-t+b_0)^s}
 \right\rfloor
 \le
 \left\lfloor
 \frac{\mu(n-s+1)^{s-1}}
      {(\mu-t+b_0)^s}
 \right\rfloor .
 }
 \label{eq:min-word-bound}
\end{equation}
Writing $\mu=R+a$, $a\ge1$, the soft form is
\begin{equation}
 L\le
 \left\lfloor
 \frac{(R+a)(n-s+1)^{s-1}}
      {(w+a+b_0)^s}
 \right\rfloor .
 \label{eq:min-word-soft}
\end{equation}
\end{corollary}

\begin{proof}
A generalized-weight minimizer is support-saturated by \cref{lem:weight-minimizer-saturated}, so the minimum-weight line admits a support-saturated flag.  Apply the all-extender form of \cref{thm:min-distance-flag-resources}.  The initial basis count is $\mu$, and $\mu-t+b_0$ is a lower bound at every extension level.  Since every remaining support layer is nonempty and $d_{s-1}\le n$, one has $d_\ell-\ell\le n-s+1$.
\end{proof}

\begin{proposition}[Exact active minimum-support transitions]
\label{prop:min-support-transitions}
At $b_0=0$, \eqref{eq:min-word-soft} pays every first-unpaid affine chart at and above the following support excesses:
\[
\begin{array}{c|r|r|r|r}
\toprule
\text{row}&s&a_{\rm paid}&\text{bound at }a_{\rm paid}&\text{bound at }a_{\rm paid}-1\\
\midrule
\text{KoalaBear MCA}&15&70832&274974324324469113&275003903412638889\\
\text{KoalaBear list}&15&70832&274974569967317804&275004149082130935\\
\text{Mersenne-31 MCA}&7&111529&16776653&16777295\\
\text{Mersenne-31 list}&7&111529&16776668&16777309\\
\bottomrule
\end{array}
\]
Thus a surviving zero-mismatch chart has $a\le70831$ on KoalaBear or $a\le111528$ on Mersenne-31.  Positive $b_0$ only shrinks this residual.
\end{proposition}

\begin{proposition}[Exact mismatch curve]
\label{prop:min-mismatch-curve}
For selected support excesses, the least outside mismatch paid by \eqref{eq:min-word-soft} is
\[
\begin{array}{c|rrrrr}
\toprule
& a=1&a=1000&a=10000&a=50000&a=70000\\
\midrule
\text{KoalaBear MCA/list}&70230&69240&60318&20659&825\\
\text{Mersenne-31 MCA/list}&108962&107987&99202&60141&40599\\
\bottomrule
\end{array}
\]
The curve is strictly decreasing in the actual direction-code support excess.
\end{proposition}

\begin{remark}[Combination with the actual first generalized weight]
Let the full support have size $d_s=R+\nu$.  Replacing the ambient lower bound in the common-zero Johnson compiler by $d_1(C')=R+a$ gives
\begin{equation}
 L\le
 \left\lfloor
 \frac{(R+\nu)(w+a+b_0)}
 {(w+\nu+b_0)^2-(R+\nu)(\nu-a)}
 \right\rfloor
 \label{eq:actual-d1-johnson}
\end{equation}
when the denominator is positive.  Combining \eqref{eq:actual-d1-johnson} with the nonuniform form of \eqref{eq:min-word-bound} gives the exact residual wedge in $(a,\nu,b_0)$.
\end{remark}

\begin{theorem}[Fixed-mismatch puncturing compiler]
\label{thm:fixed-mismatch-puncturing}
Let $\mathcal L\subseteq\mathcal A$ be a family of codewords, each agreeing with a received word $u$ on at least $m$ coordinates.  Suppose a fixed set $Z\subseteq D$, $|Z|=z$, is a mismatch set for every member of $\mathcal L$.  Then $z\le t$, and if the affine span of $\mathcal L$ has direction dimension at most $s$, then
\begin{equation}
 \boxed{
 |\mathcal L|
 \le
 \left\lfloor
 \frac{\binom{R-z+s}{s}}
      {\binom{w+s}{s}}
 \right\rfloor .
 }
 \label{eq:fixed-mismatch-puncturing}
\end{equation}
\end{theorem}

\begin{proof}
Every listed word has at most $t$ mismatches, so $z\le t$.  Puncture the coordinates in $Z$.  Since $z<t<R=n-K$, evaluation of degree-$<K$ polynomials remains injective on $D\setminus Z$, no agreement is deleted, and the affine direction dimension does not increase.  Apply the recursive affine-span compiler to the punctured code of length $n-z$.
\end{proof}

\begin{corollary}[Active puncture thresholds]
\label{cor:active-puncture-thresholds}
At the first unpaid affine dimensions, a subfamily with a common fixed mismatch set is paid at the following first integers:
\[
\begin{array}{c|r|r|r}
\toprule
\text{row}&z_{\rm paid}&\text{bound at }z_{\rm paid}&\text{bound at }z_{\rm paid}-1\\
\midrule
\text{KoalaBear MCA}&67313&274978201175970494&274982404608707323\\
\text{KoalaBear list}&67314&274978201175970494&274982404608707323\\
\text{Mersenne-31 MCA}&322321&16777133&16777295\\
\text{Mersenne-31 list}&322322&16777133&16777295\\
\bottomrule
\end{array}
\]
\end{corollary}

\begin{proposition}[Exchange defect dominates support excess]
\label{prop:exchange-defect}
In one boundary prefix fiber, let the associated codeword direction space have minimum support $R+a$.  For distinct supports $M,M'$ of size $m$, write
\[
 e=|M\setminus M'|=|M'\setminus M|=w+r.
\]
Then
\begin{equation}
 \boxed{r\ge a.}
 \label{eq:exchange-defect}
\end{equation}
\end{proposition}

\begin{proof}
The prefix-rigidity identity of CAP25 v13.2 gives
\[
 \Lambda_M-\Lambda_{M'}=G(A-B),
 \qquad \deg(A-B)\le r-1,
\]
where $G$ is the locator of $M\cap M'$ and has degree $m-e$.  The corresponding nonzero codeword difference vanishes on the $m-e$ common support coordinates, so its weight is at most
\[
 n-(m-e)=t+e=R+r.
\]
By the definition of $a$, its weight is at least $R+a$.  Hence $r\ge a$.
\end{proof}

\begin{corollary}[Top-seam localization]
\label{cor:top-seam-localization}
The top seam $e=w+1$ can occur only in the exact-minimum-distance branch $a=1$.  Every chart with $a\ge2$ has first possible exchange defect at least $a$.
\end{corollary}

\begin{theorem}[Exact primitive residual after minimum-distance payment]
\label{thm:min-distance-residual}
After all unconditional first-match payments in Parts~I--IV, an over-budget first-unpaid Q or list chart must lie in the following simultaneous wedge:
\begin{enumerate}[label=\textup{(R\arabic*)}]
\item its support excess $a$ lies below the exact curve \eqref{eq:min-word-soft};
\item its full-support nullity lies beyond both \eqref{eq:actual-d1-johnson} and the minimum-word flag bound;
\item it contains no fixed mismatch block meeting \cref{thm:fixed-mismatch-puncturing};
\item every boundary neighbour has exchange defect at least $a$, and the top seam is restricted to $a=1$;
\item it is not in an earlier quotient, occupancy, extension, common-zero, saturation, or rank-defect cell.
\end{enumerate}
\end{theorem}


\section{The quadratic mean-overlap staircase}\label{sec:mean-overlap-v3}
\begin{lemma}[Large-overlap collapse]
\label{lem:large-overlap-collapse}
Fix a received pair and an exact agreement \(a=n-r\ge k+1\).  Choose one
noncommon exact witness support \(S_z\) and explaining polynomial \(p_z\)
for every slope \(z\) in its bad set \(Z\).  If two distinct chosen
supports satisfy
\[
                    \abs{S_z\cap S_w}\ge k+r,       \tag{MO1}
\label{eq:large-overlap-collapse}
\]
then \(\abs Z\le r+1\).
\end{lemma}

\begin{proof}
Put
\[
 c_1=\frac{p_z-p_w}{z-w},\qquad c_0=p_z-zc_1.
\]
On \(S_z\cap S_w\), the two line equations imply that the received pair
equals \((c_0,c_1)\).  Let \(C_0\) be the full common support of this
codeword pair and write \(c=\abs{C_0}\ge k+r\).  For every \(u\in Z\),
\(\abs{S_u\cap C_0}\ge a+c-n\ge k\), so the explaining polynomial
\(p_u\) equals \(c_0+uc_1\).

Subtract this codeword pair from the received pair.  At every coordinate
outside \(C_0\) at least one of the two residual values is nonzero, and a
coordinate can cancel at at most one finite slope.  A noncommon witness
for \(u\) must contain at least \(\max\{1,a-c\}\) such cancellations:
the support-size condition gives \(a-c\), and one is still necessary when
\(a\le c\).  Double counting therefore gives
\[
 \abs Z\le
 \floor{\frac{n-c}{\max\{1,a-c\}}}\le r+1.        \tag{MO2}
\label{eq:large-overlap-count}
\]
Indeed, for \(c\ge a\) the numerator is at most \(r\); for \(c=a-1\)
it is \(r+1\); and, with \(x=a-c\ge2\), the quotient is
\((r+x)/x\le r+1\).
\end{proof}

\begin{theorem}[Quadratic mean-overlap theorem]
\label{thm:quadratic-mean-overlap}
Let \(C=\RS_\F(D,k)\), put \(R=n-k\), let
\(0\le r\le R-1\), and let \(\varnothing\ne\Gamma\subseteq\F\).  If
\[
 (n-r)^2\ge n(k+r),
 \qquad\text{equivalently}\qquad
 r^2\ge n(3r-R),                                  \tag{MO3}
\label{eq:quadratic-overlap-condition}
\]
then
\[
 B_{C,\Gamma}^{\rm MCA}(n-r)=\min\{\abs\Gamma,r+1\}. \tag{MO4}
\label{eq:quadratic-overlap-exact}
\]
The same upper bound holds for every linear \([n,k]\) MDS code.
\end{theorem}

\begin{proof}
The universal tangent construction imported from \cite[Sec.~13]{CAP25v132} gives the lower bound.  For the upper
bound, fix a received pair and its challenge-restricted bad set \(Z\).
If \(\abs Z\le r+1\) there is nothing to prove.  Otherwise choose
\(r+2\) slopes and distinct exact witness supports
\(S_1,\ldots,S_{r+2}\), each of size \(a=n-r\).  Suppose every pair had
intersection at most \(k+r-1\), and put
\(d_x=\#\{i:x\in S_i\}\).  Then
\[
 \sum_x\binom{d_x}{2}\le\binom{r+2}{2}(k+r-1),    \tag{MO5}
\]
whereas Cauchy--Schwarz and \(a^2/n\ge k+r\) give
\[
 \sum_x\binom{d_x}{2}
 \ge\frac{(r+2)^2a^2/n-(r+2)a}{2}
 \ge\frac{(r+2)^2(k+r)-(r+2)a}{2}.                \tag{MO6}
\]
The last lower bound exceeds the upper bound in \textup{(MO5)} by
\[
 \frac{r+2}{2}\bigl(k+3r-n+1\bigr).
\]
When \(3r>R\) this is positive.  When \(3r\le R\), the exact deep theorem
already gives the desired upper bound.  Thus in the remaining case some
pair satisfies \eqref{eq:large-overlap-collapse}, and
\cref{lem:large-overlap-collapse} contradicts \(\abs Z\ge r+2\).
The exact-support reduction and large-overlap lemma use only the MDS
information-set property, proving the stated extension.
\end{proof}

\begin{corollary}[Exact quadratic staircase]
\label{cor:mean-overlap-exact}
Put
\[
 r_{\rm quad}(n,k)=
 \floor{\frac{3n-\sqrt{n(5n+4k)}}2}.
 \tag{MO7}\label{eq:quadratic-radius}
\]
For every \(0\le r\le r_{\rm quad}(n,k)\) and every nonempty challenge
set \(\Gamma\),
\[
 B_{C,\Gamma}^{\rm MCA}(n-r)=\min\{\abs\Gamma,r+1\}.
\]
If \(k/n\to\rho\), then
\(r_{\rm quad}(n,k)/n\to
\alpha_\rho=(3-\sqrt{5+4\rho})/2\).
\end{corollary}

\begin{proof}
The smaller root of \(r^2-3nr+n(n-k)=0\) is the real number inside the
floor in \eqref{eq:quadratic-radius}.  Apply
\cref{thm:quadratic-mean-overlap}; the limiting formula follows by
division by \(n\).
\end{proof}

\begin{proposition}[Limit of pairwise-overlap information]
\label{prop:pairwise-overlap-limit}
Let \(k/n\to\rho\in(0,1)\), \(r/n\to\alpha\in(0,1)\), and suppose
\[
 \alpha^2<3\alpha-(1-\rho).
 \tag{MO8}\label{eq:pairwise-overlap-limit}
\]
For all sufficiently large \(n\), there are \(r+2\) distinct
\(r\)-subsets \(T_i\) such that, whenever
\(j=3r-(n-k)\le r\),
\(\abs{T_i\cap T_\ell}\le j-1\) for every \(i\ne\ell\).  If \(j>r\),
that intersection condition is automatic.  Hence pairwise support overlap
alone cannot extend the exact theorem past the quadratic boundary; further
progress must use the syndrome, determinant, or fiber incidence.  This is
a limitation of the method, not a construction of an MCA-bad line.
\end{proposition}

\begin{proof}
When \(j\le r\), choose \(r+2\) independent uniform \(r\)-subsets.
The intersection of a fixed pair is hypergeometric with mean
\(r^2/n=(\alpha^2+o(1))n\), whereas
\(j=(3\alpha-(1-\rho)+o(1))n\).  More explicitly,
\[
 \Pr(\abs{T_1\cap T_2}=s)
 =\frac{\binom rs\binom{n-r}{r-s}}{\binom nr}.
\]
After writing \(s=xn\), Stirling's formula gives the exponent
\[
 \alpha h(x/\alpha)+(1-\alpha)
 h((\alpha-x)/(1-\alpha))-h(\alpha),
\]
which is strictly concave and has its unique maximum zero at
\(x=\alpha^2\).  The strict gap in
\eqref{eq:pairwise-overlap-limit}, followed by summing at most \(n+1\)
terms, therefore gives probability \(e^{-\Omega(n)}\) that one fixed
intersection reaches \(j\).  A union bound over \(O(n^2)\) pairs is
smaller than one.  A
successful family is automatically distinct, because equal members have
intersection \(r\ge j\).  When \(j>r\), choose any \(r+2\) distinct
\(r\)-sets.
\end{proof}

The following complementary theorem can be stronger when the overlap
defect \(j=3r-R\) is small.  Its only new input is a packing inequality;
in particular it uses no smoothness or Fourier estimate.

\begin{theorem}[Exact overlap--packing criterion]
\label{thm:overlap-packing-exact}
Let \(C=\RS_\F(D,k)\), put \(R=n-k\), and fix an integer
\(0\le r\le R-1\).  Set \(a=n-r\) and \(j=3r-R\).  If either
\[
 j\le0,
 \tag{OP0}\label{eq:op-deep}
\]
or
\[
 1\le j\le r,
 \qquad
 \binom nj\le(r+1)\binom rj,
 \tag{OP}\label{eq:op-packing}
\]
then
\[
                    B_C^{\rm MCA}(n-r)=r+1.       \tag{OP$'$}
\label{eq:op-exact}
\]
The same upper bound holds for every linear \([n,k]\) MDS code; the
matching lower bound holds whenever the field contains at least \(r+1\)
distinct slope parameters.
\end{theorem}

\begin{proof}
Condition \eqref{eq:op-deep} is the imported deep-regime case \cite[Sec.~13]{CAP25v132}, so assume
\eqref{eq:op-packing}.  Fix an arbitrary received pair \((f,g)\), let
\(Z\) be its bad-slope set, and use the exact-agreement reduction imported from \cite[Sec.~13]{CAP25v132} to choose, for every \(z\in Z\), an
exact witness support \(S_z\), \(\abs{S_z}=a\), and an explaining
polynomial \(p_z\in\F[X]_{<k}\).  A single noncommon support cannot carry
two distinct slopes, so the supports chosen for distinct elements of
\(Z\) are distinct.

For the stated MDS extension, the same exact-support reduction is
available: interpolate the unexplained component on any \(k\)-subset of
the original witness support; because every \(k\)-subset is an information
set, nonexplanation is detected at one further coordinate, and that
\((k+1)\)-set can be enlarged to size \(a\).

First suppose that two of them have a large intersection:
\[
 \abs{S_z\cap S_w}\ge n+k-a=k+r
 \qquad(z\ne w).
 \tag{OP1}\label{eq:op-large-overlap}
\]
Then \cref{lem:large-overlap-collapse} gives \(\abs Z\le r+1\).

It remains to consider the complementary case, in which every pair of
chosen supports satisfies
\(\abs{S_z\cap S_w}\le n+k-a-1\).  Put
\(T_z=D\setminus S_z\), so \(\abs{T_z}=r\).  For distinct slopes,
\[
 \abs{T_z\cap T_w}
 =n-2a+\abs{S_z\cap S_w}
 \le3r-R-1=j-1.                                  \tag{OP3}
\label{eq:op-small-overlap}
\]
No \(j\)-subset of \(D\) can therefore lie in two different \(T_z\)'s.
Counting pairs \((z,J)\) with \(J\in\binom{T_z}{j}\) yields
\[
                 \abs Z\binom rj\le\binom nj\le
                 (r+1)\binom rj,
\]
and hence \(\abs Z\le r+1\).  The universal tangent construction gives
the reverse inequality.  The proof used only linearity, the fact that
agreement of two codewords on \(k\) coordinates forces equality, and the
exact-support reduction; these are MDS properties.  The coordinate
tangent construction gives the stated MDS lower bound.
\end{proof}

\begin{proposition}[General overlap--packing upper bound]
\label{prop:general-overlap-packing-upper}
For \(1\le j=3r-(n-k)\le r\), let \(A(n,r,j)\) be the largest size of a
family of \(r\)-subsets of an \(n\)-set whose distinct members intersect
in at most \(j-1\) points.  Then
\[
 B_{C,\Gamma}^{\rm MCA}(n-r)
 \le\min\left\{\abs\Gamma,
     \max\left(r+1,A(n,r,j)\right)\right\}          \tag{OP5}
\label{eq:general-overlap-packing-upper}
\]
and
\[
 A(n,r,j)\le
 \floor{\frac{\binom nj}{\binom rj}}.              \tag{OP6}
\label{eq:packing-number-bound}
\]
\end{proposition}

\begin{proof}
For a fixed received pair, either two chosen exact witness supports have
intersection at least \(k+r\), in which case
\cref{lem:large-overlap-collapse} gives at most \(r+1\) slopes, or every
pair of complementary \(r\)-sets intersects in at most \(j-1\) points.
The latter family has size at most \(A(n,r,j)\).  Finally, no \(j\)-subset
can lie in two members of such a packing, so counting contained
\(j\)-subsets proves \eqref{eq:packing-number-bound}.
\end{proof}


\section{Exact thresholds for explicit prize rows}\label{sec:prize-rows}

\begin{table}[H]
\centering
\caption{Explicit maximal-dimension prize-admissible rows.  In each row
\(k=2^{40}\), \(B=\floor{p/2^{128}}=r_\rho(n)+1\), and \((s,a_0)\) is part of
the Proth certificate in \cref{prop:proth-row-check}.}
\label{tab:prize-proth-rows}
\small
\begin{tabular}{@{}ccrrc@{}}
\hline
rate & length & \(B\) & bits of \(p\) & \((s,a_0)\)\\
\hline
\(1/2\)  & \(2^{41}\) & \(389500552609\)  & 167 & \((92,3)\)\\
\(1/4\)  & \(2^{42}\) & \(1210584858040\) & 169 & \((93,13)\)\\
\(1/8\)  & \(2^{43}\) & \(2879806199253\) & 170 & \((95,5)\)\\
\(1/16\) & \(2^{44}\) & \(6233898019554\) & 171 & \((97,5)\)\\
\hline
\end{tabular}
\end{table}

For completeness, the four field orders are
\begin{align*}
p_{1/2}&=132540169958804033333249306710494641010898987122689,\\
p_{1/4}&=411940680852499481698306614369841346700408394874881,\\
p_{1/8}&=979947269755402568812854322316630667196565607677953,\\
p_{1/16}&=2121285573237585848299875619011192262679065433997313.
\end{align*}
Their odd Proth coefficients \(u\), in the same order, are
\begin{align*}
u_{1/2}&=26766274163673319604503,
&u_{1/4}&=41595378994516821279015,\\
u_{1/8}&=24737346889219389259839,
&u_{1/16}&=13387194060291799253121.
\end{align*}
For \(F_{n,k}(r)=r^2-n(3r-(n-k))\), the exact boundary checks are
\[
\begin{array}{c|r|r}
 \rho&F_{n,k}(B-1)&F_{n,k}(B)\\ \hline
 1/2&5154112775168&-663955886271\\
 1/4&7590647904465&-3182321912768\\
 1/8&13908181940112&-6720484728007\\
 1/16&19335616403905&-20973145690236.
\end{array}                                                    \tag{PC0}
\label{eq:quadratic-row-certificates}
\]

\begin{proposition}[Exact arithmetic and primality certificates]
\label{prop:proth-row-check}
For the four rows in \cref{tab:prize-proth-rows}, the displayed integer
\(p\) satisfies
\[
 p=u\,2^s+1,\qquad u\text{ odd},\qquad u<2^s,\qquad
 a_0^{(p-1)/2}\equiv-1\pmod p.                  \tag{PC1}
\label{eq:proth-certificates}
\]
Consequently every \(p\) is prime.  Moreover
\[
 n\mid p-1,\qquad p<2^{256},\qquad
 B\,2^{128}\le p<(B+1)2^{128},                  \tag{PC2}
\label{eq:prize-row-arithmetic}
\]
where \(n\) and \(B\) are the corresponding entries of the table.
They also satisfy \(F_{n,k}(B-1)\ge0>F_{n,k}(B)\).  Since
\(F_{n,k}'(r)=2r-3n<0\) on \([0,n]\), these two signs locate its smaller
root and give \(B-1=r_\rho(n)=r_{\rm quad}(n,k)\).
Thus \(\F_p^*\) contains a subgroup of order \(n\), and each displayed code exists with the asserted slope budget.
\end{proposition}

\begin{proof}
All identities and inequalities in
\eqref{eq:quadratic-row-certificates}--\eqref{eq:prize-row-arithmetic}
are direct
integer calculations; the values of \((s,a_0)\), \(u\), \(p\), \(n\),
and \(B\) have all been printed above so that the calculation is
independently reproducible.  We recall why the certificate proves
primality.  If a prime \(\ell\) divides \(p\), then
\(a_0^{u2^{s-1}}\equiv-1\pmod\ell\), whereas
\(a_0^{u2^s}\equiv1\pmod\ell\).  The multiplicative order of \(a_0\)
modulo \(\ell\) therefore has exact \(2\)-part \(2^s\), so
\(2^s\mid\ell-1\).  Since \(u<2^s\), one has \(2^s>\sqrt p\); hence a
composite \(p\) would have a prime divisor both at most \(\sqrt p\) and
at least \(2^s+1\), a contradiction.  This is Proth's theorem in the
special form used here.  Finally \(s\ge\log_2 n\) gives \(n\mid p-1\),
and cyclicity of \(\F_p^*\) gives the required subgroup.
\end{proof}

\begin{theorem}[Exact first adjacent row over a separating field]
\label{thm:exact-first-adjacent-row}
Let \(C=\RS_\F(D,k)\), put \(R=n-k\) and
\(M=\binom n{R-1}=\binom n{k+1}\), and define
\[
                    Q_{\rm sep}(M)=\max\left\{M,\binom M2\right\}.
\]
If \(\abs\F>Q_{\rm sep}(M)\), then for the full-field challenge
\[
                         B_C^{\rm MCA}(k+1)=M.      \tag{AD1}
\]
If \(R\ge2\) and \(b=\floor{\eps\abs\F}\), the exact target threshold
therefore satisfies
\[
 \begin{array}{ll}
 b\ge M &\Longrightarrow\quad a^*=k+1,\\
 \displaystyle\binom n{R-2}\le b<M
   &\Longrightarrow\quad a^*=k+2.
 \end{array}                                       \tag{AD2}
\]
Thus the first finite adjacent comparison has literal binomial constants;
no \(e^{o(n)}\) or unspecified polynomial factor occurs.
\end{theorem}

\begin{proof}
The upper bound in \textup{(AD1)} is the exact support-atlas bound imported from \cite[Sec.~13]{CAP25v132}.  For the reverse inequality use the
Reed--Solomon parity check of redundancy \(R\).  For every
\((R-1)\)-set \(E\subseteq D\), its column span \(V_E\) is a hyperplane
in \(\F^R\), and distinct sets give distinct hyperplanes: otherwise the
union of two distinct \((R-1)\)-sets would contain at least \(R\) columns
in one \((R-1)\)-space, contradicting the MDS column property.

Every proper linear hyperplane in \(\F^R\) has
\(\abs{\F}^{R-1}\) points.  Since \(M<\abs\F\), the union of the \(M\)
spaces \(V_E\) has fewer than \(\abs{\F}^{R}\) points; choose \(y_1\)
outside it.  For each \(E\),
choose a nonzero linear functional \(\ell_E\) with kernel \(V_E\).
The line \(y_0+\gamma y_1\) meets \(V_E\) at
\[
                 \gamma_E=-\ell_E(y_0)/\ell_E(y_1).
\]
For \(E\ne E'\), equality \(\gamma_E=\gamma_{E'}\) cuts out a proper
hyperplane in the variable \(y_0\), because the normalized functionals
\(\ell_E/\ell_E(y_1)\) and
\(\ell_{E'}/\ell_{E'}(y_1)\) are distinct.  There are at most
\(\binom M2<\abs\F\) collision hyperplanes; their union also has fewer
than \(\abs{\F}^{R}\) points, so choose \(y_0\) outside it.  Since
\(y_1\notin V_E\), the line is not contained in \(V_E\), and every
intersection is transverse.  The \(M\) parameters \(\gamma_E\) are then
finite and pairwise distinct.  Syndrome surjectivity and
\cref{thm:syndrome-secant-exact} realize them as \(M\) bad slopes at
agreement \(k+1\), proving \textup{(AD1)}.

If \(b\ge M\), the support upper bound makes the first allowed agreement
safe.  If \(\binom n{R-2}\le b<M\), equality \textup{(AD1)} makes
\(k+1\) unsafe, whereas
the exact support-atlas bound imported from \cite[Sec.~13]{CAP25v132} makes \(k+2\) safe.  Monotonicity proves
\textup{(AD2)}.
\end{proof}


\section{Asymptotic consequences of the imported prefix floor}\label{sec:shallow-prefix-exponent}
\begin{theorem}[Unconditional shallow-prefix RS--MCA exponent]
\label{thm:unconditional-asymptotic-rs-mca}
Let \(C_n=\RS_{\F_n}(D_n,k_n)\), where
\(D_n\subseteq\B_n\subseteq\F_n\), \(\abs{D_n}=n\), and
\(k_n/n\to\rho\in(0,1)\).  Let \(a_n\ge k_n+1\), put
\(w_n=a_n-k_n-1\), and assume
\[
 \frac{a_n}{n}\to\alpha\in(0,1),\qquad
 w_n\log\abs{\B_n}=o(n),                           \tag{AS1}
\]
\(\abs{\F_n}>n\), and, for some finite \(\phi\ge0\),
\[
 \frac1n\log\abs{\F_n}\to\phi,
 \qquad
 \frac1n\log(\abs{\F_n}-n)\to\phi.               \tag{AS2}
\]
For the full-field challenge,
\[
 \frac1n\log B_{C_n}^{\rm MCA}(a_n)
   =\min\{h(\alpha),\phi\}+o(1),                   \tag{AS3}
\]
and therefore
\[
 \frac1n\log e_{\rm MCA}
   \left(C_n,1-\frac{a_n}{n}\right)
   =-\bigl(\phi-h(\alpha)\bigr)_++o(1).            \tag{AS4}
\]
Because \(\abs{\B_n}\ge n\), condition \textup{(AS1)} forces
\(\alpha=\rho\).  No smoothness, Fourier estimate, or auxiliary ray
hypothesis is assumed.
\end{theorem}

\begin{proof}
Put \(M_n=\binom n{a_n}\) and
\(L_n=\ceil{M_n\abs{\B_n}^{-w_n}}\).  Exact-agreement reduction and
fixed-support slope uniqueness give the unconditional upper bound
\[
 B_{C_n}^{\rm MCA}(a_n)\le\min\{\abs{\F_n},M_n\}.  \tag{AS5}
\]
The exact locator-prefix list and collision-aware pole compiler imported from \cite[Secs.~14--18]{CAP25v132} give
\[
 B_{C_n}^{\rm MCA}(a_n)\ge
 \left\lceil
 \frac{L_n(\abs{\F_n}-n)}
 {\abs{\F_n}-n+k_n(L_n-1)}
 \right\rceil.                                     \tag{AS6}
\]
By Stirling's formula and \textup{(AS1)},
\(\log L_n=h(\alpha)n+o(n)\).  For positive \(x,y\),
\(xy/(x+k_ny)\ge\frac12\min\{y,x/k_n\}\); hence
\textup{(AS2)}, \(\log k_n=o(n)\), and the ceiling show that the
right side of \textup{(AS6)} has logarithm
\(n\min\{h(\alpha),\phi\}+o(n)\).  The upper bound
\textup{(AS5)} has the same exponent, proving \textup{(AS3)}; subtracting
\(\log\abs{\F_n}\) proves \textup{(AS4)}.  Finally
\(w_n\log n=o(n)\), so \((a_n-k_n)/n\to0\) and
\(\alpha=\rho\).
\end{proof}

\begin{corollary}[Unconditional large-challenge RS--MCA threshold]
\label{cor:unconditional-rs-mca-capacity}
Under the row hypotheses of
\cref{thm:unconditional-asymptotic-rs-mca}, let the full-field target
satisfy \(\eps_n=\exp(-\lambda n+o(n))\) for a finite
\(\lambda\ge0\).  If
\[
                         \lambda<\phi-h(\rho),      \tag{AS7}
\]
then, for every sufficiently large \(n\),
\[
 a_n^*=k_n+1,\qquad
 \delta_n^*=1-\frac{k_n+1}{n}=1-\rho+o(1).         \tag{AS8}
\]
Conversely, if \(\lambda>(\phi-h(\rho))_+\), every agreement sequence
satisfying \((a_n-k_n-1)\log\abs{\B_n}=o(n)\) is unsafe.  The critical
boundary \(\lambda=(\phi-h(\rho))_+\) is governed by the
exact finite constants and is not decided by exponent asymptotics.  This
is a large scalar-extension regime.  If \(\abs{\F_n}\) is polynomial in
\(n\), then \(\phi=0\) and \textup{(AS7)} cannot hold; the corollary does
not settle that Proximity-Prize parameter regime.
\end{corollary}

\begin{proof}
At \(a=k_n+1\), \textup{(AS5)} is
\(B_{C_n}^{\rm MCA}(a)\le\binom n{k_n+1}
=\exp(h(\rho)n+o(n))\).  Under \textup{(AS7)},
\(\eps_n\abs{\F_n}=\exp((\phi-\lambda)n+o(n))\) has strictly larger
exponent, and taking the integer floor does not change the comparison.
Thus the first permitted agreement is safe
and \textup{(AS8)} follows.  The converse follows from \textup{(AS4)}
and the strict reverse exponential inequality.
\end{proof}


\part{Exact profile refinements and the profile envelope}
\section{Profile envelope and the sequel compiler}\label{sec:profile-envelope-compiler}
A realized first-match profile $\lambda$ has a full support slice $\Omega^0_\lambda$, a coefficient field $\B_\lambda$, a proved boundary map
$\Phi_\lambda:\Omega^0_\lambda\to\B_\lambda^{R_\lambda}$, realized image size
$L_\lambda=|\Phi_\lambda(\Omega^0_\lambda)|$, and image-normalized average
\[
 \bar N_\lambda=\frac{|\Omega^0_\lambda|}{L_\lambda}.
\]
For a received line and agreement $a$, let $\Lambda(a)$ be the realized first-match profile set and define
\begin{equation}
 \mathfrak E_n(a)=1+(n-a+1)+
 \sup_{\text{received lines}}\sum_{\lambda\in\Lambda(a)}(1+\bar N_\lambda).\label{eq:profile-envelope}
\end{equation}
The formal codomain size $|\B_\lambda|^{R_\lambda}$ may replace $L_\lambda$ only after a full-image certificate
$L_\lambda\ge e^{-o(n)}|\B_\lambda|^{R_\lambda}$ has been proved.  Otherwise the leaf stays at realized-image scale or is routed to an earlier effective-image-collapse cell.

\begin{definition}[Admissible sequel compiler]\label{def:admissible-sequence}
A row sequence is \emph{admissible} when: (i) every bad witness belongs to an ordered first-match atlas with $e^{o(n)}$ realized profiles; (ii) all quotient, planted, field-drop, tangent, rank, extension, and directly counted ray cells have proved distinct-slope payments; (iii) every remaining primitive profile has a fixed-density Boolean support slice, an exact boundary map, and image normalization; (iv) its nontrivial effective Fourier characters satisfy the certified minor/major payment of \cref{def:effective-major-minor,def:aggregate-minor-payment,def:major-arc-aggregate}, or an equivalent direct Sidon payment; and (v) every residual ray chart has either a direct slope bound or the ray compiler of \cref{hyp:ray-compiler}.  All profile fields and quotient scales are included in \eqref{eq:profile-envelope}.
\end{definition}

\begin{theorem}[Conditional profile-envelope compiler]\label{thm:main-smooth-circle}
Every admissible row sequence satisfies
\begin{equation}
 B_{C_n}^{\rm MCA}(a_n)\le e^{o(n)}\mathfrak E_n(a_n).\label{eq:conditional-numerator}
\end{equation}
For a finite row the same proof is literal after every $e^{o(n)}$ term is replaced by its exact integer constant and all first-match cells are summed once.
\end{theorem}
\begin{proof}
The first-match atlas pays the algebraic cells.  The effective Fourier/Sidon step and \cref{thm:primitive-q} pay primitive Q, while the ray compiler and \cref{prop:q-implies-sp} convert the support bounds to distinct slopes.  Summation over the realized profile set gives \eqref{eq:conditional-numerator}.  The finite statement is the same disjoint first-match count with integer constants.
\end{proof}

\begin{definition}[Complete-fiber folding map]\label{def:structured-folding}
Let $D\subseteq\B$ be finite.  A polynomial map $\phi:D\to Q=\phi(D)$ is a $c$-fold complete-fiber folding map when every fiber has exactly $c$ points.  A support is complete at this scale when it is a union of full fibers.  Multiplicative power maps and the Chebyshev twin-coset foldings of the foundation \cite[Secs.~7 and 11]{CAP25v132} are the motivating examples.
\end{definition}

\section{Quotient--remainder prefix normal form}\label{sec:quotient-remainder-profiles}
The quotient-prefix floors and descent mechanisms of CAP25 v13.2 are imported from \cite[Secs.~3, 8, and 14]{CAP25v132}.  The normal forms and obstruction theorems below refine that package for the sequel's Q compiler.

The quotient and remainder variables can be separated exactly at the
locator-prefix level.  This removes a spurious factor coming from summing
over marked remainder sets, but it also exposes a genuine scaled quotient-
coefficient-field term which prevents an unrestricted comparison with the
identity profile.

\begin{definition}[Quotient/remainder profile and field drop]
\label{def:quotient-remainder-profile}
Let \(\phi:D\to Q\) be a \(c\)-fold complete-fiber folding map in the
sense of \cref{def:structured-folding}.  A
\emph{complete-fiber-plus-remainder support} is
\[
                  S=\phi^{-1}(E)\sqcup R,
\]
where \(E\subseteq Q\) and
\(R\subseteq D\setminus\phi^{-1}(E)\).  The set \(R\) is the
\emph{support remainder}: it records the selected points in fibers that
are not declared complete.  A \emph{quotient/remainder profile} fixes
\(\phi\), its degree \(c\), the size \(r=\abs R\), the relevant
coefficient field, and any planted remainder data, while \(E\) and the
unfixed part of \(R\) range subject to the displayed disjointness.
A fixed-\(R\) profile is the subprofile in which the remainder support
itself is part of the label.  The adjective \emph{Euclidean} used below
refers to separating the locator factor of degree \(r<c\) from the
composed quotient locator; it does not mean that an arbitrary support
remainder has size less than \(c\).

We call \(\B_\phi\subseteq\B\) a \emph{scaled quotient coefficient
field} when, for some \(\eta\in\B^\times\),
\[
                         Q=\phi(D)\subseteq\eta\B_\phi .
\]
When this field is recorded in a profile, it is taken minimal among the
subfields with this property.
Then the \(j\)-th elementary quotient coefficient lies in
\(\eta^j\B_\phi\).  A \emph{field drop} is the use of a proper such
subfield.  It is a \emph{positive-rate field drop} along a sequence when
\[
       1-\frac{\log\abs{\B_\phi}}{\log\abs\B}
\]
stays bounded away from zero and a positive-density number on the
entropy scale of quotient-prefix slots is live.  The payment saved in
each live slot is the factor \(\abs\B/\abs{\B_\phi}\); this is the term
quantified in \cref{prop:identity-quotient-comparison}.
\end{definition}

\begin{definition}[Balanced quotient core]
\label{def:balanced-quotient-core}
A \emph{split locator} is a monic squarefree polynomial all of whose roots
lie in the evaluation domain \(D\).
After common locator factors, fixed quotient data, and planted remainder
points have been removed, let \(A\) and \(B\) be two monic split
locators of the same degree \(e\).  They form a
\emph{depth-\(w\) balanced core} if they have the same first \(w\)
nonleading coefficients, equivalently
\[
                         \deg(A-B)\le e-w-1.
\]
It is a \emph{balanced quotient core} when the remaining locator factors
are pullbacks through the same folding map \(\phi\).  The word
\emph{balanced} records equality of degree and leading normalization;
the \emph{core} consists of the roots still allowed to move after common
and planted factors are removed.  If the resulting projective locator
family is one-dimensional it is a split pencil
\(Q_0+\lambda Q_1\); in higher dimension this definition supplies no
distinct-ray bound, and the separate ray compiler is required.
\end{definition}

For a monic polynomial
\[
 A(X)=X^d+a_1X^{d-1}+\cdots+a_d
\]
and an integer $0\le u\le d$, write
\(\operatorname{pref}_u(A)=(a_1,\ldots,a_u)\).  Equivalently, if
\(A^\vee(Z):=Z^dA(Z^{-1})\), then
\(\operatorname{pref}_u(A)\) is the coefficient vector of
\(A^\vee(Z)\bmod Z^{u+1}\), with its constant coefficient omitted.

\begin{theorem}[Exact quotient--remainder prefix normal form]
\label{thm:exact-quotient-remainder-normal-form}
Let $\B\subseteq\F$ be finite fields, let $D\subseteq\B$ have $n$
distinct elements, and let $\phi\in\B[X]$ be monic of degree $c\ge2$.
Assume that the restriction
\[
             \phi:D\longrightarrow Q:=\phi(D)
\]
has complete fibers of size $c$; in particular, $N:=|Q|=n/c$.
Fix integers $m,r$ with $0\le r<c$ and $cm+r\le n$.  For
\[
 E\in\binom Qm,
 \qquad
 R\in\binom{D\setminus\phi^{-1}(E)}r,
\]
put
\[
 \begin{split}
 V_E(Y)&:=\prod_{y\in E}(Y-y)
       =Y^m+\sum_{j=1}^m v_j(E)Y^{m-j},\\
 P_R(X)&:=\prod_{x\in R}(X-x)
       =X^r+\sum_{i=1}^r p_i(R)X^{r-i},\\
 L_{E,R}(X)&:=P_R(X)V_E(\phi(X)).
 \end{split}                                      \tag{QR1}\label{eq:qr-locators}
\]
Then $L_{E,R}$ is the monic locator of the support
\(\phi^{-1}(E)\sqcup R\), of degree $A=cm+r$.

Let $0\le w\le A$ and put $d=\lfloor w/c\rfloor$.

\begin{enumerate}[label=\textup{(\roman*)},leftmargin=*]
\item If $w\ge r$, every nonempty fiber of
      \((E,R)\mapsto\operatorname{pref}_w(L_{E,R})\) determines $R$
      uniquely and, after $R$ is recovered, is exactly one quotient-prefix
      fiber.  More precisely, for every prefix value $z$ there are unique
      data $R_z$ and $\eta_z\in\B^d$, whenever the fiber is nonempty, such
      that
      \[
      \begin{split}
       &\{(E,R):\operatorname{pref}_w(L_{E,R})=z\}\\
       &\quad=
       \{(E,R_z): E\in\tbinom{Q\setminus\phi(R_z)}m,
           \ \operatorname{pref}_d(V_E)=\eta_z\}.
      \end{split}                                  \tag{QR2}\label{eq:qr-fiber-exact}
      \]
      Thus no factor counting possible remainder sets occurs in a fixed
      global prefix fiber.

\item If $w<r$, then necessarily $w<c$, and the quotient set $E$ is
      invisible to the first $w$ coefficients.  There is a fixed
      unitriangular bijection $T_{\phi,m,w}:\B^w\to\B^w$ such that
      \[
          \operatorname{pref}_w(L_{E,R})
          =T_{\phi,m,w}(\operatorname{pref}_w(P_R))                 \tag{QR3}
      \]
      for every admissible $(E,R)$.  Consequently, with the convention
      \(\binom uv=0\) for $v>u$,
      \[
      \begin{split}
       &\#\{(E,R):\operatorname{pref}_w(L_{E,R})=z\}\\
       &\quad=
       \sum_{\substack{R\in\binom Dr\\
          \operatorname{pref}_w(P_R)=T_{\phi,m,w}^{-1}(z)}}
          \binom{N-|\phi(R)|}{m}.
      \end{split}                                  \tag{QR4}\label{eq:qr-shallow-remainder}
      \]
      Thus this case is a remainder-prefix problem, multiplied by the
      number of quotient sets still available after $R$ is fixed.
\end{enumerate}
\end{theorem}

\begin{proof}
For a monic degree-$e$ polynomial $A$, set
\(A^\vee(Z)=Z^eA(Z^{-1})\).  Also set
\[
       \Phi(Z):=Z^c\phi(Z^{-1})
       =1+\phi_1Z+\cdots+\phi_cZ^c.
\]
Taking reciprocal polynomials in \eqref{eq:qr-locators} gives the exact
formal-power-series identity
\[
 L_{E,R}^\vee(Z)=P_R^\vee(Z)
 \left(
   \Phi(Z)^m+
   \sum_{j=1}^m v_j(E)Z^{cj}\Phi(Z)^{m-j}
 \right).                                          \tag{QR5}\label{eq:qr-infinity}
\]
All factors on the right have constant coefficient one except for the
displayed scalar coefficients $v_j(E)$.

Suppose first that $w\ge r$.  Since $r<c$, no term containing a $v_j$
occurs below degree $c$.  For $1\le i\le r$, the coefficient of $Z^i$
on the right of \eqref{eq:qr-infinity} is
\(p_i(R)\) plus an expression depending only on
\(p_1(R),\ldots,p_{i-1}(R)\), on $m$, and on $\Phi$.  The first $r$
coefficients therefore recover $p_1(R),\ldots,p_r(R)$ successively.
They recover the monic polynomial $P_R$, hence its root set $R$, uniquely.

We may now divide \(L_{E,R}^\vee\) by the unit
\(P_R^\vee\) modulo $Z^{w+1}$.  In the remaining factor of
\eqref{eq:qr-infinity}, the coefficient of $Z^{cj}$ is
\(v_j(E)\) plus an expression depending only on
\(v_1(E),\ldots,v_{j-1}(E)\), on $m$, and on $\Phi$.  Hence the
coefficients at degrees $c,2c,\ldots,dc$ recover
\(v_1(E),\ldots,v_d(E)\) successively.  Conversely, $P_R$ and these $d$
coefficients determine the right side of \eqref{eq:qr-infinity} modulo
$Z^{w+1}$, because every term containing $v_j$ with $j>d$ starts in degree
$cj>w$.  This proves \eqref{eq:qr-fiber-exact}; admissibility of $E$ is
exactly the condition $E\cap\phi(R)=\varnothing$.

If $w<r<c$, every term involving a $v_j$ starts in degree at least
$c>w$, so modulo $Z^{w+1}$ the bracket in \eqref{eq:qr-infinity} is the
fixed unit $\Phi(Z)^m$.  Multiplication by this unit sends the first $w$
coefficients of $P_R^\vee$ to the first $w$ coefficients of
$L_{E,R}^\vee$ by a unitriangular transformation.  It is independent of
$E$ and invertible.  For a fixed $R$, the admissible $E$ are precisely the
$m$-subsets of $Q\setminus\phi(R)$, of which there are
\(\binom{N-|\phi(R)|}{m}\).  Summing over the indicated remainder-prefix
fiber proves \eqref{eq:qr-shallow-remainder}.
\end{proof}

\begin{proposition}[Complete support factorization and its limitation]
\label{prop:complete-support-factorization}
Under the complete-fiber hypotheses of
\cref{thm:exact-quotient-remainder-normal-form}, every support
\(S\subseteq D\) has the unique decomposition
\[
 E(S)=\{y\in Q:\phi^{-1}(y)\subseteq S\},\qquad
 R(S)=S\setminus\phi^{-1}(E(S)),
 \tag{QR5a}\label{eq:complete-support-factorization}
\]
and
\(Q_S(X)=P_{R(S)}(X)V_{E(S)}(\phi(X))\).  The remainder meets every
nonfull \(\phi\)-fiber in at most \(c-1\) points, but its total size may
be \(\Theta(n)\).  The reciprocal identity \eqref{eq:qr-infinity} remains
valid for this arbitrary remainder.  When \(\abs{R(S)}\ge c\), however,
the quotient coefficients and the remainder coefficients begin to
interlace at degree \(c\), so the triangular recovery theorem above does
not give a comparison without an additional partial-occupancy atlas.
\end{proposition}

\begin{proof}
A fiber is put into \(E(S)\) exactly when all of its points lie in \(S\);
all other selected points form \(R(S)\).  This proves existence,
disjointness, the occupancy bound, and uniqueness.  Multiplying the
corresponding linear factors gives the locator identity.  The reciprocal
calculation used for \eqref{eq:qr-infinity} is purely multiplicative and
does not require \(\deg P_R<c\).  If \(\deg P_R\ge c\), both
\(p_c(R)\) and \(v_1(E)\) occur in degree \(c\), which proves the stated
limitation of the triangular separation.
\end{proof}

\begin{theorem}[Exact partial-occupancy disintegration]
\label{thm:exact-partial-occupancy}
Let \(D=D_0\sqcup X\), where \(\abs X=b\), and let
\(\phi:D_0\twoheadrightarrow Q\) have exactly \(N=\abs Q\) fibers, each
of cardinality \(c\).  For \(S\subseteq D\), put
\[
 X_S=S\cap X,\qquad
 E(S)=\{y\in Q:\phi^{-1}(y)\subseteq S\},\qquad
 R(S)=(S\cap D_0)\setminus\phi^{-1}(E(S)).
\]
Write
\[
 t=\abs{X_S},\qquad m=\abs{E(S)},\qquad
 p=\abs{\phi(R(S))},\qquad r=\abs{R(S)}.
\]
Then this datum is unique, every fiber meeting \(R(S)\) contains between
one and \(c-1\) points of \(R(S)\), and \(\abs S=t+cm+r\).  If
\(\Omega_{t,m,p,r}\) denotes the family of supports with these four
parameters, then
\[
 \abs{\Omega_{t,m,p,r}}
 =\binom bt\binom Np\binom{N-p}m
   [x^r]\bigl((1+x)^c-1-x^c\bigr)^p.                 \tag{PO1}
\]
Consequently the nonempty families \(\Omega_{t,m,p,r}\) partition
\(\binom Da\), and their exact add-back is
\[
 \sum_{\substack{t,m,p,r\\t+cm+r=a}}\abs{\Omega_{t,m,p,r}}
 =\binom{cN+b}{a}.                                   \tag{PO2}
\]
If \(\phi\) is the restriction of a monic degree-\(c\) polynomial and
every fiber \(\phi^{-1}(y)\) is the complete simple root set in \(D_0\)
of \(\phi(X)-y\), then their locators satisfy
\[
             Q_S=Q_{X_S}P_{R(S)}V_{E(S)}(\phi).
\]
\end{theorem}

\begin{proof}
Full fibers, partial fibers, and empty fibers are mutually exclusive, so
the displayed decomposition is canonical.  Choose the \(t\) exceptional
points, the \(p\) partial fibers, and the \(m\) full fibers in
\(\binom bt\binom Np\binom{N-p}m\) ways.  On one partial fiber the
weight enumerator of the allowed nonempty proper subsets is
\((1+x)^c-1-x^c\); hence the coefficient in \textup{(PO1)} counts the
choices having total partial weight \(r\).  Summing \textup{(PO1)} and
extracting the coefficient of \(x^a\) gives
\[
 [x^a](1+x)^b
 \left(1+x^c+\bigl((1+x)^c-1-x^c\bigr)\right)^N
 =[x^a](1+x)^{b+cN}=\binom{b+cN}{a}.
\]
Multiplying the linear factors belonging to the three disjoint parts of
the support proves the locator identity.
\end{proof}

\begin{proposition}[Partial-occupancy Fourier factorization]
\label{prop:partial-occupancy-fourier}
In the setting of \cref{thm:exact-partial-occupancy}, let \(G\) be a
finite abelian group, let \(g:D\to G\), and put
\(\Phi(S)=\sum_{x\in S}g(x)\).  For a character
\(\chi\in\widehat G\) and a regular fiber \(F_y=\phi^{-1}(y)\), define
\[
 A_y(\chi)=\prod_{x\in F_y}\chi(g(x)),\qquad
 P_y(\chi;z)=
 \prod_{x\in F_y}\bigl(1+z\chi(g(x))\bigr)-1-z^cA_y(\chi).
\]
Then
\[
 \sum_{S\in\Omega_{t,m,p,r}}\chi(\Phi(S))
 =[u^tv^m\zeta^pz^r]
 \prod_{x\in X}\bigl(1+u\chi(g(x))\bigr)
 \prod_{y\in Q}
 \bigl(1+vA_y(\chi)+\zeta P_y(\chi;z)\bigr).          \tag{PO3}
\]
Consequently, for every \(s\in G\),
\[
 \abs{\{S\in\Omega_{t,m,p,r}:\Phi(S)=s\}}
 =\frac1{\abs G}\sum_{\chi\in\widehat G}
   \overline{\chi(s)}
   \sum_{S\in\Omega_{t,m,p,r}}\chi(\Phi(S)).         \tag{PO4}
\]
At the trivial character, \textup{(PO3)} is exactly the support count
\textup{(PO1)}.  Thus support add-back is unconditional, while a finite
flatness theorem is exactly a bound for the nontrivial-character sum in
\textup{(PO4)}.
\end{proposition}

\begin{proof}
For a fixed regular fiber the three terms in the second product encode,
respectively, an empty fiber, a full fiber, and a nonempty proper subset.
The variables \(v,\zeta,z\) record the numbers \(m,p,r\), while \(u\)
records the exceptional weight \(t\).  Multiplicativity of \(\chi\) turns
the character of the support sum into the product of its selected-point
characters, proving \textup{(PO3)} by coefficient extraction.
Formula \textup{(PO4)} is character orthogonality on \(G\).
\end{proof}

\begin{remark}[Effective normalization of a partial-occupancy slice]
Fix a nonempty slice \(\Omega_\lambda=\Omega_{t,m,p,r}\), choose
\(S_0\in\Omega_\lambda\), and let
\[
 G_\lambda=
 \left\langle\Phi(S)-\Phi(S_0):S\in\Omega_\lambda\right\rangle
 \le G.
\]
The attained boundary values lie in the affine coset
\(\Phi(S_0)+G_\lambda\), and the effective inversion is
\[
 \abs{\{S\in\Omega_\lambda:\Phi(S)=s\}}
 =\frac1{\abs{G_\lambda}}
   \sum_{\chi\in\widehat{G_\lambda}}
   \overline{\chi(s-\Phi(S_0))}
   \sum_{S\in\Omega_\lambda}
      \chi(\Phi(S)-\Phi(S_0)).                       \tag{PO5}
\]
Every character of \(G_\lambda\) extends to the ambient finite abelian
group, so the inner sum may be evaluated from \textup{(PO3)} after
multiplication by the fixed factor
\(\overline{\chi(\Phi(S_0))}\).  Different extensions agree on support
differences.  Equivalently, ambient characters group by their restriction
to \(G_\lambda\), and the annihilator multiplicity cancels exactly against
the ambient denominator.  This identifies the correct Fourier denominator;
it does not assert that the realized image fills the affine group.
\end{remark}

\begin{remark}[What partial-occupancy add-back does and does not prove]
The exponentially many local proper-subset choices are already contained
inside the coarse slices \(\Omega_{t,m,p,r}\); they incur no separate
profile-count loss in \textup{(PO2)}.  The theorem does not bound the
nontrivial Fourier terms in \textup{(PO4)}, the size of a realized
boundary image, or the projection of a witness cell to distinct slopes.
Those are separate analytic and ray payments.
\end{remark}

\begin{theorem}[Canonical partial-occupancy atlas and bounded-boundary closure]
\label{thm:canonical-partial-occupancy-atlas}
Fix a received line and exact agreement \(a\).  For every nonempty
canonical slice \(\Omega_\lambda=\Omega_{t,m,p,r}\) of
\cref{thm:exact-partial-occupancy}, let \(\Ccal_\lambda\) be the witnesses
whose support lies in \(\Omega_\lambda\), and order these cells
arbitrarily.  Then they form a witness-exhaustive first-match atlas and,
with \(Z_\lambda^\circ\) its first-match slope projections,
\[
 \abs{Z_\lambda^\circ}\le\abs{\Omega_\lambda},
 \qquad
 \sum_\lambda\abs{Z_\lambda^\circ}
 \le\sum_\lambda\abs{\Omega_\lambda}=\binom na.    \tag{PO6}
\]

More generally, give each slice a boundary map
\(\Phi_\lambda:\Omega_\lambda\to G_\lambda^{\rm amb}\) into a finite
abelian group of order \(A_\lambda\), let \(N_\lambda\) be its
active-coordinate count, let \(L_\lambda\) be its realized image size,
and put
\(\bar N_\lambda=\abs{\Omega_\lambda}/L_\lambda\).  The exact finite ray
multiplier and challenge-set bound are
\[
             \abs{Z_\lambda^\circ}
             \le L_\lambda\bar N_\lambda.          \tag{PO7}
\]
\[
 B_{C,\Gamma}^{\rm MCA}(a)
 \le\min\left\{\abs\Gamma,\
       \sup_r\sum_{\lambda\in\Lambda(r;a)}
              L_\lambda\bar N_\lambda\right\}
 \le\min\left\{\abs\Gamma,\binom na\right\}.        \tag{PO7a}
\]
Define the coarse partial-occupancy envelope
\[
 \mathfrak E_n^{\rm PO}(a)
 =1+\sup_r\sum_{\lambda\in\Lambda(r;a)}(1+\bar N_\lambda).
\]
If, uniformly over received lines and the polynomially many nonempty
canonical slices,
\[
                         \log A_\lambda=o(n),       \tag{PO8}
\]
then this atlas has a closed asymptotic profile-envelope compiler:
\[
 B_{C}^{\rm MCA}(a)\le
 e^{o(n)}\mathfrak E_n^{\rm PO}(a).                \tag{PO9}
\]
Condition \textup{(PO8)} discharges atlas exhaustivity, support add-back,
and the direct alternative in \textup{(RC)}, whose exact multiplier is
\(L_\lambda\le A_\lambda\).  If, in addition,
\[
                         \log A_\lambda=o(N_\lambda) \tag{PO10}
\]
on every slice to which effective Fourier terminology is applied, then
the effective inversion \textup{(PO5)} may designate every nontrivial
effective character major: its minor sum is empty and its exact major
multiplier is at most \(A_\lambda-1\).  Thus \textup{(PO10)} verifies the
slice analogues of MI and MA rather than assuming them.  A concrete
specialization is a polynomial-size coefficient field with uniformly
bounded boundary depth and \(N_\lambda=\Theta(n)\) on every live canonical
folding profile.  This is a coarse support-paid envelope, not an
identity-scale or refined quotient-envelope comparison.  It gives a safe
target whenever the literal right side of \textup{(PO7a)} is at most
\(B^*\); a sharp frontier additionally uses the unsafe inequality in
\textup{(SB2)}.  Condition \textup{(PO8)} alone does not locate a target
frontier.
\end{theorem}

\begin{proof}
The support slices partition \(\binom Da\) by \textup{(PO2)}, hence the
corresponding witness cells partition the exact witness incidence.  For
each first-match slope choose one witness support in its cell.  A fixed
noncommon support carries at most one slope, so this choice is injective
and proves the first inequality in \textup{(PO6)}.  Exact add-back gives
the equality there, and \textup{(PO7)} is the same inequality rewritten
using \(\abs{\Omega_\lambda}=L_\lambda\bar N_\lambda\).

The exact challenge bound follows by summing \textup{(PO7)}, taking the
maximum over received lines, and using \textup{(PO2)}.  Since
\(L_\lambda\le A_\lambda\), condition \textup{(PO8)} turns
\textup{(PO7)} into
\(e^{o(n)}\bar N_\lambda\).  There are only polynomially many quadruples
\((t,m,p,r)\), so their sum proves \textup{(PO9)}.  If
\(A_\lambda\le\abs\B^{R_\lambda}\), a polynomial-size \(\B\) and
bounded \(R_\lambda\) imply \textup{(PO8)}.  Finally, the effective
difference group in \textup{(PO5)} has at most \(A_\lambda\) characters.
Under \textup{(PO10)}, designating all of its nontrivial characters major
makes the minor sum empty and bounds the normalized major aggregate by
\(A_\lambda-1=e^{o(N_\lambda)}\).
\end{proof}

\begin{remark}[Exact dimension rounding]
\label{rem:qr-rounding}
If the global locator degree is $A=cm+r$, the list-code dimension is $K$,
and $w=A-K$, then the quotient depth in
\cref{thm:exact-quotient-remainder-normal-form} is
\[
 d=\left\lfloor\frac{A-K}{c}\right\rfloor
   =m-\left\lceil\frac{K-r}{c}\right\rceil.         \tag{QR5b}\label{eq:qr-rounding}
\]
Thus the descended list dimension is exactly
\(K_c=\lceil(K-r)/c\rceil\).  More explicitly, if \(w=cu+s\) with
\(0\le s<c\), then
\[
 \abs{\B_\phi}^{-\floor{w/c}}
 =\abs{\B_\phi}^{s/c}\abs{\B_\phi}^{-w/c}.          \tag{QR5c}
\]
Thus the exact rounding multiplier is \(\abs{\B_\phi}^{s/c}\), lying
between one and \(\abs{\B_\phi}^{(c-1)/c}\).  Retaining the integer
\(\floor{w/c}\) avoids even this comparison; no polynomial or
\(e^{o(n)}\) placeholder is needed in a finite certificate.
\end{remark}

\begin{remark}[Nonmonic maps and Chebyshev maps]
\label{rem:qr-chebyshev}
If the leading coefficient of $\phi$ is $b\ne0$, replace $\phi$ by
$b^{-1}\phi$ and $Q$ by $b^{-1}Q$.  This changes quotient coefficients by
fixed nonzero scalar factors and leaves every support and every fiber
cardinality unchanged.  The theorem therefore applies to the Chebyshev
folding map $T_c$ after this normalization.  On a circle twin-coset domain
on which $T_c$ has complete fibers of size $c$, it gives the same statement
verbatim.  The fixed or ramified endpoints form an exceptional set \(X\)
of cardinality \(b\) and are included exactly by the factor \(\binom bt\)
in \textup{(PO1)}; the crude total exceptional-point multiplier is at most
\(2^b\), with \(b\le2\) for involution-fixed endpoints.  Coincident twin
cosets are not a bounded planted set: they form the separate dihedral
collapse profile of the circle edge-case audit imported from \cite[Sec.~11]{CAP25v132}.
\end{remark}

For a declared complete-fiber-plus-Euclidean-remainder profile, the exact
normal-form theorem separates algebraic bookkeeping from prefix
equidistribution only when the remainder degree is less than \(c\).  The
exact disintegration of \cref{thm:exact-partial-occupancy} exhausts
arbitrary partial-occupancy supports, but when \(\deg P_R\ge c\) the
quotient and remainder coefficients interlace.  Thus that disintegration
does not by itself prove prefix flatness on the resulting coarse slices.
Let
$\B_\phi\subseteq\B$ be a subfield and suppose
\(Q\subseteq\eta\B_\phi\) for some $\eta\in\B^\times$.  Then
\(v_j(E)\in\eta^j\B_\phi\).  Hence, in case
\textup{(i)}, the natural average scale of the quotient-prefix fiber with
fixed $R$ is
\[
 \overline N_{\phi,R}(w)
 :=\binom{N-|\phi(R)|}{m}
       |\B_\phi|^{-\lfloor w/c\rfloor}.             \tag{QR6}\label{eq:qr-natural-scale}
\]
Pigeonholing proves a fiber of size at least
\(\lceil\overline N_{\phi,R}(w)\rceil\).  An upper bound of the same order
is not a consequence of the normal form: it is precisely the uniform
prefix-flatness, or Sidon/Fourier, input on the quotient row.

\begin{proposition}[Identity--quotient exponent comparison]
\label{prop:identity-quotient-comparison}
Consider a sequence to which
\cref{thm:exact-quotient-remainder-normal-form} applies with fixed
$c\ge2$ and fixed $0\le r<c$.  Suppose
\[
 \frac an\longrightarrow\alpha\in(0,1),\qquad
 \log|\B|=o(n),\qquad
 \log\binom na-w\log|\B|=o(n),                     \tag{QR7}\label{eq:identity-crossing-qr}
\]
where $a=cm+r$, and suppose
\(Q\subseteq\eta\B_c\) for a subfield $\B_c\subseteq\B$.
Write \(\overline N_{c,r}(w)=\overline N_{\phi,R}(w)\) for the natural
scale in \eqref{eq:qr-natural-scale}, with this fixed remainder profile.
Then $w\to\infty$, so $w\ge r$ eventually, and
\[
 \log\overline N_{c,r}(w)
 =\frac1c\log\!\left(\binom na|\B|^{-w}\right)
  +\frac wc\log\frac{|\B|}{|\B_c|}+o(n),           \tag{QR8}\label{eq:qr-comparison-general}
\]
where replacing $N-|\phi(R)|$ by $N$ changes only the $o(n)$ term.
Equivalently, if
\(\lambda_c=\log|\B_c|/\log|\B|\) has a limit, then at the identity
crossing
\[
       \frac1n\log\overline N_{c,r}(w)
       =\frac{h(\alpha)}c(1-\lambda_c)+o(1),         \tag{QR9}\label{eq:qr-comparison-crossing}
\]
using the entropy convention fixed in the Introduction.
Thus the natural average scale of this bounded-degree quotient profile is
comparable with the identity scale at exponential accuracy exactly when
\(\log|\B_c|/\log|\B|=1-o(1)\).  A proper field drop of fixed relative
degree produces a positive exponential quotient term.
\end{proposition}

\begin{proof}
Stirling's formula, $a/n\to\alpha$, and fixed $c,r$ give
\[
 \log\binom{n/c+O(1)}{(a-r)/c}
 =\frac1c\log\binom na+o(n).
\]
Also
\(\lfloor w/c\rfloor\log|\B_c|
=(w/c)\log|\B_c|+o(n)\), because the rounding error is at most
\(\log|\B_c|\le\log|\B|=o(n)\).  This proves
\eqref{eq:qr-comparison-general}.  Since
\(\log\binom na=nh(\alpha)+o(n)\),
\eqref{eq:identity-crossing-qr} gives
\(w\log|\B|=nh(\alpha)+o(n)\).  In particular $w\to\infty$ because
\(\log|\B|=o(n)\).  Substitution in
\eqref{eq:qr-comparison-general} proves
\eqref{eq:qr-comparison-crossing}.
\end{proof}

\begin{remark}[Unbounded quotient degrees]
If $c=c_n\to\infty$, the complete-fiber choice contributes at most
\(2^{n/c}=e^{o(n)}\).  When $w\ge r$, the remainder is recovered exactly
and no further remainder multiplicity occurs.  When $w<r<c$, the exact
formula \eqref{eq:qr-shallow-remainder} reduces the profile to an ordinary
prefix problem for the remainder support, multiplied by $e^{o(n)}$
quotient choices.  Thus only bounded quotient degrees can create a new
positive-rate field-drop term; large degrees either cost $e^{o(n)}$ or
return to an identity/remainder prefix problem.  This does not prove
flatness of the remaining prefix problem.
\end{remark}


\section{Effective-image Fourier normalization}\label{sec:partial-occupancy}
There is, however, an exact normalization that must be performed before
declaring any character major.  Assume \(1\le m\le\abs T-1\), as on every
fixed-density primitive leaf, put \(p=\operatorname{char}\B\), and fix
\(S_0\in\binom Tm\) and \(t_0\in T\).
Then
\[
 \Span_{\F_p}\{\Psi(S)-\Psi(S_0):S\in\tbinom Tm\}
 =\Span_{\F_p}\{g(t)-g(t_0):t\in T\}.               \tag{EF0}
\]
Define the \emph{effective Fourier span}
\[
 V_g=\Span_{\F_p}\{g(t)-g(t_0):t\in T\}\subseteq\B^R,
 \qquad A_{\rm eff}=\abs{V_g}.
 \tag{EF1}\label{eq:effective-fourier-span}
\]
The fixed-weight sum map takes values in the affine space
\(m g(t_0)+V_g\).  Thus Fourier inversion belongs on the additive group
\(V_g\), not automatically on the ambient group \(\B^R\).

\begin{lemma}[Effective-span Fourier inversion]
\label{lem:effective-span-fourier}
For \(z\in V_g\), let
\[
 N_g(z)=\abs{\left\{S\in\binom Tm:
       \sum_{t\in S}g(t)=m g(t_0)+z\right\}}.
\]
Then
\[
 N_g(z)=\frac1{A_{\rm eff}}
 \sum_{\chi\in\widehat{V_g}}\overline{\chi(z)}
 e_m\bigl(\chi(g(t)-g(t_0)):t\in T\bigr).           \tag{EF2}
\]
In particular,
\[
 \max_zN_g(z)\le\frac1{A_{\rm eff}}
 \left(\binom{\abs T}m+
 \sum_{1\ne\chi\in\widehat{V_g}}
 \left|e_m\bigl(\chi(g(t)-g(t_0)):t\in T\bigr)\right|\right). \tag{EF3}
\]
Every character of \(V_g\) extends to a trace character of \(\B^R\),
and a nontrivial effective character is not constant on all the generators
\(g(t)-g(t_0)\).  Ambient characters that are trivial on \(V_g\) occur
with multiplicity \(\abs{\B}^R/A_{\rm eff}\), which cancels exactly
against the ambient Fourier normalization.
\end{lemma}

\begin{proof}
The inclusion from left to right in \textup{(EF0)} follows by expanding a
difference of two support sums.  Conversely, for distinct \(t,u\in T\),
choose an \((m-1)\)-set \(A\subseteq T\setminus\{t,u\}\); such a set
exists because \(1\le m\le\abs T-1\).  Then
\(\Psi(A\cup\{t\})-\Psi(A\cup\{u\})=g(t)-g(u)\), so the two spans agree.
The first formula is character orthogonality on the finite additive group
\(V_g\), applied after translating by \(m g(t_0)\); the second follows by
the triangle inequality.  The trace pairing on the \(\F_p\)-space
\(\B^R\) is nondegenerate, so every character of its subspace \(V_g\)
extends to the ambient space.  A character trivial on each displayed
generator is trivial on their span.  Finally, the annihilator
\(V_g^\perp\) has size \(\abs{\B}^R/A_{\rm eff}\); grouping ambient
characters by their restriction to \(V_g\) proves the cancellation claim.
\end{proof}

\begin{definition}[Certified effective major/minor partition]
\label{def:effective-major-minor}
Let
\(\operatorname{res}:\widehat{\B^R}\twoheadrightarrow\widehat{V_g}\)
be restriction.  A \emph{certified effective Fourier partition} is a
disjoint union
\[
 \widehat{V_g}\setminus\{1\}
 =\mathfrak m_{\rm eff}\sqcup\mathfrak M_{\rm eff}
\]
such that every \(\chi\in\mathfrak m_{\rm eff}\) is supplied with an
ambient lift whose rational phase satisfies the character-sum hypotheses
used to pay it.  Characters without such a certified lift may be placed in
\(\mathfrak M_{\rm eff}\), but their aggregate must then be paid by
\textup{(MA)}.  This designation is data of the chart; it is not a
property of an unspecified ambient lift.
\end{definition}

\begin{definition}[Effective Fourier payment]
\label{def:effective-fourier-payment}
The full slice has an \emph{effective Fourier payment with constant
\(\kappa\ge1\)} if
\[
 \sum_{1\ne\chi\in\widehat{V_g}}
 \left|e_m\bigl(\chi(g(t)-g(t_0)):t\in T\bigr)\right|
 \le (\kappa-1)\binom{\abs T}m.                    \tag{EFP}
\]
By \cref{lem:effective-span-fourier}, this implies the exact max-fiber
bound
\[
               \max_zN_g(z)\le
               \kappa\binom{\abs T}m/A_{\rm eff}. \tag{EF4}
\]
It also implies that the realized image contains at least
\(A_{\rm eff}/\kappa\) points.  The ambient instances of \textup{(PF)}
and \textup{(MA)} are one sufficient way to prove \textup{(EFP)}, but
\textup{(EFP)} is phase-dependent and does not pay an ambient annihilator
twice.
\end{definition}

\begin{proof}[Proof of the image-size assertion]
The fibers sum to \(\binom{\abs T}m\), and each nonempty fiber is at most
\(\kappa\binom{\abs T}m/A_{\rm eff}\) by \textup{(EF4)}.
Hence at least \(A_{\rm eff}/\kappa\) fibers are nonempty.
\end{proof}

\begin{lemma}[Constant and Artin--Schreier modes are annihilator modes]
\label{lem:constant-modes-annihilator}
Suppose an ambient trace character has phase constant on \(T\).  Then its
restriction to \(V_g\) is trivial.  In particular, if an ambient rational
phase is Artin--Schreier of the form \(G^p-G+c\) and \(G\) is defined at
the \(\B\)-rational points of \(T\), then its character is an annihilator
mode and disappears from the nontrivial sum in \textup{(EF2)}.
\end{lemma}

\begin{proof}
The quotient of the character values at \(g(t)\) and \(g(t_0)\) is one
for every \(t\), so the character is trivial on the generators of
\(V_g\).  For an Artin--Schreier phase,
\(\Tr_{\B/\F_p}(G(t)^p-G(t))=0\), hence its additive-character value is
constant on all points where it is defined.
\end{proof}

\begin{remark}[Effective normalization does not prove the remaining aggregate]
The effective-span reduction removes exactly the constant annihilator
multiplicity.  A nonconstant character may still factor through a quotient
map or have a phase-dependent estimate too weak for \textup{(EFP)}.  Such
modes require a recursive quotient estimate, the aggregate
\textup{(MA)}, or a direct full-slice bound.  Thus effective normalization
repairs the Fourier denominator but is not a proof of unrestricted deep-prefix
flatness.
\end{remark}

\begin{definition}[Major-character aggregate]
\label{def:major-arc-aggregate}
Let \(A_*\) be the additive target on which Fourier inversion is being
performed, let \(h:T\to A_*\), and let
\(\mathfrak M_*\subseteq\widehat{A_*}\setminus\{1\}\) be the designated
major characters.  The chart satisfies \textup{(MA)} on \(A_*\) if
\[
 \tag{MA}\label{eq:major-arc-aggregate}
 \sum_{\chi\in\mathfrak M_*}
 \left|e_m\bigl(\chi(h(t)):t\in T\bigr)\right|
 \le e^{o(\abs T)}\binom{\abs T}m .
\]
Unless the ambient target is explicitly named, \(A_*=V_g\) and
\(h(t)=g(t)-g(t_0)\).  For the ambient target \(A_*=\B^R\), the same
condition is equivalently
\[
 \abs\B^{-R}\sum_{\alpha\in\mathfrak M_{\rm amb}}
 \left|e_m\bigl(\psi(\alpha\cdot g(t)):t\in T\bigr)\right|
 \le e^{o(\abs T)}\barN_0.
\]
The condition is vacuous when the designated major set is empty.  It can
be proved by a recursive quotient census or replaced by a direct
full-slice max-fiber bound; first-match terminology alone does not imply
it.
\end{definition}

\begin{lemma}[Sparse effective major set]
\label{lem:sparse-major-payment}
On the effective group \(V_g\), let \(\mathfrak M_{\rm eff}\) be any set
of nontrivial characters designated major.  Its exact normalized Fourier
loss is at most \(\abs{\mathfrak M_{\rm eff}}\).  Consequently
\(\abs{\mathfrak M_{\rm eff}}=e^{o(\abs T)}\) proves the effective form
of \textup{(MA)}; for a finite row, the literal cardinality is a valid
major-character constant.
\end{lemma}

\begin{proof}
Every fixed-weight Fourier coefficient is a sum of \(\binom{\abs T}m\)
complex numbers of modulus one, and hence has absolute value at most that
binomial coefficient.  Sum this bound over the designated set.
\end{proof}

\begin{theorem}[Exact low-boundary-entropy closure]
\label{thm:small-effective-dual-closure}
Let \(\Omega=\binom Tm\), let \(W\) be a finite additive
\(\F_p\)-space, and let \(\Psi:\Omega\to W\) be a fixed-weight additive
boundary map.  Fix \(S_0\in\Omega\), put \(z_0=\Psi(S_0)\), and define
the effective difference span
\[
 V=\Span_{\F_p}\{\Psi(S)-\Psi(S_0):S\in\Omega\}\subseteq W.
\]
Then \(\Psi(\Omega)\subseteq z_0+V\).  Put
\[
       A=\abs V,\qquad L=\abs{\Psi(\Omega)},
       \qquad M=\binom{\abs T}m,\qquad \bar N=M/L.
\]
Designate every nontrivial character of \(V\) as major and no character
as minor.  Then the literal finite-row constants are
\[
 C_{\rm min}=0,\qquad C_{\rm maj}\le A-1,\qquad
 \max_z\abs{\Psi^{-1}(z)}\le L\bar N=M.            \tag{SE1}
\]
For every exact-agreement witness cell whose support projection is
contained in \(\Omega\), its first-match slope set satisfies
\[
                         \abs Z\le L\bar N=M.       \tag{SE2}
\]
Consequently, if \(\log A=o(\abs T)\), effective \textup{(MI)},
effective \textup{(MA)}, image-normalized Q, and the direct alternative
in \textup{(RC)} all hold with subexponential loss.  Taking the
support-restricted incidence over \(\Omega\) as one ordered cell is
exhaustive for that restricted incidence; it is exhaustive for the whole
exact-agreement incidence when \(\Omega=\binom Da\), or after an
exhaustive partition into such slices.  For a locator prefix with
\(\abs T=\Theta(n)\), the concrete condition
\((a-k-1)\log\abs\B=o(n)\) implies this hypothesis because
\(A\le\abs\B^{a-k-1}\).
\end{theorem}

\begin{proof}
Fourier inversion is performed on the effective group \(V\).  The minor
sum is empty.  Every fixed-weight Fourier coefficient has
modulus at most \(M\), so the normalized major loss is at most \(A-1\).
The max-fiber estimate in \textup{(SE1)} is the trivial bound by the whole
support slice, written as the exact image-normalized multiplier \(L\).
By exact-agreement reduction, every bad slope in the cell has a noncommon
\(m\)-support, and a fixed noncommon support carries at most one slope.
Choosing one support for each slope is therefore injective and proves
\textup{(SE2)}.  Since \(L\le A\), all three finite multipliers
\(A-1,L,L\) are \(e^{o(\abs T)}\) when \(\log A=o(\abs T)\).
The one-cell atlas is exhaustive for its support-restricted incidence by
construction; the stated whole-incidence alternatives give global
exhaustivity.  If it is refined by a partial-occupancy partition,
\textup{(PO1)}--\textup{(PO2)} add the same support set back exactly and
do not change the conclusion.
\end{proof}

\begin{definition}[Aggregate minor payment]
\label{def:aggregate-minor-payment}
After effective-span normalization, let \(\mathfrak m_{\rm eff}\) be the
nontrivial characters not designated major.  The chart satisfies
\textup{(MI)} if
\[
 \sum_{\chi\in\mathfrak m_{\rm eff}}
 \left|e_m\bigl(\chi(g(t)-g(t_0)):t\in T\bigr)\right|
 \le e^{o(\abs T)}\binom{\abs T}m.                 \tag{MI}
\]
For a finite row the corresponding exact loss is this sum divided by
\(\binom{\abs T}m\).  The pointwise condition \textup{(PF)} is only one
sufficient way to prove \textup{(MI)}; phase-dependent estimates may be
strictly stronger.
\end{definition}

\begin{proposition}[Effective MI plus MA]
\label{prop:effective-mi-ma-flatness}
Partition the nontrivial effective characters as
\(\widehat{V_g}\setminus\{1\}
=\mathfrak m_{\rm eff}\sqcup\mathfrak M_{\rm eff}\).  Suppose
\textup{(MI)} and the effective form of \textup{(MA)} in
\cref{def:major-arc-aggregate} hold.
Then the full slice and every first-match residual satisfy
\[
 \max_z\abs{\{S:\Psi(S)=z\}}
 \le e^{o(\abs T)}\binom{\abs T}m/A_{\rm eff}.     \tag{EF5}
\]
The realized image has size at least
\(e^{-o(\abs T)}A_{\rm eff}\), so the same assertion holds at image
normalization.  For a finite row the exact multiplier is
\[
 1+\binom{\abs T}m^{-1}
 \sum_{1\ne\chi\in\widehat{V_g}}
 \left|e_m\bigl(\chi(g(t)-g(t_0)):t\in T\bigr)\right|.
 \tag{EF6}
\]
\end{proposition}

\begin{proof}
Add the minor and major sums and apply
\cref{lem:effective-span-fourier,def:effective-fourier-payment}.  A
first-match residual only deletes supports from a full-slice fiber.
\end{proof}

\begin{corollary}[Exact finite minor and major constants]
\label{cor:exact-finite-fourier-constants}
Put \(M=\binom{\abs T}m\) and, for an arbitrary partition of the
nontrivial effective dual, define
\[
 C_{\rm min}=M^{-1}\sum_{\chi\in\mathfrak m_{\rm eff}}\abs{E_m(\chi)},
 \qquad
 C_{\rm maj}=M^{-1}\sum_{\chi\in\mathfrak M_{\rm eff}}\abs{E_m(\chi)},
\]
where
\(E_m(\chi)=e_m(\chi(g(t)-g(t_0)):t\in T)\).  Then
\[
              \max_zN_g(z)\le
              \frac{M}{A_{\rm eff}}
              (1+C_{\rm min}+C_{\rm maj}).          \tag{EF7}
\]
If every minor character has power sums bounded by \(\Lambda\) through
order \(m\), then
\[
 C_{\rm min}\le
 \frac{\abs{\mathfrak m_{\rm eff}}}{M}
 \binom{\ceil\Lambda+m-1}{m}.                      \tag{EF8}
\]
Thus \textup{(EF7)}, with the two actual finite sums, records every finite
Fourier loss in an adjacent-row certificate exactly.
\end{corollary}

\begin{proof}
Formula \textup{(EF7)} is \textup{(EF3)} after splitting the nontrivial
characters.  The cycle-index coefficient bound gives
\(\abs{E_m(\chi)}\le\binom{\ceil\Lambda+m-1}{m}\), proving
\textup{(EF8)}.
\end{proof}

\begin{definition}[Pointwise minor-arc range]\label{def:prefix-flat-range}
A primitive weighted prefix chart satisfies the pointwise minor-arc range if the characteristic is larger than \(m\), every certified minor lift has a rational phase of total pole-degree at most \(\Delta_{\rm pole}\) and satisfies the nondegeneracy hypotheses of \cref{prop:weighted-weil-minor-arcs}, and, with
\[
        \Lambda=C_0(\Delta_{\rm pole}+1)\sqrt{\abs\B}+\abs P
\]
for the absolute Weil constant \(C_0\) and planted exceptional set \(P\), one has
\[
        \tag{PF}
        R\log\abs\B+
        \log\binom{\Lambda+m-1}{m}
        -\log\binom{\abs T}{m}\le o(\abs T).
\]
For circle twin-cosets the same definition is used with \(2C_0(\Delta_{\rm pole}+1)\sqrt{\abs\B}\) in place of \(C_0(\Delta_{\rm pole}+1)\sqrt{\abs\B}\), accounting for the two inverse-coset lifts on the torus before the involution quotient.
\end{definition}


\part{Primitive Q and the residual compiler}
\section{The primitive prefix leaf}\label{sec:primitive-leaf}

After first-match removal of paid cells, a residual prefix-boundary chart
has a fixed set of active coordinates \(T\), \(\abs T=N\), and a fixed
support size \(m\).  The full profile slice is
\[
        \Omega_{T,m}\defeq\{x\in\{0,1\}^{T}:\sum_{t\in T}x_t=m\},
\]
and the primitive residual is a subset \(\Omc\subseteq\Omega_{T,m}\).
The superscript circle records a first-match residual; no Zariski openness
is asserted.  The normalization is taken from the full profile slice, not
from the possibly much smaller residual mass.  It is nevertheless the
\emph{realized image}, rather than the formal codomain, that supplies the
default number of target fibers.

The three words in \emph{primitive first-match residual} have separate
content.  ``First match'' means that a slope is removed as soon as one of
its witnesses occurs in an earlier ordered cell; ``residual'' means that
only witnesses surviving those removals remain; and ``primitive'' means
that the row-specific quotient, planted, field-descent, rank, and ray-saturation
certificates named by the atlas have all been applied.  Primitivity does
not mean merely trivial setwise stabilizer, and it does not assume the
Fourier, Sidon, max-fiber, or ray estimates that will be imposed below.

The prefix syndrome map is a linear map \(\Phi:\Z^T\to\B^R\), restricted
first to the full slice and then to \(\Omc\).  In the unprofiled locator
chart, when \(R<\operatorname{char}\B\), it is equivalent through Newton's
triangular change of coordinates to
\begin{equation}
        \Phi(x)=\sum_{t\in T}x_t(t,t^2,\ldots,t^R).
        \label{eq:exact-power-sum-map}
\end{equation}
Equivalently, the augmented columns
\(\widetilde v_t=(1,t,\ldots,t^R)\) have constant first coordinate
\(\sum_tx_t=m\), and the other \(R\) coordinates are \(\Phi\).  A residual
quotient or rational chart is allowed to replace these columns by
\(v_t=\rho(t)(1,t,\ldots,t^{R-1})\) only when that genuine weighted
Vandermonde representation, with \(\rho(t)\ne0\), has been proved for the
chart.  It is not a consequence merely of constructibility or of deleting
a Jacobian locus.  The number \(R\) is the number of independent prefix
equations; in the unprofiled MCA boundary leaf \(R=w=a-K\).

Put \(M=\abs{\Omega_{T,m}}\), and write
\begin{equation}
 \Scal=\Phi(\Omega_{T,m}),\qquad L=\abs\Scal,
 \qquad A=\abs\B^R,
 \qquad \barN^{\rm img}=M/L,
 \qquad \barN^{\rm amb}=M/A,
 \label{eq:image-ambient-scales}
\end{equation}
and for \(s\in\Scal\) put
\(F_s=\Omc\cap\Phi^{-1}(s)\) and \(f_s=\abs{F_s}\).  Unless an
ambient full-image certificate has been proved, we set
\(\barN=\barN^{\rm img}\).  The certificate is
\[
        L\ge e^{-o(N)}A.                       \tag{FI}\label{eq:full-image-certificate}
\]
The set \(\Scal\) is the \emph{effective image}: it consists of the
prefix targets actually attained by the full profile slice.  The formal
ambient target space \(\B^R\) can be much larger.  Condition \textup{(FI)},
the \emph{full-image certificate}, says that these two target counts differ
by at most a subexponential factor.  An \emph{effective-image collapse} is
the failure of this comparison; it must be retained as an effective-image
rank-collapse profile, or handled throughout with image normalization.
Under \eqref{eq:full-image-certificate}, the two scales differ by only
\(e^{o(N)}\), so either may be used at exponential accuracy.  If the
certificate fails, effective-image collapse must be retained as an
effective-image rank-collapse profile, or all estimates on this leaf must stay
at image scale.  Deleting earlier cells can only decrease every fiber.

The image-scale frontier condition is
\(\log M-\log L=o(N)\).  The familiar ambient equation
\(\log\binom Nm-R\log\abs\B=o(N)\) is interchangeable with it only under
\eqref{eq:full-image-certificate}.  This distinction is immaterial for a
surjective prefix map but essential for a collapsed primitive image.

\begin{lemma}[Exact image--ambient moment conversion]
\label{lem:image-ambient-moment-conversion}
For \(q\ge1\), extend \(f_s\) by zero from \(\Scal\) to \(\B^R\) and
define
\[
 \Gamma_q^{\rm img}=\frac1L\sum_{s\in\Scal}
       \left(\frac{f_s}{\barN^{\rm img}}\right)^q,
 \qquad
 \Gamma_q^{\rm amb}=\frac1A\sum_{s\in\B^R}
       \left(\frac{f_s}{\barN^{\rm amb}}\right)^q.
\]
Then
\begin{equation}
 \Gamma_q^{\rm amb}
   =\left(\frac{A}{L}\right)^{q-1}\Gamma_q^{\rm img}.
 \label{eq:image-ambient-moment-identity}
\end{equation}
Thus an ambient-normalized moment bound implies the corresponding
image-normalized bound.  The converse is valid with subexponential loss at
logarithmic order only under \eqref{eq:full-image-certificate}.
\end{lemma}

\begin{proof}
Since \(\barN^{\rm amb}=M/A\) and \(\barN^{\rm img}=M/L\), zero extension
and one rearrangement give
\[
 \Gamma_q^{\rm amb}
 =\frac1A\sum_{s\in\Scal}\left(\frac{Af_s}{M}\right)^q
 =\left(\frac AL\right)^{q-1}
   \frac1L\sum_{s\in\Scal}\left(\frac{Lf_s}{M}\right)^q.
\]
Because \(L\le A\), the asserted one-way implication follows.  Under
\eqref{eq:full-image-certificate},
\((q-1)\log(A/L)=o(Nq)\), which proves the reverse transfer at the scale
used here.
\end{proof}

\begin{remark}[Full-slice flatness certifies the image size]
\label{rem:flatness-certifies-image}
If a Fourier argument proves
\(\max_s\abs{\Omega_{T,m}\cap\Phi^{-1}(s)}
 \le e^{o(N)}\barN^{\rm amb}\), then averaging over the nonempty image
fibers gives \(A/L\le e^{o(N)}\).  Hence such an ambient flatness theorem
itself proves \eqref{eq:full-image-certificate}; image collapse cannot be
silently assumed away before that estimate has been established.
\end{remark}

\begin{definition}[Primitive Q]\label{def:primitive-q}
The \emph{Q branch} is the primitive max-fiber problem: it asks whether one
effective prefix target contains substantially more residual supports than
the image-normalized average.  A primitive prefix leaf satisfies Q with
loss \(e^{o(N)}\) if
\[
        \max_{s\in\Scal}f_s\le e^{o(N)}\barN^{\rm img}.
\]
A sequence of primitive leaves satisfies logarithmic-moment Q if, at every
logarithmic order \(q=q_N\to\infty\), \(q\le(\log N)^C\), used by the
argument,
\[
 \Gord_q\defeq\Gamma_q^{\rm img}
 =L^{-1}\sum_{s\in\Scal}
       \left(\frac{f_s}{\barN^{\rm img}}\right)^q
 \le e^{o(Nq)}.
\]
Here a \emph{logarithmic moment order} means
\(q=q_N\to\infty\) with \(q\le(\log N)^C\) for a fixed \(C\).  The
displayed normalized \(q\)-th moment is a R\'enyi-type probe of the fiber
distribution: letting \(q\) grow slowly makes it approximate the largest
normalized fiber without introducing an exponential moment-order cost.
An ambient formulation is permitted only when it is proved directly or
when \eqref{eq:full-image-certificate} has already been certified.
\end{definition}

The next lemma is the primitive logarithmic moment equivalence used in the earlier RS--MCA reductions.  It is included because it is the exact interface between moment estimates and max-fiber Q.

\begin{lemma}[Logarithmic moments and primitive Q]\label{lem:logmoment-q}
Assume \(q=q_N\to\infty\) and \(\log L/q=o(N)\).  Then
\(\Gord_q\le e^{o(Nq)}\) implies primitive Q.  Conversely, primitive Q
implies \(\Gord_q\le e^{o(Nq)}\).
\end{lemma}

\begin{proof}
Let \(M_*=\max_s f_s/\barN^{\rm img}\).  Since
\(M_*^q\le\sum_s(f_s/\barN^{\rm img})^q=L\Gord_q\), the moment bound gives
\(\log M_*\le(\log L+\log\Gord_q)/q=o(N)\).  Conversely,
\(\sum_sf_s\le\abs{\Omega_{T,m}}=L\barN^{\rm img}\), so
\(\Gord_q\le M_*^{q-1}\).  Primitive Q gives \(M_*=e^{o(N)}\), hence
\(\Gord_q=e^{o(Nq)}\).
\end{proof}

The genuine Vandermonde condition removes the ordinary rank-defect alternative.

\begin{lemma}[Augmented Vandermonde independence]\label{lem:vandermonde}
Let \(t_1,\ldots,t_r\in\B\) be distinct and \(r\le R+1\).  Then the
columns \(\widetilde v_{t_i}=(1,t_i,\ldots,t_i^R)\) are linearly
independent.  The same holds for \(\rho(t_i)(1,t_i,\ldots,t_i^{R-1})\)
when \(r\le R\) and every \(\rho(t_i)\ne0\).
\end{lemma}

\begin{proof}
In either case the leading \(r\times r\) minor, after removing the nonzero
column weights when present, is \((t_i^j)_{0\le j<r,1\le i\le r}\).  Its
determinant is \(\prod_{i<j}(t_j-t_i)\ne0\).
\end{proof}

In the high-energy argument below the Vandermonde lemma is not used:
Boolean-cube growth gives the contradiction directly.  The lemma certifies
only the exact locator chart, or a residual chart for which the displayed
weighted form has separately been verified; it does not classify arbitrary
rank defects as paid cells.


\section{Prefix coordinates and the Vandermonde map}\label{sec:prefix-coordinates}

This section explains why the primitive residual has the linear Vandermonde form used in \cref{sec:primitive-leaf}.  Let \(S\subseteq D\) have size \(m\), and write its monic locator as
\[
        Q_S(X)=\prod_{t\in S}(X-t)=X^m+q_1(S)X^{m-1}+\cdots+q_m(S).
\]
The first \(R\) boundary coordinates may be taken either as elementary symmetric coefficients \(q_1,\ldots,q_R\) or, when \(R<\operatorname{char}\B\), as power sums \(p_j(S)=\sum_{t\in S}t^j\), \(1\le j\le R\).  Newton identities give triangular polynomial changes of variables between these two coordinate systems; the diagonal entries are \(1,2,\ldots,R\), so the change is invertible in characteristic greater than \(R\).

For the unprofiled locator chart, fixing the support size makes the zeroth
power sum constant and Newton's identities identify the first \(R\)
locator coefficients with the power sums
\(\sum_tx_tt^j\), \(1\le j\le R\).  Thus
\eqref{eq:exact-power-sum-map} is exact.  For a general residual chart one
first fixes the planted roots, quotient data, common factors, and all
coordinates charged to earlier cells.  The fixed part contributes a
constant syndrome \(s_0\).  To use the primitive theorem one must then
exhibit the remaining map explicitly, either as the same power-sum map, as
a genuine common-weight Vandermonde map, or as a map for which the required
Sidon estimate has been proved directly.  Coordinate-dependent rational
weights do not become a Vandermonde system merely because they are
nonzero.  Here ``common-weight'' means that one scalar \(\rho(t)\)
multiplies every monomial coordinate of the column at \(t\), rather than
using a different weight for each coordinate.  This verification is part of the primitive-map requirement in
\cref{def:admissible-sequence}.

The fixed-density condition is also intrinsic.  In a prefix leaf the total
support size is fixed by the agreement threshold, and all paid planted
blocks have fixed size.  Thus the number of remaining choices is fixed.  If
a profile leaves several colors of coordinates, one obtains a product of
fixed-density slices; the proof applies to each color after encoding the
colors in a bounded alphabet.  For the main statement we present the
Boolean case because it is the MCA support case and because the
Boolean-cube difference bound applies directly.

\begin{lemma}[Newton-coordinate fiber equivalence]\label{lem:newton-equivalence}
Assume \(R<\operatorname{char}\B\).  For supports of fixed size, the fibers of the first \(R\) elementary locator coefficients are exactly the fibers of the first \(R\) power sums.  Consequently the Q max-fiber problem may be stated using either coordinate system.
\end{lemma}

\begin{proof}
Newton identities give, for \(1\le j\le R\),
\[
        p_j+q_1p_{j-1}+\cdots+q_{j-1}p_1+jq_j=0.
\]
Since \(j\) is invertible in \(\B\), \(q_j\) is determined by \(p_1,\ldots,p_j\), and conversely \(p_j\) is determined by \(q_1,\ldots,q_j\).  The transformations are triangular inverse polynomial maps on the first \(R\) coordinates.
\end{proof}

The \emph{description entropy} of a profile family is the logarithm of the
number of allowed auxiliary/profile descriptions.  It measures the census
of cells and is distinct from the algebraic description complexity of one
cell, which measures the number and degrees of its defining equations.

\begin{lemma}[Primitive profiles have subexponential multiplicity]\label{lem:profile-multiplicity}
If the cell ledger fixes only bounded-complexity quotient, planted, tangent,
extension, and split-pencil profiles, and the total description entropy of
all data outside the moving support is \(o(n)\) bits, then the number of
primitive profiles is \(e^{o(n)}\).
\end{lemma}

\begin{proof}
There are at most \(2^{o(n)}=e^{o(n)}\) descriptions by hypothesis.  In
particular, a planted or common block of size \(o(n)\) has
\(\sum_{j=o(n)}\binom nj=e^{o(n)}\) possible supports, but recording
\(s=o(n)\) freely chosen field values costs \(s\log_2\abs\B\) bits and is
subexponential only when that product is \(o(n)\).  Positive-density data
requires its own profile census and entropy payment; it cannot be discarded
as a bounded-complexity label.
\end{proof}

\begin{remark}
The subexponential profile lemma is the asymptotic counterpart of the
finite ledger census.  It is the reason a row proof must distinguish
between a genuine cell payment and a mere change of variables.  A change
of variables with exponentially many charts would consume the whole
frontier budget.
\end{remark}

\section{The Sidon split}\label{sec:sidon-split}

The Q moment can be large for two opposite reasons.  A fiber may contain
many repeated additive differences, which is the high-energy situation
seen by inverse additive combinatorics, or it may be large while its
differences almost never repeat, which is the Sidon-like situation.  The
split below separates these mechanisms before either one is estimated.

Every support is identified with its incidence vector in the torsion-free
group \(\Z^T\).  In particular, all sums and differences in this section
are integer operations, even when the Reed--Solomon coefficient field has
characteristic two.  For a finite set \(F\subseteq\{0,1\}^T\), define its
additive energy
\[
        E(F)=\abs{\{(a,b,c,d)\in F^4:a-b=c-d\}}.
\]
Let \(X,X'\) be independent uniform points of \(F\).  Then \(E(F)=\abs F^4\sum_z\Prob(X-X'=z)^2\).  It is convenient to normalize by
\[
        \Delta(F)=\frac{E(F)}{\abs F^3}.
\]
The quantity \(\Delta(F)\) is an additive closure probability:
\[
        \Delta(F)=\Prob_{a,b,c\in F}(a-b+c\in F).
\]
Indeed, once \(a,b,c\) are chosen, the fourth point in an energy quadruple is forced to be \(d=a-b+c\).

Thus additive energy counts repeated differences, and \(\Delta(F)\) is
their normalized density.  Values \(\Delta(F)=e^{-o(N)}\) indicate
substantial additive closure at exponential scale, whereas
\(\Delta(F)\le e^{-\sigma N}\) is the low-energy, Sidon-like regime.
The logarithmic moments below are the slowly growing normalized moments
from \cref{def:primitive-q}; they measure how much of the Q obstruction is
carried by fibers of either type.
Throughout these sections, \emph{positive rate} means a fixed positive
exponential rate along a subsequence: for a fiber it means an excess of
the form \(e^{cN-o(N)}\), and for a logarithmic moment it means an excess
of the form \(e^{cNq-o(Nq)}\), for some constant \(c>0\).

For the fiber \(F_s\), write \(\Delta_s=E(F_s)/f_s^3\), with \(\Delta_s=0\) if \(f_s=0\).

\begin{definition}[Sidon-heavy logarithmic moment]\label{def:sidon-heavy}
For \(\sigma>0\), the Sidon-heavy part of the normalized moment is
\[
        \Gsid_{q,\sigma}=L^{-1}\sum_{\Delta_s\le e^{-\sigma N}}
        \left(\frac{f_s}{\barN}\right)^q.
\]
A sequence has no positive-rate Sidon-heavy obstruction if \(\Gsid_{q,\sigma}\le e^{o(Nq)}\) for every fixed \(\sigma>0\) and every logarithmic \(q\).
Here ``Sidon'' refers to the exponentially small normalized energy
\(\Delta_s\), while ``heavy'' refers to the positive-exponential
contribution of those fibers to the normalized logarithmic moment; it does
not assert that every low-energy fiber is individually large.
\end{definition}

The following dichotomy is the correct replacement for the naive assertion that collision excess directly creates small doubling.

\begin{proposition}[Ordinary moment split]\label{prop:ordinary-moment-split}
Fix \(\eta>0\) and \(\sigma>0\).  If \(\Gord_q\ge e^{\eta Nq}\) and \(\Gsid_{q,\sigma}\le \frac12 e^{\eta Nq}\), then there is a syndrome \(s\) such that
\[
        f_s\ge e^{\eta N-o(N)}\barN,
        \qquad
        \Delta_s>e^{-\sigma N}.
\]
If no positive-rate Sidon-heavy obstruction is present, then any positive-rate moment failure gives, after taking \(\sigma\to0\) slowly, a fiber with \(f_s\ge e^{\eta N-o(N)}\barN\) and \(E(F_s)\ge f_s^3e^{-o(N)}\).
\end{proposition}

\begin{proof}
The high-energy part of the moment is at least \(\frac12e^{\eta Nq}\).  Since there are \(L\) syndromes, some high-energy syndrome satisfies
\[
        \left(\frac{f_s}{\barN}\right)^q
        \ge \frac{L}{2}\,e^{\eta Nq}\cdot L^{-1}
        =\frac12 e^{\eta Nq}.
\]
Thus \(f_s\ge e^{\eta N-o(N)}\barN\), and by construction \(\Delta_s>e^{-\sigma N}\).  If \(\Gsid_{q,\sigma}\le e^{o(Nq)}\) for every fixed \(\sigma\), apply the first part along a sequence \(\sigma_j\downarrow0\) and diagonalize; the normalization \(\Delta_s=E(F_s)/f_s^3\) gives the energy conclusion.
\end{proof}

The proof also explains the no-go point.  Suppose a fiber is heavy but \(\Delta_s\le e^{-\sigma N}\).  A closure-rich diagram using tests \(a-b+c\in F_s\) pays a factor \(\Delta_s\) for each nonredundant closure.  Such diagrams cannot account for the ordinary contribution \(f_s^q\) of \(q\) independent choices in \(F_s\).  Therefore Sidon-heavy fibers must be removed by a separate cell; they cannot be forced into the high-energy inverse branch.

\begin{definition}[Sidon moment payment]\label{def:sidon-paid-cell}
A primitive prefix leaf has a Sidon moment payment if
\(\Gsid_{q,\sigma}=e^{o(Nq)}\) for every fixed \(\sigma>0\) and every
logarithmic moment order used by the primitive-Q argument, uniformly over
profiles.  This is an analytic support-fiber statement.  It is not
equivalent to a first-match distinct-slope payment; the latter still
requires \textup{(RC)} or a direct ray bound.
\end{definition}

The equivalence in the definition follows from \cref{lem:logmoment-q} restricted to the low-energy set of syndromes.  Operationally, one proves payment either directly by bounding these fibers or indirectly by proving Fourier flatness on the primitive complement of algebraic major arcs.


\section{Why the Sidon cell is necessary}\label{sec:sidon-necessity}

It is tempting to try to prove primitive Q by expanding the logarithmic moment into diagrams and then showing that every primitive diagram contains many additive closure tests.  That strategy is correct only for a closure-rich submoment.  It cannot capture the ordinary moment on a heavy low-energy fiber.  The obstruction is simple enough to record as a lemma.

Let \(F=F_s\) be a fiber of size \(f\).  A closure test is an equation of the form \(a-b+c\in F\), with \(a,b,c\in F\) already chosen.  Its success probability is \(\Delta(F)=E(F)/f^3\).  A closure-rich diagram with \(\ell\) nonredundant closure tests and \(r\) free choices contributes, up to a subexponential diagram factor,
\[
        f^r\Delta(F)^\ell .
\]
After the free coordinates have been chosen, in the image normalization of a primitive prefix leaf, the remaining \(q-r\) coordinates are supplied only at the random scale \(\barN\).  Thus its contribution is at most
\[
        e^{o(Nq)}\barN^{q-r}f^r\Delta(F)^\ell .
\]

\begin{lemma}[No-go for ordinary moment capture]\label{lem:no-go}
Assume that \(f\ge e^{\eta N}\barN\), \(\Delta(F)\le e^{-\sigma N}\), and that a family of diagrams has \(r=o(q)\) free fiber choices and \(\ell\ge c q\) nonredundant closure tests.  Then the total contribution of these diagrams is smaller than the ordinary fiber contribution \(f^q\) by a factor
\[
        \exp\bigl(-(\eta+c\sigma)Nq+o(Nq)\bigr)
\]
up to changing \(\eta\) by \(o(1)\).
\end{lemma}

\begin{proof}
For one diagram the ratio to \(f^q\) is at most
\[
        e^{o(Nq)}\left(\frac{\barN}{f}\right)^{q-r}\Delta(F)^\ell
        \le e^{o(Nq)}\exp(-\eta N(q-r)-c\sigma Nq).
\]
Since \(r=o(q)\), this is \(\exp(-(\eta+c\sigma)Nq+o(Nq))\).  The number of bounded-complexity logarithmic diagrams is \(e^{O(q\log q)}=e^{o(Nq)}\), so summing over them does not change the exponential rate.
\end{proof}

The lemma has an important conceptual consequence.  A large Sidon-like fiber is not a hidden approximate group.  It is a different phenomenon: many independent choices land in the same syndrome, but their pairwise differences almost never repeat.  Additive-combinatorics inverse tools are designed for the opposite situation, where repeated differences are abundant.  Therefore the first-match ledger must contain a Sidon/Fourier cell before the primitive inverse argument begins.

The valid implication is exactly \cref{prop:ordinary-moment-split}.  Ordinary
moment mass splits into a low-energy part, which needs a separately proved
Sidon moment estimate, and a high-energy part, which is dispatched by
\cref{prop:high-energy-impossible}.

This mechanism does not replace a prefix-flatness or Sidon-payment theorem;
it identifies exactly where one is needed.  The first-match order is:
algebraic major arcs first, then a separately certified Sidon/Fourier cell,
and only then the high-energy primitive inverse step.  Naming a large
low-energy fiber a cell does not pay it.

\section{Fourier flatness as a Sidon payment}\label{sec:fourier}\label{sec:sidon-closure}

This section records the general weighted form of the Fourier/Sidon payment.  It is phrased for an arbitrary finite field \(\B\), because the primitive residual charts that come from the moduli ledger need not have unweighted power-sum coordinates.  Let \(T\) be a finite set, \(\abs T=N\), let \(g:T\to\B^R\) be any function, and set
\[
        \Omega_{T,m}=\{x\in\{0,1\}^T: \sum_{t\in T}x_t=m\},
        \qquad
        \Psi(x)=\sum_{t\in T}x_tg(t).
\]
In applications \(g(t)=(\rho_0(t),\rho_1(t)t,\ldots,\rho_{R-1}(t)t^{R-1})\) with nonzero rational weights \(\rho_i\) on the primitive chart.

Fourier inversion has three pieces.  The zero character produces the
average-fiber main term.  On the effective dual, the chart-specific
partition of \cref{def:effective-major-minor} designates certified minor
characters and the complementary major characters.  The ambient algebraic
exceptional locus of \cref{def:major-arc} helps construct that partition
but does not canonically determine it, because an effective character has
many ambient lifts.  The pointwise condition \textup{(PF)} controls every
certified minor lift, while the aggregate condition \textup{(MA)} controls
the sum of all designated major characters.  Neither estimate substitutes
for the other.

The condition \textup{(PF')} below is the raw power-sum flatness
inequality when one has a uniform bound for every nonzero character.
Condition \textup{(PF)} is its row-specific Weil specialization for minor
characters after substituting the pole-degree bound; the omitted major
characters are then handled by \textup{(MA)}.

\begin{theorem}[Prefix flatness from power-sum control]\label{thm:prefix-flatness-power-sum}
Fix a nontrivial additive character \(\psi\) of \(\B\).  For \(\alpha\in\B^R\setminus\{0\}\) and \(1\le j\le m\), put
\[
        P_j(\alpha)=\sum_{t\in T}\psi\bigl(j\alpha\cdot g(t)\bigr).
\]
Assume \(\abs{P_j(\alpha)}\le \Lambda\) for every \(\alpha\ne0\) and every \(1\le j\le m\).  Then, for every \(z\in\B^R\),
\[
        \abs{\Psi^{-1}(z)}
        \le \abs\B^{-R}\binom Nm+\binom{\Lambda+m-1}{m}.
\]
Consequently, if
\[
        \tag{PF'}
        R\log\abs\B+\log\binom{\Lambda+m-1}{m}-\log\binom Nm\le o(N),
\]
then \(\max_z\abs{\Psi^{-1}(z)}\le e^{o(N)}\abs\B^{-R}\binom Nm\).
\end{theorem}

\begin{proof}
By orthogonality,
\[
        \abs{\Psi^{-1}(z)}=\abs\B^{-R}
        \sum_{\alpha\in\B^R}\psi(-\alpha\cdot z)
        e_m\bigl(\psi(\alpha\cdot g(t)):t\in T\bigr),
\]
where \(e_m\) denotes the \(m\)-th elementary symmetric coefficient.  The term \(\alpha=0\) is the main term \(\abs\B^{-R}\binom Nm\).  For \(\alpha\ne0\), write \(a_t=\psi(\alpha\cdot g(t))\).  The generating identity
\(\sum_{r\ge0}e_r(a_t)u^r=\exp(\sum_{j\ge1}(-1)^{j-1}(\sum_ta_t^j)u^j/j)\) and the hypothesis \(\abs{\sum_ta_t^j}\le\Lambda\) for \(j\le m\) give
\(\abs{e_m(a_t:t\in T)}\le[u^m](1-u)^{-\Lambda}=\binom{\Lambda+m-1}{m}\).  Summing the nontrivial characters and using the outside factor \(\abs\B^{-R}\) gives the first bound.  The displayed condition \textup{(PF')} makes the error term subexponential relative to the random-fiber main term.
\end{proof}

\begin{remark}[Small characteristic cycles]\label{rem:small-characteristic-cycles}
If some integers \(j\le m\) vanish in the characteristic, the formal
Li--Wan cycle index \cite{LiWanSieve} separates the cycles divisible by
\(p=\operatorname{char}\B\); the corresponding coefficient majorant is
\([u^m](1-u)^{-\Lambda}(1-u^p)^{-(N-\Lambda)/p}\).  This bookkeeping
identity does not supply the missing character-sum bounds for those
cycles, exclude Artin--Schreier phases, or prove a uniform Sidon payment.
The pointwise theorem above is therefore deliberately stated in the
large-characteristic range.  A small-characteristic leaf needs a direct
cycle-by-cycle certificate or a separate image-normalized Q/Sidon theorem.
\end{remark}

\begin{proposition}[Weighted Weil minor arcs]\label{prop:weighted-weil-minor-arcs}
Let \(T\) be either a multiplicative coset in \(\B^\times\) or the \(x\)-coordinate of a torus twin coset.  Suppose that every certified minor lift \(\alpha\in\B^R\) has phase \(R_\alpha(t)=\alpha\cdot g(t)\) which is a nonconstant rational function on the corresponding genus-zero curve, is not of Artin--Schreier form, and has total pole-degree at most \(\Delta_{\rm pole}\).  If \(m<\operatorname{char}\B\), then the stated power-sum bound holds for every certified minor lift, with
\[
        \Lambda=C_0(\Delta_{\rm pole}+1)\sqrt{\abs\B}+\abs P
\]
for multiplicative cosets, and with \(2C_0(\Delta_{\rm pole}+1)\sqrt{\abs\B}+\abs P\) for circle twin-cosets.  Here \(P\) is the planted exceptional set and \(C_0\) is an absolute Weil constant.
\end{proposition}

\begin{proof}
For a multiplicative coset \(\theta H\subseteq\B^\times\), write the coset indicator as the average of the multiplicative characters of \(\B^\times/H\).  Each nontrivial power sum becomes an average of mixed multiplicative-additive character sums \(\sum_x\chi(x)\psi(jR_\alpha(\theta x))\).  Since \(1\le j<m<\operatorname{char}\B\), multiplication by \(j\) does not make the rational function constant or Artin--Schreier.  Weil's bound for one-variable mixed character sums \cite{Weil} gives \(C_0(\Delta_{\rm pole}+1)\sqrt{\abs\B}\).  Removing planted exceptional points changes the sum by at most \(\abs P\).  The circle case is the same after pulling back by \(x=(u+u^{-1})/2\) to the norm-one torus; the two torus branches give the factor two.
\end{proof}

\begin{theorem}[Unconditional shallow-prefix Fourier flatness]
\label{thm:unconditional-shallow-mi-ma}
Let \(p\to\infty\) through odd primes.  For each \(p\), let \(T_p\) be
either a multiplicative coset in \(\F_p^\times\) or a Chebyshev
twin-coset \(x\)-domain over \(\F_p\), in the latter case with
\(p\equiv3\pmod4\), and put \(N_p=\abs{T_p}\).
Fix \(0<\alpha<1/2\), choose integers \(m_p,R_p\) with
\[
 1\le R_p,\qquad
 \alpha\le \frac{m_p}{N_p}\le1-\alpha,
 \qquad g_p(t)=(t,t^2,\ldots,t^{R_p}),
 \qquad R_p\sqrt p=o(N_p),                         \tag{SFM1}
\]
and let \(\Psi_p(S)=\sum_{t\in S}g_p(t)\) on
\(\binom{T_p}{m_p}\).  Then \(V_{g_p}=\F_p^{R_p}\), and the certified effective partition may take
every nontrivial effective character to be minor and no character to be
major.  For some \(c_\alpha>0\),
\[
 \sum_{1\ne\chi\in\widehat{V_{g_p}}}
 \left|e_{m_p}\bigl(\chi(g_p(t)-g_p(t_0)):t\in T_p\bigr)\right|
 \le e^{-c_\alpha N_p}\binom{N_p}{m_p}.            \tag{SFM2}
\]
Consequently effective \textup{(MI)} holds, \textup{(MA)} is vacuous,
and
\[
 \max_z\abs{\Psi_p^{-1}(z)}
 \le\bigl(1+e^{-c_\alpha N_p}\bigr)
       p^{-R_p}\binom{N_p}{m_p}.                   \tag{SFM3}
\]
Indeed every \(z\in\F_p^{R_p}\) is attained for all sufficiently large
rows, and uniformly in \(z\),
\[
 \abs{\Psi_p^{-1}(z)}
 =p^{-R_p}\binom{N_p}{m_p}
  \left(1+O\bigl(e^{-c_\alpha N_p}\bigr)\right).   \tag{SFM4}
\]
Thus these smooth and circle prefix charts have unconditional effective
Fourier payment and image coverage in the shallow range \textup{(SFM1)}.
\end{theorem}

\begin{proof}
Condition \textup{(SFM1)} implies \(R_p<p\), \(R_p<N_p\), and
\(R_p\log p=o(N_p)\).  Moreover \(N_p\le p+1\), so the fixed upper
density bound gives \(m_p<p\) for every sufficiently large \(p\).
For a nonzero ambient parameter, let
\(d=\max\{j:a_j\ne0\}\).  Then \(1\le d\le R_p<p\), and
\(P(X)=\sum_{j=1}^{R_p}a_jX^j\) has a pole of order \(d\) at infinity.
Every pole of a nonconstant Artin--Schreier function \(G^p-G+c\) has
order divisible by \(p\), so \(P\) is non-Artin--Schreier.  In the circle
case \(P((u+u^{-1})/2)\) has poles of the same order \(d<p\) at the two
ends of the torus compactification and is likewise non-Artin--Schreier.

For a multiplicative coset, the character expansion of its indicator and
the mixed Weil bound give, uniformly for \(1\le j\le m_p<p\),
\[
 \left|\sum_{t\in T_p}\psi(jP(t))\right|=O(R_p\sqrt p)=o(N_p).
\]
For a circle domain, pullback by
\(x=(u+u^{-1})/2\) gives the same conclusion with an inessential factor
two.  Put \(\Lambda_p=O(R_p\sqrt p)=o(N_p)\).  The cycle-index estimate
therefore bounds every nontrivial fixed-weight Fourier coefficient by
\[
 \binom{\ceil{\Lambda_p}+m_p-1}{m_p}=e^{o(N_p)}.
\]
If a nonzero ambient character annihilated \(V_{g_p}\), its character
value would be constant on \(T_p\), making the power sum have modulus
\(N_p\), contrary to the preceding \(o(N_p)\) bound.  Thus the
annihilator is zero and \(V_{g_p}=\F_p^{R_p}\).  There are
\(p^{R_p}=e^{o(N_p)}\) effective characters, whereas
Stirling's formula and the density interval give
\(\binom{N_p}{m_p}\ge
e^{(h(\alpha)+o(1))N_p}\).  Summing the coefficient bounds proves
\textup{(SFM2)}, after decreasing the positive constant
\(c_\alpha\).  All nontrivial effective characters have just received a
certified minor lift, so the major set is empty.  Finally
Fourier inversion shows that the absolute nonzero-character error in each
fiber is \(e^{o(N_p)}\), whereas its main term is
\(p^{-R_p}\binom{N_p}{m_p}=e^{(h(m_p/N_p)+o(1))N_p}\).
This proves \textup{(SFM3)}--\textup{(SFM4)} and eventual surjectivity.
\end{proof}

\begin{definition}[Ambient algebraic exceptional locus]\label{def:major-arc}
An ambient phase parameter is \emph{algebraically exceptional} if its
weighted phase is constant on a routed quotient fiber, is Artin--Schreier
degenerate, has a pole or zero supported on a positive-density moving
divisor, or factors through a proper quotient, Chebyshev map,
field-descent locus, tangent rank drop, differential-locator defect, or
ray-saturation collapse.  On an effective chart this locus is used only
to construct the certified partition of
\cref{def:effective-major-minor}; it is not itself a numerical estimate.
In particular, an Artin--Schreier phase defined on the whole active set
restricts to the trivial effective character by
\cref{lem:constant-modes-annihilator}.
\end{definition}

\begin{proposition}[Major-arc algebraic interface]\label{prop:major-arcs-are-cells}
Each condition in \cref{def:major-arc} is a constructible condition on the
dual phase parameters and points to a quotient, Chebyshev, field-descent,
planted-divisor, differential, or rank-loss profile.  This classification
does not imply that the character vanishes after primal first-match
deletion and does not prove \textup{(MA)}.  A flatness argument must also
bound \eqref{eq:major-arc-aggregate}, recursively or directly.
\end{proposition}

\begin{proof}
Constancy on a quotient fiber, Chebyshev factorization,
Artin--Schreier degeneracy, a positive-density pole divisor, and a rank
drop are respectively coefficient, Frobenius, resultant, or determinantal
conditions.  They therefore identify constructible dual loci on which a
recursive estimate can be attempted.  Fourier inversion sums characters,
not primal witnesses, so no support-level payment follows from this fact.
\end{proof}

\begin{proposition}[Frontier Weil separation]\label{prop:frontier-weil-separation}
Let \(N=\abs T\), let \(m=\theta N+o(N)\) with \(0<\theta<1\), and suppose a primitive smooth or circle chart has minor-arc power-sum parameter \(\Lambda=o(N)\).  If the ambient random-fiber scale \(\barN=\abs\B^{-R}\binom Nm\) satisfies \(\log\barN=o(N)\), or more generally \(\log\barN\ge -o(N)\), then the prefix-flat inequality \textup{(PF')} holds.  In particular, minor arcs alone cannot create a positive-rate max-fiber or Sidon-heavy moment at this near-zero ambient entropy scale.
\end{proposition}

\begin{proof}
The elementary estimate \(\binom{\Lambda+m-1}{m}=\binom{\Lambda+m-1}{\Lambda-1}\le (e(m+\Lambda)/\Lambda)^{\Lambda}\) gives \(\log\binom{\Lambda+m-1}{m}=o(N)\), since \(\Lambda=o(N)\).  The relative minor-arc error exponent is
\[
        \log\binom{\Lambda+m-1}{m}-\log\barN
        =R\log\abs\B+\log\binom{\Lambda+m-1}{m}-\log\binom Nm.
\]
The assumed frontier inequality makes this at most \(o(N)\); it may be
negative by a linear amount, which is only stronger.  Thus the minor-arc
contribution to every full-profile fiber is at most
\(e^{o(N)}\barN\).  This statement concerns only minor characters; a
full flatness conclusion also requires \textup{(MA)}.
\end{proof}

\begin{theorem}[Sidon payment in the pointwise flat range]\label{thm:sidon-resolved-payment}
Let \(\Omc\subseteq\Omega_{T,m}\) be a primitive first-match residual of a
smooth or circle chart, with weighted prefix map \(\Psi\) and ambient
scale \(\barN^{\rm amb}=\abs\B^{-R}\binom Nm\).  Put
\(A_{\rm amb}=\abs\B^R\) and
\(L=\abs{\Psi(\Omega_{T,m})}\).  Suppose the pointwise minor-arc
condition \textup{(PF)} and the ambient instance of the aggregate
condition \textup{(MA)} hold.
Then, for every fixed \(\sigma>0\) and every
logarithmic \(q\le(\log N)^C\), for a fixed \(C\),
\[
        \Gsid_{q,\sigma}(\Omc)\le e^{o(Nq)},
\]
where \(\Gsid\) uses the image normalization of
\cref{def:sidon-heavy}.
Equivalently, no such pointwise-flat residual leaf supports positive-rate
low-energy heavy fibers.
\end{theorem}

\begin{proof}
It is enough to bound every residual fiber, and a residual is contained in
the full slice.  Fourier inversion on that slice splits into the zero,
minor, and major characters.  The zero character contributes
\(\barN^{\rm amb}\).
The cycle-index estimate and \cref{prop:weighted-weil-minor-arcs}, together
with \textup{(PF)}, bound the complete minor contribution by
\(e^{o(N)}\barN^{\rm amb}\).  Condition \textup{(MA)} bounds the complete major
contribution by the same quantity.  Thus every full-slice, and hence every
residual, fiber is \(e^{o(N)}\barN^{\rm amb}\).  This proves the
full-image certificate by \cref{rem:flatness-certifies-image}, and
\cref{lem:image-ambient-moment-conversion} transfers the ambient moment
bound to image normalization.
\end{proof}

\begin{remark}[Role of the Sidon-heavy cell]
The theorem implements ``pay structure, flatten Sidon, inverse high energy''
only in the range where both \textup{(PF)} and \textup{(MA)} are verified.
At entropy depth a phase
can have degree up to \(R\), and the raw Weil parameter is of order
\(R\sqrt{\abs\B}\); this is not automatically \(o(N)\) when
\(R\asymp N/\log\abs\B\).  Outside the pointwise-flat range a different
Sidon payment remains necessary.
\end{remark}

\section{External additive combinatorics}\label{sec:additive-tools}

Only two external additive-combinatorics inputs are used in the primitive high-energy proof.

\begin{theorem}[Balog--Szemeredi--Gowers, energy form]\label{thm:bsg}
There is an absolute constant \(C\) such that if \(A\) is a finite subset of an abelian group and \(E(A)\ge \abs A^3/K\), then there is \(A'\subseteq A\) with \(\abs{A'}\ge K^{-C}\abs A\) and \(\abs{A'-A'}\le K^C\abs{A'}\).
\end{theorem}

This is the usual energy form of the Balog--Szemeredi--Gowers theorem; polynomial dependence on \(K\) follows from Gowers's refinement of the Balog--Szemeredi theorem \cite{BalogSzemeredi,GowersSzemeredi}.  The difference-set form follows from the sumset form by replacing \(A\) by \(-A\) and using standard Ruzsa calculus \cite{TaoVu}, or directly by applying the graph form to differences.

\begin{theorem}[Boolean-cube difference growth]\label{thm:quasicube}
If \(U\subseteq\{0,1\}^d\), then
\[
                  \abs{U-U}\ge \abs U^{3/2}.
\]
\end{theorem}

This is the Boolean-cube specialization of a broader sumset inequality
of Green--Matolcsi--Ruzsa--Shakan--Zhelezov
\cite{GreenMatolcsiRuzsaShakanZhelezov}.  It also follows from the
induced-doubling identity of Matolcsi--Ruzsa--Shakan--Zhelezov
\cite{MatolcsiRuzsaShakanZhelezov}.  Only the displayed Boolean statement
is used below.

\begin{remark}
Entropy inverse theorems of Tao and the entropy-doubling theory of Green--Manners--Tao are compatible with this proof and are often useful after one has produced small-doubling trades \cite{TaoEntropy,GreenMannersTao}.  They are not needed as hypotheses here.  The Boolean-cube growth theorem gives a direct contradiction to a large small-doubling subset of a primitive Boolean support fiber.
\end{remark}


\section{Entropy inverse theory and its role}\label{sec:entropy-role}

The terminology ``Tao--Gowers method'' can refer to two related but
distinct steps.  The first is BSG, used directly in
\cref{prop:high-energy-impossible}.  The second is Freiman-type structural
theory, meaning inverse theorems that model a set with small sumset or
difference set by a generalized arithmetic progression or a coset
progression.  It is not needed for the Boolean contradiction, but it can suggest
candidate geometric strata in more general bounded-alphabet residuals;
those strata still require direct payment proofs.

A \emph{bounded integer trade} is a vector \(Y\in\Z^T\) whose
coordinates are bounded independently of \(N\), with \(\Phi(Y)=0\) and,
on a fixed-weight slice, \(\sum_tY_t=0\).  If \(Y\) is random, \(Y'\) is
an independent copy, and
\(\Ent(Y-Y')\le \Ent(Y)+o(N)\), then entropy BSG and entropy Freiman
theory produce, after conditioning and losing \(o(N)\) entropy, a set
model with doubling \(e^{o(N)}\).  Tao's entropy inverse theorem gives
this in \emph{coset-progression} language---a subgroup coset plus a
generalized arithmetic progression---for Shannon entropy \cite{TaoEntropy};
Green--Manners--Tao recast the theory in terms of entropic doubling and
Ruzsa distance \cite{GreenMannersTao}; Goh formulates an entropic analogue
of additive energy and the entropic BSG connection \cite{GohEnergy}.  In a
torsion model one may also use polynomial Freiman--Ruzsa technology after
reducing bounded integer trades modulo a large prime with no wraparound.

In the present paper the Boolean-cube difference theorem makes this
structural compression unnecessary for the primitive Q step.  Once BSG
finds \(A'\subseteq F_s\subseteq\{0,1\}^T\) with
\(\abs{A'-A'}\le e^{o(N)}\abs{A'}\), Boolean-cube difference growth
already forces \(\abs{A'}=e^{o(N)}\).  This contradicts the exponential
size supplied by moment excess.  Therefore there is no separate
entropy-Freiman hypothesis in the proof.

For non-Boolean residual trades, entropy inverse theory can suggest
coset-progression incidence strata or rank-defect tests.  This is a useful
heuristic for constructing a row-specific atlas, not an exhaustive
dichotomy: each suggested stratum still needs the explicit first-match census and
slope-projection payment required by \cref{def:admissible-sequence}, and arbitrary weighted charts
are not covered by \cref{lem:vandermonde}.  No such heuristic is used in
the proof of the Boolean primitive theorem.

\section{Primitive Q after the Sidon cut}\label{sec:primitive-q-proof}

The key primitive theorem is now short and self-contained modulo \cref{thm:bsg,thm:quasicube}.  It also addresses the main obstruction: ordinary \Renyi{} mass can be Sidon-heavy, but after that cell is removed it must be high-energy, and high-energy is impossible in a large Boolean fiber.

\begin{proposition}[High-energy Boolean fibers are small]\label{prop:high-energy-impossible}
Let \(F\subseteq\{0,1\}^T\), \(\abs T=N\).  If
\(E(F)\ge \abs F^3/K\), then \(\abs F\le K^{O(1)}\).  Consequently an
exponential set \(\abs F\ge e^{cN-o(N)}\) cannot have
\(E(F)\ge\abs F^3/e^{o(N)}\).
\end{proposition}

\begin{proof}
Apply \cref{thm:bsg}.  There is \(F'\subseteq F\) with \(\abs{F'}\ge K^{-C}\abs F\) and \(\abs{F'-F'}\le K^C\abs{F'}\).  Since \(F'\subseteq\{0,1\}^T\), \cref{thm:quasicube} gives \(\abs{F'-F'}\ge \abs{F'}^{3/2}\).  Hence \(\abs{F'}^{1/2}\le K^C\), so \(\abs F\le K^{3C}\) after increasing \(C\).  Taking \(K=e^{o(N)}\) rules out \(\abs F\ge e^{cN-o(N)}\).
\end{proof}

\begin{theorem}[Primitive Q after a Sidon moment payment]\label{thm:primitive-q}
Let \(\Omega_N^0\subseteq\{0,1\}^{T_N}\) be fixed-density full profile
slices, with \(\abs{T_N}=N\) and
\(m_N/N\in[\alpha,1-\alpha]\), and let
\(\Omc_N\subseteq\Omega_N^0\) be primitive first-match residuals.  Suppose
there is a logarithmic moment
order \(q_N\to\infty\) with
\(\log L_N/q_N=o(N)\), where
\(L_N=\abs{\Phi_N(\Omega_N^0)}\) and
\(\barN_N=\abs{\Omega_N^0}/L_N\), for which
\(\Gsid_{q_N,\sigma}\le e^{o(Nq_N)}\) for every fixed \(\sigma>0\).
Then the leaves satisfy primitive Q:
\[
        \max_s\abs{\Omc_N\cap\Phi_N^{-1}(s)}
        \le e^{o(N)}\barN_N .
\]
\end{theorem}

\begin{proof}
Assume primitive Q fails at positive exponential rate.  The lower half of
the moment sandwich gives
\(\Gord_{q_N}\ge e^{\eta Nq_N-o(Nq_N)}\) along a subsequence.  The Sidon
moment hypothesis and \cref{prop:ordinary-moment-split} then give a fiber
\(F_s\) with \(f_s\ge e^{\eta N-o(N)}\barN_N\) and
\(E(F_s)\ge f_s^3e^{-o(N)}\).  Since
\(\barN_N=\abs{\Omega_N^0}/L_N\ge1\), this fiber has size
\(e^{\eta N-o(N)}\), contradicting
\cref{prop:high-energy-impossible} with \(K=e^{o(N)}\).  Thus no
positive-rate max-fiber failure occurs.
\end{proof}

\section{Q, SP, and the MCA compiler}\label{sec:compiler}
This section uses the prefix rigidity, exact second moment, and slope-elimination results of \cite[Secs.~19--21 and Apps.~E--F]{CAP25v132} as imported inputs.  The new content is the occupancy/ray compiler and the precise hypotheses under which Q controls MCA slopes.

The \emph{Q branch} is the maximum-prefix-fiber problem.  The \emph{SP
branch} (support-pair, or shift-pair, branch) is its second-moment problem:
it counts ordered pairs of supports in one prefix fiber.  These names label
two counting interfaces and carry no assertion that either bound has
already been proved.  Neither count is automatically the MCA numerator: a
bad slope begins with one support, while converting support pairs to
projective slope directions requires a proved multiplicity or incidence
bound.  We first record the exact moment identities and then state that
separate ray compiler.

\begin{definition}[Q and SP ledgers]\label{def:q-sp}
For a first-match residual profile \(\lambda\), let \(\Omega^0_\lambda\)
be its full profile slice, let
\(\Omc_\lambda\subseteq\Omega^0_\lambda\) be the residual, and put
\[
 \Scal_\lambda=\Phi_\lambda(\Omega^0_\lambda),\qquad
 L_\lambda=\abs{\Scal_\lambda},\qquad
 \barN_\lambda=\abs{\Omega^0_\lambda}/L_\lambda,
 \qquad
 Q_\lambda(s)=\abs{\Omc_\lambda\cap\Phi_\lambda^{-1}(s)}.
\]
The Q ledger is discharged if
\(\max_sQ_\lambda(s)\le e^{o(n)}\barN_\lambda\) for every primitive
profile.  The support-pair SP statistic is
\(\sum_sQ_\lambda(s)^2\).  A distinct-ray ledger is discharged only by a
direct slope bound or by \textup{(RC)} below.
If the ambient target count \(\abs{\B_\lambda}^{R_\lambda}\) is used in
place of \(L_\lambda\), the ledger must also record the full-image
certificate \eqref{eq:full-image-certificate} or a direct ambient theorem.
\end{definition}


The compiler uses two elementary inequalities repeatedly.  They are included to make explicit how the logarithmic moment, max-fiber Q, and shift-pair second moment interact.

\paragraph{Imported moment inequality.}
For fiber sizes \(N(s)\), mean \(\bar N\), ratios \(R(s)=N(s)/\bar N\),
and \(\Gamma_r=L^{-1}\sum_sR(s)^r\), CAP25 v13.2
\cite[Sec.~19 and Apps.~E--F]{CAP25v132} gives the elementary bounds
\[
        \max_sR(s)\le(L\Gamma_r)^{1/r},
        \qquad
        \Gamma_r\le(\max_sR(s))^{r-1}.
        \tag{M}\label{eq:imported-moment-interface}
\]


\begin{corollary}[Finite moment floor]\label{cor:finite-moment-floor}
Suppose a finite row has remaining Q bit margin \(\Delta\), meaning that one needs \(\max_sR(s)\le2^\Delta\).  A moment proof of order \(r\) can certify this only if
\[
        \log_2\Gamma_r+\log_2 L\le r\Delta,
\]
where \(L\) is the number of realized targets in the selected
normalization.  In an ambient prefix proof this specializes to
\(\log_2\Gamma_r^{\rm amb}+R\log_2\abs\B\le r\Delta\); transferring an
image-normalized moment to that statement additionally costs
\((r-1)\log_2(\abs\B^R/L)\).
\end{corollary}

\begin{proof}
The moment sandwich gives \(\log_2\max_sR(s)\le(\log_2L+\log_2\Gamma_r)/r\).  To make this at most \(\Delta\) one needs the displayed inequality.
\end{proof}

The finite corollary explains why a small adjacent margin cannot be closed
by a fixed-order moment when \(R\log\abs\B\) is large.

\begin{definition}[Explanation image and retained-support occupancy]
\label{def:explanation-occupancy}
Fix \(a\ge k+1\), a received pair \(r=(r_0,r_1)\), and a cell
\(\Ccal\subseteq\Wcal_a(r)\).  Its \emph{explanation image} and
\emph{slope image} are
\[
 \Rcal(\Ccal)=\{(\gamma,h):\exists S\ (\gamma,S,h)\in\Ccal\},
 \qquad Z(\Ccal)=\{\gamma:\exists h\ (\gamma,h)\in\Rcal(\Ccal)\}.
\]
For \(\rho=(\gamma,h)\in\Rcal(\Ccal)\), its
\emph{retained-support occupancy} is
\[
 \nu_{\Ccal}(\rho)
   =\abs{\{S\in\tbinom Da:(\gamma,S,h)\in\Ccal\}}.
\]
Thus the witness projection factors canonically as
\[
 \Ccal\longrightarrow\Rcal(\Ccal)\longrightarrow Z(\Ccal),
 \qquad (\gamma,S,h)\longmapsto(\gamma,h)\longmapsto\gamma.
\]
The first map forgets an exact-agreement support; the second forgets the
explaining polynomial.  Distinct explanation states may have the same
slope, so neither map is silently treated as a bijection.  A
\emph{ray compiler} means a proved estimate for the final slope image,
obtained from these actual fibers or from a direct incidence bound.
\end{definition}

\begin{lemma}[Exact occupancy compiler]
\label{lem:exact-occupancy-compiler}
With the notation above,
\[
 \abs{\Ccal}=\sum_{\rho\in\Rcal(\Ccal)}\nu_{\Ccal}(\rho),
 \qquad
 \abs{\Rcal(\Ccal)}
   =\sum_{w\in\Ccal}
       \frac1{\nu_{\Ccal}(\pi_{\rm exp}(w))},
 \qquad
 \abs{Z(\Ccal)}\le\abs{\Rcal(\Ccal)}.             \tag{RC$_{\rm occ}$}
\]
Here \(\pi_{\rm exp}(\gamma,S,h)=(\gamma,h)\).  In particular, if every
retained explanation state has occupancy at least \(H\ge1\), then
\[
        \abs{Z(\Ccal)}\le
        \floor{\frac{\abs{\Ccal}}H}.                \tag{RC1}
\]
Also, because one noncommon exact-\(a\) support carries at most one slope
and the explaining polynomial on \(a\ge k+1\) points is unique,
\[
 \abs{Z(\Ccal)}\le
 \abs{\{S:\exists\gamma,h\ (\gamma,S,h)\in\Ccal\}}. \tag{RC2}
\]
All these statements are exact finite-row identities or inequalities.
\end{lemma}

\begin{proof}
Partition \(\Ccal\) by the fibers of \(\pi_{\rm exp}\).  This gives the
first identity; summing the reciprocal of the fiber size over each fiber
gives the second.  Projection \((\gamma,h)\mapsto\gamma\) can only
decrease cardinality.  If every first fiber has size at least \(H\),
\textup{(RC1)} follows.  Finally, if the same noncommon support carried two
distinct slopes, slope elimination would simultaneously explain the
received pair on that support.  For fixed support and slope, uniqueness of
interpolation gives at most one \(h\).  This proves \textup{(RC2)}.
\end{proof}

\begin{lemma}[Max-occupancy add-back]
\label{lem:max-occupancy-addback}
Let \(\Ccal\subseteq\Wcal_a(r)\) be partitioned into profile cells
\(\Ccal=\coprod_{\lambda\in\Lambda}\Ccal_\lambda\), and put
\(P=\abs\Lambda\).  For every explanation state
\(\rho=(\gamma,h)\in\Rcal(\Ccal)\), let
\[
 M(\rho)=\abs{\{S:(\gamma,S,h)\in\Ccal\}}.
\]
Assign \(\rho\) to a profile \(\lambda(\rho)\) on which
\(\nu_{\Ccal_\lambda}(\rho)\) is maximal, and let
\(\Rcal_\lambda^\star\) be the assigned states.  Then
\[
 \nu_{\Ccal_\lambda}(\rho)\ge\ceil{M(\rho)/P}
 \qquad(\rho\in\Rcal_\lambda^\star).               \tag{RC$_{\rm max}$1}
\]
Consequently, if \(M(\rho)\ge H\) for every retained state, then
\[
 \abs{Z(\Ccal)}
 \le\sum_{\lambda\in\Lambda}\abs{\Rcal_\lambda^\star}
 \le\sum_{\lambda\in\Lambda}
   \floor{\frac{\abs{\Ccal_\lambda^\star}}
                   {\ceil{H/P}}},                  \tag{RC$_{\rm max}$2}
\]
where \(\Ccal_\lambda^\star\) consists of the witnesses in
\(\Ccal_\lambda\) lying above assigned states.  Call \(\Ccal\)
\emph{explanation-state complete} if, whenever
\((\gamma,h)\in\Rcal(\Ccal)\), every exact-\(a\) noncommon witness
\((\gamma,S,h)\in\Wcal_a(r)\) also belongs to \(\Ccal\).  If \(\Ccal\)
is explanation-state complete, every retained state has exact agreement
\(b\ge a\), and the received pair is not simultaneously explained on any
\(a\)-subset of that agreement set, then one may take \(H=\binom ba\).
\end{lemma}

\begin{proof}
The \(P\) profile occupancies of a fixed state sum to \(M(\rho)\), so
their maximum is at least \(\ceil{M(\rho)/P}\).  Apply
\cref{lem:exact-occupancy-compiler} on each assigned cell and sum.  Under
the final hypotheses, explanation-state completeness retains every
\(a\)-subset of the \(b\)-point agreement set, and every such subset is
noncommon.  Hence \(M(\rho)=\binom ba\).
\end{proof}

\begin{definition}[Support interpolation functionals]
\label{def:pair-to-ray-map}
Let \(T\subseteq D\), \(\abs T=m\ge k\), put \(w=m-k\), and let
\(Q_T(X)=\prod_{x\in T}(X-x)\).  For a word \(u\), define
\begin{equation}
 \mathcal A_T(u)=
 \left(
   \sum_{x\in T}\frac{u(x)x^r}{Q_T'(x)}
 \right)_{0\le r<w}\in\F^w.
 \label{eq:interpolation-functional}
\end{equation}
\end{definition}

\paragraph{Imported slope interface.}
CAP25 v13.2 \cite[Sec.~19 and Apps.~E--F]{CAP25v132} proves that a word
\(u\) is explained by \(\RS_\F(D,k)\) on \(T\) exactly when
\(\mathcal A_T(u)=0\).  Hence
\[
 \mathcal A_T(u_0)+\gamma\mathcal A_T(u_1)=0,
 \tag{SE}\label{eq:imported-slope-elimination}
\]
and every noncommon support carries at most one finite slope.  We use this only as an imported incidence input to the occupancy compiler above.


\paragraph{Imported support-pair identity.}
For every residual prefix profile \(\lambda\), the foundational paper
\cite[Sec.~19 and Apps.~E--F]{CAP25v132} identifies the ordered
same-prefix support-pair count as
\[
        \sum_sQ_\lambda(s)^2
        \le
        \bigl(\max_sQ_\lambda(s)\bigr)\sum_sQ_\lambda(s).
        \tag{SP$_0$}\label{eq:imported-support-pair}
\]


Q therefore controls a support-pair statistic.  It does not justify dividing
that statistic by \(\abs{\Omega^0_\lambda}\): a ray can have a single
representative at exact agreement.

\begin{proposition}[Q and SP do not determine rays]
\label{prop:q-sp-no-ray}
Let \(\Phi:\Omega\to\Sigma\) be any support-profile map with
\(\Omega\ne\varnothing\), put
\(Q(s)=\abs{\Phi^{-1}(s)}\), and
\({\rm SP}=\sum_sQ(s)^2\).  Regard \(\Omega\) as an abstract set of
candidate noncommon supports, and suppose the only retained
slope-incidence axiom is that each support carries at most one slope.  For
every nonempty finite slope set \(\Gamma\) and every
\(1\le M\le\min\{\abs\Omega,\abs\Gamma\}\), within the class of abstract
incidence structures obeying these retained data there is a partial
support-to-slope map having exactly \(M\) distinct values.  There is also
a one-value map.  Hence no ray bound can be deduced from Q, SP, and that
axiom alone; an RS-specific bound must use additional algebraic incidence
information.
\end{proposition}

\begin{proof}
Choose \(M\) distinct supports and \(M\) distinct slopes and map them
bijectively, leaving all other supports unlabeled.  This changes neither
\(\Phi\), its fibers, nor SP, and it respects one-slope-per-support.
Mapping the same chosen supports to one slope gives identical Q and SP
data with a one-ray image.  The missing datum is the image incidence, not
another moment of the prefix fibers.
\end{proof}

\begin{proposition}[Exact pair-to-ray multiplicity requirement]
\label{prop:pair-ray-multiplicity}
Let
\(\mathcal P=\{(A,B)\in\Omega^2:\Phi(A)=\Phi(B)\}\), so
\(\abs{\mathcal P}={\rm SP}\).  Let \(Z\) be a bad-slope set and
\(I\subseteq Z\times\mathcal P\) an incidence relation, and let
\(H,J\ge1\).  If every
\(\gamma\in Z\) is incident with at least \(H\) pairs and every pair is
incident with at most \(J\) slopes, then
\begin{equation}
 \abs Z\le\frac{J\,{\rm SP}}H
       \le\frac{J(\max_sQ(s))\abs\Omega}{H}.
 \label{eq:pair-ray-multiplicity}
\end{equation}
Consequently this pair-count argument yields a ray estimate at the same
profile scale provided \(J\abs\Omega/H=e^{o(n)}\); alternatively one may
prove a direct transverse-secant estimate.  Merely proving that every ray
has one representing pair is insufficient.
\end{proposition}

\begin{proof}
Double counting gives
\(H\abs Z\le\abs I\le J\abs{\mathcal P}=J\,{\rm SP}\).  The second
inequality follows from
\({\rm SP}\le(\max_sQ(s))\sum_sQ(s)=(\max_sQ(s))\abs\Omega\).
\end{proof}

\begin{corollary}[Exact finite pair-to-ray constant]
\label{cor:finite-pair-ray-compiler}
For a residual profile \(\lambda\), put
\(m_\lambda=\sum_sQ_\lambda(s)\),
\(L_\lambda=\abs{\Phi_\lambda(\Omega^0_\lambda)}\), and
\(\barN_\lambda=\abs{\Omega^0_\lambda}/L_\lambda\).  For an integer
\(q\ge2\), define
\[
 \Gamma_{q,\lambda}=L_\lambda^{-1}\sum_s
 \left(\frac{Q_\lambda(s)}{\barN_\lambda}\right)^q.
\]
If the pair-to-ray incidence has lower ray degree \(H_\lambda\) and upper
pair degree \(J_\lambda\), then
\[
 \abs{Z_\lambda^\circ}\le
 \floor{\frac{J_\lambda m_\lambda}{H_\lambda}\,
  \barN_\lambda
  \bigl(L_\lambda\Gamma_{q,\lambda}\bigr)^{1/q}}.   \tag{RC$_{\rm pair}$}
\]
Thus the finite compiler uses the actual residual mass and literal
incidence degrees; it contains no unnamed polynomial or asymptotic loss.
\end{corollary}

\begin{proof}
The moment sandwich gives
\(\max_sQ_\lambda(s)\le\barN_\lambda
(L_\lambda\Gamma_{q,\lambda})^{1/q}\).  Hence
\(\sum_sQ_\lambda(s)^2\le m_\lambda\max_sQ_\lambda(s)\).
Apply \cref{prop:pair-ray-multiplicity} and take the integer floor.
\end{proof}

\begin{proposition}[Curve-degree ray compiler]
\label{prop:curve-degree-ray-compiler}
Let \(V\) be a finite-dimensional space of locator polynomials and, for
\(x\in D\), let
\(H_x=\{[L]\in\Pj(V):L(x)=0\}\) be its \emph{evaluation hyperplane}.
Suppose a residual slope set \(Z\) admits an injective assignment into a
disjoint union of projective integral curves
\[
              Z\hookrightarrow\coprod_{i=1}^sX_i(\F),
\]
where ``integral curve'' means an irreducible reduced one-dimensional
projective variety, together with locator morphisms
\(\psi_i:X_i\to\Pj(V)\).  Put
\[
 \delta_i=\deg\psi_i^*\mathcal O_{\Pj(V)}(1),
\]
the degree on \(X_i\) of the pullback of a projective hyperplane.  For
every \(\gamma\in Z\), if its assigned point is \(y\in X_i(\F)\), require
\(\psi_i(y)=[L_{\gamma,y}]\), where \(L_{\gamma,y}\) is a locator from an
actual residual witness to \(\gamma\).  For each \(i\), put
\[
 G_i=\{x\in D:\psi_i(X_i)\subseteq H_x\},
 \qquad D_i^{\rm mov}=D\setminus G_i.
\]
Assume the locator attached to every chosen representative on \(X_i\)
has at least \(h_i\ge1\) roots in \(D_i^{\rm mov}\).  Then
\[
 \abs Z\le
 \floor{\sum_{i=1}^s
       \frac{\abs{D_i^{\rm mov}}\delta_i}{h_i}}
 \le\floor{n\sum_{i=1}^s\frac{\delta_i}{h_i}}.
 \tag{RC$_{\rm curve}$}
\]
In particular, a cover of subexponential total locator degree gives a
subexponential ray budget when the \(h_i\) have no exponential loss.  The
standard degree-one pencil \([s:t]\mapsto[sA+tB]\) is the special case
\(X=\Pj^1\) and \(\delta=1\).
\end{proposition}

\begin{proof}
Choose one representative point for each slope.  For fixed \(i\) and
\(x\in D_i^{\rm mov}\), the pullback of \(H_x\) is a nonzero effective
divisor of degree \(\delta_i\), and hence contains at most \(\delta_i\)
distinct rational points.  Each chosen point on \(X_i\) supplies at least
\(h_i\) moving-root incidences.  Thus the number assigned to \(X_i\) is
at most \(\abs{D_i^{\rm mov}}\delta_i/h_i\).  Sum over \(i\) and take
the integer floor.
\end{proof}

\begin{corollary}[Automatic exact-support locator-curve compiler]
\label{cor:automatic-locator-curve-ray}
Fix a received line and exact agreement \(a\ge k+1\).  For each slope in
a residual set \(Z\), choose one noncommon witness support
\(S_\gamma\in\binom Da\) and its locator \(Q_{S_\gamma}\).  Suppose the
locator points are covered by
\(\psi_i(X_i(\F))\), where
\(\psi_i:X_i\to\Pj(V)\) are nonconstant morphisms from projective
integral curves.  Put
\[
 \delta_i=\deg\psi_i^*\mathcal O(1),\qquad
 G_i=\{x\in D:\psi_i(X_i)\subseteq H_x\},
 \qquad g_i=\abs{G_i}.
\]
Then
\[
 \abs Z\le
 \#\{i:g_i=a\}+
 \floor{\sum_{i:g_i<a}\frac{(n-g_i)\delta_i}{a-g_i}}. \tag{RC$_{\rm loc}$}
\]
No separate slope-to-curve injectivity hypothesis is required.  A
component with \(g_i=a\) contains at most one selected degree-\(a\)
locator, because its \(a\) common roots determine the monic support
locator uniquely; this accounts for the first term.  Components with
\(g_i>a\) contain no selected degree-\(a\) locator and are discarded.
\end{corollary}

\begin{proof}
Two distinct slopes cannot use the same chosen noncommon support.  Since a
monic squarefree locator determines its root set, their projective locator
points are distinct as well.  Assign each locator to the first covering
component and choose one rational preimage; this is the injective
assignment required by \cref{prop:curve-degree-ray-compiler}.  On
component \(i\), every selected degree-\(a\) locator has the \(g_i\)
common roots and at least \(a-g_i\) roots outside them.  If \(g_i=a\),
there is only one such monic locator.  On the other components apply
\textup{(RC$_{\rm curve}$)} with \(h_i=a-g_i\).
\end{proof}

\begin{corollary}[Automatic exact-support one-pencil compiler]
\label{cor:automatic-one-pencil-ray}
In the setting above, suppose all selected locators lie in the projective
pencil spanned by \(A,B\), and let \(g<a\) be the number of common roots
of \(A,B\) in \(D\).  Then
\[
                         \abs Z\le\floor{(n-g)/(a-g)}. \tag{RC$_{\rm pen}$}
\]
The same argument applies to residual locators of degree \(e\) whose roots
lie in a fixed active domain \(T\subseteq D\), provided the fixed profile
data together with the residual locator determine the exact support.  If
\(N=\abs T\) and \(g_T\) is the number of common roots of \(A,B\) in
\(T\), then
\[
                         \abs Z\le\floor{(N-g_T)/(e-g_T)}.
\]
\end{corollary}

\begin{proof}
The exact-support locator assignment is injective by the preceding proof.
Apply \cref{lem:pencil-payment}, or the curve corollary with
\(X=\Pj^1\) and \(\delta=1\).
\end{proof}

By \cref{thm:syndrome-secant-exact}, the genuinely missing invariant can
be taken to be \(\Theta_t(y_0,y_1)\).  Equivalently one may construct a
higher-dimensional split-pencil \emph{eliminant}, meaning a polynomial in
the slope parameter obtained after eliminating locator and support
variables, whose degree bounds that transverse secant incidence.  By
\cref{lem:independent-union-rays}, only dependent support unions can
support more than \(t+1\) rays.  Those unions are no longer untreated:
\cref{thm:fixed-union-ray,thm:directional-johnson-ray} bound every fixed
union, and \cref{prop:active-fixed-union-payments} gives the exact active
low-nullity ranges.  What remains is an uncontrolled family of many such
unions, a high-nullity chart, or a direction whose rank defect exceeds the
quantitative threshold of \cref{cor:quantitative-direction-defect}; a
higher-dimensional balanced core must still be routed into one of those
certified cells or paid by its own eliminant.

Here an \emph{MDS-circuit stratum} is a chart on which the relevant union
of generalized Vandermonde parity-check columns has a minimal linear
dependence.  Equivalently, deleting any one column restores independence.
The name identifies the only dependence pattern left by the MDS property;
it does not by itself place all witnesses in one fixed circuit or provide
a count for an unrestricted union of such strata.

The next hypothesis packages exactly this missing conversion.  Its number
\(H_\lambda\) is a lower bound on how many certified support pairs represent
one bad ray, while \(J_\lambda\) is an upper bound on how many bad rays one
pair can represent.  The ratio in \textup{(RC)} is what makes the resulting
double count stay at the profile scale.

\begin{hypothesis}[Ray compiler]\label{hyp:ray-compiler}
For every received line and primitive residual profile \(\lambda\), let
\(Z_\lambda^\circ\) be its actual first-match set of unpaid bad slopes and
put
\[
\Pcal_\lambda=\{(A,B)\in(\Omc_\lambda)^2:
                       \Phi_\lambda(A)=\Phi_\lambda(B)\}.
\]
Put \(m_\lambda=\abs{\Omc_\lambda}=\sum_sQ_\lambda(s)\).
The row satisfies \textup{(RC)} on this profile if either it has the direct
bound \(\abs{Z_\lambda^\circ}\le e^{o(n)}(1+\barN_\lambda)\), or there are
an incidence
\(I_\lambda\subseteq Z_\lambda^\circ\times\Pcal_\lambda\) and numbers
\(H_\lambda,J_\lambda\ge1\) such that
\[
 \deg_{I_\lambda}(\gamma)\ge H_\lambda,\qquad
 \deg_{I_\lambda}(A,B)\le J_\lambda,\qquad
 \frac{J_\lambda m_\lambda}{H_\lambda}=e^{o(n)}.
 \tag{RC}\label{eq:ray-compiler}
\]
All bounds are uniform in the received line and realized profile.  The
incidence must arise from the RS witness geometry; its existence is not a
consequence of Q or SP.  In particular, one representing pair per ray does
not supply the required lower degree.
\end{hypothesis}

\begin{proposition}[Q plus RC implies rays]\label{prop:q-implies-sp}
If primitive Q is discharged and \textup{(RC)} holds on every residual
profile, then all primitive ray ledgers are discharged
at total cost \(e^{o(n)}\mathfrak E_n\).
\end{proposition}

\begin{proof}
The boundary prefix fibers are bounded by Q.  For every remaining unpaid
slope choose one noncommon support, as permitted by
the imported slope interface \eqref{eq:imported-slope-elimination}; \textup{(RC)} bounds their distinct slope set by
\(e^{o(n)}\barN_\lambda\).  Sum over the subexponential profile set and use
\eqref{eq:profile-envelope}.
\end{proof}

The proof of \cref{thm:main-smooth-circle} is now complete: the Sidon payment and \cref{thm:primitive-q} pay primitive Q, while \cref{prop:q-implies-sp} and \textup{(RC)} convert the remaining support bounds into distinct-ray bounds.


\part{Finite deployment: Q, one-parameter BC, and the exact ledger}
\section{Finite max-fiber calibration}\label{sec:finite-calibration}
We import the four unsafe endpoints and challenge budgets from CAP25 v13.2 \cite[Cor.~14.3 and Def.~21.1]{CAP25v132}.  Only the adjacent-row quantities consumed by the new upper compiler are printed here.

\begin{table}[H]
\centering
\caption{Imported active frontier rows and the adjacent safe-side calibration.}
\label{tab:active-frontier-rows}
\small
\begin{tabular}{@{}lrrrr@{}}
\toprule
row & $a_0$ & $a_+=a_0+1$ & $w=a_+-K$ & $B^*$\\
\midrule
KoalaBear MCA & 1116047 & 1116048 & 67471 & 274980728111395087\\
KoalaBear list & 1116046 & 1116047 & 67471 & 274980728111395087\\
Mersenne-31 MCA & 1116023 & 1116024 & 67447 & 16777215\\
Mersenne-31 list & 1116022 & 1116023 & 67447 & 16777215\\
\bottomrule
\end{tabular}
\end{table}
Here $n=2^{21}$, $K=2^{20}+1$ on the MCA route, and $K=2^{20}$ on the list route.  Put
\[
 \overline F=\binom n{a_+}|\B|^{-w}.
\]
The exact ceilings of this average are, in table order,
\begin{equation}
57198030366,\quad 65065153468,\quad 1752700,\quad 1993678.\label{eq:finite-average-ceilings}
\end{equation}
Hence the full-budget Q overhead ceilings are respectively
\begin{equation}
4807520.9295,\quad 4226236.5253,\quad 9.5722,\quad 8.4152.\label{eq:finite-q-ratios}
\end{equation}
Any payment made to another cell lowers these values.

The exact moment interface imported from the foundation is
\begin{equation}
 \max_z R(z)\le (|\B|^w\Gamma_r)^{1/r},
 \qquad R(z)=\frac{|\Fib_w(z)|}{\binom n{a_+}|\B|^{-w}}.\label{eq:imported-prefix-moment-sandwich}
\end{equation}
The exact second-moment decomposition is written
\begin{equation}
 \sum_zN_w(z)^2=\binom n{a_+}+\sum_eP_e.\label{eq:imported-second-moment}
\end{equation}

\begin{theorem}[Finite moment criterion for Q]\label{thm:moment-q}
If $\Gamma_r\le e^A$, then the normalized maximum fiber satisfies
\[
 R_Q^{\max}\le\exp\!\left(\frac{A+w\log|\B|}{r}\right).
\]
Thus a finite row with remaining bit margin $\Delta$ can be proved by this route only if
\begin{equation}
 \log_2\Gamma_r+w\log_2|\B|\le r\Delta.\label{eq:finite-moment-criterion}
\end{equation}
\end{theorem}
\begin{proof}
Take logarithms in \eqref{eq:imported-prefix-moment-sandwich}.
\end{proof}

\begin{lemma}[Pencil payment]\label{lem:pencil-payment}
Let $L_{[s:t]}=sA+tB$ be a projective polynomial pencil and let $T$ be a finite set.  If every counted member has at least $h$ roots in $T$ that are not common roots of $A$ and $B$, then at most $\lfloor |T|/h\rfloor$ projective parameters are counted.
\end{lemma}
\begin{proof}
A point $x\in T$ that is not a common root solves the homogeneous linear equation $sA(x)+tB(x)=0$ for exactly one projective parameter.  Count moving-root incidences.
\end{proof}

\begin{definition}[projective locator pencil and its common \(D\)-part]\label{def:projective-locator-pencil}
Let \(D\subseteq\F\) be finite, put \(\Lambda_D=\prod_{x\in D}(X-x)\), and let \(A,B\in\F[X]\) be not both zero.  The projective locator pencil generated by \(A,B\) is
\[
        L_{[s:t]}(X)=sA(X)+tB(X),\qquad [s:t]\in\mathbb P^1(\F).
\]
Its fixed \(D\)-part is
\[
        G_D(A,B)=\gcd(A,B,\Lambda_D),
        \qquad g_D(A,B)=\deg G_D(A,B).
\]
A root of \(G_D(A,B)\) is a common root of every member of the pencil; these are precisely the roots assigned to the common-GCD branch in the first-match ledger.
\end{definition}

\begin{theorem}[moving-root bound for one-parameter split pencils]\label{thm:bc-moving-root}
Let \(D\subseteq\F\) have size \(n\), and let \(L_{[s:t]}=sA+tB\) be a projective locator pencil as in \cref{def:projective-locator-pencil}.  Put \(G=G_D(A,B)\) and \(g=\deg G\).  Let \(\mathcal Z\subseteq\mathbb P^1(\F)\), and suppose that for every \(\lambda=[s:t]\in\mathcal Z\), the polynomial \(L_\lambda/G\) has at least \(h\ge1\) distinct roots in \(D\setminus Z(G)\).  Then
\[
        |\mathcal Z|\le \left\lfloor\frac{n-g}{h}\right\rfloor .
\]
In particular, if every counted member has degree \(j\), has fixed \(D\)-part exactly \(G\), and is \(D\)-split, then
\[
        |\mathcal Z|\le
        \left\lfloor\frac{n-g}{j-g}\right\rfloor,
        \qquad j>g.
\]
The case \(j=g\) is not a moving split-pencil cell: all \(D\)-roots are fixed and the member belongs to the common-GCD branch.
\end{theorem}

\begin{proof}
Count incidences
\[
        I=\{(\lambda,x):\lambda\in\mathcal Z,
        \ x\in D\setminus Z(G),\ L_\lambda(x)=0\}.
\]
For a fixed moving root \(x\in D\setminus Z(G)\), the pair \((A(x),B(x))\) is not \((0,0)\).  Hence the homogeneous linear equation
\[
        sA(x)+tB(x)=0
\]
has exactly one solution \([s:t]\in\mathbb P^1(\F)\).  Thus each moving domain point contributes to at most one incidence, and \(|I|\le n-g\).  On the other hand, every \(\lambda\in\mathcal Z\) contributes at least \(h\) incidences by hypothesis, so \(|I|\ge h|\mathcal Z|\).  This proves the first bound.

If a counted member is \(D\)-split of degree \(j\) and has fixed \(D\)-part exactly \(G\), then exactly \(j-g\) of its \(D\)-roots are moving roots in \(D\setminus Z(G)\).  Applying the first bound with \(h=j-g\) gives the displayed inequality.  If \(j=g\), no moving root remains; the split condition has already been paid by the common-GCD branch rather than by a primitive moving pencil.
\end{proof}

\begin{corollary}[BC is settled on primitive one-parameter locator pencils]\label{cor:bc-one-pencil}
Consider an MCA split-pencil chart at agreement \(m\), and write \(\omega=n-m\).  Suppose that after the first-match tangent, common-support, quotient, extension, degree-drop, and common-GCD branches have been removed, the residual candidate error locators in this chart are represented by a projective locator pencil \(L_\lambda=sA+tB\), with the slope parameter injecting into \(\lambda\), and that every residual bad slope in the chart has at least one degree-\(\omega\) \(D\)-split locator in the pencil.  Then the number of residual bad finite slopes in this chart is at most
\[
        \left\lfloor\frac{n-g}{\omega-g}\right\rfloor,
        \qquad g=g_D(A,B)<\omega .
\]
In the primitive case \(g=0\), this is \(\lfloor n/(n-m)\rfloor\).  For the two active MCA adjacent rows in \cref{tab:active-frontier-rows},
\[
\begin{array}{c|c|c|c}
\text{row} & m=a_+ & \omega=n-m & \lfloor n/\omega\rfloor\\ \hline
\text{KoalaBear MCA} & 1116048 & 981104 & 2\\
\text{Mersenne-31 MCA} & 1116024 & 981128 & 2
\end{array}
\]
so every primitive one-parameter BC pencil contributes at most two finite bad slopes.
\end{corollary}

\begin{proof}
The first-match removals ensure that the chart is counted only through moving roots of the displayed pencil.  A residual bad slope supplies at least one degree-\(\omega\) split locator, and the slope-to-pencil parameter is injective by hypothesis; hence the slope count is bounded by the number of split parameters.  Apply \cref{thm:bc-moving-root} with \(j=\omega\).  The displayed row values use \(n=2^{21}\), \(a_+=1116048\) for KoalaBear MCA, and \(a_+=1116024\) for Mersenne-31 MCA.
\end{proof}


\begin{definition}[saturated codeword rays and support census]\label{def:saturated-rays}
Let \(C=\RS[\F,D,K]\), let \(U:D\to\F\), and put
\[
        S_c(U)=\{x\in D:U(x)=c(x)\},\qquad s_c(U)=|S_c(U)|
\]
for \(c\in C\).  The saturated ray set at agreement \(m\) is
\[
        \operatorname{Ray}(U;m)=\{c\in C:s_c(U)\ge m\}.
\]
Let \(\omega=n-m\), and define the support-locator census
\[
\begin{aligned}
\Cen(U;m)=\#\{(W,N)\in \F[X]^2:\;& W\text{ monic},\ \deg W=\omega,\ W\mid \ell_D,\\
        & WU=N\text{ on }D,\quad W\mid N,\quad \deg(N/W)<K\}.
\end{aligned}
\]
Equivalently, \(W=\ell_{D\setminus T}\) for an agreement support \(T\subseteq D\) of size \(m\).
\end{definition}

\begin{theorem}[exact saturation identity]\label{thm:saturation}
With \(C=\RS[\F,D,K]\) and \(\Cen(U;m)\) as in \cref{def:saturated-rays}, for every word \(U:D\to\F\),
\[
        \Cen(U;m)=\sum_{c\in C}\binom{s_c(U)}m .
\]
Moreover, the support locators lying above a fixed codeword ray \(c\) are exactly
\[
        (W,N)=H\,(G_c,G_cc),
        \qquad
        G_c=\ell_{D\setminus S_c(U)},
        \qquad
        H\mid\ell_{S_c(U)},
        \qquad
        \deg H=s_c(U)-m .
\]
Consequently the fiber size of the forgetful map from support locators to saturated codeword rays is \(\binom{s_c(U)}m\) above \(c\).
\end{theorem}

\begin{proof}
Let \((W,N)\) be counted by \(\Cen(U;m)\).  Since \(W\) is monic, divides \(\ell_D\), and has degree \(n-m\), there is a unique \(m\)-set \(T\subseteq D\) such that \(W=\ell_{D\setminus T}\).  Since \(W\mid N\) and \(\deg(N/W)<K\), write \(N=Wc\) for a unique codeword \(c\in C\).  On every \(x\in T\), the value \(W(x)\) is nonzero; the identity \(WU=N=Wc\) on \(D\) therefore gives \(U(x)=c(x)\).  Thus \(T\subseteq S_c(U)\).

Conversely, if \(c\in C\) and \(T\subseteq S_c(U)\) has \(|T|=m\), then \(W=\ell_{D\setminus T}\) and \(N=Wc\) satisfy all defining conditions of \(\Cen(U;m)\).  For fixed \(c\), the number of such supports is \(\binom{s_c(U)}m\), and summing over \(c\) gives the first identity.

For the fiber statement, write \(G_c=\ell_{D\setminus S_c(U)}\).  If \(T\subseteq S_c(U)\), then
\[
        \ell_{D\setminus T}
        =\ell_{D\setminus S_c(U)}\ell_{S_c(U)\setminus T}
        =G_cH,
\]
where \(H=\ell_{S_c(U)\setminus T}\), hence \(H\mid\ell_{S_c(U)}\) and \(\deg H=s_c(U)-m\).  The converse is the same construction reversed, because every monic divisor \(H\) of \(\ell_{S_c(U)}\) of degree \(s_c(U)-m\) is the locator of \(S_c(U)\setminus T\) for a unique \(m\)-set \(T\subseteq S_c(U)\).
\end{proof}

\begin{corollary}[raw support BC is not the MCA object]\label{cor:raw-bc-fails}
In the setting of \cref{thm:saturation}, if \(U\) agrees with a single codeword \(c\) on \(n-d\) positions and with no other codeword on \(m\) or more positions, then
\[
        |\operatorname{Ray}(U;m)|=1,
        \qquad
        \Cen(U;m)=\binom{n-d}{m} .
\]
Thus an exponentially large raw support-locator family may represent one MCA ray.  A split-pencil proof for MCA must therefore count saturated primitive line-rays after quotient, prefix, tangent, and common-support branches have been paid; a raw support census is a list/moment object, not the final MCA numerator.
\end{corollary}

\begin{proof}
The ray statement is the hypothesis.  The support-census identity is \cref{thm:saturation} with one nonzero summand \(\binom{s_c(U)}m=\binom{n-d}m\).  The final sentence is just the same equality interpreted in MCA terms: MCA counts a bad slope once, whereas \(\Cen(U;m)\) counts all size-\(m\) supports below the same saturated agreement ray.
\end{proof}

\begin{definition}[line rays]\label{def:line-rays}
In the setting of \cref{def:saturated-rays}, for an affine line \(U_Z=u+Zv\) and a set of finite slopes \(E\subseteq\F\), define
\[
        \LineRay_E(u,v;m)
        =
        \{(z,c)\in E\times C:\ |\{x\in D:u(x)+zv(x)=c(x)\}|\ge m\}.
\]
Write \(s_{z,c}=|\{x\in D:u(x)+zv(x)=c(x)\}|\).
\end{definition}

\begin{proposition}[line-ray saturation identity]\label{prop:line-ray-saturation}
With notation as in \cref{def:line-rays}, for every \(E\subseteq\F\),
\[
        \sum_{z\in E}\Cen(U_z;m)
        =
        \sum_{(z,c)\in\LineRay_E(u,v;m)}\binom{s_{z,c}}m .
\]
Consequently
\[
        |\{z\in E:\exists c\ (z,c)\in\LineRay_E(u,v;m)\}|
        \le
        |\LineRay_E(u,v;m)|
        \le
        \sum_{z\in E}\Cen(U_z;m).
\]
The two inequalities are the exact deduplication losses between raw support locators, saturated codeword rays, and slope numerators.
\end{proposition}

\begin{proof}
Apply \cref{thm:saturation} separately to the word \(U_z=u+zv\) for each \(z\in E\), and then sum over \(z\).  The first inequality forgets the codeword coordinate of a line ray, and the second follows because each line ray has \(s_{z,c}\ge m\), hence contributes at least one support to the right-hand side of the identity.  \Cref{cor:raw-bc-fails} shows that the second inequality can be exponentially loose for a single slope and a single codeword ray, so an MCA upper ledger must either count slopes directly or control this saturation factor.
\end{proof}

\begin{theorem}[Q discharges the shift-pair ledger]\label{thm:q-implies-sp}
Let \(N_w(z)=|\Fib_w(z)|\), let
\[
        \overline F=\binom nm |\B|^{-w},
\]
and let \(P_e\) denote the full off-diagonal shift-pair stratum in the imported decomposition \eqref{eq:imported-second-moment}, namely
\[
        P_e=\sum_{\substack{R\subseteq D\\ |R|=m-e}}\operatorname{sp}_w(e;D\setminus R).
\]
If Q holds in max-fiber form with constant \(\kappa\),
\[
        \max_z N_w(z)\le \kappa\overline F,
\]
then necessarily \(\kappa\overline F\ge1\) whenever \(\binom Dm\) is nonempty, and
\[
        \sum_{e=w+1}^{\min(m,n-m)}P_e
        \le
        \left(\kappa-\frac{1}{\overline F}\right)\binom nm\,\overline F .
\]
Every quotient-removed or primitive shift-pair subledger satisfies the same upper bound.  Hence SP is not an independent residual input in any closing theorem that already proves Q as a max-fiber theorem.
\end{theorem}

\begin{proof}
The imported exact second-moment identity \eqref{eq:imported-second-moment} says
\[
        \sum_z N_w(z)^2
        =\binom nm+
        \sum_{e=w+1}^{\min(m,n-m)}P_e .
\]
On the other hand,
\[
        \sum_zN_w(z)^2
        \le (\max_zN_w(z))\sum_zN_w(z)
        \le \kappa\overline F\binom nm,
\]
because \(\sum_zN_w(z)=\binom nm\).  Since at least one fiber is nonempty when \(\binom Dm\ne\varnothing\), the hypothesis also implies \(\kappa\overline F\ge1\), so the displayed right-hand side is nonnegative in the nonvacuous case.  Subtracting the diagonal term \(\binom nm\) gives the displayed inequality.  Any primitive or quotient-removed ledger is a sub-sum of the full off-diagonal ledger, so it is bounded a fortiori.
\end{proof}


\begin{theorem}[SP: unconditional primitive shift-pair bound]\label{thm:sp-proper}
Use the notation of the imported exact second-moment ledger.  For \(w+1\le e\le\min(m,n-m)\), put
\[
        T_e(n,m)=\binom n{m-e}\binom{n-m+e}{e}\binom{n-m}{e}.
\]
Then the total off-diagonal contribution at distance \(e\) satisfies
\[
        P_e^{\rm quot}+P_e^{\rm prim}\le T_e(n,m).
\]
If \(D\subseteq\B\), then the top stratum \(e=w+1\) also satisfies the sharper constant-shift bound
\[
        P_{w+1}^{\rm quot}+P_{w+1}^{\rm prim}
        \le
        (|\B|-1)\binom n{m-w-1}\binom{n-m+w+1}{w+1}.
\]
Consequently
\[
        \Gamma_2
        \le
        \frac{1}{\overline F}
        +\frac{\sum_{e=w+1}^{\min(m,n-m)}T_e(n,m)}{\binom{n}{m}\overline F},
\]
with the option of replacing the \(e=w+1\) term by the sharper displayed top-stratum bound.
\end{theorem}

\begin{proof}
Fix \(e\).  In the imported exact second-moment decomposition \eqref{eq:imported-second-moment}, an off-diagonal pair is encoded by a common part \(R\subseteq D\) of size \(m-e\), then by two disjoint degree-\(e\) root sets inside \(D\setminus R\).  There are \(\binom n{m-e}\) choices for \(R\).  Once \(R\) is fixed, the ambient set \(D\setminus R\) has size \(n-m+e\); the first root set has at most \(\binom{n-m+e}{e}\) choices and the second, disjoint from it, has at most \(\binom{n-m}{e}\) choices.  This ignores the depth equation and the quotient/primitive first-match assignment, so it is an upper bound for \(P_e^{\rm quot}+P_e^{\rm prim}\).

When \(e=w+1\), the imported decomposition \eqref{eq:imported-second-moment} shows that every top-stratum shift pair has the form \(A-B=c\in\F^\times\).  If both polynomials split over \(D\subseteq\B\), then both have coefficients in \(\B\), hence \(c\in\B^\times\).  For each fixed \(R\), choose \(A\) by choosing its \(e\)-element root set in \(D\setminus R\), and choose \(c\in\B^\times\); the polynomial \(B=A-c\) is then determined.  Requiring \(B\) to split over \(D\setminus R\), to be disjoint from \(A\), and then assigning the pair to either the quotient or primitive ledger can only reduce the count.  This gives the top-stratum bound.  The displayed \(\Gamma_2\) inequality is the imported exact second-moment ledger with the total quotient-plus-primitive off-diagonal contribution in every stratum replaced by the just-proved upper bound.
\end{proof}


\begin{proposition}[size of the proper Q bound at the active adjacent rows]\label{prop:proper-q-gap}
At the four active adjacent agreements, the imported Q packing estimate \cite[Secs.~19--21]{CAP25v132} gives the following one-term upper bounds for the remaining normalized prefix overhead:
\[
\begin{array}{c|ccc}
\text{row} & w & \log_2 R_Q^{\rm pack}\text{ allowed by the bound} & \text{spare margin}\\
\hline
\text{KoalaBear MCA} & 67471 & 1660789.024 & 22.1969\\
\text{KoalaBear list} & 67471 & 1660789.017 & 22.0109\\
\text{Mersenne-31 MCA} & 67447 & 1660926.120 & 3.2589\\
\text{Mersenne-31 list} & 67447 & 1660926.114 & 3.0730 .
\end{array}
\]
Thus the unconditional packing proof is a real theorem but is far too weak to supply the adjacent safe ledger.
\end{proposition}

\begin{proof}
For each row, \(n=2097152\), \(K=k+1=1048577\) on the MCA route, \(K=k=1048576\) on the list route, and \(w=a_+-K\).  The base-field orders are \(2^{31}-2^{24}+1\) for KoalaBear and \(2^{31}-1\) for Mersenne-31.  The imported Q packing inequality \cite[Secs.~19--21]{CAP25v132} gives
\[
        \log_2 R_Q^{\rm pack}
        \le
        w\log_2|\B|-\log_2\binom{a_+}{\lfloor w/2\rfloor}
             -\log_2\binom{n-a_+}{\lfloor w/2\rfloor}.
\]
Substituting the exact field orders and
\[
        a_+\in\{1116048,1116047,1116024,1116023\}
\]
gives the displayed values.  For these four rows the Johnson-ball summand
\[
        J_i=\binom{a_+}{i}\binom{n-a_+}{i}
\]
is increasing on \(0\le i\le t\).  Indeed
\[
        \frac{J_i}{J_{i-1}}
        =
        \frac{(a_+-i+1)(n-a_+-i+1)}{i^2},
\]
and at \(i=t\) this ratio is already \(>900\) in all four rows, while the ratio is larger for smaller \(i\).  Hence \(V_t(n,a_+)\le(t+1)J_t\), so replacing the one-term denominator by the full ball can lower the displayed logarithmic upper bound by less than \(\log_2(t+1)<16\) bits.  This is negligible compared with the gap between the displayed million-bit overhead and the row margins.
\end{proof}


\begin{definition}[row-sharp Q atom bound]\label{def:q-row-atom}
Fix a deployed adjacent row, let \(a_+\) be its adjacent agreement, let \(K\) be the list dimension used on that route, and put \(w=a_+-K\).  For a first-match residual prefix-boundary family \(\mathcal P_Q(z)\subseteq\Fib_w(z)\), define
\[
        R_Q^{\max}
        =
        |\B|^w\max_z\frac{|\mathcal P_Q(z)|}{\binom n{a_+}}.
\]
The row-sharp Q atom bound with bit margin \(\Delta_Q\) is
\[
        R_Q^{\max}\le 2^{\Delta_Q}.
\]
Equivalently, if the entire row budget is hypothetically assigned to Q, the sufficient integer form is
\[
        \max_z|\mathcal P_Q(z)|\le B^* .
\]
Any non-Q cell paid in the same first-match ledger only decreases the available \(\Delta_Q\).
\end{definition}

\begin{proposition}[exact finite Q atom target at the four adjacent rows]\label{prop:q-exact-target}
At the active adjacent agreements, the full-budget Q atom target is as follows:
\[
\begin{array}{c|r|r|r|r}
\text{row} & w & \left\lceil \binom n{a_+}|\B|^{-w}\right\rceil & B^* & B^*/\left\lceil \binom n{a_+}|\B|^{-w}\right\rceil\\
\hline
\text{KoalaBear MCA} & 67471 & 57198030366 & 274980728111395087 & 4807520.9295\\
\text{KoalaBear list} & 67471 & 65065153468 & 274980728111395087 & 4226236.5253\\
\text{Mersenne-31 MCA} & 67447 & 1752700 & 16777215 & 9.5722\\
\text{Mersenne-31 list} & 67447 & 1993678 & 16777215 & 8.4152
\end{array}
\]
In bits, these are exactly the spare margins already printed in the row table:
\[
        22.1969,
        \quad 22.0109,
        \quad 3.2589,
        \quad 3.0730 .
\]
Thus the binding Q problem is the Mersenne-31 list row: after all first-match payments, the remaining primitive prefix-boundary atom must be at most about \(8.42\) times its average size, and less if any other cell consumes budget.
\end{proposition}

\begin{proof}
For each row, substitute \(n=2^{21}\), \(k=2^{20}\), the displayed adjacent agreement \(a_+\), and \(K=k+1\) on the MCA route and \(K=k\) on the list route.  The average prefix fiber is \(\binom n{a_+}|\B|^{-w}\), and the integer lower-route replay uses its ceiling.  The budgets are
\[
        B^*_{\rm KB}=\left\lfloor\frac{(2^{31}-2^{24}+1)^6}{2^{128}}\right\rfloor,
        \qquad
        B^*_{\rm M31}=\left\lfloor\frac{(2^{31}-1)^4}{2^{100}}\right\rfloor .
\]
The displayed ratios are the exact integer quotients evaluated as decimals.  Taking base-two logarithms of \(B^*/(\binom n{a_+}|\B|^{-w})\) gives the four printed bit margins.
\end{proof}


\begin{proposition}[Exact fixed-union payments at the active rows]
\label{prop:active-fixed-union-payments}
Let \(t=n-a_+\).

\textup{(a) Fixed-word Q and list charts.}
Let \(K=k+1\) on an MCA route and \(K=k\) on a list route, put
\(R_Q=n-K\), and let \(\mathcal A\) be a subfamily of one boundary-Q
prefix fiber, or of one arbitrary-word list, whose error supports are all
contained in a fixed union \(U\) of size \(R_Q+\nu\).  Then
\cref{thm:fixed-union-list-johnson} gives a full-budget payment for every
integer nullity in the ranges below:
\[
\begin{array}{c|r|r|r|r}
\text{row} & R_Q & R_Q-t & 0\le\nu\le &
 \max\limits_{\nu\text{ in range}}|\mathcal A|\\ \hline
\text{KoalaBear MCA}  &1048575&67471&4983&94008\\
\text{KoalaBear list} &1048576&67471&4983&94632\\
\text{Mersenne-31 MCA}&1048575&67447&4980&3555853\\
\text{Mersenne-31 list}&1048576&67447&4980&4735771
\end{array}
\tag{FU-Q/list}\label{eq:active-fixed-union-list}
\]
Every number in the last column is below the corresponding challenge
budget \(B^*\) of \cref{prop:q-exact-target}.  In particular, a global Q
counterexample at the binding Mersenne-31 list row cannot be contained in
one support union of nullity at most \(4980\).

\textup{(b) Fixed-line MCA charts.}
For the actual MCA code, \(R=n-k=1048576\).  The unconditional
multiplicity compiler \cref{thm:fixed-union-ray} pays the following fixed
unions within the full row budget:
\[
\begin{array}{c|r|r|r}
\text{row}&0<\nu\le&
 \max\limits_{\nu\text{ in range}}|Z|&B^*\\ \hline
\text{KoalaBear MCA}&2&2847909263951&274980728111395087\\
\text{Mersenne-31 MCA}&1&8150817&16777215
\end{array}
\tag{FU-MCA}\label{eq:active-fixed-union-mca}
\]
If the direction is rank-regular on its union, so that
\(d_U(y_1)=R\), the directional compiler
\cref{cor:rank-regular-fixed-union-ray} pays the much larger ranges
\[
\begin{array}{c|r|r}
\text{row}&0\le\nu\le&
 \max\limits_{\nu\text{ in range}}|Z|\\ \hline
\text{KoalaBear MCA}&4982&94008\\
\text{Mersenne-31 MCA}&4979&3555853
\end{array}
\tag{FU-MCA-reg}\label{eq:active-rank-regular-mca}
\]
The direction-defect threshold of
\cref{cor:quantitative-direction-defect} pays more than the rank-regular
case.  With the full row budget, selected exact thresholds are
\[
\begin{array}{c|rrrrrrr}
\text{KoalaBear }\nu&0&1000&4000&4500&4979&4980&4982\\ \hline
J_{B^*}(\nu)&4341&3466&853&418&3&2&0
\end{array}
\tag{FU-KB-def}\label{eq:active-kb-direction-defect}
\]
and
\[
\begin{array}{c|rrrrr}
\text{Mersenne-31 }\nu&0&1000&4000&4500&4979\\ \hline
J_{B^*}(\nu)&4338&3463&849&415&0
\end{array}
\tag{FU-M31-def}\label{eq:active-m31-direction-defect}
\]
Thus an unpaid fixed-union MCA chart is either above the displayed
nullity range or has defect \(j>J_{B^*}(\nu)\); it need not be declared
exceptional merely because \(j>0\).  At the Mersenne row, two disjoint
one-circuit charts cost at most
\(2\cdot8150817=16301634<16777215\), leaving \(475581\) units of the
full budget.
\end{proposition}

\begin{proof}
For a boundary-Q fiber, the imported exact witness-list dictionary
\cite[Secs.~18--21]{CAP25v132} associates to every support a distinct
codeword of \(\RS[\F,D,K]\) agreeing with one received word on that
support.  Hence a subfamily whose error supports lie in \(U\) is a subset
of the affine syndrome fiber \(\mathcal L_U(y,t)\) in
\cref{thm:fixed-union-list-johnson}.  The same statement is tautological
for an arbitrary-word list.  Substituting \(N=R_Q+\nu\) and
\(h=N-t\) into \eqref{eq:fixed-union-list-johnson}, and evaluating every
integer in the displayed ranges, gives \eqref{eq:active-fixed-union-list}.
At the next nullity the Johnson denominator is nonpositive in every row:
\(\nu=4984\) for KoalaBear and \(\nu=4981\) for Mersenne-31.

For MCA, substitute \(R=1048576\) and the two deployed values of \(t\)
into \eqref{eq:fixed-union-ray}.  The largest full-budget nullities are
\(2\) and \(1\), giving \eqref{eq:active-fixed-union-mca}.  Under rank
regularity, substitute the same row data into
\eqref{eq:rank-regular-fixed-union-ray}; its denominator is positive
through \(\nu=4982\) and \(4979\), respectively, and the displayed maxima
occur at the final integer in each range.  Finally substitute the row
budgets into \eqref{eq:direction-defect-threshold} to obtain
\cref{eq:active-kb-direction-defect,eq:active-m31-direction-defect}.
All comparisons are exact integer comparisons and are reproduced by the
release verifier.
\end{proof}

\begin{proposition}[Exact recursive affine-span payments at the active rows]
\label{prop:active-affine-span-payments}
In the boundary-Q witness list, or in an arbitrary-word list for the same
route, suppose the associated codewords lie in one affine subspace of
\(\RS[\F,D,K]\) of dimension \(s\).  The recursive compiler
\cref{thm:affine-span-list} gives the following full-budget payments:
\[
\begin{array}{c|r|r|r}
\text{row}&0\le s\le&
 \max\limits_{s\text{ in range}}|\mathcal L_{\mathcal A}|&
 \text{compiler value at the next }s\\ \hline
\text{KoalaBear MCA}&14&47876303026096432&743896698428332665\\
\text{KoalaBear list}&14&47876942243704279&743907339955962922\\
\text{Mersenne-31 MCA}&6&14115447&219426634\\
\text{Mersenne-31 list}&6&14115528&219428099
\end{array}
\tag{AF-Q/list}\label{eq:active-affine-span-list}
\]
The middle column is below the corresponding full challenge budget in
every row.  The final column only records that this compiler no longer
fits at dimension \(15\) or \(7\); it is not a lower bound on an actual
list.  Consequently a Q counterexample surviving the full-budget payments
must have affine codeword-span dimension at least \(15\) on a KoalaBear
row and at least \(7\) on a Mersenne-31 row.
\end{proposition}

\begin{proof}
Apply \cref{thm:affine-span-list} with \(m=a_+\), with \(K=k+1\) on an
MCA boundary route and \(K=k\) on a list route.  Evaluate the single
rational expression in \eqref{eq:affine-span-list} by exact integer
products.  The value is increasing over each displayed range, so the
maximum occurs at its last dimension.  Comparison with the exact budgets
in \cref{prop:q-exact-target} gives the claim; the next-dimension values
mark the limit of the recursive compiler.
\end{proof}

\begin{remark}[Retraction of the former active affine-span MCA payments]
The former KoalaBear and Mersenne-31 MCA affine-span payments used the
refuted ordered-basis denominator of
\cref{prop:affine-span-mca-counterexample}.  Their falling- and
rising-product evaluations were arithmetically correct evaluations of that
expression, but they are not upper bounds for MCA slope families.  They
therefore imply no affine-rank lower wall for a surviving MCA chart.  The
ordinary affine-span list payments immediately above are unaffected.  The
next proposition gives a distinct conditional wall from the repaired
support-transverse theorem, with its deleted stratum and local margin stated
explicitly.
\end{remark}

\begin{proposition}[Conditional post-deletion affine-error router on KoalaBear]
\label{prop:kb-post-deletion-affine-error-router}
Use the KoalaBear parameters
\[
 (n,k,m,w,B_*)=
 (2097152,1048576,1116048,67472,274980728111395087).
\]
On one received line, let \(Z_{\rm bad}\) be the finite support-wise
MCA-bad slopes and define the intrinsic near-rational stratum
\[
 N=\{\gamma\in Z_{\rm bad}:\min_{c\in C}\wt(r_\gamma-c)\le w\},
 \qquad Z=Z_{\rm bad}\setminus N.
\]
Suppose \(|N|\le2w=134944\), as supplied under the printed KoalaBear
hypotheses by the separately pinned two-anchor near-rational theorem.  Its
guards hold here because $w\ge1$ and
$3w=202416\le1048576=n-k$.
For every \(\gamma\in Z\), choose one exact same-support-noncontained
\(m\)-record, and let \(a\) be the affine rank of the resulting error
words.  If \(a\le9\), then
\[
 |Z_{\rm bad}|\le
 110390969172173096+134944
 =110390969172308040<B_*,
\]
with slack \(164589758939087047\).

For \(|Z|\ge2\), apply \cref{thm:mca-error-rank-gauge} and recompute the
margin \(\theta\) of \eqref{eq:support-local-mca-theta} in the translated
explanation direction space.  At error ranks \(a=10,11,12\), the line is
also paid whenever respectively
\[
 \theta\ge13,\qquad \theta\ge388,
 \qquad \theta\ge12050.
\]
Consequently an over-budget line forces \(a\ge10\) for every selected
post-deletion family.  At ranks \(10,11,12\), every such gauge emits an
actual selected support and a codeword direction agreeing with the original
\(r_1\) on all but at most respectively \(12,387,12049\) support
coordinates.
\end{proposition}

\begin{proof}
The exact-record selection is legitimate.  Given a bad witness
\((S,c)\) with \(|S|\ge m\), suppose every \(m\)-subset of \(S\) were
pair-contained.  Adjacent \(m\)-subsets overlap in
\(m-1\ge K\) coordinates, so Reed--Solomon uniqueness makes their
explaining pairs identical.  Connectedness of the \(m\)-subset Johnson
graph would then glue one pair explaining all of \(S\), contradicting
badness.  Hence some exact \(m\)-subwitness is pair-noncontained.

If \(|Z|\le1\), the displayed bound is immediate.  Otherwise
\cref{thm:mca-error-rank-gauge} gives a translated explanation family of
affine dimension \(s=a-1\).  For \(2\le a\le9\), apply
\cref{thm:support-transverse-affine-span-mca} with \(\theta\ge1\).  The
largest cap occurs at \(s=8\) and is \(110390969172173096\).

If \(a=1\), all translated explanations coincide.  Same-support
noncontainment forces each selected support to contain a coordinate where
the translated line direction is nonzero; at such a coordinate the
agreement equation determines at most one finite slope.  Hence \(|Z|\le n\),
which is smaller than the displayed cap.  The case \(a=0\) has at most one
slope, since the injective map in the proof of
\cref{thm:mca-error-rank-gauge} has nonzero slope projection.

For \(a=10,11,12\), substitute \(s=9,10,11\) in
\eqref{eq:support-local-proper-subspace-mca}.  Exact integer arithmetic
shows that the least paying margins after the \(2w\) add-back are
\(13,388,12050\).  At the preceding integers the totals are respectively
\[
 285894151677963688,\qquad
 275500176064828033,\qquad
 274992929018868606,
\]
all strictly above \(B_*\).  Negating the paying condition gives the three
same-record direction-exception terminals.  Explicitly, if $b_0\in C$
is the gauge codeword from \cref{thm:mca-error-rank-gauge} and failure of
the relevant margin produces $b_1$ in the translated explanation
direction space, then $r'_1=r_1-b_0$ agrees with $b_1$ off the printed
exception set.  Thus the inverse gauge transports the terminal to the
original record as $r_1=b_0+b_1$ on the identical support outside that
same set.
\end{proof}


\begin{theorem}[Sub-square-root common-support interleaving collapse]
\label{thm:subsquare-interleaving-collapse}
Let \(C\leq\F^D\) be an \(\F\)-linear code, put \(q=|\F|\), and fix an
agreement threshold \(A\).  Let \(L\) be the largest ordinary list size
at agreement \(A\), over all received words.  For \(\mu\geq1\), let
\(L_\mu\) be the corresponding largest list of \(\mu\)-tuples in
\(C^\mu\) agreeing with one received \(\mu\)-tuple on one common support
of size at least \(A\).  If \(1\leq L<q\), then
\[
 L\leq L_\mu\leq
 \left\lfloor\frac{L(q-1)}{q-L}\right\rfloor .
\]
In particular, if \(L^2<q\), then \(L_\mu=L\) for every \(\mu\geq1\).
\end{theorem}

\begin{proof}
The case \(\mu=1\) is the definition of \(L\), so assume \(\mu\geq2\).
The diagonal embedding gives \(L_\mu\geq L\).  For the reverse bound fix
a common-support list \(X\subseteq C^\mu\), \(N=|X|\), and for
\(\alpha=(\alpha_2,\ldots,\alpha_\mu)\in\F^{\mu-1}\) project by
\[
 \Phi_\alpha(c_1,\ldots,c_\mu)
 =c_1+\sum_{i=2}^\mu\alpha_i c_i.
\]
The projected set is an ordinary list at the identical agreement
threshold, so it has at most \(L\) images.  Cauchy--Schwarz gives at least
\((N^2/L-N)/2\) unordered colliding pairs for every \(\alpha\).  A fixed
distinct pair collides for at most \(q^{\mu-2}\) choices of \(\alpha\): if
some difference after the first component is nonzero, fix all other
coefficients and solve for that coefficient; if all those differences
vanish, the nonzero first difference prevents collision.  Summing over
all \(q^{\mu-1}\) projections yields
\[
 q^{\mu-1}(N^2/L-N)\leq N(N-1)q^{\mu-2},
\]
and hence \(N\leq L(q-1)/(q-L)\).  Finally
\[
 \frac{L(q-1)}{q-L}-L=\frac{L(L-1)}{q-L}<1
\]
when \(L^2<q\); integrality and the diagonal lower bound give equality.
\end{proof}

\begin{theorem}[Margin/interleaving split for a gauged MCA family]
\label{thm:mca-margin-interleaving-split}
Let \(C=\RS[\F,D,K]\) have length \(n\), put \(m=K+w\leq n\), and
let \(Z\) be a finite family of distinct slopes with exact
same-support-pair-noncontained \(m\)-records.  Suppose their explanations
lie in one affine flat \(c_0+C'\subseteq C\), where
\(s=\dim C'\).  For each \(\gamma\in Z\), define
\[
 \theta_\gamma=
 \min\left\{w+1,
 \min_{b\in C'}|\{x\in S_\gamma:r_1(x)\ne b(x)\}|\right\}.
\]
Fix an integer \(T\) with \(2\leq T\leq w\), put
\[
 A=m-T+1>K,
 \qquad
 M_s(T)=
 \left\lfloor
 \frac{\binom{n-K+s}{s}}{\binom{w-T+1+s}{s}}
 \right\rfloor,
\]
and assume \(M_s(T)^2<|\F|\).  Define
\[
 H_0(T)=\left\lfloor\frac nT\right\rfloor
\]
and, for \(1\leq j\leq s\),
\[
 H_j(T)=\left\lfloor\max\left\{
 \frac{n^{\underline{j+1}}}
 {mT(w+1)^{\overline{j-1}}},
 \frac{(n-K+j)^{\underline{j+1}}}
 {T(w+1)^{\overline j}}
 \right\}\right\rfloor .
\]
Then
\begin{equation}
 |Z|\leq
 \max_{0\leq j\leq s}H_j(T)+(n-A)M_s(T).
 \label{eq:mca-margin-interleaving-split}
\end{equation}
\end{theorem}

\begin{proof}
Partition \(Z=Z_{\rm hi}\sqcup Z_{\rm lo}\) according as
\(\theta_\gamma\geq T\) or \(\theta_\gamma\leq T-1\).
If the affine explanation span of \(Z_{\rm hi}\) has direction dimension
\(j\geq1\), its direction space is contained in \(C'\); minimizing over
that smaller space cannot decrease the margins.  Thus
\cref{thm:support-transverse-affine-span-mca} gives
\(|Z_{\rm hi}|\leq H_j(T)\).  If \(j=0\), every explanation equals one
word.  Each support contains at least \(T\) coordinates where the line
direction is nonzero, and at each coordinate the agreement equation
determines at most one slope.  Incidence counting gives
\(|Z_{\rm hi}|\leq H_0(T)\).

For every \(\gamma\in Z_{\rm lo}\), choose \(b_\gamma\in C'\) with at
most \(T-1\) mismatches on \(S_\gamma\), and put
\(a_\gamma=h_\gamma-\gamma b_\gamma\in c_0+C'\).  The pair
\((a_\gamma,b_\gamma)\) agrees with \((r_0,r_1)\) on at least \(A\)
common coordinates.  Translating the first component by \(c_0\) places
all such pairs in the two-fold common-support interleaving of \(C'\).
The ordinary affine-span list compiler
\cref{thm:affine-span-list} bounds each projected ordinary list by
\(M_s(T)\).  The hypothesis \(M_s(T)^2<|\F|\) and
\cref{thm:subsquare-interleaving-collapse} therefore bound the number of
distinct pairs by \(M_s(T)\).

For one fixed pair \((a,b)\), let
\(H_{a,b}=\{x:r_0(x)=a(x),\ r_1(x)=b(x)\}\); then
\(|H_{a,b}|\geq A\).  Pair noncontainment gives every owning slope an
actual \(x\in S_\gamma\setminus H_{a,b}\).  At such an \(x\),
\[
 \gamma=-\frac{r_0(x)-a(x)}{r_1(x)-b(x)}.
\]
Choosing one exception coordinate per slope is injective, so the pair
owns at most \(n-|H_{a,b}|\leq n-A\) slopes.  Adding the disjoint high and
low bounds proves \eqref{eq:mca-margin-interleaving-split}.
\end{proof}

\begin{proposition}[KoalaBear affine-error rank-ten payment]
\label{prop:kb-affine-error-rank-ten-payment}
Retain the intrinsic near-rational deletion and exact post-deletion record
selection of \cref{prop:kb-post-deletion-affine-error-router}.  If the
selected error words have affine rank \(a=10\), and the received line is
over the deployed sextic field
\[
 |\F|=2130706433^6,
\]
then
\[
 |Z_{\rm bad}|\leq
 2w+5143522968716559+56727790457914040
 =61871313426765543<B_*.
\]
The slack is \(213109414684629544\).  Thus the direct KoalaBear
post-near affine-error-rank-ten branch is paid.  The same printed split
does not pay affine error ranks \(11\) or \(12\).
\end{proposition}

\begin{proof}
The reversible gauge \cref{thm:mca-error-rank-gauge} converts error rank
\(a=10\) into an explanation flat of dimension \(s=9\), preserving the
slopes, errors, exact supports, and same-support pair noncontainment.  Use
\cref{thm:mca-margin-interleaving-split} with \(T=667\).  Exact integer
arithmetic gives
\[
 A=1115382,\qquad
 M_9(667)=57781140652,\qquad
 M_9(667)^2<2130706433^6=|\F|,
\]
\[
 \max_{0\leq j\leq9}H_j(667)=5143522968716559,\qquad
 (n-A)M_9(667)=56727790457914040.
\]
Adding the disjoint near charge \(2w=134944\) proves the displayed bound.
An exhaustive exact scan of \(2\leq T\leq w\) shows that \(T=16\) is the
first paying threshold and \(T=667\) is the unique minimizer for \(s=9\).
For \(s=10\), corresponding to error rank \(11\), the minimum of this
formula is \(1040506078215897711>B_*\), attained at \(T=876\); the
minima at \(s=11,12\) are larger still.
\end{proof}

\begin{theorem}[Nonuniform support-margin resource]
\label{thm:mca-nonuniform-support-margin}
Retain the hypotheses and notation of
\cref{thm:support-transverse-affine-span-mca}, but define the individual
support margins
\[
 \theta_\gamma=\min\left\{w+1,
  \min_{b\in C'}|\{x\in S_\gamma:r_1(x)\ne b(x)\}|\right\}.
\]
Then
\begin{equation}
 \sum_{\gamma\in Z}\theta_\gamma\le C_s,
 \qquad
 C_s=\left\lfloor\max\left\{
 \frac{n^{\underline{s+1}}}{m(w+1)^{\overline{s-1}}},
 \frac{(n-K+s)^{\underline{s+1}}}{(w+1)^{\overline s}}
 \right\}\right\rfloor .
 \label{eq:mca-nonuniform-support-margin}
\end{equation}
\end{theorem}

\begin{proof}
Use the ordered-normal-basis proof of
\cref{thm:support-transverse-affine-span-mca}, before replacing the
individual final-extension counts by their minimum.  With the proof's
global zero-normal parameters \((z,g)\), the parameter point belonging to
\(\gamma\) owns at least
\[
 (m-g)\theta_\gamma(w+1)^{\overline{s-1}}
\]
ordered independent coordinate-normal tuples.  One independent
\((s+1)\)-tuple determines at most one parameter point.  Summing first gives
\[
 (m-g)(w+1)^{\overline{s-1}}
 \sum_\gamma\theta_\gamma
 \le (n-z)^{\underline{s+1}}.
\]
Here \((z,g)\) are global: a zero normal's affine incidence equation is either
identically satisfied on the parameter flat or never satisfied.  Now
\(0\le g\le z\le K-s\).  For fixed \(z\), the ratio is largest at \(g=z\),
and the successive-ratio calculation in the cited proof puts its maximum at
\(z=0\) or \(z=K-s\).  These are the two terms in
\eqref{eq:mca-nonuniform-support-margin}.
\end{proof}

\begin{theorem}[Weighted minimizing-pair core route cut]
\label{thm:mca-weighted-pair-core-route-cut}
Retain the hypotheses of
\cref{thm:mca-nonuniform-support-margin}.  Fix
\(1\le\tau\le w-1\), and put
\[
 Q_s(\tau)=\left\lfloor
  \frac{\binom{n-K+s}{s}}{\binom{w-\tau+s}{s}}
 \right\rfloor ,
 \qquad Q_s(\tau)^2<|\F|.
\]
For every \(\gamma\) with \(\theta_\gamma\le\tau\), fix one minimizing
\(b_\gamma\in C'\), and put
\[
 a_\gamma=c_\gamma-\gamma b_\gamma,
 \qquad \lambda_\gamma=\tau+1-\theta_\gamma,
\]
and group records by their ordered pair \((a_\gamma,b_\gamma)\).  If
\(|Z|\ge E\) and \((\tau+1)E>C_s\), some fixed pair has total weight at
least
\begin{equation}
 L_s(\tau,E)=\left\lceil
 \frac{\max\{0,(\tau+1)E-C_s\}}{Q_s(\tau)}
 \right\rceil .
 \label{eq:mca-weighted-pair-load}
\end{equation}
For a fixed pair \((a,b)\), write
\[
 H_{a,b}=\{x\in D:r_0(x)=a(x),\ r_1(x)=b(x)\},
 \qquad \delta=\max\{1,m-|H_{a,b}|\}.
\]
The pair owns at most
\begin{equation}
 c_\delta=\left\lfloor\frac{n-m+\delta}{\delta}\right\rfloor
 \label{eq:mca-pair-core-multiplicity}
\end{equation}
records, and its weight is at most
\(c_\delta(\tau+1-\delta)\).
\end{theorem}

\begin{proof}
The positive-part identity and
\eqref{eq:mca-nonuniform-support-margin} give
\[
 \sum_{\theta_\gamma\le\tau}
   (\tau+1-\theta_\gamma)
 =\sum_{\gamma\in Z}(\tau+1-\theta_\gamma)_+
 \ge (\tau+1)|Z|-C_s.
\]
The selected pair \((a_\gamma,b_\gamma)\) agrees with the received pair on
at least \(m-\tau>K\) coordinates.  The ordinary affine-span list theorem
and the sub-square interleaving collapse
\cref{thm:subsquare-interleaving-collapse} bound the number of distinct
ordered pairs by \(Q_s(\tau)\).  Pigeonholing proves
\eqref{eq:mca-weighted-pair-load}; there is no bipartite factor two because
the interleaving theorem counts complete ordered pairs.

Fix one pair and an owning record.  On \(S_\gamma\), the identity
\[
 r_0-a=-\gamma(r_1-b)
\]
shows that
\(S_\gamma\cap H_{a,b}=\{x\in S_\gamma:r_1(x)=b(x)\}\).  Since the margin
cap is inactive, \(\theta_\gamma=|S_\gamma\setminus H_{a,b}|\ge\delta\).
For distinct finite slopes these exception sets are disjoint: at a point
outside \(H_{a,b}\), the displayed equation determines \(\gamma\) uniquely.
Their union lies in \(D\setminus H_{a,b}\).  Division by \(\delta\), with
the harmless overestimate \(n-|H_{a,b}|\le n-m+1\) when
\(|H_{a,b}|\ge m\), proves \eqref{eq:mca-pair-core-multiplicity}.  Each
edge weight is at most \(\tau+1-\delta\), proving the final assertion.
\end{proof}

\begin{proposition}[KoalaBear rank-eleven pair-core terminal]
\label{prop:kb-rank-eleven-pair-core-terminal}
Retain the intrinsic near deletion, exact record selection, and reversible
gauge of \cref{prop:kb-post-deletion-affine-error-router}.  Suppose the
selected error words have affine rank \(11\), so the gauged explanations
lie in an affine ten-flat over the deployed field
\(|\F|=2130706433^6\).  If the received line has more than \(B_*\) bad
slopes, then for each row below there exists a fixed actual minimizing pair
satisfying that row.  The pairs supplied by the two independently optimized
cutoffs need not coincide.
\[
\begin{array}{c|r|r|r}
 \tau&\text{pair weight}&\text{owned slopes}&\delta\le\ \\ \hline
 6486&743449148&114624&8\\
 1795&360132809&200632&4
\end{array}
\tag{PC11}\label{eq:kb-rank-eleven-pair-core-terminal}
\]
Parallel records are allowed; the owned-slope column is not a
distinct-neighbor degree.

Moreover, the complete one-cutoff upper bound obtained from the printed
high-margin theorem, all cumulative interleaved pair caps, and the exact
core-deficiency multiplicities has unique minimum
\[
 813929118931913384
 =134944+5401690553097387+808527428378681053
 >B_*
\]
at cutoff \(J=19737\).  Thus those inputs do not pay rank eleven.  The
remaining theorem must couple different fixed pairs or route the dense
parallel pair cores to a chronology-correct owner.
\end{proposition}

\begin{proof}
An over-budget line leaves at least
\(E=B_*-2w+1\) post-near slopes.  Apply
\cref{thm:mca-weighted-pair-core-route-cut} with \(s=10\).  Exact integer
optimization of \eqref{eq:mca-weighted-pair-load} gives the two rows of
\eqref{eq:kb-rank-eleven-pair-core-terminal}; comparison with
\(c_\delta(\tau+1-\delta)\) gives the displayed deficiency walls.

For the final assertion, first fix one minimizing \(b_\gamma\), and hence
one pair type \((a_\gamma,b_\gamma)\), for every selected record, once and
for all cutoffs.  Let \(G(t)\) count distinct selected pair types
whose deficiency is at most \(t\).  Such a pair has common agreement at
least \(m-t\), so interleaving gives \(G(t)\le Q_{10}(t)\).  Since
\(c_\delta\) decreases, summation by parts bounds the low records through
\(J\) by
\[
 Q_{10}(1)c_1+
 \sum_{\delta=2}^{J}
   \bigl(Q_{10}(\delta)-Q_{10}(\delta-1)\bigr)c_\delta.
\]
The records with margin at least \(J+1\) are bounded by the maximum of the
support-transverse theorem over affine subranks \(0\le j\le10\).  Scanning
all \(1\le J\le w-1\), subject to the exact sub-square guard, gives the
unique displayed minimum.  The release certificate recomputes every term,
the adjacent cutoffs, and hostile mutations.
\end{proof}

\begin{remark}[Scope of the rank-eleven route cut]
The proposition moves no active-v4 atom and does not close KoalaBear.  A
small constant-code star shows that the fixed-pair multiplicity and parallel
edges are genuine under the local hypotheses, so no simple-graph expansion
or smaller per-pair multiplier may be inferred.  The missing result is an
actual-record cross-pair compatibility or owner theorem, not another scalar
threshold optimization.
\end{remark}

\begin{proposition}[Two-anchor reserve repricing]
\label{prop:mca-two-anchor-reserve-repricing}
Suppose the near-rational first-match set has size at most $2w$, the
residual exception set has size at most $31$, and a large rational owner
has size $g>2m-K$.  If the active S/A/E assembly is used after those two
sets, then its contained-owner input must be required at the target
\begin{equation}
 B_{\rm owner}^{(2w)}(g)
 \le B_*-(2w+31)-(n-g).
 \label{eq:mca-two-anchor-repriced-owner}
\end{equation}
With this target the four charges sum exactly to $B_*$.  On the deployed
KoalaBear row, the targets at $g=2m-K+1$ and $g=n$ are respectively
\[
 274980728110346481,
 \qquad
 274980728111260112.
\]
On the Mersenne-31 arithmetic stress row they are respectively
\[
 15728609,
 \qquad
 16642288.
\]
The corresponding full-owner targets are at least $4807520$ and $9$
times the proved exact average ceilings.
\end{proposition}

\begin{proof}
Write $N$ for the near-rational set and $E$ for the residual exception
set.  Since $2w$ equals $134944$ on KoalaBear and $134896$ on the
Mersenne-31 stress row, neither bound can be absorbed into the $31$-slope
bound without an additional containment theorem.  Charging the two sets in
first-match order leaves at most $n-g$ crossing slopes, so the remaining
target is forced by
\[
 2w+31+(n-g)+\bigl(B_*-(2w+31)-(n-g)\bigr)=B_*.
\]
Substitution at the two endpoints gives the displayed integers.  Exact
division of the full-owner targets by the average ceilings gives quotients
$4807520$ and $9$, with positive remainders.  This proves arithmetic
viability, not the contained-owner bound in
\eqref{eq:mca-two-anchor-repriced-owner}.
\end{proof}

\begin{theorem}[Degree-guarded effective-shift witness adapter]
\label{thm:mca-degree-guarded-shift-adapter}
Let $D$ be an $n$-point Reed--Solomon evaluation domain, let
$k<m\le n$, put $\omega=n-m$, and for a word $U$ define
\[
 \mathcal M_U=\{(W,N):W(x)U(x)=N(x)\text{ for every }x\in D\}.
\]
For a nonzero vector define
\[
 s_k(W,N)=\max\{\deg W,\deg N-(k-1)\},\qquad
 s_{k+1}(W,N)=\max\{\deg W,\deg N-k\}.
\]
Then
\[
 s_{k+1}(W,N)\le s_k(W,N)\le s_{k+1}(W,N)+1,
\]
and the two lattice minima obey the same inequalities.

Suppose $W$ is the monic split squarefree locator of $D\setminus T$,
where $|T|=m$, and $N=Wc$.  Under the effective envelope
$s_{k+1}(W,N)\le\omega$, the following are equivalent:
\begin{equation}
 \deg c<k,
 \qquad \deg N\le\omega+k-1,
 \qquad s_k(W,N)\le\omega.
 \label{eq:mca-effective-shift-degree-guard}
\end{equation}
They hold exactly when the lattice vector reconstructs a degree-below-$k$
explanation of $U$ on the identical support $T$.  Conversely every such
explanation has a unique guarded representation.

For a received line $U=u+\gamma v$, interpolate $u|_T$ and $v|_T$ to
their unique degree-below-$m$ polynomials.  The pair is simultaneously
code-explained on $T$ if and only if both interpolants have degree below
$k$.  Thus a guarded explanation together with failure of this pair test
is an actual support-wise MCA-bad witness on the same line, slope, and
support.
\end{theorem}

\begin{proof}
Put $a=\deg W$ and $b=\deg N-k$, with the usual convention
$\deg 0=-\infty$.  Then
\[
 s_{k+1}=\max\{a,b\},\qquad s_k=\max\{a,b+1\},
\]
which proves the one-step comparison, including after taking minima.
Since $\deg W=\omega$ and $N=Wc$, the effective cap permits
$\deg c\le k$, while each condition in
\eqref{eq:mca-effective-shift-degree-guard} removes exactly the degree-$k$
coefficient.

At every point of $T$, $W$ is nonzero, so the lattice identity gives
$U=c$.  Conversely an explanation $h$ on $T$ gives
\[
 W=\prod_{x\in D\setminus T}(X-x),\qquad N=Wh.
\]
The monic locator and $h$ are unique because $m\ge k$.  Finally, unique
degree-below-$m$ interpolation on $T$ turns simultaneous pair
explainability into the two stated degree tests.
\end{proof}

\begin{proposition}[Transport mutation and typed pole-line control]
\label{prop:mca-effective-shift-controls}
On the deployed KoalaBear row let $D=\langle\zeta\rangle$ have order
$2^{21}$, let $E$ be the first $67473$ subgroup powers, and let $S$
be the following $1116048$ powers.

First, for
\[
 u=\mathbf 1_E,\qquad v=X^k,\qquad\gamma=0,
\]
the identical slope/support record is support-wise MCA-bad for
$RS[\F,D,k]$ but not for $RS[\F,D,k+1]$.  Hence silent dimension
transport does not preserve badness or first-owner semantics.

Second, over the deployed sextic extension, let $\alpha$ be a root of
$X^6+X+6$ and set
\[
 v(x)=-\frac1{x-\alpha},\qquad
 u(x)=\mathbf 1_E(x)+\frac{\alpha}{x-\alpha},\qquad
 \gamma=\alpha.
\]
This is an actual support-wise MCA-bad witness on $S$, it passes
\eqref{eq:mca-effective-shift-degree-guard} with quotient zero, and its
lattice minimum is exactly $67473$ under both shifts.  Its Q, BC, and
new-owner fields remain unassigned.
\end{proposition}

\begin{proof}
For the first record, zero explains $u$ on $S$.  A degree-below-$k$
polynomial agreeing with $X^k$ on $m>k$ points would make a nonzero
degree-$k$ polynomial have more than $k$ roots.  At dimension $k+1$,
the pair $(0,X^k)$ is a simultaneous explanation on the same support.

For the pole line, direct cancellation gives
$u+\alpha v=\mathbf 1_E$, again explained by zero on $S$.  If $g$ of
degree below $k+1$ agreed with $v$ there, then
$(X-\alpha)g+1$ would have $m>k+1$ roots and degree at most $k+1$,
which is impossible.  The complement locator with numerator zero therefore
passes the guarded adapter.

The prefix locator gives shifted degree $67473$.  If a nonzero lattice
vector had shifted degree at most $67472$, its numerator would vanish on
all $2097152-67473=2029679$ points outside $E$, while its degree would be
at most $1116048$ under either shift.  Thus the numerator is zero.  The
denominator would then vanish at all $67473$ points of $E$ despite
degree at most $67472$, another contradiction.  No owner inference occurs
in this argument.
\end{proof}

\begin{remark}[Acceptance contract for the rank-eleven owner route]
\label{rem:mca-dense-core-owner-acceptance-contract}
The pair-core terminal of
\cref{prop:kb-rank-eleven-pair-core-terminal} forces a possibly fixed pair
of deficiency at most $4$ owning at least $200632$ slopes.  A
chronology-correct owner theorem for that terminal may use the preceding
results as an interface: if it reuses S/A/E it must prove the repriced bound
\eqref{eq:mca-two-anchor-repriced-owner}; if it computes certificates in the
effective $k+1$ envelope it must impose
\eqref{eq:mca-effective-shift-degree-guard} and exact-support pair
noncontainment; and it must preserve the actual line, support, slope, and
owner chronology.  The preceding results prove witness soundness and
arithmetic consistency only.  They do not construct the owner, pay rank
eleven, move an active-v4 atom, or close KoalaBear.
\end{remark}


\begin{proposition}[Exact rank-flat and common-zero reductions at the active rows]
\label{prop:active-rank-flat-common-zero}
Consider the first affine dimensions not paid by
\cref{prop:active-affine-span-payments}: \(s=15\) on KoalaBear and
\(s=7\) on Mersenne-31.  If such an affine Q/list chart exceeds the full
row budget, \cref{cor:rank-flat-concentration} forces a proper direction
subcode with the following coarse support ceiling:
\[
\begin{array}{c|r|r|r|r}
\text{row}&s&\Theta^{\rm coarse}_{s,B^*}&d_j\le&d_j-(n-K)\le\\ \hline
\text{KoalaBear MCA}&15&152252&1133356&84781\\
\text{KoalaBear list}&15&152252&1133357&84781\\
\text{Mersenne-31 MCA}&7&232423&1213551&164976\\
\text{Mersenne-31 list}&7&232423&1213552&164976
\end{array}
\tag{RF-active}\label{eq:active-rank-flat}
\]
Thus an over-budget first-unpaid affine chart is not generalized-weight
uniform: it contains a proper direction subcode supported on a union much
smaller than the full domain.

There is also an exact common-zero refinement beyond the former
fixed-union boundary.  Write the common-zero set of the affine direction
space as \(z=K-\nu\), and let \(b\) be the number of those coordinates on
which the common affine value disagrees with the received word.  At the
first nullity not paid by \cref{prop:active-fixed-union-payments}, the
common-zero Johnson compiler gives
\[
\begin{array}{c|r|r|r|c}
\text{row}&\nu&\min b&\text{bound}&\text{residual }b\\ \hline
\text{KoalaBear MCA}&4984&2&537111&b\le1\\
\text{KoalaBear list}&4984&2&558125&b\le1\\
\text{Mersenne-31 MCA}&4981&7&590532&b\le6\\
\text{Mersenne-31 list}&4981&7&616024&b\le6
\end{array}
\tag{CZ-active}\label{eq:active-common-zero}
\]
At nullity \(5000\), the corresponding paid mismatch thresholds are
\(b\ge102\) on the two KoalaBear rows and \(b\ge126\) on the two
Mersenne-31 rows.
\end{proposition}

\begin{proof}
For \eqref{eq:active-rank-flat}, use the coarse form of
\cref{cor:rank-flat-concentration}.  The integer roots are
\[
 \left\lfloor
 \left(
 \frac{n^{\underline{15}}}{(B^*_{\rm KB}+1)(67471+15)}
 \right)^{1/14}
 \right\rfloor=152252
\]
and
\[
 \left\lfloor
 \left(
 \frac{n^{\underline 7}}{(B^*_{\rm M31}+1)(67447+7)}
 \right)^{1/6}
 \right\rfloor=232423.
\]
Adding the exact radii \(t=n-a_+\) gives the support and nullity ceilings
in the table.

For \eqref{eq:active-common-zero}, substitute the row data into
\eqref{eq:common-zero-nu} and choose the least nonnegative integer \(b\)
for which the denominator is positive and the integer floor is at most the
row budget.  All displayed comparisons, including the nullity-
\(5000\) thresholds, are exact and are reproduced by the verifier.
\end{proof}

\begin{proposition}[Exact codimension-one recursion at the active rows]
\label{prop:active-codim-one}
Consider a first-unpaid affine Q/list chart: dimension \(15\) on a
KoalaBear route and dimension \(7\) on a Mersenne-31 route.  Suppose the
low-support direction subcode forced by
\cref{prop:active-rank-flat-common-zero} has support closure of codimension
one in the chart direction space.  Write \(j=14\) or \(j=6\) for the
dimension of that support-saturated hyperplane, \(d\) for its support
size, \(e\) for the support increment to the whole chart, and \(b_0\) for
the common mismatch count outside the whole support.  Then the exact
full-budget payments are
\[
\begin{array}{c|r|r|r|r|r}
\text{row}&j&d-(n-K)\le&e\text{ paid from}&b_0\text{ paid from}
 &\text{largest residual nullity}\\ \hline
\text{KoalaBear MCA}&14&84781&8565&9813&93345\\
\text{KoalaBear list}&14&84781&8565&9813&93345\\
\text{Mersenne-31 MCA}&6&164976&73266&40074&238241\\
\text{Mersenne-31 list}&6&164976&73267&40074&238242
\end{array}
\tag{CR-active}\label{eq:active-codim-one-ranges}
\]
Here ``\(e\) paid from'' means that every chart with at least that support
increment is below the full row budget, and similarly for \(b_0\).  The
largest residual nullity is \(d+e-(n-K)\) after both payments have failed.
More explicitly, every surviving codimension-one branch lies in
\begin{align}
 e&\le8564, & b_0&\le9812
 &&\text{on either KoalaBear route},
 \tag{CR-KB}\label{eq:active-codim-one-kb}\\
 e&\le73265, & b_0&\le40073
 &&\text{on the Mersenne-31 MCA route},
 \tag{CR-MCA}\label{eq:active-codim-one-mca}\\
 e&\le73266, & b_0&\le40073
 &&\text{on the Mersenne-31 list route}.
 \tag{CR-list}\label{eq:active-codim-one-list}
\end{align}
The corresponding certified maxima at the first paid integers are
\begingroup\scriptsize
\[
\begin{array}{c|r|r|r}
\text{row}&C_e&C_{b_0}&B^*\\ \hline
\text{KoalaBear MCA}&274959863953460332&274954310643125130&274980728111395087\\
\text{KoalaBear list}&274963447501329540&274957707108580869&274980728111395087\\
\text{Mersenne-31 MCA}&16777189&16776155&16777215\\
\text{Mersenne-31 list}&16777174&16776251&16777215
\end{array}
\tag{CR-budget}\label{eq:active-codim-one-budgets}
\]
\endgroup
where $C_e$ is the closed bound at the first paid support increment and
$C_{b_0}$ is the all-extender bound at the first paid mismatch count.
The immediately preceding integer fails for the displayed MDS-soft
certificate in every row; this records the sharp integer transition of
this compiler, not a lower bound on an actual list.
\end{proposition}

\begin{proof}
The support closure has the same support as the direction subcode supplied
by \cref{prop:active-rank-flat-common-zero}; hence its support size is in
the exact ranges
\[
\begin{array}{c|c}
\text{row}&d\text{ range}\\ \hline
\text{KoalaBear MCA}&1048589\le d\le1133356\\
\text{KoalaBear list}&1048590\le d\le1133357\\
\text{Mersenne-31 MCA}&1048581\le d\le1213551\\
\text{Mersenne-31 list}&1048582\le d\le1213552.
\end{array}
\]
For the MDS-soft profile \(\Pi\) of
\cref{cor:codim-one-mds-soft}, the largest initial factor is
\(A=d-t+b_0\) and \(Q=w+1+b_0\).  Uniformly over the displayed ranges,
\[
 (j-1)Q-A\ge
 \begin{cases}
 13\cdot67472-152252=724884,&\text{KoalaBear},\\
 5\cdot67448-232423=104817,&\text{Mersenne-31}.
 \end{cases}
\]
The final term grows by \((j-2)b_0\), so
\cref{lem:codim-one-profile-interpolation} proves profile interpolation
for every \(e\) and \(b_0\) in the active ranges.

The closed bound \eqref{eq:codim-one-closed} decreases with \(e\) and
with \(b_0\), while the all-extender bound
\eqref{eq:codim-one-all-extenders-mds} decreases with \(b_0\).  It
therefore suffices to use \(b_0=0\) for the support-increment payment and
to test the displayed mismatch threshold for the all-extender payment.
Exact rational comparison over every integer \(d\) in the four finite
ranges gives the maxima in \eqref{eq:active-codim-one-budgets}.  At the
preceding support increments the maxima are respectively
\[
274997419229066860,
\quad275001003273987385,
\quad16777303,
\quad16777288,
\]
and at the preceding mismatch counts they are
\[
275005816112956831,
\quad275009213214651183,
\quad16777247,
\quad16777343.
\]
Each exceeds its row budget.  Finally combine the support ceiling with
\(e\le e_{\rm paid}-1\) to obtain the residual nullities in
\eqref{eq:active-codim-one-ranges}.  The release verifier performs all
comparisons with exact integers.
\end{proof}




\begin{proposition}[Active higher-codimension flag payment]
\label{prop:active-support-flag}
Consider a first-unpaid affine Q/list chart: dimension $15$ on a
KoalaBear route and dimension $7$ on a Mersenne-31 route.  If the chart is
over the full row budget, choose a generalized-weight minimizer supplied
by \cref{cor:rank-flat-concentration}.  By
\cref{lem:weight-minimizer-saturated,lem:support-saturated-flag}, its
support closure admits the exact coset-free flag certificate of
\cref{thm:support-flag-recursion} in every codimension.

For support-closure codimension at least two, the coarse all-extender
certificate already pays every chart with
\begin{equation}
 b_0\ge70230
 \quad\text{on either KoalaBear route},
 \qquad
 b_0\ge108962
 \quad\text{on either Mersenne-31 route}.
 \label{eq:active-flag-mismatch}
\end{equation}
At the first paid integers, the worst exact bounds are
\[
\begin{array}{c|r|r|r}
\text{row}&b_0&\text{certified maximum}&B^*\\ \hline
\text{KoalaBear MCA}&70230&274974976450914526&274980728111395087\\
\text{KoalaBear list}&70230&274975238687487221&274980728111395087\\
\text{Mersenne-31 MCA}&108962&16776934&16777215\\
\text{Mersenne-31 list}&108962&16776950&16777215
\end{array}
\tag{HF-active}\label{eq:active-flag-budgets}
\]
The preceding mismatch integer fails the same uniform certificate, with
values respectively
\[
 275004931457252515,
 \quad275005193722392530,
 \quad16777600,
 \quad16777616.
\]
These are transitions of the all-extender flag certificate, not lower
bounds on actual lists.  Below \eqref{eq:active-flag-mismatch}, the exact
binary-pattern linear program \eqref{eq:flag-dual-bound} and the closed
tensor certificate \eqref{eq:flag-tensor-bound} remain available for the
actual support-layer profile.
\end{proposition}

\begin{proof}
The rank-flat conclusion supplies a generalized-weight minimizer of some
dimension $j<s$.  It is support-saturated by
\cref{lem:weight-minimizer-saturated}, and therefore has a
support-saturated flag to the whole chart by
\cref{lem:support-saturated-flag}.  For codimension at least two, use
\eqref{eq:flag-all-extender}, replace the flag-support product by
$(n-j)^{\underline{s-j}}$, and use the exact support ceilings in
\eqref{eq:active-rank-flat}.  The resulting bound is
\[
 \frac{d^{\underline j}(n-j)^{\underline{s-j}}}
 {(d-t+b_0)
  \prod_{i=1}^{j-1}(w+i+b_0)
  (w+1+b_0)^{s-j}}.
\]
For fixed $j$ its maximum over the allowed support interval occurs at an
endpoint: the ratio of the $d$-dependent factor at $d+1$ and $d$ changes
sign at most once, since the numerator of the comparison is
\[
 (j-1)d-j(t-b_0)+j-1.
\]
Checking the two endpoints for $1\le j\le13$ on KoalaBear and
$1\le j\le5$ on Mersenne-31 gives the exact transitions displayed above.
The independent verifier performs these integer comparisons without
floating point.
\end{proof}


\begin{proposition}[Exact paving-basis MCA payments at the active rows]
\label{prop:active-paving-payments}
Let \(j=R-d_U(y_1)\) be the direction defect of a separated fixed-union
MCA chart.  With the full row budget, the paving threshold of
\cref{cor:paving-threshold} is
\[
\begin{array}{c|rrrrr}
\text{KoalaBear }\nu&0\text{--}9&10&11&12&13\\ \hline
J_{B^*}^{\rm pav}(\nu)&67471&67468&67423&66715&55724
\end{array}
\tag{PAV-KB}\label{eq:active-paving-kb}
\]
and
\[
\begin{array}{c|rrrrr}
\text{Mersenne-31 }\nu&0\text{--}1&2&3&4&5\\ \hline
J_{B^*}^{\rm pav}(\nu)&67447&67432&67213&63797&10700.
\end{array}
\tag{PAV-M31}\label{eq:active-paving-m31}
\]
Consequently every separated KoalaBear direction defect is paid through
nullity \(9\).  At Mersenne nullity \(2\), only the top fifteen possible
defects \(67433\le j\le67447\) remain outside this compiler.  At every
nullity the certified range is the maximum of the paving threshold and the
directional-Johnson threshold in
\cref{eq:active-kb-direction-defect,eq:active-m31-direction-defect}.
\end{proposition}

\begin{proof}
Substitute \(R=1048576\), the deployed radii, and the exact row budgets
into \eqref{eq:paving-threshold}.  At the last paid defects in the nontrivial
columns, the paving bounds are respectively
\[
\begin{array}{c|r|r}
\nu&J_{B^*}^{\rm pav}(\nu)&\text{bound at }J_{B^*}^{\rm pav}(\nu)\\ \hline
10&67468&215264294199835875\\
11&67423&273052400452931367\\
12&66715&274631447525088415\\
13&55724&274966795861149003
\end{array}
\]
and
\[
\begin{array}{c|r|r}
\nu&J_{B^*}^{\rm pav}(\nu)&\text{bound at }J_{B^*}^{\rm pav}(\nu)\\ \hline
2&67432&15838845\\
3&67213&16764415\\
4&63797&16774561\\
5&10700&16776956
\end{array}
\]
The next defect exceeds the budget in every displayed nonterminal column.
All calculations use exact binomial integers.
\end{proof}

\begin{proposition}[moment-order floor for a finite Q proof]\label{prop:q-moment-order-floor}
Let
\[
        \tau=\frac{\sum_z|\mathcal P_Q(z)|}{\binom nm}\in(0,1]
\]
be the mass of the pruned residual relative to the full support slice.  A proof based only on the moment criterion of \cref{thm:moment-q} must satisfy the necessary condition
\[
        r\bigl(\Delta_Q-\log_2\tau\bigr)
        \ge w\log_2|\B| .
\]
In the full-mass case \(\tau=1\), this reduces to
\[
        r\ge\left\lceil\frac{w\log_2|\B|}{\Delta_Q}\right\rceil .
\]
For the active rows the corresponding full-mass floors are
\[
\begin{array}{c|r|r|r}
\text{row} & w & \Delta_Q\text{ bits} & \text{minimum }r\\
\hline
\text{KoalaBear MCA} & 67471 & 22.1969 & 94196\\
\text{KoalaBear list} & 67471 & 22.0109 & 94991\\
\text{Mersenne-31 MCA} & 67447 & 3.2589 & 641593\\
\text{Mersenne-31 list} & 67447 & 3.0730 & 680397
\end{array}
\]
Thus fixed moments cannot prove the finite adjacent rows in the full-mass model.  After first-match pruning, the displayed row floors remain applicable only when \(\tau=1\); otherwise the residual mass must be carried explicitly, or one must use an off-diagonal or falling-factorial moment or isolate sparse-heavy fibers as a separate cell.
\end{proposition}

\begin{proof}
Normalize the residual masses by the full slice and put
\(\mu(z)=|\mathcal P_Q(z)|/\binom nm\), so \(\sum_z\mu(z)=\tau\).  Convexity over at most \(|\B|^w\) prefix values gives
\[
        |\B|^{w(r-1)}\sum_z\mu(z)^r\ge\tau^r.
\]
In the notation of \cref{thm:moment-q}, this is \(\Gamma_r\ge\tau^r\).  The finite criterion requires
\[
        \log_2\Gamma_r+w\log_2|\B|\le r\Delta_Q .
\]
Substituting the lower bound and rearranging gives the stated necessary condition.  Setting \(\tau=1\) yields the table.  The numerical entries use the exact row prime in \(\log_2|\B|\).
\end{proof}

\begin{proposition}[the BC moving-root theorem does not pay boundary Q]\label{prop:bc-not-q}
In the boundary-Q profile of the imported boundary-Q dictionary \cite[Secs.~18--21]{CAP25v132}, the unknown divisor is
\[
        Q_z+A,
        \qquad
        \deg A\le \omega-w-1,
        \qquad
        \omega=n-a_+ .
\]
Thus the coefficient parameter space has dimension \(\omega-w\).  At the active rows this dimension is
\[
\begin{array}{c|r|r|r}
\text{row} & \omega=n-a_+ & w & \omega-w\\
\hline
\text{KoalaBear MCA} & 981104 & 67471 & 913633\\
\text{KoalaBear list} & 981105 & 67471 & 913634\\
\text{Mersenne-31 MCA} & 981128 & 67447 & 913681\\
\text{Mersenne-31 list} & 981129 & 67447 & 913682
\end{array}
\]
Therefore \cref{thm:bc-moving-root,cor:bc-one-pencil} cannot be applied directly to boundary Q: those theorems count split parameters on a projective one-parameter locator pencil, while Q is a high-dimensional affine divisor slice.  To route Q through BC one would need an additional theorem decomposing this affine slice into paid one-parameter pencils with a first-match bound on the number of relevant pencils.  No such theorem is present in this manuscript.
\end{proposition}

\begin{proof}
The imported boundary-Q dictionary \cite[Secs.~18--21]{CAP25v132} states exactly that a prefix fiber is counted by the divisibility condition \(Q_z+A\mid X^n-\beta\) with \(\deg A\le\omega-w-1\).  A polynomial of degree at most \(\omega-w-1\) has \(\omega-w\) free coefficients.  Substituting the four values of \(a_+\), \(K\), and \(w=a_+-K\) gives the table.  Since \(\omega-w\) is larger than nine hundred thousand in all four cases, the family is not a one-parameter pencil.  The moving-root theorem is still valid for every line inside this affine space, but a line-by-line decomposition without a bound on the number of lines gives no row budget.
\end{proof}



\part{Order-\texorpdfstring{$32$}{32} line rank and rational-atom classification}
\label{part:rank-atoms}
The rest of the paper changes the MCA completion mechanism.  Instead of adding many local cell budgets, it seeks a maximum-type classification: every large bad-slope set is either one global affine codeword line, one coherent rational atom, or one primitive spread family.  This part proves the exact rank dictionary, the dyadic quotient branch, low-degree and low-complexity atom extraction, and atom coherence.

\section{Rank-theoretic setup}
Let $D\subseteq\F$ consist of $n$ distinct evaluation points and let
\[
 C=\RS[\F,D,k]
 =\{(h(x))_{x\in D}:h\in\F[X],\ \deg h<k\}.
\]
Fix distinct finite slopes $\gamma_1,\ldots,\gamma_L$ and a received line $r_\gamma=r_0+\gamma r_1$.  Suppose $h_i\in\F[X]_{<k}$ agrees with $r_{\gamma_i}$ on a support $S_i$ of size $m$; shrinking supports to exact size $m$ does not affect the rank calculation.  Put
\[
 I_x=\{i:x\in S_i\}.
\]
The tuple is globally affine if $h_i=A+\gamma_iB$ for some $A,B\in\F[X]_{<k}$.

\section{The rank-two reduced intersection matrix}

\subsection{Local relation spaces}

For $I\subseteq[L]$, define
\[
 \cK_\gamma(I)
 :=\left\{c\in\F^L:
 \supp(c)\subseteq I,\quad
 \sum_{i=1}^Lc_i=0,\quad
 \sum_{i=1}^Lc_i\gamma_i=0
 \right\}.
\]
Since the slopes are distinct,
\[
 \dim\cK_\gamma(I)=(|I|-2)_+.
\]
Let $v_k(x)=\vander{x}{k}\in\F^k$.

\begin{definition}[rank-two reduced intersection matrix]\label{def:r2rim}
Choose any basis $\mathcal B_x$ of $\cK_\gamma(I_x)$ for every $x\in D$.  The matrix
\[
 \cM^{(2)}(D,\mathcal S,\gamma)
\]
has $Lk$ columns, grouped into $L$ blocks of length $k$, and one row
\[
 c\otimes v_k(x)
\]
for every $x\in D$ and $c\in\mathcal B_x$.
\end{definition}

Changing the local bases left-multiplies the corresponding row blocks by invertible matrices, so the kernel and rank are intrinsic.

\begin{theorem}[exact kernel dictionary]\label{thm:kernel-dictionary}
For $h_i\in\F[X]_{<k}$, let $\mathbf h$ be the concatenated coefficient vector.  Then
\[
 \mathbf h\in\ker\cM^{(2)}(D,\mathcal S,\gamma)
\]
if and only if, for every $x\in D$, there exist scalars $a_x,b_x\in\F$ such that
\[
 h_i(x)=a_x+\gamma_i b_x
 \qquad(i\in I_x).
\]
In particular, every family of codeword explanations of one received line belongs to the kernel.
\end{theorem}

\begin{proof}
For fixed $x$, the local equations are
\[
 \sum_{i\in I_x}c_i h_i(x)=0
 \qquad(c\in\cK_\gamma(I_x)).
\]
The orthogonal complement of $\cK_\gamma(I_x)$ in $\F^{I_x}$ is exactly
\[
 \Span\{\one_{I_x},(\gamma_i)_{i\in I_x}\}.
\]
Thus the local value vector is affine in the slope exactly when all local rows annihilate it.  Applying this independently at each coordinate proves the claim.  For a received line, take $a_x=r_0(x)$ and $b_x=r_1(x)$.
\end{proof}

\begin{proposition}[row count and trivial kernel]\label{prop:row-count}
The matrix has
\[
 \sum_{x\in D}(r_x-2)_+
 \ge Lm-2n
\]
rows.  It always contains the $2k$-dimensional kernel
\[
 h_i=A+\gamma_iB,
 \qquad A,B\in\F[X]_{<k}.
\]
Consequently its maximum useful rank is $(L-2)k$, and the equation surplus over the nontrivial unknown dimension is
\[
 \boxed{\ L(m-k)-2(n-k).\ }
\]
\end{proposition}

\begin{proof}
Since $(r-2)_+\ge r-2$,
\[
 \sum_x(r_x-2)_+
 \ge\sum_xr_x-2n
 =\sum_i|S_i|-2n
 =Lm-2n.
\]
A globally affine tuple satisfies every local relation.  Distinct slopes make the map $(A,B)\mapsto(A+\gamma_iB)_i$ injective, so this kernel has dimension $2k$.  Subtracting $(L-2)k$ from the row lower bound gives the displayed surplus.
\end{proof}

\section{Why the deployed critical order is exactly 32}

The actual Reed--Solomon dimension on both active MCA rows is
\[
 n=2^{21}=2097152,
 \qquad k=2^{20}=1048576.
\]
The adjacent safe agreements are
\[
 m_{\rm KB}=1116048,
 \qquad
 m_{\rm M31}=1116024.
\]

\begin{proposition}[critical order]\label{prop:critical-32}
For either deployed row,
\[
 L_*:=\left\lfloor\frac{2(n-k)}{m-k}\right\rfloor+1=32.
\]
More precisely,
\[
\begin{array}{c|r|r|r}
\toprule
\text{row}&m-k&31(m-k)-2(n-k)&32(m-k)-2(n-k)\\
\midrule
\text{KoalaBear MCA}&67472&-5520&61952\\
\text{Mersenne-31 MCA}&67448&-6264&61184\\
\bottomrule
\end{array}
\]
Thus order $31$ can remain underdetermined for dimension reasons, whereas order $32$ has a positive row surplus.
\end{proposition}

\begin{proof}
Direct substitution into \cref{prop:row-count}.
\end{proof}

\section{What full relative rank would buy}

\begin{theorem}[one global code-line block]\label{thm:global-block}
Let $Z$ be a set of bad slopes whose explaining codewords lie on one global codeword line,
\[
 h_\gamma=A+\gamma B.
\]
Assume the received line is support-wise MCA-bad at agreement $m$.  Then
\[
 \boxed{|Z|\le n-m+1.}
\]
\end{theorem}

\begin{proof}
Let
\[
 C_0:=\{x\in D:r_0(x)=A(x),\ r_1(x)=B(x)\},
 \qquad c:=|C_0|.
\]
MCA-badness gives $c<m$.  Outside $C_0$, the equation
\[
 r_0(x)-A(x)+\gamma(r_1(x)-B(x))=0
\]
has at most one finite solution $\gamma$.  Counting agreements over all $\gamma\in Z$ therefore yields
\[
 |Z|m\le |Z|c+(n-c).
\]
Hence
\[
 |Z|\le\frac{n-c}{m-c}\le n-m+1,
\]
where the last expression is the maximum over integers $0\le c\le m-1$.
\end{proof}

\begin{corollary}[conditional generic closure]\label{cor:generic-closure}
Suppose every $32$-subfamily of a set of bad explanations has full relative rank, meaning that the kernel of its rank-two matrix is exactly the global affine kernel.  Then all explanations lie on one global codeword line and
\[
 |Z|\le n-m+1.
\]
At the deployed rows the bounds are
\[
 981105
 \qquad\text{and}\qquad
 981129,
\]
respectively.
\end{corollary}

\begin{proof}
Any $32$ explanations are globally affine.  If there are more than $32$, fix two slopes and place any third slope in a $32$-subset containing the fixed pair; the unique codeword line through the fixed pair contains the third.  Apply \cref{thm:global-block}.
\end{proof}

The Mersenne-31 challenge numerator is $16777215$, so the generic order-$32$ branch would fit with more than fifteen million units of margin.  The difficult issue is therefore not the generic branch.  It is classification and exclusive payment of specialized rank defects.

\section{A necessary partition inequality}

The following obstruction identifies the correct rank-two analogue of the partition conditions in higher-order MDS theory.

For a partition $\cP$ of $[L]$, put
\[
 c(\cP):=\sum_{P\in\cP}\min\{2,|P|\},
\]
and for $I\subseteq[L]$ define
\[
 w_\cP(I)
 :=\sum_{P\in\cP}\min\{2,|I\cap P|\}-\min\{2,|I|\}.
\]

\begin{proposition}[partition obstruction]\label{prop:partition-obstruction}
If
\[
 \sum_{x\in D}w_\cP(I_x)<k\bigl(c(\cP)-2\bigr)
\]
for some partition $\cP$, then the rank-two matrix has kernel dimension strictly larger than $2k$.
Consequently full relative rank requires
\[
 \boxed{
 \sum_{x\in D}w_\cP(I_x)
 \ge k\bigl(c(\cP)-2\bigr)
 \quad\text{for every partition }\cP.
 }
\]
\end{proposition}

\begin{proof}
Consider tuples that are independently affine on each block of $\cP$: a singleton block contributes one free degree-$<k$ polynomial, and a block of size at least two contributes two.  This space has dimension $kc(\cP)$.  At coordinate $x$, forcing the blockwise-affine local values on $I_x$ to lie on one affine slope line imposes at most $w_\cP(I_x)$ scalar conditions.  Hence the intersection with the rank-two kernel has dimension at least
\[
 kc(\cP)-\sum_xw_\cP(I_x).
\]
Under the strict inequality this exceeds $2k$.
\end{proof}

For generic evaluation points, higher-order MDS and row-span constrained local-property machinery make closely related partition inequalities sufficient for the expected rank and for curve-decoding phenomena \cite{BGM23,GGSW26}.  The deployed domains are not generic.  The point of the next sections is to identify what specializations must be removed before a sufficiency theorem can be true.  Recent work explicitly notes that identifying deterministic Reed--Solomon evaluation sets with the random-evaluation curve-decoding behavior remains open \cite{GGSW26}.

\section{Exact dyadic quotient descent}

\begin{theorem}[pair descent on a multiplicative dyadic domain]\label{thm:pair-descent}
Assume $\operatorname{char}\F\ne2$ and
\[
 D=D_+\sqcup(-D_+).
\]
Suppose the incidence pattern is pair-invariant:
\[
 I_x=I_{-x}\qquad(x\in D_+).
\]
Write uniquely
\[
 h_i(X)=u_i(X^2)+Xv_i(X^2),
\]
where
\[
 \deg u_i<\left\lceil\frac{k}{2}\right\rceil,
 \qquad
 \deg v_i<\left\lfloor\frac{k}{2}\right\rfloor.
\]
Then the rank-two kernel on $D$ is naturally the direct sum of two rank-two kernels on the squared domain
\[
 D^{(2)}:=\{x^2:x\in D_+\},
\]
with the common incidence sets $I_x=I_{-x}$ and degree bounds $\lceil k/2\rceil$ and $\lfloor k/2\rfloor$.
In particular, every nontrivial pair-invariant defect descends to a nontrivial defect on the half-size quotient in at least one parity component.
\end{theorem}

\begin{proof}
Fix $x\in D_+$ and $c\in\cK_\gamma(I_x)=\cK_\gamma(I_{-x})$.  The two local equations are
\[
 \sum_i c_i\bigl(u_i(x^2)+xv_i(x^2)\bigr)=0,
 \qquad
 \sum_i c_i\bigl(u_i(x^2)-xv_i(x^2)\bigr)=0.
\]
Because $2x\ne0$, adding and subtracting gives
\[
 \sum_i c_i u_i(x^2)=0,
 \qquad
 \sum_i c_i v_i(x^2)=0.
\]
The converse follows by reversing the calculation.  The even/odd polynomial decomposition is a vector-space isomorphism, so the full kernel is the direct sum claimed.  The global affine kernels descend componentwise; hence any excess kernel dimension appears in at least one component.
\end{proof}

Iterating \cref{thm:pair-descent} proves exact descent for every fully periodic dyadic incidence atom.  This is the rigorous quotient branch of the proposed specialization theorem.  What fails is the assertion that every non-descending atom must be a simple pole.

\section{Rational atoms and profile depth}

Let
\[
 \Lambda_i(X):=\prod_{x\in S_i}(X-x)
\]
be the monic support locator of degree $m$.

\begin{definition}[scalar-locator rational atom]\label{def:rational-atom}
A family is a scalar-locator rational atom if there exist
\[
 Q,A,B\in\F[X],
 \qquad c_i=c_0+c_1\gamma_i,
\]
with $Q(x)\ne0$ for all $x\in D$, such that
\[
 \boxed{
 Qh_i+c_i\Lambda_i=A+\gamma_iB
 \qquad(i\in[L]).
 }
\]
\end{definition}

\begin{proposition}[rational atoms are MCA explanations]\label{prop:rational-to-mca}
Every atom of \cref{def:rational-atom} is explained by the rational received line
\[
 r_\gamma(x)=\frac{A(x)+\gamma B(x)}{Q(x)}.
\]
The codeword $h_i$ agrees with $r_{\gamma_i}$ on all roots of $\Lambda_i$.
\end{proposition}

\begin{proof}
At $x\in S_i$, the locator term vanishes and division by $Q(x)$ is valid.
\end{proof}

\begin{proposition}[converse under the scalar degree profile]\label{prop:mca-to-rational}
Put $m=k+d$.  Suppose
\[
 r_\gamma(x)=\frac{A(x)+\gamma B(x)}{Q(x)}
\]
on $D$, where $Q$ has no root in $D$, $\deg Q=s\le d$, and $\deg A,\deg B\le m$.  If $h_i\in\F[X]_{<k}$ agrees with $r_{\gamma_i}$ on the $m$ roots of $\Lambda_i$, then
\[
 A+\gamma_iB-Qh_i=c_i\Lambda_i
\]
for a scalar $c_i$, and $c_i$ is affine in $\gamma_i$.
\end{proposition}

\begin{proof}
The left side has degree at most $m$ and vanishes at the $m$ distinct roots of $\Lambda_i$, so it is a scalar multiple of $\Lambda_i$.  Since $\deg(Qh_i)\le k+s-1\le m-1$, its degree-$m$ coefficient is zero.  Comparing degree-$m$ coefficients gives
\[
 c_i=[X^m]A+\gamma_i[X^m]B.
\]
\end{proof}

\begin{corollary}[scaled-locator prefix depth]\label{cor:profile-depth}
Under the hypotheses of \cref{prop:mca-to-rational}, the coefficients of
\[
 c_i\Lambda_i
\]
in degrees
\[
 k+s,k+s+1,\ldots,m
\]
are affine functions of $\gamma_i$.  Apart from the leading scaling coefficient, this is a depth-$(d-s)$ affine prefix constraint.
\end{corollary}

\begin{proof}
The product $Qh_i$ has degree at most $k+s-1$, so the displayed high coefficients of $c_i\Lambda_i$ equal those of $A+\gamma_iB$.
\end{proof}

Thus a simple pole, $s=1$, is only one endpoint of a profile continuum.  At the opposite endpoint $s=d$, only the leading scaling is forced, while the locator residues modulo $Q$ must lie on a projective affine line.  Any correct rank-defect classification must account for the entire continuum or prove that the already-declared first-match cells remove all but selected endpoints.

\section{An exact degree-two rank defect over \texorpdfstring{$\F_{17}$}{F17}}

We now give a completely explicit obstruction to an unqualified ``quotient or simple pole'' theorem.

\begin{theorem}[degree-two support-wise MCA atom]\label{thm:f17-atom}
Let
\[
 \F=\F_{17},\qquad
 D=\mu_8=\{1,9,13,15,16,8,4,2\},
 \qquad k=4,\quad m=6.
\]
Take slopes
\[
 (\gamma_0,\gamma_1,\gamma_2,\gamma_3,\gamma_4)=(0,1,4,7,2)
\]
and supports
\[
\begin{array}{c|l}
 i&S_i\\
\midrule
0&\{1,15,16,8,4,2\}\\
1&\{1,9,13,15,16,8\}\\
2&\{1,13,15,16,4,2\}\\
3&\{1,9,15,16,4,2\}\\
4&\{1,9,13,15,8,4\}.
\end{array}
\]
Let
\[
\begin{aligned}
 h_0&=h_1=0,\\
 h_2&=9+13X+8X^2+4X^3,\\
 h_3&=3+10X+14X^2+7X^3,\\
 h_4&=11+8X+9X^2+6X^3.
\end{aligned}
\]
Define a received line by the following values, in the displayed order of $D$:
\[
\begin{array}{c|rrrrrrrr}
 x&1&9&13&15&16&8&4&2\\
\midrule
r_0(x)&0&16&6&0&0&0&0&0\\
 r_1(x)&0&1&11&0&0&0&5&12.
\end{array}
\]
Then:
\begin{enumerate}[label=\textup{(\roman*)}]
\item $h_i$ agrees with $r_0+\gamma_i r_1$ on $S_i$ for every $i$;
\item the received line is support-wise MCA-bad at agreement $m=6$;
\item the rank-two matrix has rank $11$, whereas full relative rank is $12$;
\item the defect is not first-level quotient-periodic;
\item it has no scalar-locator rational representation with $\deg Q\le1$, but it has one with $\deg Q=2$.
\end{enumerate}
\end{theorem}

\begin{proof}
Direct substitution proves (i).  For (ii), observe that $r_0$ is zero at exactly the six points $S_0$.  On any other six-point set, a degree-$<4$ polynomial agreeing with $r_0$ would have at least four zero roots and also take a prescribed nonzero value, which is impossible; with five zero roots the same conclusion is immediate.  Hence the only possible six-point common explanation set for the first coefficient is $S_0$.  On $S_0$, the word $r_1$ is zero at $1,15,16,8$ and nonzero at $4,2$, so no degree-$<4$ polynomial can agree with it there.  Thus no pair of codeword polynomials simultaneously explains the received line on six coordinates.

Part (i) and \cref{thm:kernel-dictionary} give a nontrivial kernel vector.  It is not globally affine because $h_0=h_1=0$ at slopes $0$ and $1$, while $h_2\ne0$.  Exact row reduction gives rank $11$; one $11\times11$ pivot minor has determinant $9\in\F_{17}^\times$.  This proves (iii).

For (iv), the incidence sets at $1$ and $-1=16$ are different: the former contains all five slopes, while the latter omits slope $\gamma_4$.  Thus the hypothesis of \cref{thm:pair-descent} fails already at the first dyadic level.

For (v), let
\[
\begin{aligned}
 Q&=12+3X+5X^2,\\
 A&=13+10X+7X^2+13X^3+5X^4+11X^5+9X^6,\\
 B&=7+5X+6X^2+14X^3+3X^4+15X^5+X^6,
\end{aligned}
\]
and $c_i=9+\gamma_i$.  The polynomial $Q$ is nonzero on every point of $D$, and exact multiplication gives
\[
 Qh_i+c_i\Lambda_i=A+\gamma_iB
 \qquad(i=0,\ldots,4).
\]
For a denominator bound $s$, the identities of \cref{def:rational-atom} form a homogeneous coefficient system in $Q,A,B,c_0,c_1$.  At $s=0$ and $s=1$ the matrices have full column ranks $17$ and $18$, respectively.  At $s=2$ the $35\times19$ system has rank $18$ and its one-dimensional kernel is generated by the displayed certificate.  Hence degree two is minimal within this atom model.
\end{proof}

\begin{remark}[the common-support qualifier is essential]\label{rem:common-factor}
The example has the common support factor
\[
 H=(X-1)(X-15)=X^2+X+15.
\]
It divides every locator and also $A,B,h_0,\ldots,h_4$.  Therefore \cref{thm:f17-atom} refutes an \emph{unqualified} quotient-or-pole theorem, but it does not refute a theorem stated only after the common-support first-match branch has been removed.  Instead, it proves that the qualifier cannot be omitted and that rational denominator profiles can be hidden inside common-support reductions.
\end{remark}

The incidence pattern also satisfies every partition inequality of \cref{prop:partition-obstruction}, with minimum slack $1$.  Thus the defect is a genuine specialization phenomenon rather than a failure of the obvious partition count.  This finite claim is included in the verifier.

\section{Maximal supports and the slope-degree barrier}

For each explanation, take its \emph{maximal} agreement support
\[
 \widehat S_i
 :=\{x\in D:h_i(x)=r_0(x)+\gamma_i r_1(x)\},
 \qquad |\widehat S_i|\ge m.
\]
Put
\[
 \widehat I_x:=\{i:x\in\widehat S_i\},
 \qquad
 C_\cap:=\bigcap_{i=1}^{L}\widehat S_i,
 \qquad c:=|C_\cap|.
\]

Let $R_0,R_1\in\F[X]_{<n}$ be the interpolation polynomials of $r_0,r_1$ on $D$.  Interpolating the codeword polynomials coefficientwise in the slope gives the unique polynomial
\[
 H(X,Z)\in\F[X,Z],
 \qquad
 \deg_X H<k,
 \qquad
 \deg_Z H<L,
\]
with
\[
 H(X,\gamma_i)=h_i(X)
 \qquad(i\in[L]).
\]
Define the slope-error polynomial
\[
 E(X,Z):=H(X,Z)-R_0(X)-ZR_1(X).
\]

\begin{lemma}[zero slope fibers are exactly common agreements]\label{lem:zero-fibers}
For $x\in D$,
\[
 E(x,Z)\equiv0
 \quad\Longleftrightarrow\quad
 x\in C_\cap.
\]
\end{lemma}

\begin{proof}
If $x\in C_\cap$, then $E(x,\gamma_i)=0$ for all $L$ distinct slopes.  Since $\deg_ZE(x,Z)<L$, the polynomial is zero.  Conversely, if $E(x,Z)$ is zero, then evaluating at every $\gamma_i$ gives
\[
 h_i(x)=R_0(x)+\gamma_iR_1(x)=r_0(x)+\gamma_i r_1(x),
\]
so $x$ belongs to every maximal support.
\end{proof}

\begin{theorem}[maximal-support slope-degree barrier]\label{thm:slope-degree}
Assume the explanations are not globally affine: there do not exist $A,B\in\F[X]_{<k}$ with
\[
 h_i=A+\gamma_iB
 \qquad(i\in[L]).
\]
Then $E\ne0$.  Writing
\[
 r:=\deg_ZE,
\]
one has $c<k$ and
\[
 \boxed{
 Lm\le Lc+r(n-c).
 }
\]
Equivalently,
\[
 \boxed{
 r\ge
 \ceil{\frac{L(m-c)}{n-c}}.
 }
\]
In particular, after the common-support branch is removed ($c=0$),
\[
 r\ge\ceil{\frac{Lm}{n}}.
\]
\end{theorem}

\begin{proof}
If $c\ge k$, choose two slopes and let $A,B\in\F[X]_{<k}$ be the unique codeword line through their polynomials.  For every $i$, the polynomial $h_i-A-\gamma_iB$ has degree $<k$ and vanishes on $C_\cap$, hence is zero.  This contradicts non-affinity, so $c<k$.

For $x\notin C_\cap$, \cref{lem:zero-fibers} says $E(x,Z)$ is nonzero.  It has the $|\widehat I_x|$ distinct roots $\gamma_i$ indexed by agreements at $x$, whence
\[
 |\widehat I_x|\le r.
\]
Counting maximal-support incidences gives
\[
 Lm
 \le\sum_i|\widehat S_i|
 =\sum_{x\in C_\cap}L+\sum_{x\notin C_\cap}|\widehat I_x|
 \le Lc+r(n-c).
\]
It remains only to justify that $E\ne0$.  If it were zero, then $H$ would be affine in $Z$ and all $h_i=H(X,\gamma_i)$ would lie on one global codeword line.
\end{proof}

\begin{corollary}[the active primitive degree is at least eighteen]\label{cor:degree-18}
For $L=32$ and either deployed adjacent row, a non-affine explanation with no common agreement coordinate satisfies
\[
 \boxed{18\le r\le31.}
\]
Indeed,
\[
 \ceil{\frac{32m_{\rm KB}}{n}}
 =
 \ceil{\frac{32m_{\rm M31}}{n}}
 =18.
\]
\end{corollary}

Thus every slope-degree-$\le17$ explanation is already global-affine or common-support.  This removes the entire low-degree branch before any determinant specialization argument begins.

\section{Local incidence complexity}

For the rational extraction theorem, choose exact $m$-subsupports
\[
 S_i\subseteq\widehat S_i,
 \qquad |S_i|=m,
\]
and set
\[
 I_x:=\{i:x\in S_i\},
 \qquad r_x:=|I_x|.
\]
Let $N_a$ count coordinates of incidence multiplicity $a$, let
\[
 N_{\ge2}:=|\{x:r_x\ge2\}|,
 \qquad
 U:=\bigcup_iS_i,
 \qquad u:=|U|.
\]

\begin{definition}[two-cover incidence complexity]\label{def:chi}
Define
\[
 \chi(\cS)
 :=\sum_{x\in D}\min\{2,r_x\}
 =N_1+2N_{\ge2}.
\]
Equivalently,
\[
 \boxed{
 \chi(\cS)=2u-N_1=2n-2N_0-N_1.
 }
\]
\end{definition}

\begin{proposition}[exact relation to the rank-two row count]\label{prop:chi-rows}
For exact $m$-supports,
\[
 \sum_{x\in D}(r_x-2)_+
 =Lm-\chi(\cS).
\]
\end{proposition}

\begin{proof}
At multiplicity zero or one, both sides contribute zero after subtracting $\min\{2,r_x\}$.  At multiplicity at least two, the difference is $r_x-2$.  Sum over $x$ and use $\sum_xr_x=Lm$.
\end{proof}

The quantity $\chi$ is the number of scalar conditions needed to force one common rational received line on the occupied coordinates: two conditions at a multiply occupied coordinate and one at a singleton.

\section{Support-collapsed rational interpolation}

Let
\[
 \Lambda_i(X):=\prod_{x\in S_i}(X-x)
\]
be the monic locator of $S_i$.

\begin{definition}[scalar-locator certificate]\label{def:certificate}
A scalar-locator certificate of denominator degree at most $s$ is a tuple
\[
 (Q,A,B,c_0,c_1)
\]
with
\[
 \deg Q\le s,
 \qquad
 \deg A,\deg B\le m,
\]
not all zero, such that, with $c_i=c_0+c_1\gamma_i$,
\[
 \boxed{
 Qh_i+c_i\Lambda_i=A+\gamma_iB
 \qquad(i\in[L]).
 }
\]
It is a rational atom when $Q\ne0$ and $(c_0,c_1)\ne(0,0)$, a pure locator pencil when $Q=0$, and globally affine when $(c_0,c_1)=(0,0)$.
\end{definition}

When $Q$ has no root on $D$, a rational atom is explained by the received line $(A+\gamma B)/Q$.  A root of $Q$ on $D$ is retained as a separate denominator-root first-match branch rather than being silently divided away.

\begin{theorem}[low-complexity atom extraction]\label{thm:atom-extraction}
Assume $L\ge2$.  Let
\[
 d=m-k
\]
and fix an integer $0\le s\le d$.  If
\[
 \boxed{
 \chi(\cS)<2m+s+3,
 }
\]
then the explanations admit a scalar-locator certificate of denominator degree at most $s$.

More precisely, one of the following holds:
\begin{enumerate}[label=\textup{(\roman*)}]
\item the codeword explanations are globally affine;
\item they lie in a pure locator pencil;
\item they form a rational atom, possibly with a denominator-root degeneracy on $D$.
\end{enumerate}
\end{theorem}

\begin{proof}
Use as unknowns the coefficients of
\[
 Q\in\F[X]_{\le s},
 \qquad
 A,B\in\F[X]_{\le m}.
\]
There are
\[
 (s+1)+2(m+1)=2m+s+3
\]
unknowns.

At a coordinate $x$ with $r_x\ge2$, impose the two homogeneous linear equations
\[
 A(x)=Q(x)r_0(x),
 \qquad
 B(x)=Q(x)r_1(x).
\]
At a singleton coordinate, with $I_x=\{i\}$, impose
\[
 A(x)+\gamma_iB(x)=Q(x)h_i(x).
\]
No equation is needed at an unoccupied coordinate.  The total number of equations is at most $\chi(\cS)$.  Under the strict inequality, the homogeneous system has a nonzero solution $(Q,A,B)$.

For every $i$ and every $x\in S_i$, the imposed equations give
\[
 A(x)+\gamma_iB(x)=Q(x)h_i(x).
\]
The polynomial
\[
 P_i:=A+\gamma_iB-Qh_i
\]
has degree at most $m$ because
\[
 \deg(Qh_i)\le s+k-1\le m-1.
\]
It vanishes on the $m$ distinct roots of $\Lambda_i$, so
\[
 P_i=c_i\Lambda_i
\]
for a scalar $c_i$.  Comparing coefficients of $X^m$ gives
\[
 c_i=[X^m]A+\gamma_i[X^m]B,
\]
which is affine in $\gamma_i$.  This is the required certificate.

If $Q=0$, nontriviality and two distinct slopes rule out $c_0=c_1=0$, so the certificate is a pure locator pencil.  If $Q\ne0$ but $c_0=c_1=0$, then
\[
 Qh_i=A+\gamma_iB.
\]
Using two slopes shows $Q$ divides both $A$ and $B$, and cancellation gives a global affine codeword line.  The remaining case is a rational atom.
\end{proof}

\begin{remark}[relation to rational-function interpolation]
The proof is a support-collapsed Berlekamp--Welch argument: it interpolates only the two received-line coefficients on multiply occupied coordinates, rather than imposing a separate condition for every support incidence.  This is compatible with the rational-function and curve-decoding viewpoint in recent proximity-gap work \cite{BSCHKS25,GG26}, but the theorem above is elementary and finite; no random-evaluation or asymptotic result is imported.
\end{remark}

\begin{corollary}[full denominator profile]\label{cor:full-profile}
Taking $s=d=m-k$, every explanation with
\[
 \boxed{
 \chi(\cS)<3m-k+3
 }
\]
is globally affine, a pure locator pencil, or a scalar-locator rational atom of denominator degree at most $m-k$.
\end{corollary}

\begin{corollary}[support-union form]\label{cor:union-form}
If
\[
 2u-N_1<3m-k+3,
\]
then the same alternatives hold.  In particular, it suffices that
\[
 u\le\floor{\frac{3m-k+2}{2}}.
\]
\end{corollary}

\begin{corollary}[exact deployed extraction thresholds]\label{cor:active-extraction}
For the active $32$-tuples, put
\[
 \Xi:=3m-k+3.
\]
Then
\[
\begin{array}{c|r|r|r|r}
\toprule
\text{row}&\Xi&2n-\Xi&\text{automatic union ceiling}&\text{union nullity ceiling}\\
\midrule
\text{KoalaBear MCA}&2299571&1894733&1149785&101209\\
\text{Mersenne-31 MCA}&2299499&1894805&1149749&101173\\
\bottomrule
\end{array}
\]
where the last column subtracts $n-k=1048576$ from the union ceiling.

Equivalently, extraction is automatic if
\[
 2N_0+N_1\ge1894734
\]
on KoalaBear, or if
\[
 2N_0+N_1\ge1894806
\]
on Mersenne-31.
\end{corollary}

These are classification thresholds, not challenge-numerator payments.  Their use is to show that a large low-complexity support family belongs to one atom instead of being charged as many unrelated cells.

\section{The degree-two counterexample lies at the extraction threshold}

The degree-two $\F_{17}$ example of \cref{thm:f17-atom} has
\[
 n=8,
 \qquad k=4,
 \qquad m=6,
 \qquad L=5.
\]
Its five listed supports are maximal.  Their incidence multiplicities are
\[
 (5,3,3,5,4,3,4,3),
\]
so
\[
 \chi=16
 \quad\text{while}\quad
 3m-k+3=17.
\]
Thus \cref{thm:atom-extraction} forces a denominator of degree at most
\[
 d=m-k=2,
\]
which is exactly the minimal degree established by \cref{thm:f17-atom}.

The common support has size $c=2$.  The coefficientwise slope interpolation of the five codewords has degree $r=4$, and
\[
 \ceil{\frac{5(6-2)}{8-2}}=4.
\]
Hence the example also attains equality in \cref{thm:slope-degree}.  The counterexample is therefore not an unexplained accident: it sits simultaneously on the common-support slope-degree boundary and the low-complexity rational-extraction boundary.

\section{Collision pencils for two atom certificates}

Suppose the same indexed explanations
\[
 (\gamma_i,h_i,\Lambda_i),
 \qquad i\in I,
\]
satisfy two nontrivial scalar-locator certificates
\[
 Qh_i+(c_0+c_1\gamma_i)\Lambda_i=A+\gamma_iB,
\]
\[
 Q'h_i+(c'_0+c'_1\gamma_i)\Lambda_i=A'+\gamma_iB'.
\]
Assume $|I|\ge2$, $Q,Q'\ne0$, and both scalar pairs are nonzero.  Define
\[
\begin{aligned}
 P_0&:=Q'c_0-Qc'_0,
 &P_1&:=Q'c_1-Qc'_1,\\
 R_0&:=Q'A-QA',
 &R_1&:=Q'B-QB'.
\end{aligned}
\]
Eliminating $h_i$ gives the collision identity
\[
 \boxed{
 (P_0+\gamma_iP_1)\Lambda_i
 =R_0+\gamma_iR_1.
 }
\]

Let
\[
 G:=\{x\in D:R_0(x)=R_1(x)=0\},
 \qquad g:=|G|,
\]
and
\[
 H:=\{x\in G:P_0(x)=P_1(x)=0\},
 \qquad h:=|H|.
\]

\begin{theorem}[atom collision rigidity]\label{thm:collision-rigidity}
If $R_0=R_1=0$ and the certificates share at least two distinct slopes, then the certificates are projectively identical.

Otherwise, if $g<m$, then
\[
 \boxed{
 |I|\le\floor{\frac{n-g}{m-g}}.
 }
\]
Moreover, every coordinate in $G\setminus H$ belongs to at least $|I|-1$ of the supports $S_i$.

If both denominators have degree at most $d$, then two distinct certificates satisfy
\[
 h\le d.
\]
\end{theorem}

\begin{proof}
First suppose $R_0=R_1=0$.  Then
\[
 (P_0+\gamma_iP_1)\Lambda_i=0
\]
as a polynomial.  Since $\Lambda_i\ne0$, two distinct slopes imply $P_0=P_1=0$.  Thus
\[
 Q'(c_0,c_1)=Q(c'_0,c'_1).
\]
Choose a nonzero component of $(c_0,c_1)$.  Because $Q,Q'\ne0$, this identity gives a scalar $\lambda\in\F^\times$ with
\[
 Q'=\lambda Q,
 \qquad
 (c'_0,c'_1)=\lambda(c_0,c_1).
\]
The equations $R_0=R_1=0$ then give
\[
 A'=\lambda A,
 \qquad B'=\lambda B.
\]
So the certificates are projectively identical.

Now assume $(R_0,R_1)\ne(0,0)$.  At every root of $\Lambda_i$,
\[
 R_0(x)+\gamma_iR_1(x)=0.
\]
After removing the at most $g$ roots in $G$, each $S_i$ contributes at least $m-g$ incidences.  At a coordinate outside $G$, the affine equation
\[
 R_0(x)+\gamma R_1(x)=0
\]
has at most one slope solution.  Hence
\[
 |I|(m-g)\le n-g,
\]
which proves the bound.

For $x\in G\setminus H$, the polynomial
\[
 P_0(x)+\gamma P_1(x)
\]
is a nonzero affine function of $\gamma$ and therefore vanishes at at most one of the shared slopes.  The collision identity then forces $\Lambda_i(x)=0$ for all other $i$, proving the near-sunflower assertion.

Finally, if the certificates are distinct, then $(P_0,P_1)\ne(0,0)$ by the first part.  Both have degree at most $d$, so their common root set in $D$ has size at most $d$.
\end{proof}

\begin{corollary}[primitive atom uniqueness]\label{cor:primitive-unique}
If $G=\varnothing$, two distinct atom certificates share at most
\[
 \floor{\frac nm}=1
\]
explanation at either deployed row.  Thus two primitive certificates sharing two slopes are projectively identical.
\end{corollary}

\section{Affine lines versus rational atoms}

A similar argument prevents a primitive rational atom from masquerading as a global affine codeword line.

\begin{theorem}[affine--atom collision]\label{thm:affine-atom}
Suppose $h_i=A_0+\gamma_iB_0$ for $i\in I$, and the same explanations satisfy a nontrivial atom certificate
\[
 Qh_i+(c_0+c_1\gamma_i)\Lambda_i=A+\gamma_iB,
 \qquad (c_0,c_1)\ne(0,0).
\]
Put
\[
 E_0:=A-QA_0,
 \qquad
 E_1:=B-QB_0,
\]
and let
\[
 G_{\rm aff}:=\{x\in D:E_0(x)=E_1(x)=0\},
 \qquad g_{\rm aff}:=|G_{\rm aff}|.
\]
If $|I|\ge2$, then $(E_0,E_1)\ne(0,0)$, and if $g_{\rm aff}<m$,
\[
 \boxed{
 |I|\le\floor{\frac{n-g_{\rm aff}}{m-g_{\rm aff}}}.
 }
\]
Every point of $G_{\rm aff}$ belongs to at least $|I|-1$ supports.
In particular, a primitive atom and a global affine line share at most one explanation.
\end{theorem}

\begin{proof}
Substitution gives
\[
 (c_0+c_1\gamma_i)\Lambda_i=E_0+\gamma_iE_1.
\]
If $E_0=E_1=0$, then $c_0+c_1\gamma_i=0$ for every $i$, since $\Lambda_i\ne0$.  A nonzero affine function of the slope has at most one root, contradicting $|I|\ge2$.

The root-incidence count for the pencil $E_0+\gamma E_1$ is identical to the proof of \cref{thm:collision-rigidity}.  At a common fixed root, all supports except possibly the unique slope with $c_0+c_1\gamma=0$ contain the coordinate.
\end{proof}

\section{Coherence or a quantified near-sunflower}

The graph whose vertices are the $32$-subsets of a fixed slope set and whose edges join subsets with $31$ common slopes is connected.  Collision rigidity can therefore propagate a local classification globally.

\begin{theorem}[order-$32$ coherence]\label{thm:coherence}
Let $Z$ be a set of at least $32$ bad slopes.  Suppose every $32$-subset of the corresponding explanations is classified as either
\begin{enumerate}[label=\textup{(\alph*)}]
\item one global affine codeword line, or
\item one nontrivial scalar-locator rational atom of denominator degree at most $d=m-k$.
\end{enumerate}
Assume all cross-fixed sets arising between distinct local classifications have already been removed by the first-match rule.  Then the entire set $Z$ lies on one global affine codeword line or in one projective rational atom certificate.
\end{theorem}

\begin{proof}
Adjacent $32$-subsets share $31$ explanations.  Two affine codeword lines sharing two distinct slopes are identical.  By \cref{cor:primitive-unique}, two primitive atom certificates sharing two explanations are identical.  By \cref{thm:affine-atom}, a primitive atom and an affine line cannot share two explanations.  Therefore adjacent vertices of the $32$-subset graph receive the same type and the same structure.  Connectedness propagates that structure to every $32$-subset and hence to every explanation in $Z$.
\end{proof}

Without the primitive hypothesis, failure of coherence is not unstructured.

\begin{corollary}[exact near-sunflower alternative on a $31$-overlap]\label{cor:near-sunflower}
Suppose two distinct local classifications share $q=31$ explanations.  Then at least one classification is an atom, and their cross-fixed set has size at least
\[
 g_{31}:=\ceil{\frac{31m-n}{30}}.
\]
For two rational atoms, at least $g_{31}-d$ of those coordinates belong to at least $30$ of the $31$ supports.  For an affine-line/atom collision, all $g_{31}$ fixed coordinates have that property.

At the deployed rows,
\[
\begin{array}{c|r|r|r}
\toprule
\text{row}&g_{31}&d=m-k&g_{31}-d\\
\midrule
\text{KoalaBear MCA}&1083345&67472&1015873\\
\text{Mersenne-31 MCA}&1083320&67448&1015872\\
\bottomrule
\end{array}
\]
Thus every failure of local atom coherence produces a $30$-of-$31$ near-sunflower on more than one million domain coordinates.
\end{corollary}

\begin{proof}
Rearrange the collision inequality
\[
 31(m-g)\le n-g
\]
to obtain
\[
 g\ge\ceil{\frac{31m-n}{30}}.
\]
For two atoms, \cref{thm:collision-rigidity} gives a collision core $H$ of size at most $d$; every point outside it lies in at least thirty supports.  The affine--atom statement follows from \cref{thm:affine-atom}.
\end{proof}

For reference, if two classifications share all $32$ explanations, the corresponding fixed-root thresholds are
\[
\begin{array}{c|r|r}
\toprule
\text{row}&g_{32}=\ceil{(32m-n)/31}&g_{32}-d\\
\midrule
\text{KoalaBear MCA}&1084400&1016928\\
\text{Mersenne-31 MCA}&1084375&1016927\\
\bottomrule
\end{array}
\]

\section{The proved part of the relative order-\texorpdfstring{$32$}{32} theorem}

The results can now be stated as one relative classification theorem for the actual explanation vector, which is weaker than classifying every vector in the rank-two kernel but is sufficient for MCA.

\begin{theorem}[partial relative order-\texorpdfstring{$32$}{32} explanation theorem]\label{thm:partial-relative}
Fix either deployed adjacent MCA row and $32$ distinct bad slopes.  After maximalizing supports, removing their common-support branch, and selecting exact $m$-subsupports, any non-global-affine explanation outside the pure-locator, denominator-root, and scalar-locator rational-profile cells must satisfy both
\[
 \boxed{18\le\deg_ZE\le31}
\]
and
\[
 \boxed{\chi(\cS)\ge3m-k+3.}
\]
Numerically, the second inequality is
\[
 \chi(\cS)\ge2299571
 \quad\text{on KoalaBear},
\]
and
\[
 \chi(\cS)\ge2299499
 \quad\text{on Mersenne-31}.
\]

Moreover, if every residual $32$-subset is classified as affine or rational, then either all bad slopes belong to one coherent structure or a $31$-overlap enters the quantified near-sunflower branch of \cref{cor:near-sunflower}.
\end{theorem}

\begin{proof}
The slope-degree statement is \cref{cor:degree-18}.  The incidence-complexity statement is the contrapositive of \cref{cor:full-profile}.  The coherence alternative follows from \cref{thm:coherence,cor:near-sunflower}.
\end{proof}

This proves the relative theorem in two substantial regimes:
\begin{itemize}[leftmargin=2em]
\item low slope degree ($r\le17$) is impossible outside global-affine/common-support structure;
\item low two-cover complexity ($\chi<3m-k+3$) automatically produces a rational or pure-locator atom.
\end{itemize}
It also discharges the separate atom-coherence obligation once the near-sunflower first-match cell is applied.


\part{Primitive spread abundance and the exact MCA endgame}
\label{part:spread-endgame}
The rank and atom theorems reduce a large bad-slope set to one coherent global structure, except for a quantified exceptional set.  This part proves that universal spread exclusion is false, identifies the correct spread-abundance target, localizes rational atoms to a canonical owner, and bounds broad classes of spread-core corrections.

\section{A critical spread core exists}

A universal theorem excluding primitive spread critical tuples on dyadic domains would be stronger than the deployed abundance statement.  The following exact example shows that such a universal prohibition is false.

\begin{proposition}[exact dyadic critical spread core]\label{prop:f17-spread}
Let
\[
 \F=\F_{17},
 \qquad
 D=\F_{17}^{\times}=\langle3\rangle,
 \qquad
 n=16,
 \quad k=8,
 \quad m=10.
\]
The critical order is
\[
 L_*=\floor{\frac{2(n-k)}{m-k}}+1=9.
\]
There exist nine distinct slopes, nine codeword explanations, and exact size-$10$ witness supports with all of the following properties.
\begin{enumerate}[label=\textup{(\roman*)}]
\item The received pair is support-wise MCA-bad at agreement $10$.
\item The maximal agreement supports have empty common intersection.
\item The rank-two reduced intersection matrix has rank $55$, whereas full relative rank is $(9-2)8=56$.
\item The two-cover complexity is $32$, above the extraction threshold of \cref{cor:full-profile} $3m-k+3=25$.
\item The incidence pattern is not periodic at any nontrivial level of the dyadic tower.
\item There is no scalar-locator certificate with denominator degree $0$, $1$, or $2=m-k$.
\item The coefficientwise slope interpolation has degree $8$, above the common-support-free lower barrier $\ceil{9m/n}=6$.
\item The direction code contains a word of weight $9=n-k+1$, so the minimum-support excess is exactly one.
\end{enumerate}
Moreover, for this same received line every slope in $\F_{17}$ has a degree-$<8$ codeword agreeing on at least ten coordinates.
\end{proposition}

\begin{proof}[Exact certificate]
Order the domain as
\[
 (1,3,9,10,13,5,15,11,16,14,8,7,4,12,2,6).
\]
Use slopes $\gamma_i=i$ for $0\le i\le8$ and the following exact size-$10$ supports, written by domain indices:
\[
\begin{array}{c|l}
\toprule
 i&S_i\\
\midrule
0&0,3,4,5,7,9,10,11,12,13\\
1&1,3,5,6,7,10,11,12,13,14\\
2&1,4,6,7,9,10,11,12,13,15\\
3&0,1,2,4,5,6,7,9,14,15\\
4&0,2,3,4,5,6,7,9,12,13\\
5&0,3,4,5,6,7,8,9,10,14\\
6&1,2,3,5,6,8,9,11,13,15\\
7&0,3,4,5,6,7,11,12,14,15\\
8&2,3,4,5,6,7,8,9,10,14\\
\bottomrule
\end{array}
\]
The codeword coefficient vectors, constant coefficient first, are
\[
\begin{array}{c|rrrrrrrr}
\toprule
 i&1&X&X^2&X^3&X^4&X^5&X^6&X^7\\
\midrule
0&1&13&16&16&9&16&7&12\\
1&13&9&14&12&6&4&2&13\\
2&2&2&2&4&2&12&0&1\\
3&13&3&3&4&0&13&16&5\\
4&1&4&11&4&2&16&9&16\\
5&16&13&15&13&2&15&5&7\\
6&3&3&0&13&0&8&12&4\\
7&3&2&6&8&10&13&10&12\\
8&1&12&2&14&9&10&9&4\\
\bottomrule
\end{array}
\]
and the received coefficient words on the ordered domain are
\[
\begin{aligned}
 r_0&=(5,16,0,0,0,0,0,0,5,7,8,0,8,10,0,6),\\
 r_1&=(6,14,15,0,0,0,2,13,14,12,6,3,15,13,6,15).
\end{aligned}
\]
Direct substitution verifies every claimed support.  Maximalizing the nine supports gives sizes $10$ or $11$ and empty total intersection.

For the rank statement, form the rank-two matrix of \cref{def:r2rim}.  It has $58$ rows and $72$ columns.  The displayed explanation vector supplies one kernel direction independent of the $16$-dimensional global-affine kernel, so its rank is at most $55$; exact Gaussian elimination over $\F_{17}$ gives rank $55$.  The scalar-locator systems at denominator degrees $0,1,2$ have column counts and ranks
\[
 (25,25),\qquad(26,26),\qquad(27,27),
\]
respectively.  Exact interpolation gives slope degree $8$, and direct evaluation of all pairwise codeword differences gives minimum weight $9$.

Finally, exhaustive interpolation over the $\binom{16}{10}$ candidate common supports finds no ten-coordinate support on which both $r_0$ and $r_1$ are degree-$<8$ evaluations.  Enumerating the $\binom{16}{8}$ interpolation bases at each of the seventeen slopes finds at least one agreement-$10$ codeword for every slope.  The accompanying verifier repeats all calculations in exact arithmetic.
\end{proof}

\begin{remark}[what the example does and does not refute]\label{rem:counter-scope}
The example refutes the \emph{universal} prohibition of primitive spread critical tuples.  It does not refute a row-sharp theorem for the two deployed length-$2^{21}$ domains.  It also shows why the minimum-distance refinement is relevant but not by itself decisive: the hard spread core lies exactly in the support-excess-$1$ branch, the same branch in which the boundary top seam survives by \cref{cor:top-seam-localization}.
\end{remark}

\section{The deployed spread-abundance target}

\begin{definition}[first-match primitive spread core]\label{def:spread-core}
A set of $32$ explanations is a \emph{first-match primitive spread core} if it satisfies the residual conditions of \cref{thm:partial-relative}: no global affine codeword line, common-support factor, exact dyadic quotient descent, scalar-locator certificate of denominator degree at most $m-k$, denominator-root, extension, tangent, field-drop, or quantified near-sunflower atom; slope degree lies between $18$ and $31$; and two-cover complexity is at least $3m-k+3$.
\end{definition}

A universal nonexistence theorem is false by \cref{prop:f17-spread}.  The finite MCA theorem requires the weaker and more natural abundance statement.

\begin{problem}[deployed spread-abundance theorem]\label{prob:spread-abundance}
\emph{Intermediate status.}
This is the line-level abundance consequence required from the final
spread-component routing problem \cref{prob:mca-spread-routing}; it is not an
additional terminal input to the completion certificate.

For either deployed domain, prove that if a support-wise MCA-bad received line contains one first-match primitive spread core, then its total number of bad challenge slopes is at most the row budget $B^*$.

A stronger useful form would bound the number of slopes extending a fixed core by a quantity below $B^*-32$, with every failure routed to one canonical quotient, extension, rational-owner, clone, or near-sunflower certificate.
\end{problem}

\begin{remark}
The distinction is essential.  The finite theorem does not need local higher-order-MDS rank at every critical tuple.  It needs a global statement that a line carrying one genuine spread defect cannot carry too many further slopes.  The exact example in \cref{prop:f17-spread} shows that any proof must use the deployed domain and field geometry, not only the abstract critical-order dimension count.
\end{remark}

\section{Rational atoms have a canonical owner}

Consider one coherent rational atom certificate
\begin{equation}\label{eq:atom}
 Qh_i+(c_0+c_1\gamma_i)\Lambda_i=A+\gamma_iB,
 \qquad i\in I,
\end{equation}
where $\Lambda_i$ is the locator of a witness support $S_i$, $|S_i|\ge m$, and $Q$ has no root on $D$.  Put
\[
 \widetilde r_0(x)=\frac{A(x)}{Q(x)},
 \qquad
 \widetilde r_1(x)=\frac{B(x)}{Q(x)}.
\]
The line
\[
 \widetilde r_\gamma=\widetilde r_0+\gamma\widetilde r_1
\]
is the \emph{owner line} of the atom.

\begin{definition}[fixed owner set]\label{def:owner-set}
Define
\[
 G=\{x\in D:r_0(x)=\widetilde r_0(x),\ r_1(x)=\widetilde r_1(x)\},
 \qquad g=|G|.
\]
\end{definition}

\begin{theorem}[exclusive rational-owner localization]\label{thm:owner-localization}
Retain \eqref{eq:atom} and assume every scalar $c_0+c_1\gamma_i$ is nonzero; zero-scalar slopes are a separate global-affine degeneracy and contribute at most one slope.  Then:
\begin{enumerate}[label=\textup{(\alph*)}]
\item for every $x\in D\setminus G$, at most one support $S_i$ contains $x$;
\item if $g<m$, then
\[
 \boxed{
 |I|\le\floor{\frac{n-g}{m-g}}\le n-m+1;
 }
\]
\item if $g\ge m$, at most $n-g$ supports meet $D\setminus G$, while the supports contained in $G$ form a support-wise MCA-bad instance for the punctured code
\[
 C|_G\cong\RS[\F,G,k].
\]
Consequently
\[
 \boxed{
 |I|\le n-g+\Bbad(\RS[\F,G,k],m).
 }
\]
\end{enumerate}
\end{theorem}

\begin{proof}
On $S_i$, the locator vanishes, and \eqref{eq:atom} gives
\[
 h_i(x)=\widetilde r_0(x)+\gamma_i\widetilde r_1(x).
\]
Since $S_i$ is also an agreement support for the original received line,
\[
 (r_0-\widetilde r_0)(x)
 +\gamma_i(r_1-\widetilde r_1)(x)=0.
\]
For $x\notin G$ this is a nonzero affine equation in $\gamma_i$, hence has at most one solution.  This proves (a).

If $g<m$, every support uses at least $m-g$ coordinates outside $G$.  By (a), those outside incidences are disjoint across slopes, so
\[
 |I|(m-g)\le n-g.
\]
The quotient $(n-g)/(m-g)$ is increasing for $g<m$, and at $g=m-1$ equals $n-m+1$.  This proves (b).

Now assume $g\ge m$.  Distinct supports meeting $D\setminus G$ use disjoint nonempty outside-coordinate sets, so there are at most $n-g$ of them.  Every remaining support lies in $G$ and witnesses the punctured owner line there.  Since $g\ge m>k$, evaluation on $G$ is injective and the punctured code is Reed--Solomon of dimension $k$.  If the punctured pair admitted a common size-$m$ explaining support, the same support would explain the original pair on $D$, contradicting support-wise MCA-badness.  The contained slopes are therefore bounded by the punctured MCA numerator.
\end{proof}

\begin{corollary}[all small-owner atoms are paid]\label{cor:small-owner}
For either deployed row, every coherent rational atom with
\[
 g\le2m-k
\]
has at most $n$ bad slopes.  Explicitly, the paid owner ranges are
\[
\begin{array}{c|r|r}
\toprule
\text{row}&2m-k&\text{atom bound}\\
\midrule
\text{KoalaBear MCA}&1183520&2097152\\
\text{Mersenne-31 MCA}&1183472&2097152\\
\bottomrule
\end{array}
\]
Both fit their challenge budgets; on Mersenne-31 the exact residual margin is
\[
 16777215-2097152=14680063.
\]
\end{corollary}

\begin{proof}
The case $g<m$ is \cref{thm:owner-localization}(b).  For $m\le g\le2m-k$, the punctured error count is $g-m$ and
\[
 2(g-m)\le g-k.
\]
CAP25 v13.2 proves, by its half-distance pincer and unique-decoding correlated-agreement input, that the punctured MCA numerator is at most $g$ in this range \cite{CAP25v132}.  Apply \cref{thm:owner-localization}(c):
\[
 |I|\le(n-g)+g=n.
\]
\end{proof}

\section{A self-contained minimum-distance owner payment}

The corollary uses the half-distance proximity theorem of CAP25 v13.2.  The next statement is elementary and retains the actual minimum support of the punctured direction code.

\begin{theorem}[triple-collapse from direction minimum support]\label{thm:triple-collapse}
Assume $g\ge m$ and restrict to the atom slopes whose supports lie in $G$.  Let
\[
 C'_G=\Span\{h_i-h_j:i,j\in I\}\subseteq\RS[\F,G,k]
\]
be their direction code.  If it is nonzero, write
\[
 d_1(C'_G)=g-k+a,
 \qquad a\ge1.
\]
If
\[
 \boxed{
 3m-2g\ge k-a+1,
 }
\]
then all contained codeword explanations lie on one global codeword line.  Hence the full atom, including supports meeting $D\setminus G$, satisfies
\[
 \boxed{|I|\le n-m+1.}
\]
Equivalently, this payment applies whenever
\[
 g\le\floor{\frac{3m-k+a-1}{2}}.
\]
\end{theorem}

\begin{proof}
For three distinct slopes define the affine deviation
\[
 p_{ijk}
 =(
 \gamma_j-\gamma_k)h_i
 +(\gamma_k-\gamma_i)h_j
 +(\gamma_i-\gamma_j)h_k.
\]
It belongs to $C'_G$.  On $S_i\cap S_j\cap S_k$, all three codewords equal the owner line at their slopes, so $p_{ijk}$ vanishes.  Inclusion--exclusion gives
\[
 |S_i\cap S_j\cap S_k|\ge3m-2g.
\]
A nonzero word of $C'_G$ has at most
\[
 g-d_1(C'_G)=k-a
\]
zeros.  Under the displayed hypothesis, $p_{ijk}$ must therefore be zero for every triple.

Fix two slopes.  The vanishing of every triple deviation shows that every $h_i$ lies on the same global codeword line $H_0+\gamma_iH_1$.  Let $C_0\subseteq G$ be the coordinate set on which this codeword line equals the owner line coefficientwise.  Support-wise MCA-badness gives $|C_0|<m$.  Outside $C_0$, equality of the two affine lines determines at most one slope.  Thus the contained slopes satisfy
\[
 |I_{\rm in}|(m-|C_0|)\le g-|C_0|,
 \qquad
 |I_{\rm in}|\le g-m+1.
\]
There are at most $n-g$ crossing supports by \cref{thm:owner-localization}, and hence
\[
 |I|\le(n-g)+(g-m+1)=n-m+1.
\]
\end{proof}

\begin{corollary}[exact support-excess-one thresholds]\label{cor:triple-active}
At $a=1$, \cref{thm:triple-collapse} pays owners through
\[
\begin{array}{c|r}
\toprule
\text{row}&\floor{(3m-k)/2}\\
\midrule
\text{KoalaBear MCA}&1149784\\
\text{Mersenne-31 MCA}&1149748\\
\bottomrule
\end{array}
\]
Larger direction support excess increases the threshold by one owner coordinate for every two units of excess.
\end{corollary}

\begin{remark}[connection with the minimum-distance sharpening of \cref{sec:min-distance-sharpening}]
For the boundary prefix dictionary, \cref{prop:exchange-defect} proves that a shift pair of exchange defect $r$ can occur only when $r\ge a$, where $a$ is the direction-code support excess; the top seam $r=1$ is therefore confined to $a=1$ .  The counterexample in \cref{prop:f17-spread} also has $a=1$.  Thus the difficult seam is now consistently localized: large $a$ is paid by the sharpened flag and owner bounds, while the exact-minimum-distance branch remains the genuine image-counting problem.
\end{remark}

\section{Owner pencils and a seven-point rich-incidence bound}

The exclusive atom theorem also requires control of competing owner certificates.  \Cref{thm:collision-rigidity} gives uniqueness once two slopes are shared.  The following theorem gives a complementary geometric bound when owners vary in a projective pencil.

Let
\[
 (Q_\tau,A_\tau,B_\tau)
 =(Q_0,A_0,B_0)+\tau(Q_1,A_1,B_1),
 \qquad \tau\in\mathbb P^1(\F),
\]
be a projective pencil of owner triples.  At a coordinate $x\in D$, set
\[
 F_x(\gamma,\tau)
 =A_\tau(x)+\gamma B_\tau(x)
 -Q_\tau(x)(r_0(x)+\gamma r_1(x)).
\]
This has bidegree at most $(1,1)$ on $\mathbb P^1_\gamma\times\mathbb P^1_\tau$.

\begin{theorem}[component-free owner-pencil rich points]\label{thm:owner-pencil}
Assume no $F_x$ is identically zero and no two coordinate curves $F_x=0$ share an irreducible component.  Then the number of parameter pairs $(\gamma,\tau)$ for which at least $m$ coordinate equations vanish is at most
\[
 \boxed{
 \floor{\frac{2\binom n2}{\binom m2}}
 =\floor{\frac{2n(n-1)}{m(m-1)}}.
 }
\]
For both deployed MCA rows this number is exactly $7$.
\end{theorem}

\begin{proof}
Two bidegree-$(1,1)$ curves with no common component have intersection number at most
\[
 (1,1)\cdot(1,1)=2
\]
on $\mathbb P^1\times\mathbb P^1$.  A parameter pair lying on at least $m$ coordinate curves contributes at least $\binom m2$ incidences with unordered coordinate-curve pairs.  Summing first over parameter pairs and then over the $\binom n2$ curve pairs gives
\[
 N_{\rm rich}\binom m2\le2\binom n2.
\]
The deployed integer evaluations are direct.
\end{proof}

\begin{remark}[the shared-component branches are explicit]
A nonzero bidegree-$(1,1)$ polynomial is either irreducible or factors into one vertical and one horizontal linear factor.  Hence a shared component is one of:
\begin{enumerate}[label=\textup{(\roman*)}]
\item a fixed-slope vertical component;
\item a fixed-owner horizontal component;
\item a common irreducible $(1,1)$ component, in which case the two coordinate equations are proportional and form a coordinate-clone cell.
\end{enumerate}
Identically zero $F_x$ are fixed owner coordinates.  Thus every failure of \cref{thm:owner-pencil} is already a named structural branch rather than an unpriced multiplicity.
\end{remark}

\section{The exact large-owner atom target}

After \cref{cor:small-owner}, only
\[
 g\ge2m-k+1
\]
remains for a coherent rational atom.  If at most $31$ exceptional slopes are reserved for the local-to-global first-match transition, \cref{thm:owner-localization} shows that the required bound for the owner-contained slopes is
\begin{equation}\label{eq:owner-target}
 \boxed{
 B_{\rm owner}(g)
 \le B^*-31-(n-g).
 }
\end{equation}

At the first unproved owner size and at $g=n$, the targets are
\[
\begin{array}{c|r|r|r}
\toprule
\text{row}&g_{\min}=2m-k+1&\text{target at }g_{\min}&\text{target at }g=n\\
\midrule
\text{KoalaBear MCA}&1183521&274980728110481425&274980728111395056\\
\text{Mersenne-31 MCA}&1183473&15863505&16777184\\
\bottomrule
\end{array}
\]
The Mersenne-31 row is binding.

For the full-owner boundary-prefix atom, use the CAP25 convention $K=k+1$ and $w=m-K$.  The average prefix-fiber sizes are
\[
 \overline F=\binom nm|\mathbb B|^{-w}.
\]
Their high-precision values and the full-owner multiplicative targets are
\[
\begin{array}{c|r|r|r}
\toprule
\text{row}&w&\log_2\overline F&(B^*-31)/\overline F\\
\midrule
\text{KoalaBear MCA}&67471&35.7352464170&4807520.929566\\
\text{Mersenne-31 MCA}&67447&20.7411470346&9.57219783037\\
\bottomrule
\end{array}
\]
Thus the exact remaining Mersenne theorem is a factor-$9.5722$ maximum-image or maximum-fiber statement before any additional cell consumes budget.  Owner localization has removed every smaller owner and all multi-owner component-free pencils, but it does not prove this full-owner prefix bound.

\begin{problem}[exclusive large-owner atom image bound]\label{prob:large-owner}
For every first-match coherent rational atom on either deployed domain with $g>2m-k$, prove \eqref{eq:owner-target}.  In the full-owner simple-pole branch it suffices to prove the row-sharp challenge-slope image bound
\[
 |\{U_z(\alpha)-\Lambda_M(\alpha):M\in\operatorname{Fib}_w(z)\}\cap\Gamma|
 \le B^*-31.
\]
For Mersenne-31, the stronger support-fiber estimate need only lose a factor below $9.57219783037$ from the exact average.
\end{problem}

\section{A conditional exact adjacent-row theorem}

We can now state the final certificate without reverting to a sum of many unrelated local budgets.

\begin{definition}[exception routing input]\label{def:exception-routing}
Say that exception routing holds if, for every bad received line with more than $31$ slopes, there is a set $E$ of at most $31$ slopes such that every remaining local order-$32$ classification is primitive and \cref{thm:coherence} applies; every denominator-root, extension, coordinate-clone, imperfect-quotient, and quantified near-sunflower failure is either included in $E$ or assigned to the spread-abundance or large-owner input with the same received-line owner retained.
\end{definition}

\begin{theorem}[exact conditional adjacent-row closure]\label{thm:conditional-final}
Fix one of the two deployed MCA rows.  Assume the following three inputs.
\begin{description}[leftmargin=3.4em,labelwidth=2.6em,style=nextline]
\item[\textup{(S)}] The deployed spread-abundance theorem, \cref{prob:spread-abundance}.
\item[\textup{(A)}] The exclusive large-owner atom image bound, \cref{prob:large-owner}.
\item[\textup{(E)}] Exception routing in the sense of \cref{def:exception-routing}.
\end{description}
Then
\[
 \boxed{\Bbad(C,m)\le B^*.}
\]
Together with the unconditional CAP25 v13.2 witness at agreement $m-1$, this proves that the row's exact adjacent MCA threshold is $m$.
\end{theorem}

\begin{proof}
Let $Z$ be the set of bad challenge slopes.  If $|Z|\le31$, the claim is immediate.  If $Z$ contains a primitive spread core, assumption (S) pays the whole line.

Otherwise remove the exceptional set $E$ supplied by exception routing.  Every $32$-subset of $Z\setminus E$ is classified by \cref{thm:partial-relative} as a global affine codeword line or one primitive rational atom.  The overlap graph of $32$-subsets is connected, and \cref{thm:collision-rigidity} forces one coherent global structure.

If that structure is a global affine codeword line, the standard coordinate-injection argument gives at most $n-m+1$ slopes.  Adding $|E|\le31$ remains far below either budget.

Suppose instead that it is one rational atom with owner-set size $g$.  For $g<m$, \cref{thm:owner-localization} gives at most $n-m+1$ atom slopes.  For $m\le g\le2m-k$, \cref{cor:small-owner} gives at most $n$.  Again the additional $31$ slopes fit both budgets.  Finally, for $g>2m-k$, there are at most $n-g$ crossing slopes and assumption (A), namely \eqref{eq:owner-target}, pays the contained slopes; adding the reserved $31$ gives exactly at most $B^*$.

CAP25 v13.2 already proves an unsafe witness at $m-1$, so safety at $m$ identifies the adjacent threshold.
\end{proof}

\section{Core interpolant and correction setup}

Let
\[
 C=\RS[\F,D,k],\qquad |D|=n,
\]
and let
\[
 r_\gamma=r_0+\gamma r_1\in\F^D
\]
be support-wise MCA-bad at agreement $m>k$.  Put
\[
 d=m-k,\qquad R=n-k,\qquad t=n-m.
\]
Fix one first-match primitive spread core in the sense of \cref{prob:spread-abundance}
\[
 (\gamma_i,h_i,S_i)_{i=1}^{32},
\]
where the slopes are distinct, $h_i\in\F[X]_{<k}$, and $S_i$ is the \emph{maximal} agreement support of $h_i$ with $r_{\gamma_i}$.  The common-support branch has already been removed, so
\begin{equation}\label{eq:no-common}
 \bigcap_{i=1}^{32}S_i=\varnothing.
\end{equation}

There is a unique coefficientwise interpolant
\begin{equation}\label{eq:H}
 H(X,Z)=\sum_{j=0}^{31}H_j(X)Z^j,
 \qquad H(X,\gamma_i)=h_i(X),
 \qquad \deg H_j<k.
\end{equation}
For $x\in D$ define
\begin{equation}\label{eq:Ex}
 E_x(Z)=H(x,Z)-r_0(x)-Zr_1(x)
\end{equation}
and
\[
 q_x=|\{i:x\in S_i\}|.
\]
Every core agreement contributes a root of $E_x$, so
\begin{equation}\label{eq:qsum}
 \sum_{x\in D}q_x=\sum_{i=1}^{32}|S_i|\ge32m.
\end{equation}
Moreover, \eqref{eq:no-common} and maximality imply
\begin{equation}\label{eq:nonzero-E}
 E_x\ne0\quad\text{for every }x\in D.
\end{equation}
Indeed, $E_x\equiv0$ would put $x$ in every maximal core support.

The deployed values are
\[
\begin{array}{c|r|r|r|r}
\toprule
\text{row}&m&d&n-m&B^*\\
\midrule
\text{KoalaBear MCA}&1116048&67472&981104&274980728111395087\\
\text{Mersenne-31 MCA}&1116024&67448&981128&16777215\\
\bottomrule
\end{array}
\]
with $n=2^{21}$ and $k=2^{20}$.

\section{Only fifty-eight slopes can use the core interpolant}

\begin{theorem}[core-compatible abundance bound]\label{thm:core-compatible}
Let $Z_0$ be the set of slopes $\gamma\in\F$ for which the codeword $H(X,\gamma)$ agrees with $r_\gamma$ on at least $m$ coordinates.  Then
\begin{equation}\label{eq:58}
 \boxed{
 |Z_0|\le\floor{\frac{31n}{m}}=58.
 }
\end{equation}
More precisely, outside the $32$ core slopes there are at most
\begin{equation}\label{eq:26}
 \boxed{
 \floor{\frac{31n-32m}{m}}=26
 }
\end{equation}
additional core-compatible slopes on either deployed row.
\end{theorem}

\begin{proof}
By \eqref{eq:nonzero-E}, every $E_x$ has at most $31$ roots.  A slope in $Z_0$ contributes at least $m$ coordinate roots, hence
\[
 |Z_0|m\le\sum_{x\in D}\deg E_x\le31n,
\]
which proves \eqref{eq:58}.

For a non-core slope, every agreement coordinate consumes a root distinct from the $q_x$ core roots already present at $x$.  The remaining root capacity is at most
\[
 \sum_x(31-q_x)\le31n-32m
\]
by \eqref{eq:qsum}.  Dividing by $m$ gives \eqref{eq:26}.
\end{proof}

\begin{remark}
The theorem uses maximal supports only to exclude identically zero coordinate error polynomials.  It is independent of quotient, rational-owner, or minimum-distance estimates.  Thus any over-budget spread line must obtain essentially all of its extra slopes from nonzero codeword corrections to $H$.
\end{remark}

\section{An unconditional theorem for one complete correction ray}

Fix a nonzero $P\in\F[X]_{<k}$.  Consider all parameter pairs $(\gamma,c)$ for which
\begin{equation}\label{eq:rayword}
 h_{\gamma,c}(X)=H(X,\gamma)+cP(X)
\end{equation}
agrees with $r_\gamma$ on at least $m$ coordinates.  At $x\in D$ the agreement equation is
\begin{equation}\label{eq:raycurve}
 E_x(\gamma)+cP(x)=0.
\end{equation}
No component-freeness assumption will be imposed.

For $P(x)\ne0$, put
\[
 f_x(Z)=\frac{E_x(Z)}{P(x)}\in\F[Z]_{\le31}.
\]
Coordinates with the same $f_x$ form a \emph{clone class}.  Coordinates with $P(x)=0$ form the vertical class.

\begin{lemma}[clone-class dichotomy]\label{lem:clone}
Let $C_f=\{x:P(x)\ne0,\ f_x=f\}$.  Then exactly one of the following holds.
\begin{enumerate}[label=\textup{(\roman*)}]
\item $|C_f|\le k-1$;
\item the codeword curve
\[
 G_f(X,Z)=H(X,Z)-f(Z)P(X)
\]
has slope degree at most one.  In this case the graph $c=-f(\gamma)$ is a global affine codeword line.
\end{enumerate}
At most two clone classes of type \textup{(ii)} can have size at least $k$.
\end{lemma}

\begin{proof}
For every $x\in C_f$,
\[
 G_f(x,Z)=r_0(x)+Zr_1(x)
\]
as a polynomial identity in $Z$.  If $G_f$ has a nonzero coefficient of $Z^j$ for some $j\ge2$, that coefficient is a nonzero degree-$<k$ polynomial vanishing on $C_f$, so $|C_f|\le k-1$.  Otherwise $G_f=A_f+ZB_f$ is a global affine codeword line.

Clone classes are disjoint subsets of $D$.  A class of type (ii) not covered by (i) has at least $k$ coordinates.  Since $n=2k$, there can be at most two.
\end{proof}

\begin{lemma}[global affine-line block]\label{lem:affine-block}
For any fixed global codeword line $A+\gamma B$, the number of slopes for which it agrees with $r_\gamma$ on at least $m$ coordinates is at most
\begin{equation}\label{eq:affinecost}
 n-m+1.
\end{equation}
\end{lemma}

\begin{proof}
Let $C_0$ be the coordinate set on which $A=r_0$ and $B=r_1$.  Support-wise MCA-badness gives $|C_0|<m$.  Outside $C_0$, agreement of the two affine lines determines at most one slope.  Hence
\[
 N(m-|C_0|)\le n-|C_0|,
\]
and the right quotient is maximized at $|C_0|=m-1$.
\end{proof}

\begin{theorem}[complete correction-ray bound]\label{thm:ray}
For either deployed row, the number of rich parameter pairs $(\gamma,c)$ satisfying \eqref{eq:rayword} is at most
\begin{equation}\label{eq:raybound}
 \boxed{
 2(n-m+1)
 +\floor{\frac{31\binom n2}{(k-1)(m-k+1)}}.
 }
\end{equation}
The final fraction equals $963$ on both rows.  Thus one full correction ray contributes at most
\[
\begin{array}{c|r|r}
\toprule
\text{row}&\text{ray bound}&B^*-\text{ray bound}\\
\midrule
\text{KoalaBear MCA}&1963173&274980728109431914\\
\text{Mersenne-31 MCA}&1963221&14813994\\
\bottomrule
\end{array}
\]
In particular, a single correction ray can never be the residual spread-abundance obstruction.
\end{theorem}

\begin{proof}
Remove all rich points lying on a type-(ii) clone graph whose class has at least $k$ coordinates.  By \cref{lem:clone,lem:affine-block}, there are at most two such graphs and each contributes at most $n-m+1$ slopes.

Consider a remaining rich point and choose exactly $m$ of its agreement coordinates.  Partition them into the vertical class and the clone classes meeting the point.  The vertical class has size at most $k-1$, because a nonzero degree-$<k$ polynomial has at most $k-1$ zeros.  Every clone class present at the point has size at most $k-1$: type-(i) classes have this bound, while every larger type-(ii) graph was removed.

Since $m<2(k-1)$ on both rows, a partition of $m$ elements into parts of size at most $k-1$ has at least
\begin{equation}\label{eq:crosspairs}
 (k-1)(m-k+1)
\end{equation}
heterogeneous unordered pairs.  This is the minimum obtained from the two-part partition $(k-1,m-k+1)$.

A heterogeneous pair contributes to at most $31$ rich parameter points.  For two graph coordinates from different clone classes, the condition is
\[
 f_x(\gamma)=f_y(\gamma),
\]
a nonzero polynomial equation of degree at most $31$, after which $c$ is determined.  For one vertical and one graph coordinate, $E_x(\gamma)=0$ has at most $31$ roots and again the graph equation determines $c$.  We never count vertical--vertical pairs or pairs inside one clone class.

Double counting rich-point/heterogeneous-pair incidences gives
\[
 N_{\rm res}(k-1)(m-k+1)\le31\binom n2.
\]
Adding the at most two affine-line blocks proves \eqref{eq:raybound}.  Counting parameter pairs only overestimates the number of distinct slopes.
\end{proof}

\begin{remark}[minimum-support refinement]\label{rem:a-refine}
Suppose every nonzero codeword in the span of $P$ and the high slope coefficients $H_2,\ldots,H_{31}$ has support at least $R+a$.  Then $k-1$ in the proof can be replaced by $k-a$.  Writing $q=k-a$, the exact cross-pair denominator is
\[
 \Xi(m,q)=\binom m2-u\binom q2-\binom r2,
 \qquad m=uq+r,\quad0\le r<q,
\]
and there are at most $\floor{n/(q+1)}$ large affine clone classes.  In the active residual wedge $m\le2q$, so $\Xi(m,q)=q(m-q)$.  The exact-minimum-distance case $a=1$ is the worst case in the range left by the minimum-distance sharpening of \cref{sec:min-distance-sharpening}; the same section also proves that the boundary top seam occurs only at $a=1$.
\end{remark}

\section{Higher-dimensional correction spaces}

Let $W\le C$ have dimension $s$, and consider all rich points
\[
 (\gamma,P)\in\F\times W
\]
for which $H(X,\gamma)+P(X)$ agrees with $r_\gamma$ on at least $m$ coordinates.  Choose a basis $P_1,\ldots,P_s$ of $W$ and write $P=\sum c_jP_j$.  The coordinate equation
\begin{equation}\label{eq:hypersurface}
 F_x(Z,\mathbf c)=E_x(Z)+\sum_{j=1}^s c_jP_j(x)=0
\end{equation}
has bidegree at most $(31,1)$ on $\mathbb P^1_Z\times\mathbb P^s_{\mathbf c}$.

\subsection{The proper-intersection compiler}

\begin{definition}[proper correction space]
The correction space $W$ is \emph{proper for the core} if every $s+1$ coordinate hypersurfaces from \eqref{eq:hypersurface} have empty or zero-dimensional common intersection in $\mathbb P^1\times\mathbb P^s$.
\end{definition}

\begin{theorem}[proper-intersection correction compiler]\label{thm:proper}
If $W$ is proper for the core, then the number $N_W$ of rich parameter points satisfies
\begin{equation}\label{eq:proper}
 \boxed{
 N_W\le
 \floor{
 31(s+1)\frac{\binom n{s+1}}{\binom m{s+1}}
 }.
 }
\end{equation}
At the deployed rows this pays every proper correction space through
\[
\begin{array}{c|r|r|r}
\toprule
\text{row}&s\le&\text{bound at final }s&\text{bound at }s+1\\
\midrule
\text{KoalaBear MCA}&50&148068539552473273&283695391248936790\\
\text{Mersenne-31 MCA}&15&11989622&23938324\\
\bottomrule
\end{array}
\]
The next values only mark the limit of this compiler.
\end{theorem}

\begin{proof}
The divisor class of each projective hypersurface is at most
\[
 31H_Z+H_W.
\]
In the Chow ring of $\mathbb P^1\times\mathbb P^s$,
\[
 (31H_Z+H_W)^{s+1}=31(s+1)H_ZH_W^s,
\]
because $H_Z^2=0$ and $H_W^{s+1}=0$.  Hence every proper intersection of $s+1$ coordinate hypersurfaces has degree at most $31(s+1)$.

A rich point lies on at least $m$ hypersurfaces and therefore contributes at least $\binom m{s+1}$ incidences with coordinate $(s+1)$-subsets.  Summing first over rich points and then over the $\binom n{s+1}$ coordinate subsets proves \eqref{eq:proper}.
\end{proof}

\begin{lemma}[what a nonproper tuple means]\label{lem:nonproper}
Let $T$ be a set of $s+1$ coordinates.
\begin{enumerate}[label=\textup{(\roman*)}]
\item If the evaluation functionals of $T$ on $W$ have rank below $s$, then a nonzero word of $W$ vanishes on a subset spanning the same evaluation flat.
\item If their rank is $s$ and the common intersection is positive-dimensional, choose any evaluation basis $B\subset T$.  There is a unique polynomial curve $P_B(Z)\in W\otimes\F[Z]_{\le31}$ solving the equations on $B$, and every equation in $T$ vanishes identically on that curve.
\end{enumerate}
Thus every failure of properness is either an evaluation rank-flat or an exact polynomial clone component.
\end{lemma}

\begin{proof}
Part (i) is the annihilator statement for the column matroid of the evaluation map.  For (ii), the constant evaluation matrix on $B$ is invertible, so solving its equations gives $P_B(Z)$ coefficientwise and with degree at most $31$.  Substitution into the remaining equation gives a univariate polynomial.  A positive-dimensional intersection forces that polynomial to vanish identically.
\end{proof}

\subsection{A clone-tolerant basis-curve compiler}

Let
\[
 V=\Span(W,H_2,\ldots,H_{31})\le C
\]
and write
\begin{equation}\label{eq:a}
 d_1(V)=R+a,
 \qquad a\ge1.
\end{equation}
We say that $W$ \emph{absorbs the high core} if $H_j\in W$ for every $j\ge2$.

\begin{theorem}[clone-tolerant basis-curve bound]\label{thm:clone-tolerant}
Assume $s\ge1$ and that $W$ does not absorb the high core.  Put
\begin{equation}\label{eq:Ma}
 M_a=\floor{\frac{31(R+a)}{d+a}}.
\end{equation}
Then
\begin{equation}\label{eq:clone-tolerant}
 \boxed{
 N_W\le
 \floor{
 M_a\frac{\fall ns}{\rise{(d+a)}s}
 }.
 }
\end{equation}
No component-freeness assumption is required.
\end{theorem}

\begin{proof}
Fix a rich point and an exact size-$m$ agreement support $S$.  For every $j$-dimensional subcode $U\le W$, strict increase of generalized Hamming weights and \eqref{eq:a} give
\[
 |\supp_D(U)|\ge R+a+j-1.
\]
At most $t=n-m$ of these support coordinates lie outside $S$, so the restricted code $W|_S$ has
\[
 d_j(W|_S)\ge d+a+j-1.
\]
Greedily choosing independent evaluation columns therefore gives at least
\begin{equation}\label{eq:bases}
 \frac{\rise{(d+a)}s}{s!}
\end{equation}
unordered evaluation bases $B\subset S$.

For each such $B$, solve the basis equations to obtain a unique polynomial correction curve
\[
 P_B(Z)\in W\otimes\F[Z]_{\le31}.
\]
Let $C_B$ be the coordinates whose equations vanish identically after this substitution, and put $c_B=|C_B|$.  Since $W$ does not absorb all high core coefficients, the codeword curve $H+P_B$ has a nonzero coefficient of $Z^j$ for some $j\ge2$.  That coefficient lies in $V$, so it has at most $k-a$ zeros.  Hence
\begin{equation}\label{eq:clone-size}
 c_B\le k-a.
\end{equation}

Every rich point on the curve $P=P_B(\gamma)$ agrees automatically on $C_B$ and must obtain at least $m-c_B$ further agreements outside $C_B$.  Each outside coordinate supplies a nonzero polynomial equation of degree at most $31$.  Root-capacity counting on this one curve gives
\[
 N(B)(m-c_B)\le31(n-c_B).
\]
The quotient $(n-c)/(m-c)$ is increasing for $c<m$, and \eqref{eq:clone-size} yields
\[
 N(B)\le
 \floor{\frac{31(R+a)}{d+a}}=M_a.
\]

Each rich point is incident to at least \eqref{eq:bases} basis curves.  Globally there are at most $\binom ns$ evaluation bases, and each basis curve contains at most $M_a$ rich points.  Therefore
\[
 N_W\frac{\rise{(d+a)}s}{s!}
 \le M_a\binom ns
 =M_a\frac{\fall ns}{s!},
\]
which is \eqref{eq:clone-tolerant}.
\end{proof}

\begin{corollary}[exact deployed transitions]\label{cor:transitions}
At exact minimum support $a=1$, \cref{thm:clone-tolerant} pays through
\[
\begin{array}{c|r|r|r}
\toprule
\text{row}&s\le&\text{bound at final }s&\text{bound at }s+1\\
\midrule
\text{KoalaBear MCA}&9&13013823503882165&404431535289439486\\
\text{Mersenne-31 MCA}&3&14457337&449492838\\
\bottomrule
\end{array}
\]
The actual minimum-support excess enlarges the range.  The first exact paid transitions beyond those dimensions are
\[
\begin{array}{c|r|r|r|r}
\toprule
\text{row}&s&a_{\rm paid}&\text{bound at }a_{\rm paid}&\text{bound at }a_{\rm paid}-1\\
\midrule
\text{KoalaBear MCA}&10&2434&274948660791683460&274987992550612247\\
&11&25508&274953206296382843&274985734988552690\\
&12&50875&274977972290567960&275005854396152321\\
\text{Mersenne-31 MCA}&4&64258&16777099&16777608\\
\bottomrule
\end{array}
\]
All comparisons are exact integer comparisons with the row budgets.
\end{corollary}

\begin{remark}[relation to the minimum-distance wedge]
\Cref{prop:min-support-transitions} proves that the first-unpaid charts already satisfy $a\le70831$ on KoalaBear and $a\le111528$ on Mersenne-31 at zero outside mismatch, with positive mismatch shrinking the residual further.  It also confines the boundary top seam to $a=1$.  Thus \cref{cor:transitions} is not merely an abstract monotonicity statement: it removes correction dimensions $10$--$12$ on substantial parts of the KoalaBear wedge and dimension $4$ on a substantial part of the Mersenne wedge, while the genuine top-seam endpoint remains the exact-minimum-distance case.
\end{remark}

\section{A rank-eleven dense-locator component route}
\label{sec:rank11-dense-locator-route}

We record a local theorem packet for the post-near rank-eleven route.  Its
input is the following pinned setup, supplied by the public prize-DAG source
listed in the companion certificate.  Let
\[
 R=1048576,\qquad d=67472,\qquad
 n'=R+K',\qquad m'=d+K',\qquad 10\le K'\le1048576.
\]
There is a fixed ten-dimensional correction space $V'$, a degree-at-most
$31$ interpolant $H(X,Z)$, and one affine pair $(a'_0,b'_0)$.  The deviation
\[
 D_H(X,Z)=H(X,Z)-a'_0(X)-Zb'_0(X)
\]
vanishes at eighteen distinct dense-pair anchor slopes, while ten other
values of $D_H$ span $V'$.  Every coefficient of $D_H$ lies in $V'$.  The
post-near family has distinct slopes and, if the original line is unsafe,
contains at least
\[
 N_{\min}=274980728111260126
\]
non-dense records after the exact $134944$ near charge and the eighteen
dense anchors.  These hypotheses are the scope of this section; no statement
here silently promotes an arbitrary rank-eleven family into this setup.

\begin{proposition}[Dense-root high-coefficient saturation]
\label{prop:kb-rank11-dense-root-saturation}
Let
\[
 q(Z)=\prod_{i=1}^{18}(Z-\gamma_i)
\]
be the monic locator of the dense anchor slopes.  Then
\[
 D_H=qG,\qquad \deg_ZG\le13,
 \qquad
 \Span\{H_j:j\ge2\}=V'.
\]
Consequently, if a correction space $W\le V'$ absorbs every coefficient
$H_j$ with $j\ge2$, then $W=V'$ and $\dim W=10$.
\end{proposition}

\begin{proof}
The eighteen vector-valued zeroes give the coefficientwise divisibility
$q\mid D_H$.  Write $q=Z^{18}+q_{17}Z^{17}+\cdots+q_0$ and
$G=\sum_{j=0}^{13}G_jZ^j$.  Descending from degree $31$ to degree $18$, the
coefficients of $D_H$ recover $G_{13},G_{12},\ldots,G_0$ by a triangular
system with diagonal one.  Hence
\[
 \Span\{(D_H)_{18},\ldots,(D_H)_{31}\}
 =\Span\{G_0,\ldots,G_{13}\}.
\]
The ten nonzero anchor values span $V'$, while every coefficient of $D_H$
lies in $V'$, so both displayed spans equal $V'$.  Subtracting the affine
pair changes only coefficients zero and one, hence $(D_H)_j=H_j$ for
$j\ge2$.  If $W$ absorbs those coefficients, then $V'\le W\le V'$.
\end{proof}

For a non-dense record of slope $\gamma$, put
\[
 R_\gamma=
 \frac{h'_\gamma-a'_0-\gamma b'_0}{q(\gamma)}\in V'.
\]
This is well-defined because the record slopes are distinct from the roots
of $q$.  At a coordinate $x$, rich agreement is the equation
\begin{equation}
 (a'_0-r'_0)(x)+Z(b'_0-r'_1)(x)+q(Z)R(x)=0.
 \label{eq:kb-rank11-dense-locator-coordinate}
\end{equation}

\begin{proposition}[Dense-locator component-incidence bound]
\label{prop:kb-rank11-component-incidence}
For every eleven-coordinate tuple, the isolated solutions of
\eqref{eq:kb-rank11-dense-locator-coordinate} in
$\mathbf P^1_Z\times\mathbf P^{10}_R$, counted with local multiplicity, have
total multiplicity at most $198$.  Uniformly in $K'$,
\[
 A_{\rm iso}(K')=
 \ceil{198\frac{\binom{n'}{11}}{\binom{m'}{11}}}
 \le2526815879272440.
\]
Thus at least $990810934$ parts per billion of all pairs consisting of a
non-dense record and an eleven-subset of its support lie on a
positive-dimensional component through that record.

At such a point, the component is either a full-evaluation-rank affine-owner
component or a rank-deficient kernel component.  At least one of these two
lanes carries $495405467$ parts per billion of all incidences.
\end{proposition}

\begin{proof}
After bihomogenization, each coordinate equation has divisor class at most
$18H_Z+H_R$.  Generic independent perturbations in these basepoint-free
classes make eleven equations intersect properly.  Every isolated point of
the original intersection persists with at least its local intersection
multiplicity, while the total generic intersection number is
\[
 (18H_Z+H_R)^{11}=18\cdot11H_ZH_R^{10}=198H_ZH_R^{10}.
\]
Summing this bound over all eleven-coordinate tuples and dividing by the
$\binom{m'}{11}$ tuples contributed by each record gives $A_{\rm iso}(K')$.
The binomial ratio is the product, for $0\le i\le10$, of
\[
 \frac{1048576+K'-i}{67472+K'-i}.
\]
Every factor decreases with $K'$, so the maximum occurs at $K'=10$ and the
displayed integer follows by exact evaluation.  Dividing it by
$N_{\min}$ gives an isolated-incidence ceiling of $9189066$ parts per
billion, leaving the stated component floor.

For the classification, inspect evaluation of $V'$ on the eleven-tuple.  At
rank ten, ten independent coordinate equations determine one rational
correction curve; an extra equation vanishes identically precisely on an
affine-owner component.  At rank at most nine, the compatible fiber contains
a positive-dimensional evaluation kernel.  These alternatives are
exhaustive, and halving the total component floor gives the one-lane floor.
\end{proof}

\begin{proposition}[Rank-stratified isolated-record bound]
\label{prop:kb-rank11-rank-stratified-isolated-incidence}
Fix an eleven-set $T$ and consider only retained record points
$(\gamma,R_\gamma)$ satisfying
\eqref{eq:kb-rank11-dense-locator-coordinate} on every coordinate of $T$.
At most one such point is isolated.  Consequently, if $N$ retained records
have exact support size $m'$ in a residual domain of size $n'$, at least
\[
 N\binom{m'}{11}-\binom{n'}{11}
\]
record/eleven-set incidences lie on positive-dimensional components.
This counts actual retained records and does not replace the valid generic
multiplicity bound in Proposition~\ref{prop:kb-rank11-component-incidence}.
\end{proposition}

\begin{proof}
Write $a_x=(a'_0-r'_0)(x)$ and $b_x=(b'_0-r'_1)(x)$.
Put $r=\rank(\operatorname{ev}_T:V'\to\F^T)$.  If $r\le9$, choose
$0\ne u\in\ker(\operatorname{ev}_T)$.  Every compatible point then lies on
the affine kernel line $(\gamma,R_\gamma+tu)$ and is not isolated.

Suppose $r=10$, and choose a ten-set $B\subset T$ on which evaluation is an
isomorphism.  There are unique $U_0,U_1\in V'$ whose values on $B$ cancel
the constant and linear terms in
\eqref{eq:kb-rank11-dense-locator-coordinate}.  On $q(Z)\ne0$, the ten
equations indexed by $B$ are equivalent to
\[
 q(Z)R=U_0+ZU_1.
\]
Substitution in the remaining coordinate $y\in T\setminus B$ leaves one
equation
\[
 (a_y+U_0(y))+Z(b_y+U_1(y))=0.
\]
If it is an identity, every retained solution belongs to a one-parameter
affine-owner component.  Otherwise it has at most one root.  Retained slopes
are distinct and avoid the roots of $q$, so at most one retained record point
is isolated on $T$.  Summing over the $\binom{n'}{11}$ choices of $T$ gives
the assertion.
\end{proof}

\begin{corollary}[Component-star owner-pencil trichotomy]
\label{cor:kb-rank11-component-star}
At least $540546700$ parts per billion of the non-dense records have at
least $98$ percent of their eleven-subsets on a positive-dimensional
component.  In particular, there are at least
\[
 148639925144138894
\]
such records.  For each one there is a ten-subset $B$ of its support with at
least
\[
 E(K')=\ceil{\frac{98(m'-10)}{100}}
\]
component extensions, and exactly one of the following holds:
\begin{enumerate}[label=\textup{(\roman*)}]
\item $\operatorname{rank}(\operatorname{ev}_B)=10$, and one affine owner
has support deficiency at most $22320$;
\item $\operatorname{rank}(\operatorname{ev}_B)=9$, and at least $45153$
full-rank extensions lie on one owner pencil with direction
$\beta(-\gamma u,u)$, where $\ker(\operatorname{ev}_B)=\F u$;
\item $\operatorname{rank}(\operatorname{ev}_B)\le8$, and the evaluation
kernel has dimension at least two.
\end{enumerate}
\end{corollary}

\begin{proof}
If $\alpha$ is the fraction of records above the $98$-percent threshold,
then Proposition~\ref{prop:kb-rank11-component-incidence} gives
\[
 \frac{990810934}{10^9}\le
 \alpha+(1-\alpha)\frac{98}{100},
\]
so $\alpha\ge540546700/10^9$.  Multiplication by $N_{\min}$ and integer
rounding gives the absolute floor.  Double-counting pairs $(B,x)$ inside one
support gives a ten-subset with $E(K')$ extensions.

At full rank those equations determine one owner.  At rank nine, a kernel
word $u$ already has ten roots and has at most $K'-11$ further roots, leaving
at least
\[
 E(K')-(K'-11)\ge45153
\]
full-rank extensions.  Any two corresponding owners differ by
$\beta(-\gamma u,u)$.  Rank at most eight gives kernel dimension at least
two by rank--nullity.  The owner deficiency in the full-rank case is at most
$m'-10-E(K')\le22320$.
\end{proof}

\begin{proposition}[Record-intrinsic full-rank star owner]
\label{prop:kb-rank11-component-star-owner-unique}
In the full-rank branch of
Corollary~\ref{cor:kb-rank11-component-star}, the owner of support deficiency
at most $22320$ is unique for its record.  In particular it is independent
of the ten-subset used to expose it.
\end{proposition}

\begin{proof}
Suppose two owner pairs have within-support cores $A_p,A_q$, each of size at
least $m'-22320$.  Inclusion--exclusion inside the exact support gives
\[
 |A_p\cap A_q|\ge m'-2\cdot22320=K'+22832.
\]
If the pairs are distinct, at least one component of their difference is a
nonzero RS polynomial of degree below $K'$.  Every coordinate in
$A_p\cap A_q$ is a root of that component, so the RS root bound gives
$|A_p\cap A_q|\le K'-1$, a contradiction.  The gap is uniformly $22833$.
\end{proof}

\begin{proposition}[One fixed rank-nine split-pencil cell]
\label{prop:kb-rank11-split-pencil-cell}
Fix a rank-nine ten-subset $B$ and let $g$ records be routed to this cell.
All owners agreeing with the received pair on $B$ form an affine plane
$\Pi_B$ with directions $(u,0)$ and $(0,u)$.  Each record slope $\gamma$
defines a line of direction $(-\gamma u,u)$ in $\Pi_B$.  If $t_p$ is the
number of record lines through owner point $p$, then
\[
 \sum_p\binom{t_p}{2}=\binom g2.
\]
Moreover, if $J_B$ is the common owner-plane core on $Z(u)$ and
$C_p=J_B\mathbin{\dot\cup}P_p$, then the petals $P_p$ are pairwise disjoint
and
\[
 45153g\le\sum_pt_p|P_p|
 \le981105(2097152-10)=2057516501910.
\]
Consequently
\[
 \boxed{g\le
 \floor{2057516501910/45153}=45567658}.
\]
\end{proposition}

\begin{proof}
All agreeing owner pairs differ by $(\alpha u,\beta u)$.  Ownership at slope
$\gamma$ cuts out one affine line with direction $(-\gamma u,u)$, and
distinct slopes give distinct directions.  Every pair of record lines has
one intersection point, proving the block-design identity.

On $Z(u)$ every point of $\Pi_B$ has the same pair value.  Away from $Z(u)$,
evaluation identifies the parameter plane bijectively with the pair-value
plane, so at most one owner agrees at a coordinate.  This proves the common
core and disjoint-petal decomposition.  Every record supplies at least
$45153$ full-rank extension incidences in petals on its line.

For a fixed owner with core size $c_p<m'$, the outside-core agreement sets of
its distinct slopes are disjoint.  Hence
\[
 t_p(m'-c_p)\le n'-c_p,
 \qquad t_p\le n'-m'+1=981105.
\]
Finally $|J_B|\ge10$, so $\sum_p|P_p|\le2097152-10$, giving the displayed
weighted cap and its integer floor quotient by $45153$.
\end{proof}

\begin{remark}[Integer rounding repair]
The public source node recorded the valid weaker bound $45567659$ by taking
a ceiling.  Since $g$ is an integer and $45153g\le2057516501910$, the sharp
consequence is the floor $45567658$ printed above.  Both values and the
rounding direction are retained in the companion certificate.
\end{remark}

\begin{proposition}[Split-pencil pair-core dichotomy]
\label{prop:kb-rank11-split-pencil-paircore}
Lift one fixed residual rank-nine owner plane to the original KoalaBear row.
Let $J$ be the coordinates where every owner in the plane equals the received
pair, let $j=|J|\ge10$, and let $g$ be the number of record lines in the
plane.  Then either
\[
 g\le1434405,
\]
or the entire owner plane has a common received-pair core of size
\[
 j\ge134944.
\]
\end{proposition}

\begin{proof}
Write $C_p$ for the pair core of owner $p$ and
$P_p=C_p\setminus J$.  Proposition~\ref{prop:kb-rank11-split-pencil-cell}
gives pairwise disjoint petals $P_p$ and
$\sum_p\binom{t_p}{2}=\binom g2$.

If two distinct record lines meet at $p$, their selected original supports,
each of size $m=1116048$, intersect in at least
\[
 2m-n=2(1116048)-2097152=134944
\]
coordinates.  The two distinct slope equalities recover both received
columns there, so the intersection lies in $C_p$.

Assume $j\le134943$.  Put $D=n-m=981104$,
$x=m-|C_p|$, and $y=t_p-1$ at a point with $t_p\ge2$.
Pair noncontainment gives $x\ge1$, while fixed-owner exception
disjointness gives $yx\le D$.  Since $m-j\ge D+1$ and
\[
 D+1-x-y=(D-yx)+(y-1)(x-1)\ge0,
\]
we have $|P_p|=|C_p|-j\ge t_p-1$ and $t_p\le D+1$.
Consequently
\[
\begin{split}
 g(g-1)
  &=\sum_p t_p(t_p-1)\\
  &\le981105\sum_p|P_p|\\
  &\le981105(2097152-10)=2057516501910.
\end{split}
\]
Now
\[
 1434405\cdot1434404\le2057516501910
 <1434406\cdot1434405,
\]
so $g\le1434405$.  Thus any larger plane lies in the alternative
$j\ge134944$.
\end{proof}

\begin{proposition}[Fixed nine-subset component target]
\label{prop:kb-rank11-component-ninesubset-target}
One of the affine-owner and rank-deficient component lanes admits a fixed
nine-coordinate subset $B$ carrying at least
\[
 2578110
\]
distinct records in that same lane.  Moreover, one of the following holds.
\begin{enumerate}[label=\textup{(N\arabic*)}]
\item every retained record has a rank-deficient component extension whose
kernel lies in the same nonzero $\ker(\operatorname{ev}_B)$;
\item $\operatorname{rank}(\operatorname{ev}_B)=9$, and after lifting to the
original row the resulting owner plane has a common received-pair core of
size at least $134944$;
\item $\operatorname{rank}(\operatorname{ev}_B)=8$, the retained owners lie
in one affine $U^2$ flat with $\dim U=2$, and the retained errors have affine
rank at most three.
\end{enumerate}
\end{proposition}

\begin{proof}
Choose the heavier component lane.  It carries at least
$\lambda=495405467/10^9$ of all record/eleven-subset incidences.  Mark all
$\binom{11}{9}=55$ nine-subsets of every such eleven-subset.  A fixed
record/nine-subset pair is marked at most $\binom{m'-9}{2}$ times, while
\[
 55\binom{m'}{11}=\binom{m'}9\binom{m'-9}2.
\]
Averaging over the domain nine-subsets gives one $B$ carrying at least
\[
 \ceil{\lambda\,274980728111260126
        \frac{\binom{m'}9}{\binom{n'}9}}
 \ge2578110.
\]
The nine factors in the binomial ratio increase with $K'$, and exact
evaluation at $K'=10$ gives the displayed endpoint.

In the rank-deficient lane, restriction from an extending eleven-subset
puts every component kernel inside the fixed nonzero
$\ker(\operatorname{ev}_B)$, proving (N1).  In the affine-owner lane, an
evaluation-rank-ten eleven-subset loses at most two ranks on restriction to
$B$, so the rank on $B$ is eight or nine.

At rank nine the owners form one affine plane.  The nine-coordinate version
of Proposition~\ref{prop:kb-rank11-split-pencil-paircore} has ordered-pair
resource
\[
 981105(2097152-9)=2057517483015
 <1434406\cdot1434405,
\]
so its low-core cap remains $1434405$.  Since $2578110>1434405$, the
reversible original-row lift forces the shared core in (N2).

At rank eight put $U=\ker(\operatorname{ev}_B)$, so $\dim U=2$, and anchor
one retained owner $(A_*,B_*)$.  Every retained explanation has the form
\[
 h_\gamma=A_*+\gamma B_*+v_\gamma,\qquad v_\gamma\in U.
\]
Hence anchored error differences lie in
$\operatorname{span}(U,r_1-B_*)$, of dimension at most three.  Exact locator
multiplication in the reverse common-core lift is injective, proving (N3).
\end{proof}

\begin{proposition}[A fixed-chart local-cap fence]
\label{prop:kb-rank11-fixed-chart-local-cap-fence}
The local rank-nine conclusions of
Proposition~\ref{prop:kb-rank11-component-ninesubset-target}, considered
without the weighted component incidences of their ancestors, admit an
official-row family of
\[
 4070408>2578110
\]
distinct pair-noncontained exact-support slopes.  The family has common
owner-plane core $K-1$, every record has a rank-ten eleven-subset on an
affine-owner component, and its errors have affine rank at most two.
\end{proposition}

\begin{proof}
Choose $J\subset D$ with $|J|=K-1$ and put
\[
 u(X)=\prod_{x\in J}(X-x).
\]
Fix a nine-subset $B\subset J$.  Nine interpolation polynomials of degree at
most eight together with $u$ span a ten-dimensional RS space whose
evaluation on $B$ has rank nine and kernel $Fu$.

Outside $J$, the coordinate and rich-support weights are
\[
 N=n-(K-1)=1048577,\qquad L=m-(K-1)=67473.
\]
Give eight owner points $P_i=(c+iM,0)$ weight $L-1=67472$, where
$M=508801$, and give the remaining $M$ coordinates weight one at
$Q_j=(j,1)$, $0\le j<M$.  The values
\[
 \gamma_{i,j}=c+iM-j
\]
are distinct because their eight difference intervals are consecutive and
disjoint.  There are $8M=4070408$ of them.  Since the KoalaBear base prime
is greater than $18\cdot4070408$, the common translate $c$ may also avoid
any prescribed eighteen slopes.

Label an outside coordinate by $(\alpha_x,\beta_x)$ and define
\[
 r_0(x)=\alpha_xu(x),\qquad r_1(x)=\beta_xu(x),
\]
with both columns zero on $J$.  At slope $\gamma_{i,j}$ use the codeword
$h_{i,j}=(c+iM)u$.  The line
$\alpha+\gamma_{i,j}\beta=c+iM$ contains exactly $P_i$ and $Q_j$ among the
assigned owner points.  Its support therefore has size
$(K-1)+(L-1)+1=m$.

A containing polynomial pair would have its second component equal to zero
on $J$ and the heavy fibre, altogether $m-1>K-1$ coordinates, and hence
zero identically by the RS root bound.  At $Q_j$ it would instead have to
equal the nonzero value of $u$, a contradiction.  Two coordinates in the
heavy fibre adjoined to $B$ give evaluation rank ten and the fixed-owner
component through $P_i$.  Finally
\[
 e_{i,j}=r_0+\gamma_{i,j}r_1-(c+iM)u,
\]
so all anchored error differences lie in $\operatorname{span}(r_1,u)$.
\end{proof}

\begin{proposition}[Weighted high-row exclusion of the rank-nine fixed target]
\label{prop:kb-rank11-rank9-weighted-exclusion}
In the fixed target of
Proposition~\ref{prop:kb-rank11-component-ninesubset-target}, alternative
\textup{(N2)} is empty for every residual dimension
$20618\le K'\le1048576$.  Before deduplication, one fixed nine-subset carries
marked component-extension weight at least
\[
 \ceil{\frac{495405467}{10^9}\,274980728111260126
        \frac{\binom{m'}9\binom{m'-9}2}{\binom{n'}9}}.
 \tag{W9}\label{eq:kb-rank11-weighted-nine-target}
\]
Every rank-nine chart has marked weight at most
\[
 981105(m'-10)n'.
 \tag{C9}\label{eq:kb-rank11-weighted-nine-cap}
\]
At the exact first weighted crossing the adjacent integer bounds are
\[
 \begin{array}{c|rr}
 K'&\text{demand}&\text{cap}\\ \hline
 20617&92386821615379573&92394042904582935\\
 20618&92397581841774591&92395178310909600.
 \end{array}
\]
Consequently the fixed target reduces to \textup{(N1)} or \textup{(N3)}
only from $K'=20618$ onward.  The rank-nine branch remains open below that
row.
\end{proposition}

\begin{proof}
The marking argument in the proof of
Proposition~\ref{prop:kb-rank11-component-ninesubset-target} produces
$55I$ marked triples $(\gamma,T,B)$.  Average over $B$ without dividing by
$\binom{m'-9}{2}$ and use
\[
 55\binom{m'}{11}=\binom{m'}9\binom{m'-9}2
\]
to obtain \eqref{eq:kb-rank11-weighted-nine-target}.  At $K'=10$ its
right side is $5868470021012020$; dividing by the maximum extension
multiplicity recovers the $2578110$ distinct records used above.

Now fix a rank-nine chart and let $u$ span its evaluation kernel on $B$.
Write $C_p=J\sqcup P_p$ for the pair core of an owner point in the affine
owner plane.  The petals $P_p$ are pairwise disjoint.  A rank-ten extension
$T=B\cup\{x,y\}$ determines one owner point $p$, both added coordinates
lie in $C_p$, and at least one lies outside $Z(u)$, hence in $P_p$.
Therefore one owner point contributes at most
\[
 |P_p|\,|C_p\setminus B|\le |P_p|(m'-10)
\]
extensions per owned record.  Pair noncontainment gives $|C_p|<m'$, while
fixed-owner exception disjointness gives at most
$n'-m'+1=981105$ owned records.  Summing the disjoint petals proves
\eqref{eq:kb-rank11-weighted-nine-cap}.

Exact cross-multiplication is negative at $K'=20617$ and positive at
$K'=20618$; after rounding the first strict gap is $2403530864991$.
After cancelling
$\binom{m'-9}{2}=(m'-9)(m'-10)/2$, the ratio of the unrounded lower bound
to \eqref{eq:kb-rank11-weighted-nine-cap} is a positive constant times
\[
 \frac{\binom{m'}9}{\binom{n'}9}\frac{m'-9}{n'}.
\]
Every factor increases with $K'$ because $n'-m'=981104$.  The boundary
contradiction therefore persists through the deployed endpoint.

The former low-row argument applied the original-row intersection floor
$2m-n=134944$ after reversing the shortening and then compared that lifted
core with residual $m'$.  This mixes rows: the reverse lift already inserts
$1048576-K'$ deleted locator coordinates into every owner core.  Hence an
original-row core of size $134944$ need not contain any residual core of
that size.  No conclusion for $10\le K'\le20617$ is claimed here.
\end{proof}

\begin{proposition}[Residual petal-pair refinement of the rank-nine target]
\label{prop:kb-rank11-rank9-residual-petal-capacity}
In the fixed target of
Proposition~\ref{prop:kb-rank11-component-ninesubset-target}, alternative
\textup{(N2)} is empty for every
\[
 15635\le K'\le20617.
\]
More precisely, let $J$ be the common core of one residual rank-nine owner
plane and put $j=|J|$. Then $9\le j\le K'-1$, and its marked
component-extension weight satisfies
\[
 W_B\le
 \floor{\frac{981105(n'-j)(m'+j-20)}2}.
 \tag{RP9}\label{eq:kb-rank11-residual-petal-cap}
\]
The exact adjacent integer bounds at the first crossing are
\[
 \begin{array}{c|rr}
 K'&\text{demand}&\text{cap}\\ \hline
 15634&50777401704768572&50779283449126807\\
 15635&50783693985583057&50780312213264392.
 \end{array}
\]
Together with
Proposition~\ref{prop:kb-rank11-rank9-weighted-exclusion}, this removes the
rank-nine target for every $K'\ge15635$.  It remains open for
$10\le K'\le15634$.
\end{proposition}

\begin{proof}
Let $u$ span the evaluation kernel on the fixed nine-set $B$.  Owners
agreeing on $B$ form one affine plane.  If $J$ is the set of coordinates
where the complete plane agrees with the received pair, then
$B\subseteq J\subseteq Z(u)$.  Since $u$ is nonzero of degree below $K'$,
\[
 9\le j\le K'-1.                                      \tag{1}
\]
Write the core of owner $p$ as $C_p=J\sqcup P_p$ and put
$s_p=|P_p|$.  The petals $P_p$ are pairwise disjoint, and support-wise
pair noncontainment gives $s_p\le m'-j-1$.

A marked extension $B\cup\{x,y\}$ has at least one coordinate outside
$Z(u)$, hence in $P_p$ for its owner.  The number of available unordered
pairs for $p$ is therefore at most
\[
 q_p=s_p(j-9)+\binom{s_p}{2}.
\]
Owners with $s_p=0$ contribute no marked extension.  For every remaining
owner,
\[
 \frac{q_p}{s_p}
 \le j-9+\frac{m'-j-2}{2}
 =\frac{m'+j-20}{2}.                                 \tag{2}
\]
At most $981105$ selected records have the same owner.  Summing (2) and
using $\sum_p s_p\le n'-j$ proves \eqref{eq:kb-rank11-residual-petal-cap};
the floor is valid because $W_B$ is integral.

The unfloored envelope $F(j)=(n'-j)(m'+j-20)$ has forward difference
\[
 F(j+1)-F(j)=981104-2j+19>0
\]
through the claimed interval.  Its maximum is therefore at $j=K'-1$.
Exact cross-multiplication against \eqref{eq:kb-rank11-weighted-nine-target}
is negative at $K'=15634$ and positive at $K'=15635$.  The ratio of the
lower bound to the unfloored cap is a positive constant times nine factors
$(m'-i)/(n'-i)$, $0\le i\le8$, and
\[
 \frac{(K'+67463)(K'+67462)}{2K'+67451}.
\]
The nine factors increase because $n'-m'=981104$.  After cancelling the
positive common factor $K'+67463$, the last factor has forward
cross-difference $2(K'-11)>0$.  Thus the contradiction persists through
$K'=20617$.  No original-row $134944$ core floor is used.
\end{proof}

\begin{proposition}[Exact residual petal-partition refinement]
\label{prop:kb-rank11-rank9-exact-petal-partition}
In the fixed target of
Proposition~\ref{prop:kb-rank11-component-ninesubset-target}, alternative
\textup{(N2)} is empty for every
\[
 15529\le K'\le15634.
\]
More precisely, exact integer packing of the disjoint residual owner petals
gives
\begin{equation}
 W_B\le981105\bigl(1048577K'+34798536326\bigr).
 \label{eq:kb-rank11-exact-petal-partition-cap}
\end{equation}
The exact adjacent integer bounds at the first crossing are
\[
 \begin{array}{c|rr}
 K'&\text{demand}&\text{cap}\\ \hline
 15528&50114371326035640&50115667510540110\\
 15529&50120589875892136&50116696274677695.
 \end{array}
\]
Together with Propositions~\ref{prop:kb-rank11-rank9-residual-petal-capacity}
and~\ref{prop:kb-rank11-rank9-weighted-exclusion}, this removes the
rank-nine target for every $K'\ge15529$.  It remains open for
$10\le K'\le15528$.
\end{proposition}

\begin{proof}
Retain the notation of
Proposition~\ref{prop:kb-rank11-rank9-residual-petal-capacity} and put
\[
 D_1=981105,\qquad a=m'-j-1.
\]
Then $67472\le a\le67462+K'$, every petal size satisfies $0\le s_p\le a$,
and $\sum_p s_p\le D_1+a$.  For
$q_j(s)=s(j-9)+\binom{s}{2}$, the two convex transfers are
\[
 q_j(x+y)-q_j(x)-q_j(y)=xy\quad(x+y\le a)
\]
and
\[
 q_j(a)+q_j(x+y-a)-q_j(x)-q_j(y)=(a-x)(a-y)
 \quad(x,y<a<x+y).
\]
Both are nonnegative.  Saturating the relaxed total budget and repeating
these transfers leaves
\[
 r=1+\floor{D_1/a}
\]
full petals and one remainder $b=D_1\bmod a$.  Thus
\[
 \sum_p q_j(s_p)\le r q_j(a)+q_j(b).                \tag{1}
\]
This remains an upper bound if the affine owner plane has fewer petal
slots, since (1) optimizes over a larger partition class.

Substitute $j=K'+67471-a$.  For fixed $a$, the right side of (1) is the
line
\[
 G_a(K')=(D_1+a)K'+(D_1+a)(67462-a)
 +\frac{r a(a-1)+b(b-1)}2.                          \tag{2}
\]
On $10\le K'\le15634$, the quotient $\floor{D_1/a}$ has four blocks:
\[
 \begin{array}{c|c}
 14&67472\le a\le70078\\
 13&70079\le a\le75469\\
 12&75470\le a\le81758\\
 11&81759\le a\le83096.
 \end{array}
\]
Within a block of quotient $q$, (2) is a convex quadratic in $a$ with
leading coefficient $(q^2+q-1)/2$.  It therefore suffices to compare the
eight block endpoints.  At $K'=15634$, their gaps below $G_{67472}$ are
\[
 \begin{array}{c|rrrrr}
 a&67472&70078&70079&75469&75470\\ \hline
 G_{67472}-G_a&0&676268727&676325879&3265037774&3265407519
 \end{array}
\]
and
\[
 \begin{array}{c|rrr}
 a&81758&81759&83096\\ \hline
 G_{67472}-G_a&6322001175&6322245154&7515065748.
 \end{array}
\]
Every competing line has larger slope than $G_{67472}$, so these
nonnegative comparisons at the largest row imply the same comparisons at
every earlier admissible row.  For $a=67472$, there are fifteen full
petals and remainder $36497$, and
\[
 G_{67472}(K')=1048577K'+34798536326.
\]
Multiplying by the fixed-owner record cap $981105$ proves
\eqref{eq:kb-rank11-exact-petal-partition-cap}.

Exact cross-multiplication against
\eqref{eq:kb-rank11-weighted-nine-target} is negative at $K'=15528$ and
positive at $K'=15529$.  The demand/cap ratio is a positive constant times
the nine increasing factors $(m'-i)/(n'-i)$, $0\le i\le8$, and
\[
 H(K')=\frac{(K'+67463)(K'+67462)}
              {1048577K'+34798536326}.
\]
After clearing positive denominators, the numerator of
$H(K'+1)-H(K')$ is
\[
 1048577K'^2+69598121229K'-77044697164886.
\]
Writing $K'=15529+x$ turns this into
\[
 1048577x^2+102164825695x+1256608704226512>0
\]
for $x\ge0$.  The contradiction therefore persists through $K'=15634$;
the preceding two propositions cover all larger rows.
\end{proof}

\begin{proposition}[Weighted selected-support split-pencil capacity]
\label{prop:kb-rank11-weighted-selected-support-split-pencil}
Let $A\ge3$, and let finitely many affine owner points $p$ have integral
weights
\[
 0\le s_p\le A-1,qquad \sum_p s_p\le S.
\]
For each line $L$ in a family of distinct affine lines, let integral
selected masses $x_{L,p}$, supported on the owner points of $L$, satisfy
\[
 0\le x_{L,p}\le s_p,qquad \sum_p x_{L,p}=A.
\]
Put
\[
 Q_L=\sum_p\binom{x_{L,p}}2,qquad
 h=\floor{\frac{S}{\floor{A/2}+1}}.
\]
Then
\begin{equation}
 \sum_L Q_L\le
 \floor{\frac{(A-2)S^2}{8}}+\binom S2
       +\binom h2\binom{A-1}2.
 \label{eq:kb-rank11-weighted-selected-support-cap}
\end{equation}
\end{proposition}

\begin{proof}
Write $q(x)=\binom x2$.  For one line abbreviate $x_p=x_{L,p}$ and set
\[
 X_L=\sum_{p<r}x_px_r.
\]
Since $\sum_p x_p=A$, one has $Q_L+X_L=\binom A2$.
Call an owner globally heavy when $s_p>A/2$; there are at most $h$ such
owners.

First suppose $\max_p x_p\le A/2$.  Then
\[
 \sum_p x_p^2\le \frac{A^2}{2},qquad
 2Q_L=\sum_p x_p^2-A,qquad
 2X_L=A^2-\sum_p x_p^2,
\]
so $Q_L\le X_L$.  A cross-owner coordinate pair determines its two owner
points, and those points determine a unique affine line.  Distinct lines
therefore give
\begin{equation}
 \sum_{L\ {
m balanced}}Q_L
 \le\sum_{L\ {
m balanced}}X_L
 \le\sum_{p<r}s_ps_r\le\binom S2.                  \tag{1}
\end{equation}

Every remaining line has a unique selected dominant owner.  Set aside the
lines containing at least two globally heavy owners.  Every such line
contains a pair of heavy points, each pair determines at most one line, and
convexity under $x_p\le A-1$ gives $Q_L\le\binom{A-1}2$.  These lines
contribute at most
\begin{equation}
 \binom h2\binom{A-1}2.                            \tag{2}
\end{equation}

It remains to count clean dominant lines.  If $p$ is their unique globally
heavy point, put $s=s_p$ and $d=A-s$.  Thus $A/2<s\le A-1$ and
$1\le d<A/2$.  If the selected dominant mass is $x$, superadditivity of
$q$ on the complementary masses and monotonicity away from $A/2$ give
\[
 Q_L\le q(x)+q(A-x)\le q(s)+q(d).                 \tag{3}
\]
Distinct lines through $p$ contain disjoint sets of other owner points, and
each uses at least $A-x\ge d$ light mass.  If $H$ and $L_0$ are the total
heavy and light masses, respectively, there are at most $L_0/d$ clean lines
through $p$.  The exact slack identity
\[
 ds(A-2)-\{s(s-1)+d(d-1)\}
 =(d-1)(d+s)(s-1)\ge0
\]
shows that
\[
 \frac{q(s)+q(d)}d\le\frac{s(A-2)}2.
\]
Summing first through each heavy owner and then over heavy owners yields
\begin{equation}
 \sum_{L\ {
m clean}}Q_L
 \le\frac{(A-2)H L_0}{2}
 \le\frac{(A-2)S^2}{8}.                            \tag{4}
\end{equation}
The left side is integral.  Taking the floor in (4) and adding (1)--(2)
proves the claim.
\end{proof}

\begin{proposition}[Minimal-row rank-nine split-pencil payment]
\label{prop:kb-rank11-rank9-minimal-split-pencil-payment}
In the fixed target of
Proposition~\ref{prop:kb-rank11-component-ninesubset-target}, the rank-nine
alternative is empty at $K'=10$.  In fact its full-density marked demand
and selected-support capacity satisfy
\[
 11736940042024039>9274769506943785,
\]
with gap $2462170535080254$.  Together with dimension equality, there is no
rank-eleven component target on this minimal residual row.
\end{proposition}

\begin{proof}
At $K'=10$, Proposition~\ref{prop:kb-rank11-dense-root-saturation} places
the ten-dimensional correction space $V'$ inside the ten-dimensional
space $\mathbb F[X]_{<10}$, so equality holds.  Vandermonde interpolation
gives rank ten on every eleven-set and rank nine on every nine-set.  Thus
every positive-dimensional incidence of
Proposition~\ref{prop:kb-rank11-component-incidence} is in the full-rank
owner lane; the kernel and rank-eight lanes are empty.

Write
\[
 N_{\min}=274980728111260126,qquad
 \lambda=\frac{990810934}{10^9},qquad
 (n',m')=(1048586,67482).
\]
Mark all $\binom{11}{9}$ nine-subsets of every component incidence and
average over the $\binom{n'}9$ nine-sets.  One fixed $B$ has marked weight
\begin{align*}
 W_B&\ge\ceil{\lambda N_{\min}
       \frac{\binom{m'}9\binom{m'-9}2}{\binom{n'}9}}\\
    &=11736940042024039.                           \tag{1}
\end{align*}
There is no half-lane loss because the two competing component lanes are
already empty.

The evaluation kernel on $B$ is spanned by its degree-nine locator $u$, so
$Z_D(u)=B$.  Hence the common owner core is exactly $B$.  Off $B$, the
owner petals have sizes
\[
 s_p\le m'-10=67472,qquad \sum_p s_p\le n'-9=1048577.
\]
For every retained record line, intersect its selected support with the
owner petals.  The resulting selected masses obey
\[
 0\le x_{L,p}\le s_p,qquad \sum_p x_{L,p}=m'-9=67473.
\]
Both new coordinates in each marked extension lie outside $Z_D(u)=B$ and
in the same owner petal, so that record contributes at most
$\sum_p\binom{x_{L,p}}2$.  This uses only its selected support and makes no
claim about accidental agreements outside it.

Apply Proposition~\ref{prop:kb-rank11-weighted-selected-support-split-pencil}
with
\[
 A=67473,qquad S=1048577,qquad
 h=\floor{1048577/33737}=31.
\]
The clean-dominant, balanced, and heavy-collision terms are respectively
\[
 9273161316835569,qquad 549756338176,qquad 1058433770040.
\]
Their sum is $9274769506943785$, contradicting (1).  Rank nine is therefore
empty.  The initial Vandermonde argument excludes the other component
lanes.
\end{proof}

\begin{proposition}[Selected-support split-pencil with a common-core offset]
\label{prop:kb-rank11-weighted-selected-support-core-offset}
Let distinct affine lines carry selected masses $x_{L,p}$ at owner points
$p$, where
\[
 \sum_p x_{L,p}=P,\qquad x_{L,p}\le s_p\le P-1,
 \qquad \sum_p s_p\le S.
\]
For an integer $r\ge0$, put
\[
 R_L=\sum_p\binom{x_{L,p}}2+rP,
 \quad h=\floor{\frac{S}{\floor{P/2}+1}},
 \quad M=\floor{\frac{P^2}{4}}.
\]
Then
\begin{align}
 \sum_LR_L\le{}&
 \max_{0\le\ell\le S}
 \floor{\ell\left(\frac{(P-2)(S-\ell)}2+hrP\right)}
 \notag\\
 &+\floor{\frac{M+rP}{M}\binom S2}
 +\binom h2\left(\binom{P-1}2+rP\right).       \tag{SPO}
\end{align}
\end{proposition}

\begin{proof}
Write $Q_L=\sum_p\binom{x_{L,p}}2$ and
$X_L=\sum_{p<a}x_{L,p}x_{L,a}$, so
$Q_L+X_L=\binom P2$.  If no selected owner exceeds $P/2$, convexity gives
$Q_L\le X_L$ and $X_L\ge M$.  Cross-owner pairs inject globally because
two owner points determine one affine line.  Thus the balanced lines have
total charge at most
\[
 \floor{\frac{M+rP}{M}\binom S2}.
\]

There are at most $h$ globally heavy owners.  A line containing two of them
is determined by that pair, and convexity gives $Q_L\le\binom{P-1}2$.
These collision lines contribute at most the third term of (SPO).

Every remaining nonbalanced line has one dominant globally heavy owner of
weight $s$.  Put $d=P-s$.  Its selected dominant mass is at most $s$, so
\[
 Q_L\le\binom s2+\binom d2.
\]
Distinct clean lines through this owner use disjoint light mass and each
uses at least $d$ of it.  The exact identity
\[
 ds(P-2)-\{s(s-1)+d(d-1)\}
 =(d-1)(d+s)(s-1)\ge0
\]
therefore gives
\[
 \frac{\binom s2+\binom d2+rP}{d}
 \le \frac{s(P-2)}2+rP.
\]
If the total light mass is $\ell$, summing over the at most $h$ heavy
owners bounds all clean lines by
\[
 \ell\left(\frac{(P-2)(S-\ell)}2+hrP\right).
\]
Maximizing over integral $\ell$ and adding the three classes proves (SPO).
\end{proof}

\begin{proposition}[$K'=11$ circuit-shadow split-pencil payment]
\label{prop:kb-rank11-k11-circuit-split-pencil-payment}
In the fixed target of
Proposition~\ref{prop:kb-rank11-component-ninesubset-target}, every
rank-eleven component alternative is empty at $K'=11$.  At the residual
record floor, full component demand is
\[
 901408286315387898338134887980054663001598216883356906995509296,
\]
whereas the combined high- and low-circuit capacity is
\[
 870719390190680409022824387604193486699840723094988553120053384.
\]
The strict gap is
\[
 30688896124707489315310500375861176301757493788368353875455912.
\]
\end{proposition}

\begin{proof}
At $K'=11$, the correction space $V'$ is a ten-dimensional hyperplane in
$P_{11}=\mathbb F[X]_{<11}$.  Evaluation of $P_{11}$ on every eleven-set is
an isomorphism, so every component incidence has affine-owner rank ten.
Fix the nonzero global functional $\lambda$ with $V'=\ker\lambda$.

For a rank-nine nine-shadow $B$, let $J$ be the common core of its owner
plane.  The residual root bound gives $j=|J|\in\{9,10\}$.  Outside $J$,
the selected petal mass and total owner mass are
\[
 P=m'-j,\qquad S=n'-j,
\]
and a marked extension contributes at most
\[
 \sum_p\binom{x_{L,p}}2+(j-9)P.
\]
Proposition~\ref{prop:kb-rank11-weighted-selected-support-core-offset}
gives the exact capacities
\[
 \begin{array}{c|c|c|c|r}
 j&P&S&r&\text{capacity}\\ \hline
 9&67474&1048578&0&9274924665987729\\
 10&67473&1048577&1&9275866238180030.
 \end{array}
\]
Thus every rank-nine shadow has marked load at most
$C_*=9275866238180030$, and all such shadows together have capacity
\[
 G=\binom{1048587}{9}C_*.                          \tag{1}
\]

For one component eleven-set $T$, the hyperplane
$\operatorname{ev}_T(V')$ has one circuit support $C_T$, equivalently the
unique representation
\[
 \lambda=\sum_{a\in C_T}\mu_a\operatorname{ev}_a,
 \qquad \mu_a\ne0.                                \tag{2}
\]
If $c=|C_T|$, exactly
\[
 55-\binom{11-c}{2}                                \tag{3}
\]
of its nine-shadows have rank nine.  Hence every $c\ge6$ incidence uses at
least 45 marks, and (1) bounds all such incidences by
\[
 \floor{G/45}
 =870571094319564771339440637927219938149933838522556716983865974.
                                                               \tag{4}
\]

Now suppose $c\le5$.  Subtracting two representations (2) gives a relation
among evaluation functionals on at most ten distinct points.  These are
Vandermonde-independent on $P_{11}$, so the representations are identical.
All low-circuit incidences therefore contain one fixed support $C_*'$ of
size at most five.  One record supplies at most
\[
 \binom{m'-|C_*'|}{11-|C_*'|}\le\binom{m'-1}{10}
\]
such eleven-sets.  At
$N_{\min}=274980728111260126$, the low-circuit capacity is
\[
 148295871115637683383749676973548549906884572431836136187410.
                                                               \tag{5}
\]
Adding (4)--(5) gives the asserted capacity.  The dense-locator theorem
requires
\[
 \ceil{\frac{990810934}{10^9}N_{\min}\binom{67483}{11}},
\]
which is the displayed demand.  Moreover the demand coefficient minus
$\binom{67482}{10}$ has positive cross numerator
\[
 3277538744750188045701745970410737654908675596614324570,
\]
so the contradiction persists for every allowed larger record count.
\end{proof}

\begin{proposition}[Codimension-two quotient-line sparse-circuit capacity]
\label{prop:kb-rank11-quotient-line-sparse-circuit}
Let $V\le P_{12}=\mathbb F[X]_{<12}$ have dimension ten, put
$\Lambda=V^\perp$, and let $S$ be a set of $m$ distinct field points.  For
each eleven-set $T\subseteq S$ on which $V$ has evaluation rank ten, let
$[\lambda_T]$ be the unique point of
$\mathbb P(\Lambda\cap\operatorname{span}\{\operatorname{ev}_x:x\in T\})$
and let $C_T$ be the support of its unique representation on $T$.

Put $Q_1(m)=2$, and for $2\le c\le5$ put
\[
 Q_c(m)=\max\left(\{c+1\}\cup
 \left\{c+\floor{\frac{e(m-g)}{c-g}}:
       1\le e\le c,\ 0\le g<c\right\}\right).\qquad\textup{(QL1)}
\]
Then
\[
 \#\{T:|C_T|\le5\}
 \le \sum_{c=1}^{5}Q_c(m)\binom{m-c}{11-c}.\qquad\textup{(QL2)}
\]
At $m=67484$, the five label caps are
\[
 2,\quad134968,\quad202449,\quad269928,\quad337405,
\]
and the right side of \textup{(QL2)} is
\[
 11868577829520852215896202871552159662636920.
\]
\end{proposition}

\begin{proof}
Evaluation functionals at at most twelve distinct points are independent on
$P_{12}$.  Since $\dim\Lambda=2$, annihilator duality gives the asserted
unique projective label and representation for every rank-ten $T$.

A support-one label is proportional to one evaluation functional.  Three
such labels would put three independent evaluation functionals in
$\Lambda$, so there are at most two.

Fix $2\le c\le5$.  For $\lambda\in\Lambda$ form the catalecticant
\[
 H_c(\lambda)=\bigl(\lambda(X^{i+j})\bigr)_
 {0\le i\le11-c,\ 0\le j\le c}.
\]
For a support-$c$ label this matrix has rank $c$, and its right kernel is
generated by the circuit locator $\prod_{a\in C_T}(X-a)$.  Parameterize
$\mathbb P(\Lambda)$ by $(U:V)$.  If one maximal minor of $H_c(U,V)$ is
nonzero, its degree $c+1$ bounds all support-$c$ labels by $c+1$.

Otherwise choose a $c$-row chart which has rank $c$ at one support-$c$
label.  Its signed maximal cofactors form a kernel polynomial
\[
 Q(U,V;X)=\sum_{j=0}^{c}Q_j(U,V)X^j.
\]
After cancelling the common parameter factor, the coefficient forms have
one degree $e\le c$.  At most $c$ parameter values are outside the chart.
If $e=0$, every good label has one fixed circuit support.  One full-rank
eleven-set cannot carry two such labels, since then
$\Lambda\le\operatorname{span}\{\operatorname{ev}_x:x\in T\}$ and the
evaluation rank would be at most nine.  Including the bad parameters gives
the $c+1$ alternative in \textup{(QL1)}.

Suppose $e\ge1$, and let $g$ domain points be roots of $Q(U,V;X)$ for every
parameter.  Then $g<c$.  Each good support-$c$ label visible in $S$
contributes $c-g$ nonfixed root incidences, while for each $x\in S$ the
nonzero homogeneous polynomial $Q(U,V;x)$ has at most $e$ projective roots.
Thus there are at most
$\floor{e(m-g)/(c-g)}$ good visible labels.  Adding the at most $c$ bad
parameters proves \textup{(QL1)}.  Each label
contributes at most $\binom{m-c}{11-c}$ eleven-sets, proving
\textup{(QL2)}.  Notice that this retains
moving kernel pencils and uses no classification of lines contained in a
secant variety.
\end{proof}

\begin{proposition}[$K'=12$ quotient-line circuit payment]
\label{prop:kb-rank11-k12-quotient-line-payment}
In the fixed target of
Proposition~\ref{prop:kb-rank11-component-ninesubset-target}, every
rank-eleven component alternative is empty at $K'=12$.  At the residual
record floor, the complete kernel, high-circuit, and low-circuit capacity is
\[
 873945204333998831582903951502910514268526233054054867526472861,
\]
whereas full component demand is
\[
 901555241262544083284435178226046105523688795046262319915891531.
\]
The strict gap is
\[
 27610036928545251701531226723135591255162561992207452389418670.
\]
\end{proposition}

\begin{proof}
Put $\Lambda=(V')^\perp\le P_{12}^*$.  First retain the rank-deficient
component incidences.  At shortening excess $K'-10=2$, the canonical-basis
extension factor is $\binom2{d+1}$.  Only corank $d=1$ remains, with record
cap $M_1(12)=16295594$, so all rank-deficient incidences have capacity
\[
 K_{\rm cap}=\binom{1048588}{9}16295594.
\]

For a rank-ten eleven-set $T$, let $C_T$ be its quotient-line circuit.  A
rank-nine shadow has common-core size $j\in\{9,10,11\}$.  Applying
Proposition~\ref{prop:kb-rank11-weighted-selected-support-core-offset} with
$P=m'-j$, $S=n'-j$, and offset $j-9$ gives the three chart caps
\[
 9275079825564861,\qquad9276021414531691,\qquad9276963034268184.
\]
Thus $C_*=9276963034268184$ bounds every rank-nine chart.

Exactly $55-\binom{11-|C_T|}{2}$ nine-shadows of $T$ have rank nine.  Every
$|C_T|\ge6$ incidence therefore uses at least 45 marks, and all such
incidences have capacity
\[
 H_{\rm cap}=\floor{\binom{1048588}{9}C_*/45}.
\]
Proposition~\ref{prop:kb-rank11-quotient-line-sparse-circuit} bounds the
remaining $|C_T|\le5$ incidences by
\[
 N_{\rm actual}\,
 11868577829520852215896202871552159662636920.
\]

The dense-locator theorem requires at least
\[
 \ceil{\frac{990810934}{10^9}N_{\rm actual}\binom{67484}{11}}
\]
component incidences.  Its coefficient in $N_{\rm actual}$ minus the
coefficient in (3) has positive cross numerator
\[
 3266743881505321144251837279559822170446618230146889560.
\]
The smallest gap is therefore at
$N_{\min}=274980728111260126$.  Adding the three capacities there gives the
displayed total, strictly below the displayed demand.
\end{proof}

\begin{proposition}[Codimension-three sparse-circuit completion]
\label{prop:kb-rank11-codimension-three-sparse-circuit}
Let $V\le P_{13}=\mathbb F[X]_{<13}$ have dimension ten and have no common
zero on a set $S$ of $m$ distinct field points.  Put $\Lambda=V^\perp$.
For each eleven-set $T\subseteq S$ on which $V$ has evaluation rank ten,
let $C_T$ be the support of the unique nonzero label in
\[
 \Lambda\cap\operatorname{span}\{\operatorname{ev}_x:x\in T\}.
\]
For the circuits with $2\le |C_T|\le5$, one of the following holds:
\begin{enumerate}
\item every such circuit is contained in one set $U$ with $|U|\le7$;
\item every independent $(c-1)$-set has at most two completions to a
support-$c$ circuit, for every $2\le c\le5$.
\end{enumerate}
Consequently
\begin{align*}
 \#\{T:|C_T|\le5\}\le\max\biggl\{&
 \sum_{c=2}^5\binom{7}{c}\binom{m-c}{11-c},\\
 &\sum_{c=2}^5\floor{\frac{2\binom m{c-1}
              \binom{m-c-1}{11-c}}c}\biggr\}.
\end{align*}
At $m=67485$, the two displayed caps are respectively
\[
 1679076702065233864778823429158845084750
\]
and
\[
 99254447944649683780146155758753837527116020.
\]
\end{proposition}

\begin{proof}
Write $E_A=\operatorname{span}\{\operatorname{ev}_x:x\in A\}$.  Evaluation
functionals at any at most thirteen distinct points are independent on
$P_{13}$.  If $T$ has evaluation rank ten, annihilator duality gives
\[
 \dim(\Lambda\cap E_T)=11-\operatorname{rank}(\operatorname{ev}_T|_V)=1.
\]
The unique label has a unique representation on $T$, and its nonzero support
$C_T$ is a circuit of the projected evaluations on $V$.  The common-zero
hypothesis excludes support one.

Fix $2\le c\le5$ and an independent $(c-1)$-set $A$.  Put
\[
 H_A=\{f\in V:f|_A=0\}.
\]
Then $\dim H_A=11-c$.  Every point completing $A$ to a circuit is a common
zero of $H_A$.  If $z$ distinct points are common zeros, then $H_A$ is
contained in their degree-$z$ locator times $P_{13-z}$, so
$z\le13-\dim H_A=c+2$.  The points of $A$ account for $c-1$ of these zeros;
hence $A$ has at most three circuit completions.

Suppose $A$ has three completions $x,y,z$.  Their three circuit labels are
independent: the evaluation functionals on $A\cup\{x,y,z\}$ are independent,
and each label has one private nonzero coordinate in that basis.  Put
\[
 U=A\cup\{x,y,z\},\qquad |U|=c+2\le7.
\]
They span the three-dimensional $\Lambda$, so $\Lambda\le E_U$.  If
$\mu\in\Lambda$ has another circuit support $D$ of size at most five, its
representations on $U$ and $D$ agree because $|U\cup D|\le12$ and those
evaluation functionals are independent.  Thus $D\subseteq U$, proving the
first alternative and its displayed count.

Otherwise let $b_A\le2$ be the number of completions of $A$.  If $x,y$ are
distinct completions, their labels are independent.  A rank-ten eleven-set
containing $A\cup\{x\}$ therefore cannot contain $y$, since its intersection
with $\Lambda$ is one-dimensional.  The eleven-sets charged to the
completions of $A$ are at most
\[
 b_A\binom{m-c+1-b_A}{11-c}
 \le2\binom{m-c-1}{11-c}
\]
at the stated $m$.  Every support-$c$ circuit is charged exactly $c$ times,
once after deletion of each point.  Summing over the $\binom m{c-1}$ choices
of $A$ and dividing by $c$ proves the second count.  Exact integer
evaluation gives the two printed constants.
\end{proof}

\begin{proposition}[$K'=13$ sparse-circuit completion payment]
\label{prop:kb-rank11-k13-sparse-circuit-payment}
In the fixed target of
Proposition~\ref{prop:kb-rank11-component-ninesubset-target}, every
rank-eleven component alternative is empty at $K'=13$.  At the residual
record floor, complete kernel, high-circuit, and low-circuit capacity is
\[
 898085191110430398284744062896212914931984716650701254999384513,
\]
whereas full component demand is
\[
 901702217989192688449626641411280218028664942551160634607759137.
\]
The strict gap is
\[
 3617026878762290164882578515067303096680225900459379608374624.
\]
\end{proposition}

\begin{proof}
At shortening excess three, the canonical-basis extension factors are
$\binom3{d+1}$.  Coranks one and two have record caps $16295594$ and
$253241283$, so every rank-deficient component incidence has capacity
\begin{align*}
 K_{\rm cap}
 &=3\binom{1048589}{9}16295594
   +\binom{1048589}{8}253241283\\
 &=206481189843433295842936213010503229833431068859362597823.
\end{align*}

For full-rank eleven-sets, apply
Proposition~\ref{prop:kb-rank11-weighted-selected-support-core-offset} at
common-core sizes $j=9,10,11,12$.  The four chart caps are
\[
 9275234985700485,\quad9276176591408201,\quad
 9277118227920090,\quad9278059895199813.
\]
Thus $C_*=9278059895199813$.  A circuit of size at least six creates at
least 45 rank-nine shadows, giving high-circuit capacity
\[
 H_{\rm cap}
 =870791924265139618716231673259817164224620222733319378834968170.
\]
Proposition~\ref{prop:kb-rank11-codimension-three-sparse-circuit} bounds the
remaining support-at-most-five incidences by
\[
 N_{\rm actual}\,
 99254447944649683780146155758753837527116020.
\]

The dense-locator theorem requires at least
\[
 \ceil{\frac{990810934}{10^9}N_{\rm actual}\binom{67485}{11}}
\]
component incidences.  The demand coefficient minus the sparse coefficient
has positive cross numerator
\[
 3179892509671792384744396543086450191690406044908895900.
\]
The smallest gap is therefore at
$N_{\min}=274980728111260126$.  Adding the three capacities there gives the
displayed total, strictly below the displayed demand.
\end{proof}

\begin{proposition}[Sparse-circuit completion ladder]
\label{prop:kb-rank11-sparse-circuit-completion-ladder}
Let $K\ge13$, put $q=K-10$, and let
$V\le P_K=\F[X]_{<K}$ have dimension ten and no common zero on a set $S$
of $m$ distinct field points.  Every eleven-set $T\subseteq S$ on which
$V$ has evaluation rank ten selects one minimal circuit support $C_T$.
For the circuits with $2\le c=|C_T|\le5$, one of the following holds:
\begin{enumerate}
\item every such circuit lies in one carrier of size at most $q+4$;
\item every independent $(c-1)$-set has at most $q-1$ completions to a
support-$c$ circuit, for every $2\le c\le5$.
\end{enumerate}
Consequently, if $I_c$ counts the eleven-sets selecting a support-$c$
circuit, then
\begin{align*}
 I_c\le\max\{S_c,U_c\},\qquad
 S_c&=\binom{q+4}{c}\binom{m-c}{11-c},\\
 U_c&=\floor{\frac{\binom m{c-1}}c
       \max_{0\le b\le q-1}b\binom{m-c+1-b}{11-c}}.
\end{align*}
On every specialization $14\le K'\le21$, the maximum in $U_c$ occurs at
$b=q-1$.
\end{proposition}

\begin{proof}
Put $\Lambda=V^\perp$, so $\dim\Lambda=q$.  Annihilator duality gives
$\dim(\Lambda\cap E_T)=1$ for every selected eleven-set.  Fix an
independent deletion $A$ of size $c-1$.  The space
$H_A=\{f\in V:f|_A=0\}$ has dimension $11-c$.  If its common-zero set has
size $z$, division by its locator embeds $H_A$ in $P_{K-z}$, whence
$z\le K+c-11$.  The points of $A$ already account for $c-1$ zeros, so
$A$ has at most $q$ circuit completions.

If one deletion has all $q$ completions, their private completion
coordinates make their labels independent.  They span $\Lambda$ and have
carrier $U$ of size $c-1+q\le q+4$.  Any other sparse label has the same
representations on its support and $U$; their union has size at most
$q+9=K-1$, so Vandermonde independence forces its support into $U$.  This
gives $S_c$.

Otherwise let $b\le q-1$ be the completions of one deletion.  A selected
eleven-set cannot retain two of their labels, since its intersection with
$\Lambda$ is one-dimensional.  Thus that deletion contributes at most
$b\binom{m-c+1-b}{11-c}$ eleven-sets.  Summing over all deletions and
dividing by the $c$ charges of each circuit gives $U_c$.  Finally,
$b\binom{m-c+1-b}{11-c}$ is increasing through $b=q-1$ on the eight stated
rows; the exact ratio test has threshold greater than $6700$, whereas
$q-1\le10$.
\end{proof}

\begin{proposition}[Joint sparse/high rank-nine-shadow ledger]
\label{prop:kb-rank11-joint-sparse-shadow-ledger}
Suppose every full-rank eleven-set has one circuit.  A support-$c$ circuit
has exactly
\[
 q_c=55-\binom{11-c}{2}
\]
rank-nine shadows.  Hence $q_c=19,27,34,40$ for $c=2,3,4,5$, while
$q_c\ge45$ for $c\ge6$.  If all rank-nine marks have total capacity $G$
and $R$ records have per-record sparse bounds $L_c$, then complete
full-rank incidence is at most
\[
 \floor{\frac{G+R\sum_{c=2}^5(45-q_c)L_c}{45}}.
\]
If each record lies in one of several structural branches with cap vectors
$L_c^{(a)}$, the sum is replaced by its maximum over $a$.
\end{proposition}

\begin{proof}
Omitting a pair disjoint from the circuit leaves rank eight; all other
omitted pairs leave rank nine.  This proves the formula for $q_c$.  If
$I_c$ counts sparse incidences and $H$ the incidences of support at least
six, the one global shadow resource gives
\[
 45H+\sum_{c=2}^5q_cI_c\le G.
\]
Adding $\sum(45-q_c)I_c$ and using $I_c\le RL_c$ proves the claim.  For
multiple branches, maximize the premium record by record.  In particular,
the sparse strata consume their ordinary shadows before receiving the
printed premiums $26,18,11,5$.
\end{proof}

\begin{proposition}[$K'=14$ through $K'=21$ joint sparse-shadow payment]
\label{prop:kb-rank11-k14-k21-joint-sparse-shadow-payment}
In the fixed target of
Proposition~\ref{prop:kb-rank11-component-ninesubset-target}, every
rank-eleven component alternative is empty for $14\le K'\le21$.  The exact
demand-minus-capacity gaps are
\[
\begin{array}{c@{\quad}r}
K'&\text{gap}\\ \hline
14&21650768172043394032492459946400263590275020562192971022865376\\
15&18588028475285812265606143363906671348737995144736255147834120\\
16&15524989535492821727324938882048706164611583704587475891361669\\
17&12461651333490539307570308245171118440833885273861127661424680\\
18& 9398013850103444858371174848818672822082629043827750046751608\\
19& 6334077066154944217959550399858581606171131782792291719264278\\
20& 3269840962466993657882795068981405767940223879562014346498774\\
21&  205305519860193617784849691734671763401656917434567909452790
\end{array}
\]
The identical payment first fails at $K'=22$, where capacity exceeds
demand by
\[
 2859529280846211417198922209345618432657212793529140162369036.
\]
\end{proposition}

\begin{proof}
For one row put $n=1048576+K'$, $m=67472+K'$, and $q=K'-10$.  Retain every
nonzero canonical-basis corank term:
\[
 K_{\rm cap}=\sum_{d=1}^{\min(9,q-1)}
   \binom n{10-d}M_d(K')\binom q{d+1}.
\]
For every core size $9\le j<K'$, specialize
Proposition~\ref{prop:kb-rank11-weighted-selected-support-core-offset} at
$P=m-j$, $S=n-j$, and offset $j-9$.  Let $C_*$ be the largest resulting
chart cap and put $G=\binom n9C_*$.  No rank-deficient incidence enters
this full-rank resource.

Proposition~\ref{prop:kb-rank11-sparse-circuit-completion-ladder} supplies
the two per-record vectors
\begin{align*}
 L_c^{\rm S}&=\binom{q+4}{c}\binom{m-c}{11-c},\\
 L_c^{\rm U}&=\floor{\frac{\binom m{c-1}(q-1)
                 \binom{m-c+2-q}{11-c}}c}.
\end{align*}
By Proposition~\ref{prop:kb-rank11-joint-sparse-shadow-ledger}, full-rank
capacity is
\[
 F_{\rm cap}=\floor{\frac{G+R_{\rm actual}
  \max_{a\in\{\rm S,U\}}\sum_{c=2}^5(45-q_c)L_c^a}{45}}.
\]
The dense-locator demand is
\[
 \ceil{\frac{990810934}{10^9}R_{\rm actual}\binom m{11}}.
\]
The exact replay scans every displayed $j$ and both completion branches.
On all eight rows, the demand coefficient after clearing $45\cdot10^9$ is
positive, and the full cross-product is already positive at
$R_{\min}=274980728111260126$.  Therefore demand exceeds
$K_{\rm cap}+F_{\rm cap}$ for every admissible record count.  Integer
evaluation gives the displayed gaps.  At $K'=22$, the same formulas give
capacity
\[
905885518366475292751564400874300832826807604203127204847344067
\]
and demand
\[
903025989085629081334365478664955214394150391409598064684975031,
\]
which proves only the stated method wall.
\end{proof}

\begin{proposition}[Integral heavy-owner split-pencil cap]
\label{prop:kb-rank11-integral-heavy-owner-cap}
In the notation of
Proposition~\ref{prop:kb-rank11-weighted-selected-support-core-offset},
put $P=m-j$, $S=n-j$, $r=j-9$,
$a=\lfloor P/2\rfloor+1$, $b=P-1$, and
$C_0=\binom P2+rP$.  If the globally heavy owners have selected weights
$s_i\in[a,b]$ and the selected light mass is $\ell$, then their total
clean-line charge is at most
\[
 \ell\sum_i\left(\frac{C_0}{P-s_i}-s_i\right).
\]
For fixed $\ell$ and owner count, the maximum in this bound has all but at most one $s_i$ at
an endpoint $a$ or $b$.  Hence the exact integer maximum is obtained by a
finite scan over the owner count, the number of $b$-weights, and one
residual weight.

At $K'=22$, exact evaluation for every $9\le j\le21$ gives uniform chart
cap
\[
 C_*=9269974099565290,
\]
attained at $j=21$.  This improves the preceding uniform cap by
$17960461975558$.
\end{proposition}

\begin{proof}
A clean line through a heavy owner of weight $s$ consumes at least $P-s$
light mass.  Distinct such lines use disjoint light points, and convexity of
the selected partition bounds one line by
$\binom s2+\binom{P-s}2+rP$.  This gives the displayed bound.  Its summand
\[
 \phi(s)=\frac{C_0}{P-s}-s
\]
is increasing and convex on $[a,b]$.  An exchange toward the endpoints
therefore cannot decrease the objective, proving the asserted extremal
shape.  On each residual-weight interval the objective is a rational
quadratic with second derivative changing sign at most once.  Endpoints and
the unique possible derivative crossing give an exact finite evaluation.
The certificate checks 271 intervals for each of the thirteen core offsets;
all maxima have eight owners of weight $b$, and direct substitution gives
the displayed cap and saving.
\end{proof}

\begin{proposition}[Near-saturation carrier and the $K'=22$ payment]
\label{prop:kb-rank11-k22-near-saturation-payment}
Let $\Lambda=V^\perp$ have dimension $q$, and let $2\le c\le4$.  If one
independent $(c-1)$-set has $q-1$ support-$c$ completions, then every
support-$c$ annihilator label is carried by at most $q+2c-2$ coordinates.
Otherwise every such deletion has at most $q-2$ completions.

For the fixed component target at $K'=22$, combining this dichotomy with
Proposition~\ref{prop:kb-rank11-integral-heavy-owner-cap}, the exact
support-five deletion cap, all nine nonzero kernel coranks, and the uniform
corank-one record cap $8147918$ gives
\[
 \operatorname{Cap}=901793257329000893449088123410358112576738959572328805741503920,
\]
whereas required incidence is
\[
 903025989085629081334365478664955214394150391409598064684975031.
\]
Thus $K'=22$ has no rank-nine component target; the strict gap is
\[
1232731756628187885277355254597101817411431837269258943471111.
\]
\end{proposition}

\begin{proof}
The $q-1$ completion labels have private completion coordinates, so they
span a hyperplane $\Lambda_0\le\Lambda$ on a carrier of size $q+c-2$.
Every label in $\Lambda_0$ is supported there by Vandermonde independence.
If a support-$c$ label lies outside $\Lambda_0$, adjoining its support spans
$\Lambda$ on at most $q+2c-2$ coordinates.  Comparing this representation
with any circuit representation uses at most $q+3c-2\le q+10=K$
coordinates, and Vandermonde independence forces every circuit into the
same carrier.  The complementary deletion count is the proof of
Proposition~\ref{prop:kb-rank11-sparse-circuit-completion-ladder} with
$b\le q-2$.  The exact integer ledger then gives the displayed capacity;
its cleared record coefficient and its value at the minimum admissible
record count are both positive.
\end{proof}

\begin{proposition}[Completion-defect hierarchy and the $K'=23$ payment]
\label{prop:kb-rank11-completion-defect-hierarchy}
Suppose an independent $(c-1)$-set has $q-s$ support-$c$ completions.  If
\begin{equation*}
 (s+2)c-s-1\le10.                                    \tag{DH}
\end{equation*}
then every support-$c$ label is carried by at most
$q+(s+1)(c-1)$ coordinates.  Consequently the maximal useful defect depths
for $c=2,3,4,5$ are respectively
\[
 d_c=7,2,1,0.
\]
If no carrier branch occurs through depth $d_c$, deletion counting may use
the ceiling $q-d_c-1$.

At $K'=23$, the complete defect ledger has capacity
\[
901468921726260997936972918059547622477700558490567171211759513
\]
against required incidence
\[
903173183767034183579263202870178889744727352679017565288251877.
\]
The strict gap is
\[
1704262040773185642290284810631267267026794188450394076492364.
\]
\end{proposition}

\begin{proof}
The $q-s$ completion labels span a space on a carrier of size
$q+c-1-s$.  At most $s$ further support-$c$ labels extend this to the full
support-$c$ label span, giving a carrier of size at most
$q+(s+1)(c-1)$.  Comparing an arbitrary circuit representation to a
carrier representation uses at most $q+(s+2)c-s-1$ coordinates.  Condition
(DH) and Vandermonde independence force the circuit support into the
carrier.  Maximizing the complementary deletion count over
$0\le b\le q-d_c-1$ proves the hierarchy.  Solving (DH), and checking that
the next depth fails it, gives $7,2,1,0$.  The finite ledger retains all
kernel and core-offset terms and gives the displayed strict comparison.
\end{proof}

\begin{proposition}[Universal completion incidence and full shadow deficits]
\label{prop:kb-rank11-universal-completion-full-deficit}
For every support $2\le c\le9$, the number $I_c$ of selected eleven-sets
whose circuit has support $c$ satisfies
\begin{equation*}
 I_c\le
 \floor{\frac{\binom m{c-1}}c
 \max_{0\le b\le q}b\binom{m-c+1-b}{11-c}}.
 \tag{UC}
\end{equation*}
If $G$ is the single global rank-nine-mark capacity and $R$ records have
recordwise caps $L_c^{(a)}$ in structural branch $a$, then complete
full-rank incidence is at most
\begin{equation*}
 \floor{\frac{G+R\max_a\sum_{c=2}^9
     \binom{11-c}{2}L_c^{(a)}}{55}}.
 \tag{FD}
\end{equation*}
The deficit weights in (FD), for supports two through nine, are
$36,28,21,15,10,6,3,1$; supports ten and eleven have zero deficit.
\end{proposition}

\begin{proof}
For an independent deletion $A$, the space of members of $V$ vanishing on
$A$ has dimension $11-c$.  If $A$ has $b$ completions, locator division
embeds this space in $P_{K-c+1-b}$, so $b\le K-10=q$.  The completion labels
have private nonzero coordinates and any two are independent, whereas the
annihilator of a selected rank-ten eleven-set is a line.  Such a set retains
at most one completion.  Counting the remaining coordinates, summing over
deletions, and dividing by the $c$ charges of each circuit proves (UC).

A support-$c$ circuit has exactly
$55-\binom{11-c}{2}$ rank-nine shadows.  Sum these marks in the one global
capacity, add the exact missing number of shadows to both sides, apply the
recordwise caps, and divide by 55.  This proves (FD) without spending any
rank-nine mark twice.
\end{proof}

\begin{proposition}[$K'=24$ through $K'=40$ full-deficit payment]
\label{prop:kb-rank11-k24-k40-full-deficit-payment}
In the fixed component target, every rank-nine alternative is empty for
$24\le K'\le40$.  The exact ledger uses
Proposition~\ref{prop:kb-rank11-completion-defect-hierarchy} on supports two
through five, Proposition~\ref{prop:kb-rank11-universal-completion-full-deficit}
on supports six through nine, and the full 55-shadow denominator.  Its
smallest demand-minus-capacity gap is attained at $K'=40$ and equals
\[
2272401814108959137912675549447888006236817090602808413697595.
\]
The same ledger first fails at $K'=41$, where capacity exceeds demand by
\[
4398836630793080990004182400858693750491819390616783425932508.
\]
\end{proposition}

\begin{proof}
For each row retain every nonzero refined kernel-corank term and evaluate
the integral chart of
Proposition~\ref{prop:kb-rank11-integral-heavy-owner-cap} for every
$9\le j<K'$.  In the full-rank lane maximize the complete weighted sum,
not each support separately, over the structured-carrier and
completion-defect branches, then apply (FD).  Exact integer replay checks all
423 core charts on the seventeen claimed rows and the adjacent wall.  On
$24\le K'\le40$ both the cleared record coefficient and its minimum-record
cross-product are positive.  The gaps decrease to the displayed $K'=40$
value.  At $K'=41$ both the gap and the cross-product are negative, giving
the stated honest method wall.
\end{proof}

\begin{proposition}[$K'=41$ sharp isolated-incidence payment]
\label{prop:kb-rank11-k41-sharp-isolated-payment}
In the fixed component target, every rank-nine alternative is empty at
$K'=41$.  Retain every kernel-corank term, integral core chart, completion
premium, and full 55-shadow capacity from
Proposition~\ref{prop:kb-rank11-k24-k40-full-deficit-payment}, and replace
only the generic isolated-point allowance by
Proposition~\ref{prop:kb-rank11-rank-stratified-isolated-incidence}.
Then exact component demand exceeds complete capacity by
\[
3959829848992990899082071934034620604165114037293042026746826.
\]
The same payment first fails at $K'=42$, where capacity exceeds demand by
\[
2710771376158610722953158157862051010402433288229120154217278.
\]
\end{proposition}

\begin{proof}
At $K'=41$, write $n'=1048617$, $m'=67513$, and
$N=N_{\min}=274980728111260126$.  The rank-stratified demand is exactly
\[
N\binom{m'}{11}-\binom{n'}{11}
=914185087092839732068202094579173634339667328332842235009045152.
\]
The unchanged kernel and full-rank capacities are respectively
\[
16046971808807649741426607721480807522683424615973152263500
\]
and
\[
910209210272037933519378596037417532927979530870933219830034826,
\]
whose sum is
\[
910225257243846741169120022645139013735502214295549192982298326.
\]
Subtraction gives the displayed positive gap.  After clearing the
55-shadow denominator, the coefficient of $N$ is
\[
143297102916558158042999955475306246365479723103>0,
\]
and the minimum-record cross-product is
\[
217790641694614499449513956371904133229081272051117311471075418>0.
\]
Thus the contradiction persists for every $N\ge N_{\min}$.  Exact replay
at $K'=42$ gives the displayed capacity excess, so no later row is claimed.
\end{proof}

\begin{proposition}[Cross-support completion-defect carrier]
\label{prop:kb-rank11-cross-support-defect-carrier}
Use the ten-dimensional sparse-circuit setup, put $q=K'-10$, and let
$2\le c,d\le9$.  Suppose one independent $(c-1)$-deletion has exactly
$q-s$ circuit completions, where $0\le s\le q$.  If
\[
 c+(s+1)d-s-1\le10,\qquad\text{(XC)}
\]
then every support-$d$ circuit lies in one carrier of size at most
\[
 q+c-1+s(d-1).
\]
Consequently its selected eleven-set incidence is at most
\[
 \binom{q+c-1+s(d-1)}d\binom{m'-d}{11-d}.
\]
For source support $c=5$ and defects $s=0,1,2,3,4$, the controlled target
supports are respectively
\[
 \{2,3,4,5,6\},\quad\{2,3\},\quad\{2\},\quad\{2\},\quad\{2\}.
\]
\end{proposition}

\begin{proof}
Let $A$ be the source deletion.  Normalize each of its $q-s$ completion
labels at its private completion coordinate.  These labels are linearly
independent and span a space $\Lambda_0$ of dimension $q-s$ on a set $U$
of size
\[
 |U|=q+c-1-s.
\]
Let $W_d$ be the span of all support-$d$ circuit labels.  Since the full
annihilator has dimension $q$, at most $s$ support-$d$ labels span
$(W_d+\Lambda_0)/\Lambda_0$.  Adjoining their supports to $U$ gives a
carrier $B$ with
\[
 |B|\le q+c-1-s+sd=q+c-1+s(d-1).
\]
Every support-$d$ label now has one representation on $B$ and its minimal
representation on its circuit support $D$.  Their union has size at most
\[
 q+c-1+s(d-1)+d\le q+10=K'
\]
by (XC).  Vandermonde independence through $K'$ points forces the two
representations to agree coordinatewise, hence $D\subseteq B$.  Counting
the possible $d$-subsets of $B$ and their eleven-set extensions gives the
incidence cap.  Substitution of $c=5$ gives the displayed target sets.
\end{proof}

\begin{proposition}[$K'=42$ cross-support completion-defect payment]
\label{prop:kb-rank11-k42-cross-support-payment}
In the fixed component target, every rank-nine alternative is empty at
$K'=42$.  Partition by the maximum number of completions of a support-five
deletion.  Maxima $q-s$ for $0\le s\le4$ give the five carrier branches of
Proposition~\ref{prop:kb-rank11-cross-support-defect-carrier}; the
complementary branch has at most $q-5$ completions on every deletion.  The
six exact weighted premiums are
\[
\begin{array}{c|r}
s&\text{premium}\\ \hline
0&3982795567806168516229316108688140450163630384\\
1&32010916243694499800320717073630461749362242674\\
2&38248686795246707324552098975684990817633239548\\
3&37909212899522784820182461169606692027883637848\\
4&37569681406601825198675563099275312923276993433\\
\text{fallback}&39561073029598078809344868550502487135515187669.
\end{array}
\]
The fallback is worst.  Exact component demand exceeds complete capacity by
\[
4081031051590194485758587836050845115467905186032497191061176.
\]
The same six-branch payment first fails at $K'=43$, where capacity exceeds
demand by
\[
2590504432899371163130658487199612335023802688487478696166262.
\]
\end{proposition}

\begin{proof}
For each branch, intersect the new source and cross-support caps with every
preceding supportwise completion cap, then weight support $d$ by its exact
55-shadow deficit $\binom{11-d}{2}$.  This gives the displayed exhaustive
premiums.  At $K'=42$, the all-core chart is maximized at core 41 with value
$9275468231667667$.  The kernel capacity, rank-nine marks, full-rank
capacity, and total capacity are
\[
\begin{split}
K_{\rm cap}&=17118437512304869669174162767031044097864380317587792249408,\\
G&=39184442698700870613155525220929985265240350768921233121073054790,\\
F_{\rm cap}&=910235915731681560754887897376310981744678105658368636509811656,\\
\operatorname{Cap}&=910253034169193865624557071539078012788775970038686224302061064.
\end{split}
\]
The sharp isolated-incidence demand is
\[
D_{42}=N_{\min}\binom{67514}{11}-\binom{1048618}{11},
\]
with exact value
\[
914334065220784060110315659375128857904243875224718721493122240.
\]
Subtraction gives the positive gap.  The coefficient of the record count
after clearing the 55-shadow denominator is
\[
143327002558240027109345932105933013269257449851>0,
\]
and the exact minimum-record cross-product is
\[
224456707837460696716722330982796481350734785231787345508364676>0.
\]
Thus the contradiction persists above the record floor.  Independent exact
replay of all six branches at $K'=43$ gives the displayed wall.
\end{proof}

\begin{proposition}[Descending-support completion ladder]
\label{prop:kb-rank11-descending-support-completion-ladder}
In the setup of
Proposition~\ref{prop:kb-rank11-cross-support-defect-carrier}, let $M_c$ be
the maximum number of completions of an independent deletion among
support-$c$ circuits.  Inspect source supports in the order
\[
5,4,3,2.
\]
At source $c$, terminate in defect branch $s$ when
$M_c=q-s$ for $0\le s\le9-c$; otherwise retain $M_c\le q-(10-c)$ and
continue to the next source.  These alternatives partition all completion
maxima into 26 terminal branches and one all-fallback branch.  Every
terminal branch may use the carrier cap of
Proposition~\ref{prop:kb-rank11-cross-support-defect-carrier}, and every
later branch retains all earlier fallback ceilings.
\end{proposition}

\begin{proof}
For fixed $c$, the integer $M_c$ lies in $[0,q]$.  The values
$q,q-1,\ldots,q-(9-c)$ are the terminal defects; their complement is
exactly $M_c\le q-(10-c)$.  Thus each stage is a disjoint exhaustive
partition.  Applying the four stages successively gives
\[
(10-5)+(10-4)+(10-3)+(10-2)+1=27
\]
leaves, of which 26 are terminal.  On a terminal leaf, the carrier theorem
applies with the recorded source and defect.  On a fallback leaf, the
recorded completion ceiling holds by construction.  Intersecting each new
cap with all inherited supportwise caps preserves every preceding bound.
\end{proof}

\begin{proposition}[$K'=43$ descending-support payment]
\label{prop:kb-rank11-k43-descending-support-payment}
In the fixed component target, every rank-nine alternative is empty at
$K'=43$.  Weighting all 27 leaves of
Proposition~\ref{prop:kb-rank11-descending-support-completion-ladder} by the
full shadow deficits $\binom{11-d}{2}$, the worst leaf is
\texttt{c5\_defect\_2}, with premium
\[
39510045591272162389536743615445318852720199164.
\]
Exact component demand exceeds complete capacity by
\[
4456829341030748859349785682161828589177837723939653522311506.
\]
The same 27-leaf payment first fails at $K'=44$, where capacity exceeds
demand by
\[
1729575114830772639212937201922834766923716849296098844264209.
\]
\end{proposition}

\begin{proof}
For each leaf, retain every inherited completion cap and intersect it with
the source and target carrier caps supplied by the recorded terminal
defect.  Sum the eight support contributions with weights
$36,28,21,15,10,6,3,1$, maximize once over the 27 leaves, and combine that
premium with the unchanged integral core chart, all kernel-corank terms,
and the sharp isolated-incidence demand.  Exact integer replay gives the
displayed active leaf and positive gap.  Repeating the same exhaustive
ledger at $K'=44$ gives the displayed negative gap.
\end{proof}

\begin{proposition}[Completion branch-lattice refinement]
\label{prop:kb-rank11-completion-branch-lattice}
Let a previously certified branch carry supportwise incidence caps, and
choose any source support $2\le c\le9$.  Replacing that branch by the
terminal defects $0\le s\le9-c$ and the fallback
$M_c\le q-(10-c)$ gives $11-c$ disjoint exhaustive children.  Every child
inherits all parent caps; terminal children intersect them with the
cross-support carrier caps, and the fallback child intersects the source
cap with its new completion ceiling.  The operation may be repeated on any
selected children.
\end{proposition}

\begin{proof}
The integer partition is the one-stage argument in
Proposition~\ref{prop:kb-rank11-descending-support-completion-ladder}.
Inheritance and intersection can only lower a valid parent cap.  The
carrier theorem validates every additional terminal cap.  Since the
children partition the parent branch, replacing any collection of parent
leaves preserves exhaustiveness and disjointness.
\end{proof}

\begin{proposition}[$K'=44$ completion branch-lattice payment]
\label{prop:kb-rank11-k44-branch-lattice-payment}
In the fixed component target, every rank-nine alternative is empty at
$K'=44$.  Refine the two parent leaves
\texttt{c5\_defect\_2} and \texttt{c5\_defect\_3} from
Proposition~\ref{prop:kb-rank11-k43-descending-support-payment} at source
support six.  Each is replaced by four terminal defects and one fallback,
so the ledger has
\[
27-2+2\cdot5=35
\]
leaves.  Its worst leaf is
\texttt{c5\_defect\_2\_\_c6\_defect\_2}, with premium
\[
40318474413130846902399237147930487840413149400.
\]
Exact demand and capacity are respectively
\[
\begin{split}
D_{44}&=914632087688377144021446114681200227193194740473705158570542108,\\
\operatorname{Cap}_{44}
&=914096453278432212124943531506920798021655606781504193772003859,
\end{split}
\]
and hence demand exceeds capacity by
\[
535634409944931896502583174279429171539133692200964798538249.
\]
The same 35-leaf payment first fails at $K'=45$, where capacity exceeds
demand by
\[
5651502053446174523626296867091469400380654135040887972894842.
\]
\end{proposition}

\begin{proof}
Proposition~\ref{prop:kb-rank11-completion-branch-lattice} proves that the
35 leaves are exhaustive and that every inherited and new cap is valid.
Weight each leaf by the full shadow deficits, maximize once over the leaf
premiums, and retain every preceding chart, kernel, and isolated-demand
term.  Exact integer replay yields the displayed active leaf, demand,
capacity, and gap.  The adjacent replay at $K'=45$ yields the stated method
wall.
\end{proof}

\begin{proposition}[Joint support-four/support-five zero carrier]
\label{prop:kb-rank11-support45-joint-zero-carrier}
Put $q=K'-10$.  Suppose independent support-four and support-five
deletions have respectively $q-s_4$ and $q-s_5$ circuit completions, with
$q>s_4+s_5$.  Let $U_4,U_5$ be their completion carriers and let
$H_4,H_5\le V'$ be the subspaces vanishing on those carriers.  Then
\[
 \dim(H_4\cap H_5)\ge4,
 \qquad |U_4\cup U_5|\le q+6.                       \tag{JZ1}
\]
Close $H_4\cap H_5$ under all its common zeros in the residual evaluation
set.  If $B$ is that zero set and
\[
 H_B=\{f\in V':f|_B=0\},\qquad t=\dim H_B,
 \qquad \delta=K'-t-|B|,
\]
then
\[
 4\le t\le6,
 \qquad 0\le\delta\le\min(s_4,s_5).                \tag{JZ2}
\]
For every independent three-set $A$, the union of $A$ and all its
support-four circuit completions has at most $\delta+3$ points outside
$B$.
\end{proposition}

\begin{proof}
For $c\in\{4,5\}$ write $A_c$ for the independent $(c-1)$-deletion and
$Z_c$ for its $q-s_c$ completions.  The private completion coordinates give
a $(q-s_c)$-dimensional label space supported on
$U_c=A_c\cup Z_c$, so
\[
 |U_4|=q+3-s_4,qquad |U_5|=q+4-s_5,qquad
 \dim H_4=7,qquad\dim H_5=6.                       \tag{1}
\]
Grassmann gives $\dim(H_4\cap H_5)\ge3$.  If equality held, then
$H_4+H_5=V'$.  The polynomial common-root bound applied to the
three-dimensional intersection would give
$|U_4\cup U_5|\le K'-3=q+7$.  Equation~(1) would then force
\[
 |U_4\cap U_5|\ge q-s_4-s_5>0.
\]
At a point of this overlap every polynomial in $H_4+H_5=V'$ vanishes,
contrary to the empty common zero set of $V'$.  Thus the intersection has
dimension at least four, and the same root bound proves (JZ1).

The zero closure $B$ contains both carriers and satisfies
$H_4\cap H_5\subseteq H_B\subseteq H_5$, proving $4\le t\le6$.
The root bound gives $|B|\le K'-t$, hence $\delta\ge0$.  The evaluation
functionals on $B$ are independent before restriction to $V'$, and their
rank on $V'$ is $10-t$.  Therefore the label subspace supported on $B$ has
dimension
\[
 |B|-(10-t)=q-\delta.
\]
It contains both private completion-label spaces, of dimensions
$q-s_4$ and $q-s_5$, so $\delta\le\min(s_4,s_5)$.

Finally let $U_A$ be the carrier of an arbitrary independent three-set and
let $H_A\le V'$ be its seven-dimensional vanishing space.  Then
$\dim(H_A\cap H_B)\ge t-3$.  The root bound gives
\[
 |U_A\cup B|\le K'-(t-3).
\]
Since $|B|=K'-t-\delta$, subtraction yields
$|U_A\setminus B|\le\delta+3$.
\end{proof}

\begin{proposition}[Support-four external-carrier charge]
\label{prop:kb-rank11-support4-external-charge}
Under Proposition~\ref{prop:kb-rank11-support45-joint-zero-carrier}, put
$b=K'-t-\delta$ and $N=m'-b$.  Define
\[
 C_4(t,\delta)=
 \begin{cases}
 \binom b4,&\delta=0,\\[2mm]
 \displaystyle \binom b4+
 \sum_{j=1}^4\floor{
  \frac{\binom b{4-j}\binom N{j-1}(\delta+4-j)}j},&\delta>0.
 \end{cases}                                      \tag{EC1}
\]
The number of support-four circuit supports is at most $C_4(t,\delta)$,
and their selected eleven-set incidence is at most
\[
 C_4(t,\delta)\binom{m'-4}{7}.                     \tag{EC2}
\]
A branch with terminal defects $s_4,s_5$ may therefore use the maximum of
(EC2) over $4\le t\le6$ and
$0\le\delta\le\min(s_4,s_5)$.
\end{proposition}

\begin{proof}
If $\delta=0$, every support-four label has a representation on $B$.
Comparing it with its circuit representation uses at most
$|B|+4=K'-t+4\le K'$ evaluation functionals, so Vandermonde independence
puts the circuit support inside $B$.

Suppose $\delta>0$, and fix a circuit with exactly $j$ points outside $B$.
Deleting one outside point leaves an independent three-set with $j-1$
external points.  Proposition~\ref{prop:kb-rank11-support45-joint-zero-carrier}
allows at most $\delta+4-j$ external completions.  There are at most
$\binom b{4-j}\binom N{j-1}$ possible deleted three-sets, and each circuit
is charged exactly $j$ times.  Division by $j$, integer flooring, and
summation over $j=1,\ldots,4$ prove (EC1).  Each circuit extends to at most
$\binom{m'-4}{7}$ selected eleven-sets, proving (EC2).
\end{proof}

\begin{proposition}[$K'=45$ full completion-product payment]
\label{prop:kb-rank11-k45-full-completion-product-payment}
In the fixed component target, every rank-nine alternative is empty at
$K'=45$.  Refine the completion maximum independently at every source
support $c=2,\ldots,9$.  The resulting Cartesian product has
\[
 \prod_{c=2}^9(11-c)=9!=362880
\]
disjoint exhaustive leaves.  On the $259200$ leaves where supports four
and five are both terminal, also intersect the support-four incidence cap
with Proposition~\ref{prop:kb-rank11-support4-external-charge}.

The unique worst final leaf is the all-fallback leaf, with premium
\[
40126324034612056409620566967689123241580103372.
\]
This is below the exact safe premium ceiling by
\[
323417025195949241219620287627046670164885155.
\]
Complete component demand exceeds complete capacity by
\[
1616971801308361526826641488053709685917408248376428345137933.
\]
The same payment first fails at $K'=46$, where complete capacity exceeds
demand by
\[
5057508862309072579343840146913199075599800084788396842011438.
\]
\end{proposition}

\begin{proof}
At source support $c$, the branch-lattice theorem partitions the completion
maximum into terminal values $q-s$, $0\le s\le9-c$, and the fallback
$M_c\le q-(10-c)$.  Their Cartesian product is exhaustive, and every child
inherits all preceding support caps.  Exactly
\[
 6\cdot5\cdot(9\cdot8\cdot5\cdot4\cdot3\cdot2)=259200
\]
leaves have terminal support-four and support-five defects.  On each,
$q>s_4+s_5$, so the joint zero-carrier and external-charge propositions
apply before the exact deficits $\binom{11-c}{2}$ are assigned.

Streaming exact enumeration of all leaves gives the displayed maximum.  At
$K'=45$, retaining every kernel corank and the all-core rank-nine chart gives
\[
\begin{split}
K_{\rm cap}&=
20541224524206770168358957675733727352177985505266924004331,\\
G&=39188933245978335442844817333731704607806583940216121605190822065,\\
F_{\rm cap}&=
913143619008078468726311141809692801590446007571848493844592126,\\
\operatorname{Cap}_{45}&=
913164160232602675496479500767368535317798185557353760768596457.
\end{split}
\]
Sharp isolated incidence gives demand
\[
D_{45}=
914781132033911037023306142255422245003715593805730189113734390.
\]
Their difference is the positive gap above.  The cleared record coefficient
and floor-record cross are positive, respectively
\[
142851171448156502676296279588073248696830521058
\]
and
\[
88933449071959883975465281842954032725457453660703558982586308,
\]
so the contradiction persists above the record floor.  Independent replay
at $K'=46$ again selects the all-fallback leaf and gives the stated wall.
\end{proof}

\begin{proposition}[Exact deep support-four/support-five defect partition]
\label{prop:kb-rank11-support45-deep-defect-partition}
For one residual record, let $M_c$ be the maximum number of circuit
completions of an independent deletion at support $c\in\{4,5\}$, with
$M_c=0$ if that circuit stratum is empty.  Put
\[
 s_c=q-M_c.
\]
Then $0\le s_4,s_5\le q$, and the $(q+1)^2$ exact pairs
$(s_4,s_5)$ form a disjoint exhaustive partition.  On each pair, every
inherited deletion and cross-support cap remains available, the source-$c$
deletion ceiling is $q-s_c$, and the joint zero-carrier and external
support-four charge apply whenever
\[
 s_4+s_5<q.                                           \tag{ED1}
\]
Moreover, for nonnegative deficit weights, any cap vector componentwise
dominated by another cap vector may be discarded when maximizing the
premium.
\end{proposition}

\begin{proof}
On a nonempty stratum, finiteness gives an attaining deletion and the
universal completion theorem gives $1\le M_c\le q$; the empty-stratum
convention gives $M_c=0$.  Thus $s_c=q-M_c$ is a unique integer in
$[0,q]$, proving the exact product partition.  The source incidence cap is
\[
 \left\lfloor \frac{\binom{m'}{c-1}}{c}
 \max_{0\le b\le q-s_c} b\binom{m'-c+1-b}{11-c}\right\rfloor,
\]
and intersecting it and every valid cross-support cap with the parent cap
loses no inherited resource.

Under (ED1), both maxima are positive, so attaining deletions exist, and
$q>s_4+s_5$ is precisely the hypothesis of
Propositions~\ref{prop:kb-rank11-support45-joint-zero-carrier} and
\ref{prop:kb-rank11-support4-external-charge}.  Finally, if
$a_c\le b_c$ for all $c$ and $w_c\ge0$, then for every later cap vector
$u$,
\[
 \sum_c w_c\min(a_c,u_c)\le
 \sum_c w_c\min(b_c,u_c).
\]
Hence a dominated vector cannot maximize while its dominator remains.
\end{proof}

\begin{proposition}[$K'=46$ through $K'=53$ deep-joint payment]
\label{prop:kb-rank11-k46-k53-deep-joint-payment}
In the fixed component target, every rank-nine alternative is empty for
\[
 46\le K'\le53.
\]
Use Proposition~\ref{prop:kb-rank11-support45-deep-defect-partition} at
supports four and five.  The option product at supports
$2,3,6,7,8,9$ has $8640$ raw vectors, $1182$ distinct vectors, and exactly
nine componentwise maximal vectors.  Exhausting every exact defect pair
against those nine vectors selects
\[
 s_4=s_5=\left\lfloor\frac{q-1}{2}\right\rfloor,
 \qquad\hbox{all other supports fallback}.             \tag{ED2}
\]
The minimum demand-capacity gap is attained at $K'=53$ and equals
\[
2503373059664320603163477388007627909210651834842589498907998.
\]
The same payment first fails at $K'=54$, where capacity exceeds demand by
\[
2477882110233058360154706764229180240778698202487636349407165.
\]
\end{proposition}

\begin{proof}
The other-support option counts are $9,8,5,4,3,2$.  Deduplication and the
componentwise comparison in the preceding proposition leave nine maximal
vectors on each row.  Exact enumeration of the $(q+1)^2$ defect pairs gives
(ED2).  At the last safe row the premium is
\[
40141995068282471040632276636969794558475239350,
\]
which lies below the exact safe premium ceiling by
\[
500709701466164598631229888808678512345658964.
\]
Combining this premium with the unchanged all-core rank-nine marks, every
kernel-corank capacity, and the sharp isolated demand gives the displayed
positive gap.  The cleared record coefficient and floor-record cross are
positive on every row $46\le K'\le53$, so the contradiction persists above
the record floor.

At $K'=54$, (ED2) has $s_4=s_5=21$, but its premium exceeds the exact safe
ceiling by
\[
495611154275787830253977941644262122450512788.
\]
The resulting exact capacity excess is the one displayed above.
\end{proof}

\begin{proposition}[Small-support same-source collision charge]
\label{prop:kb-rank11-small-support-self-collision-charge}
Fix a support $2\le c\le5$.  Let $M_c=q-s$ be the maximum number of
support-$c$ circuit completions of an independent $(c-1)$-deletion, with
$M_c=0$ and $s=q$ when the circuit stratum is empty.  If the stratum is
nonempty, choose an attaining deletion and let $B$ be its completion
carrier.  Then
\[
 b=|B|=q+c-1-s,
\]
and the carrier $U_A$ of every independent $(c-1)$-deletion $A$ satisfies
\[
 |U_A\setminus B|\le s+c-1.                         \tag{SC1}
\]
Writing $N=m'-b$, the number of support-$c$ circuit supports is at most
\[
 \begin{cases}
 \binom bc,&s=0,\\[2mm]
 \displaystyle \binom bc+
 \sum_{j=1}^{c}\left\lfloor
 \frac{\binom b{c-j}\binom N{j-1}(s+c-j)}{j}
 \right\rfloor,&0<s<q,\\[4mm]
 0,&s=q.
 \end{cases}                                        \tag{SC2}
\]
Its selected eleven-set incidence is at most (SC2) multiplied by
$\binom{m'-c}{11-c}$.
\end{proposition}

\begin{proof}
If a support-$c$ circuit exists, deleting one point gives an independent
set with at least that point as a completion.  Thus $M_c=0$ is equivalent
to an empty stratum.  In the nonempty case, for every independent deletion
$A$ put
\[
 H_A=\{f\in V':f|_A=0\}.
\]
Independent evaluation gives $\dim H_A=11-c$, and $H_A$ vanishes on the
complete carrier $U_A$.  For the attaining deletion $A_0$, Grassmann's
inequality gives
\[
 \dim(H_A\cap H_{A_0})\ge2(11-c)-10=12-2c>0.
\]
This intersection vanishes on $U_A\cup B$.  The common-root bound and
$|B|=q+c-1-s$ therefore give
\[
 |U_A\setminus B|
 \le K'-(12-2c)-|B|=s+c-1,
\]
proving (SC1).

If $s=0$, the attaining completion labels span the full $q$-dimensional
annihilator.  Every support-$c$ circuit label has a representation on $B$.
Comparing it with the minimal circuit support uses at most
$|B|+c=q+2c-1\le q+9<K'+1$ coordinates, so Vandermonde independence forces
that circuit into $B$.

Now suppose $0<s<q$.  A circuit with exactly $j$ points outside $B$ is
exposed by deleting each of those $j$ points.  The remaining deletion has
$j-1$ outside points, and (SC1) permits at most $s+c-j$ outside
completions.  There are at most
$\binom b{c-j}\binom N{j-1}$ such deletions.  Dividing by the exact
multiplicity $j$, flooring, and summing proves (SC2).  Finally, each circuit
extends to at most $\binom{m'-c}{11-c}$ selected eleven-sets.
\end{proof}

\begin{proposition}[$K'=54$ through $K'=59$ small-support collision payment]
\label{prop:kb-rank11-k54-k59-small-support-collision-payment}
In the fixed component target, every rank-nine alternative is empty for
\[
 54\le K'\le59.
\]
Refine the completion maxima at supports $2,3,4,5$ to every exact defect,
apply Proposition~\ref{prop:kb-rank11-small-support-self-collision-charge},
retain the joint support-four charge when $s_4+s_5<q$, and retain every
support-$6$ through support-$9$ terminal/fallback branch.  Pareto compression
of the groups
\[
 (2,3),\qquad(4,5),\qquad(6,7,8,9)
\]
leaves respectively $1$, $1$, and $7$ maximal vectors.  The active branch
on every row $54\le K'\le60$ is
\[
 s_2=s_3=s_4=s_5=\left\lfloor\frac q2\right\rfloor,
 \qquad c6F/c7F/c8F/c9F.                             \tag{SC3}
\]
The minimum demand-capacity gap is attained at $K'=59$ and equals
\[
2662571195028360324230500777441238424043251068116179184680206.
\]
The same payment first fails at $K'=60$, where capacity exceeds demand by
\[
3672733965923291717387950853821894967875078243379846951201638.
\]
\end{proposition}

\begin{proof}
The first two groups contain $(q+1)^2$ exact choices each; the final group
contains $5\cdot4\cdot3\cdot2=120$ terminal/fallback choices.  Duplicate
cap vectors may be identified, and a componentwise dominated vector may be
discarded because every deficit weight is nonnegative and every later
operation only intersects caps.  Exact enumeration leaves the three
frontiers stated above and selects (SC3).

For each row, combine the resulting premium $P$ with the unchanged
rank-nine marks $G$ and residual record floor $R$ as
\[
 \operatorname{Cap}_{\rm full}=\left\lfloor\frac{G+RP}{55}\right\rfloor.
\]
Adding every kernel-corank capacity and comparing with
$R\binom{m'}{11}-\binom{n'}{11}$ gives positive gaps on $54\le K'\le59$,
with the displayed minimum at the last row.  The cleared record coefficient
and floor-record cross are positive on all six rows, so the contradiction
persists above the record floor.  At $K'=60$, exact replay again selects
(SC3), now with all four defects equal to $25$, and gives the displayed
capacity excess.
\end{proof}

\begin{proposition}[Small-source cross-support collision charge]
\label{prop:kb-rank11-cross-support-collision-charge}
Fix source support $2\le c\le5$ and target support $2\le d\le9$ with
$c+d\le11$.  Suppose the exact source completion maximum is
$M_c=q-s>0$, and let $B$ be an attaining source carrier.  Then
\[
 |B|=b=q+c-1-s,
\]
and the carrier $U_A$ of every independent target $(d-1)$-deletion satisfies
\[
 |U_A\setminus B|\le s+d-1.                         \tag{XC1}
\]
Writing $N=m'-b$, the number of target support-$d$ circuits is at most
\[
 \begin{cases}
 \binom bd,&s=0,\\[2mm]
 \displaystyle \binom bd+
 \sum_{j=1}^{d}\left\lfloor
 \frac{\binom b{d-j}\binom N{j-1}(s+d-j)}{j}
 \right\rfloor,&0<s<q.
 \end{cases}                                        \tag{XC2}
\]
Their selected eleven-set incidence is at most (XC2) multiplied by
$\binom{m'-d}{11-d}$.  No target cap is inferred from an empty source
stratum $s=q$.
\end{proposition}

\begin{proof}
Let $A_0$ be the attaining source deletion.  For any independent target
deletion $A$, the two deletion-vanishing spaces have dimensions $11-c$ and
$11-d$.  Grassmann gives intersection dimension at least
\[
 (11-c)+(11-d)-10=12-c-d>0.
\]
The intersection vanishes on $B\cup U_A$.  The common-root bound gives
\[
 |U_A\setminus B|
 \le K'-(12-c-d)-(q+c-1-s)=s+d-1,
\]
proving (XC1).

At $s=0$, the source-completion labels span the full $q$-dimensional
annihilator.  Every target circuit label has a representation on $B$.
Comparing it with its minimal target support uses at most
$|B|+d=q+c+d-1\le K'$ coordinates, so Vandermonde independence forces the
target circuit into $B$.

For $0<s<q$, a target circuit with $j$ points outside $B$ is exposed by
each of its $j$ outside deletions.  The remaining deletion already contains
$j-1$ outside points, so (XC1) permits at most $s+d-j$ outside completions.
There are at most $\binom b{d-j}\binom N{j-1}$ such deletions.  Division by
the exact exposure multiplicity $j$, flooring, and summation prove (XC2).
The final selected-set multiplier is immediate.
\end{proof}

\begin{proposition}[$K'=60$ through $K'=70$ cross-support collision payment]
\label{prop:kb-rank11-k60-k70-cross-support-collision-payment}
In the fixed component target, every rank-nine alternative is empty for
\[
 60\le K'\le70.
\]
Refine source supports $2,3,4,5$ to every exact defect.  On each nonempty
source stratum apply
Proposition~\ref{prop:kb-rank11-cross-support-collision-charge} to all
targets with $c+d\le11$.  Retain the same-support caps, the joint
support-four charge when $s_4+s_5<q$, and every support-$6$ through
support-$9$ terminal/fallback branch.  Grouped Pareto compression selects
\[
 s_2=s_3=s_4=s_5=\left\lceil\frac q2\right\rceil,
 \qquad c6F/c7F/c8F/c9F                              \tag{XC3}
\]
on every row $60\le K'\le71$.  The minimum demand-capacity gap is attained
at $K'=70$ and equals
\[
854274172985042754802177028749324962520517760595473749602211.
\]
The same payment first fails at $K'=71$, where capacity exceeds demand by
\[
824875968499878215752683873455674299360608616555107905777434.
\]
\end{proposition}

\begin{proof}
The exact defects at supports $2$ through $5$, the existing joint charge,
and all $5\cdot4\cdot3\cdot2=120$ higher-support choices are exhaustive.
No mixed-support cap is applied at an empty source value.  Duplicate vectors
may be identified, and componentwise dominated vectors discarded, because
all deficit weights are nonnegative and every later operation intersects
caps.  Exact enumeration gives (XC3).

Combine its premium $P$ with the unchanged rank-nine marks $G$ and residual
record floor $R$ as
\[
 \operatorname{Cap}_{\rm full}=\left\lfloor\frac{G+RP}{55}\right\rfloor.
\]
Adding every kernel-corank capacity and comparing with
$R\binom{m'}{11}-\binom{n'}{11}$ gives positive gaps on $60\le K'\le70$,
with the displayed minimum at the last safe row.  The cleared record
coefficient and floor-record cross are positive on all eleven rows.  At
$K'=71$, exact replay again selects (XC3), now with all four defects equal
to $31$, and gives the displayed capacity excess.
\end{proof}

\begin{proposition}[Fixed-union multicarrier collision charge]
\label{prop:kb-rank11-multicarrier-collision-charge}
Let $D$ be a fixed set of $u$ domain points and let $W\le V'$ have
dimension $g$, with every polynomial in $W$ vanishing on $D$.  For target
support $2\le d\le9$, put
\[
 r_d=g+1-d,\qquad R_d=K'-r_d-u.
\]
Whenever $r_d>0$, every independent target $(d-1)$-deletion carrier has at
most $R_d$ points outside $D$.  Consequently the number of target
support-$d$ circuit supports is at most
\begin{equation}
 \binom ud+\sum_{j=1}^{d}\left\lfloor
 \frac{\binom u{d-j}\binom{m'-u}{j-1}
       \max\{0,R_d-j+1\}}{j}
 \right\rfloor .                                    \tag{MC1}
\end{equation}
Their selected eleven-set incidence is at most (MC1) multiplied by
$\binom{m'-d}{11-d}$.
\end{proposition}

\begin{proof}
For an independent target deletion $A$, let
$H_A=\{f\in V':f|_A=0\}$.  Then $\dim H_A=11-d$, so Grassmann gives
\[
 \dim(W\cap H_A)\ge g+(11-d)-10=r_d.
\]
This intersection vanishes on $D$ and on the complete target carrier.
The common-root bound therefore leaves at most $K'-r_d-u=R_d$ carrier
points outside $D$.

There are at most $\binom ud$ target circuits wholly inside $D$.  A circuit
with exactly $j\ge1$ points outside $D$ is exposed by each of its $j$
outside deletions.  Such a deletion already contains $j-1$ outside points,
leaving at most $\max\{0,R_d-j+1\}$ outside completions.  There are at most
$\binom u{d-j}\binom{m'-u}{j-1}$ such deletions.  Division by the exact
exposure multiplicity $j$, flooring, and summation prove (MC1).  The final
selected-set multiplier is immediate.
\end{proof}

\begin{proposition}[Support-two/three/four carrier-position trichotomy]
\label{prop:kb-rank11-k71-carrier-position-trichotomy}
Let $M_c=q-s_c>0$ be exact support-$c$ completion maxima, with attaining
carriers $B_c$.  The support-two carrier is one full nonzero parallel class
of size $b_2=M_2+1$.

For an attaining support-$c$ deletion, $c=3,4$, the projective point of
$B_2$ is either transverse to the deletion span, contained in the span of a
proper deletion subset, or contained only in the full deletion span.  In the
last case every point of $B_2$ is an exact support-$c$ completion, so
$M_c\ge M_2+1$.

Hence, when $M_3\le M_2$, only the transverse and proper-span positions
occur.  Their fixed unions have respectively
\[
 (u,g)=(b_2+b_3,7),\qquad (b_2+b_3-1,8),             \tag{PT1}
\]
and $s_2+s_3<q$ is impossible.  When
$M_3=M_4=M_2+1$, the remaining completion position is exhausted by
\[
\begin{array}{c|cc}
 &u&g\\ \hline
T23&2M_2+4&7\\ A23&2M_2+3&8\\
T24&2M_2+5&6\\ A24&2M_2+4&7\\
N34&M_2+6&6\\ N34A&M_2+5&7.
\end{array}                                           \tag{PT2}
\]
Every row of (PT1) and (PT2) feeds
Proposition~\ref{prop:kb-rank11-multicarrier-collision-charge}.
\end{proposition}

\begin{proof}
Write $v_x$ for evaluation at a domain point as a functional on $V'$.  The
empty common-zero hypothesis makes every $v_x$ nonzero.  An attaining
support-two deletion is one point $a$, and its completions are exactly the
domain evaluations parallel to $v_a$.  Thus $B_2$ is a full projective
parallel class.

Let $A_c$ be an attaining independent deletion and
$F_c=\langle v_x:x\in A_c\rangle$.  If the projective point of $B_2$ is
outside $F_c$, then $B_2\cap B_c=\varnothing$ and the two vanishing spaces
intersect in dimension at least $10-c$.  If it lies in a proper deletion
span, then $H_c\le H_2$; the carriers share one anchor in rank one and are
otherwise disjoint.  If it lies in no proper deletion span, every point of
$B_2$ completes $A_c$ minimally, forcing $M_c\ge|B_2|=M_2+1$.  This proves
the three positions and (PT1).  Applying the common-root bound to either
fixed union in (PT1) also gives $s_2+s_3\ge q$.

Suppose $M_3=M_4=M_2+1$.  If support three is transverse or anchored, its
two alternatives are $T23$ and $A23$.  Otherwise
$B_3=B_2\cup A_3$.  Applying the same analysis to support four gives $T24$
and $A24$, unless $B_4=B_2\cup A_4$.  In that final case $H_3,H_4\le H_2$
and $\dim(H_3\cap H_4)\ge6$.  The anchor sets cannot share both points of
$A_3$, since that would put the support-two projective point in a proper
span of $A_4$.  Disjoint anchors give $N34$; one shared anchor gives
$H_4\le H_3$ and $N34A$.  The six alternatives are exhaustive.
\end{proof}

\begin{proposition}[$K'=71$ carrier-trichotomy payment]
\label{prop:kb-rank11-k71-carrier-trichotomy-payment}
In the fixed component target, every rank-nine alternative is empty at
$K'=71$.  Refine every exact support-two/support-three defect pair by
Proposition~\ref{prop:kb-rank11-k71-carrier-position-trichotomy}; retain all
six cases in (PT2) when $M_3=M_4=M_2+1$, together with every inherited cap
and support-six through support-nine branch.  Exact replay selects
\[
 s_2=33,\qquad s_3=s_4=s_5=31,\qquad c6F/c7F/c8F/c9F,
\]
with completion premium
\[
41052480732722315950912559282994987357541498029.
\]
Demand exceeds complete capacity by
\[
118872281099445772155993127155914865045379156488810154591370.
\]
The same complete replay first fails at $K'=72$, where capacity exceeds
demand by
\[
4821537739796415753639473905341364357966460110033651367468100.
\]
\end{proposition}

\begin{proof}
For $M_3\le M_2$, retain both cases in (PT1) and remove the impossible
pairs $s_2+s_3<q$; at $K'=71$ this excludes exactly $961$ defect pairs.
Pairs with $M_3=M_2+1$ retain their position provenance until support four
is fixed.  When also $M_4=M_2+1$, replace the formerly uncharged leaf by all
six alternatives in (PT2).  Supports four and five retain every exact
defect pair and the joint support-four charge when $s_4+s_5<q$.  All 120
support-six through support-nine terminal/fallback choices remain explicit.
Componentwise dominated vectors may be removed because every deficit weight
is nonnegative and all later operations take componentwise minima.

For each surviving leaf, form
\[
 P=\sum_{d=2}^{9}\binom{11-d}{2}\operatorname{cap}_d,
 \qquad
 \operatorname{Cap}_{\rm full}=\left\lfloor\frac{G+RP}{55}\right\rfloor,
\]
and add every kernel-corank capacity.  The displayed $K'=71$ branch is the
exact maximum.  Its premium lies below the safe ceiling by
$23776122440930417094576446937038395558574009$, yielding the stated
positive integral gap.  At $K'=72$ the exact maximizer has
$M_3=M_2+2$, outside this one-step trichotomy, and gives the stated capacity
excess.  This is a method wall, not a counterexample.
\end{proof}

\begin{proposition}[Full-completion pairwise carrier atlas]
\label{prop:kb-rank11-full-completion-pairwise-carrier-atlas}
Retain the notation of
Proposition~\ref{prop:kb-rank11-k71-carrier-position-trichotomy}.  For an
attaining support-$c$ carrier put $b_c=M_c+c-1$.  Relative to the
support-two projective point, every support-$c$ carrier is exactly one of
the following:
\[
 T_c:(u,g)=(b_2+b_c,10-c),\qquad
 A_c:(u,g)=(b_2+b_c-1,11-c),                 \tag{FA1}
\]
or it is in full-completion position.

Suppose support three and $d\in\{4,5\}$ are both in full-completion
position.  Put
\[
 r_3=M_3-M_2+1,\qquad r_d=M_d-M_2+d-2.
\]
Then their residual overlap $t$ has the exact exhaustive range
\[
 0\le t\le\min\{r_3,M_d-M_2\},               \tag{FA2}
\]
and gives a fixed union
\[
 u=b_2+r_3+r_d-t,\qquad
 g=\begin{cases}10-d,&t=0,\\11-d,&t>0.\end{cases}       \tag{FA3}
\]
The support-four and support-five alternatives combine by Cartesian
product; no synchronization between their residual overlaps is assumed.
\end{proposition}

\begin{proof}
Let $F_c$ be the span of an attaining support-$c$ deletion.  If the
support-two projective point is outside $F_c$, or belongs to a proper
deletion span, the argument of
Proposition~\ref{prop:kb-rank11-k71-carrier-position-trichotomy} gives
$T_c$ or $A_c$ in (FA1).  Otherwise every support-two carrier point is a
minimal completion, which is the full-completion position.

In the simultaneous full case, quotient by the support-two projective
point.  The support-three carrier has $r_3$ residual points and the
support-$d$ carrier has $r_d$ residual points.  Minimality of the latter
deletion forces at least $d-2$ of its anchors outside the support-three
rank-two flat.  Therefore at most $M_d-M_2$ residual points can overlap,
while the support-three side supplies at most $r_3$.  Every integer in
(FA2) is the corresponding intersection-size case.  Rank submodularity
gives dimension $10-d$ for disjoint residuals and one additional common
vanishing dimension for positive overlap, proving (FA3).  Taking the
independent support-four and support-five position lists gives the stated
Cartesian product.
\end{proof}

\begin{proposition}[Base-field-normalized flat-coupled support-four/five census]
\label{prop:kb-rank11-flat-coupled-support45-charge}
Let a $g$-dimensional subspace $W\le V'$ vanish on a fixed union $D$ of
size $u$, where $g\ge5$, and put
\[
 R=K'-u-g,\qquad B=R+3,\qquad N=m'-u.          \tag{FC1}
\]
Assume $R\ge0$, $N>B$, and $K'\ge g+5$.
Let $L_4,L_5$ be the exact support-four and support-five circuit counts in
which at least one point lies in $D$, obtained by outside-deletion exposure:
\[
 L_d=\binom ud+
 \sum_{j=1}^{d-1}\left\lfloor
 \frac{\binom u{d-j}\binom N{j-1}R}{j}
 \right\rfloor.                               \tag{FC2}
\]
For the all-outside strata put
\[
 X_4=\min\left\{
 \left\lfloor\frac{R\binom N3}{4}\right\rfloor,
 \left\lfloor\frac{R\binom N4}{N-B}\right\rfloor
 \right\},                                    \tag{FC3}
\]
\[
 X_5=\left\lfloor
 \frac{R\binom N4-(N-B)X_4}{5}
 \right\rfloor.                               \tag{FC4}
\]
Whenever $K'\ge g+5$, the total selected support-four/support-five
incidence is at most
\[
 21(L_4+X_4)\binom{m'-4}{7}
 +15(L_5+X_5)\binom{m'-5}{6}.                 \tag{FC5}
\]
This is a joint charge in the original outside evaluation matroid, not in
the evaluation matroid of $W$.
\end{proposition}

\begin{proof}
For an independent outside triple $A$, the space $W\cap H_A$ has dimension
at least $g-3$.  Divide its polynomials by the locator of $D$.  The
common-root degree bound shows that the rank-three flat generated by $A$
has size at most $B$ in the original outside evaluation matroid.  The same
argument for an independent four-set gives rank-four flat size at most
$B+1$.  Exact exposure of $j$ outside points gives (FC2).

Write $C_4,C_5$ for the all-outside circuit counts.  Triple exposure and
the rank-three flat bound give
\[
 4C_4\le R\binom N3,
 \qquad (N-B)C_4\le R\binom N4.                \tag{FC6}
\]
Expose an outside point from each four-set.  A support-five circuit spends
five units, while each support-four circuit consumes at least $N-B$ of the
same completion resource, so
\[
 5C_5\le R\binom N4-(N-B)C_4.                 \tag{FC7}
\]
The weighted objective in (FC5) is nondecreasing in $C_4$ under (FC7)
when $K'\ge g+5$; direct cancellation of adjacent binomial factors gives
this condition.  Hence its endpoint is (FC3)--(FC4).  Adding the exact
lower strata proves (FC5).
\end{proof}

\begin{proposition}[$K'=72,73$ full-carrier-atlas payment]
\label{prop:kb-rank11-k72-k73-full-carrier-atlas-payment}
In the fixed component target, every rank-nine alternative is empty at
$K'=72$ and $K'=73$.  Refine every exact support-two through support-five
defect cell by
Proposition~\ref{prop:kb-rank11-full-completion-pairwise-carrier-atlas},
apply Proposition~\ref{prop:kb-rank11-flat-coupled-support45-charge} to
every eligible fixed union, and retain every inherited cap and
support-six through support-nine terminal/fallback branch.

At $K'=72$ the global completion premium is
\[
41089877204729279662874647920595743958596178333,
\]
below its safe ceiling by
$10440735269654784698417860383073117137496667$, and demand exceeds
complete capacity by
\[
52200017935756118667066163970702686810349690944821612538425.
\]
At $K'=73$ the global premium is
\[
41123952182016259764480350052978913220338995240,
\]
below its safe ceiling by
$424186681733896065660200756292351785430362$, and demand exceeds complete
capacity by
\[
2120784774514292837614781442321448802184060878375874298355.
\]
\end{proposition}

\begin{proof}
At $K'=72$, the conservative stream contains $7{,}991{,}221$ leaves and
exactly $36$ above-ceiling defect tuples.  Their full-atlas reroute contains
$8{,}057$ exact charged leaves, all safe.  Seven disjoint geometry lanes
contain another $113{,}124{,}235$ leaves, all safe.  The largest rerouted
premium is
$31430973522237705446704558558267588754434271701$, and the largest
geometry premium is
$35086385884924358787108064569034674887022445357$.  Both are below the
largest already-safe conservative premium printed in the proposition.

At $K'=73$, the corresponding counts are $8{,}551{,}382$ conservative
leaves, $218$ exceptional tuples, $71{,}806$ exact reroutes, and
$118{,}892{,}669$ geometry leaves.  The largest rerouted and geometry
premiums are respectively
$32133901221158725309935103349312670983455197672$ and
$35688968442860327556985962346044983398767741600$, again below the
printed conservative premium.

For each row let $R_0=274980728111260126$ be the inherited record floor,
$G$ the unchanged rank-nine marks, and $K_{\rm ern}$ the inherited kernel
capacity.  Exact integer arithmetic uses
\[
 \operatorname{Cap}_{\rm full}
 =\left\lfloor\frac{G+R_0P}{55}\right\rfloor,
 \qquad
 \operatorname{Demand}=R_0\binom{m'}{11}-\binom{n'}{11}.
\]
Subtracting $K_{\rm ern}+\operatorname{Cap}_{\rm full}$ gives the two
positive gaps displayed above.
\end{proof}

\begin{proposition}[$K'=74,\ldots,78$ full-carrier-atlas payment]
\label{prop:kb-rank11-k74-k78-full-carrier-atlas-payment}
In the fixed component target, every rank-nine alternative is empty for
$74\le K'\le78$.  The exact global completion premiums are
\[
\begin{aligned}
P_{74}&=41148280931947468743645570894078252803553423792,\\
P_{75}&=41172442942616752083301734067206968949181736144,\\
P_{76}&=41196532360070121067065901849561255059392646057,\\
P_{77}&=41220567597231178491653121647453619450645179178,\\
P_{78}&=41244614753758628801860143341643171244170576450.
\end{aligned}
\]
Their exact positive margins below the safe premium ceilings are
\[
\begin{aligned}
M_{74}&=151207872362754098743792508682842947287405,\\
M_{75}&=42309306941100164131542711491454099991233,\\
M_{76}&=3343498479116787180719767118068070898738,\\
M_{77}&=15896082187058434500851976235556574069096,\\
M_{78}&=13865957965092782628483685711822437366639.
\end{aligned}
\]
After unchanged kernel capacity is added, demand exceeds complete capacity
by
\[
\begin{aligned}
G_{74}&=755986378881174710705888877305550381612101047062792678211,\\
G_{75}&=211531709609936637728192389252911579074455707755248457894,\\
G_{76}&=16716320840480454509867840664707020900684182599020170034,\\
G_{77}&=79474840980250192811328546134901531018600026108436267508,\\
G_{78}&=69324931221842548818784073274469783546449231843657033123.
\end{aligned}
\]
\end{proposition}

\begin{proof}
Every row uses the complete pairwise carrier atlas of
Proposition~\ref{prop:kb-rank11-full-completion-pairwise-carrier-atlas}, the
base-field-normalized support-four/support-five charge of
Proposition~\ref{prop:kb-rank11-flat-coupled-support45-charge}, every
inherited cap, and every support-six through support-nine branch.  The exact
coverage ledger is
\[
\begin{array}{c|r|r|r|r}
K'&\text{plain leaves}&\text{exceptional tuples}&\text{reroute leaves}&
\text{geometry leaves}\\ \hline
74& 8{,}869{,}588&   729&   338{,}149&124{,}851{,}888\\
75& 9{,}479{,}358& 1{,}995& 1{,}228{,}878&131{,}010{,}586\\
76& 9{,}817{,}234& 3{,}800& 2{,}689{,}092&137{,}366{,}383\\
77&10{,}475{,}101& 7{,}657& 6{,}623{,}568&143{,}928{,}183\\
78&10{,}837{,}645&11{,}552&10{,}115{,}441&150{,}693{,}396
\end{array}
\]
The exceptional tuples are canonically sorted and SHA-256 pinned in the
supplemental manifest.  Their complete reroute maxima are respectively
\[
\begin{aligned}
&32837046420997427924790365894433685328482435320,\\
&33540409167198343349228340842477255779539196790,\\
&34322333430853669387933267678391560032676518774,\\
&35026060643237953572770728698158639005273906176,\\
&36102481454819719165765129089266727844130971230,
\end{aligned}
\]
all below the corresponding safe ceiling.  The largest geometry premiums
are respectively
\[
\begin{aligned}
&36187398164184684606907419896584782044225756904,\\
&36818718942605286540367916391174777496642019789,\\
&37421755192298946544884408866447573734456950687,\\
&37927152249618150254202496789638286615058058210,\\
&38552393090078383126619984181387841064702317594,
\end{aligned}
\]
each below $P_{K'}$.  Finally, applying the same floor identity as in the
preceding proposition gives $G_{K'}>0$ row by row.
\end{proof}

\begin{proposition}[$K'=79,\ldots,82$ full-carrier-atlas payment and the
$K'=83$ pairwise wall]
\label{prop:kb-rank11-k79-k82-full-carrier-atlas-payment}
In the fixed component target, every rank-nine alternative is empty for
$79\le K'\le82$.  The exact global completion premiums are
\[
\begin{aligned}
P_{79}&=41268581039515451359223235239395447056496741638,\\
P_{80}&=41292698225299493655203544545324133071167827372,\\
P_{81}&=41316738803727121977844753592626079710298860916,\\
P_{82}&=41340768193812712537232048213849077199458005267.
\end{aligned}
\]
Their exact positive margins below the safe premium ceilings are
\[
\begin{aligned}
M_{79}&=90041821230008310313867459260514662563833,\\
M_{80}&=12651432337345356124303501333284331916464,\\
M_{81}&=9200732869233857603364615518182428928433,\\
M_{82}&=14269266235451710755032134164017670137549.
\end{aligned}
\]
After unchanged kernel capacity is added, demand exceeds complete capacity
by
\[
\begin{aligned}
G_{79}&=450177555678029181663225127809721029656401221977920938493,\\
G_{80}&=63252728650428501280250074863279553490282355009311230439,\\
G_{81}&=46000440427984175715360686273569212546246763017890480411,\\
G_{82}&=71341331255235112573601253115194194916813205966998398786.
\end{aligned}
\]

The same pairwise atlas first fails at $K'=83$.  Its active cell has defects
$(50,49,49,48)$, completion maxima $(23,24,24,25)$, and premium
\[
41411584407693108041789796771180703922717609427,
\]
which exceeds the safe ceiling by
\[
46770156546844646871611081711174519620031307.
\]
The two active $(29,6)$ fixed-union charges force a four-dimensional
intersection on a union of size $29,30,31$, or $32$, but every corresponding
ordinary fixed-union cap leaves the displayed premium unchanged.
\end{proposition}

\begin{proof}
Every closed row uses the same complete pairwise carrier atlas and
base-field-normalized support-four/support-five charge as
Proposition~\ref{prop:kb-rank11-k74-k78-full-carrier-atlas-payment}.  The
exact coverage ledger is
\[
\begin{array}{c|r|r|r|r}
K'&\text{plain leaves}&\text{exceptional tuples}&\text{reroute leaves}&
\text{geometry leaves}\\ \hline
79&11{,}546{,}087&19{,}406&19{,}114{,}557&157{,}671{,}136\\
80&11{,}929{,}729&26{,}104&24{,}962{,}791&164{,}858{,}603\\
81&12{,}695{,}949&39{,}570&40{,}569{,}326&172{,}265{,}121\\
82&13{,}101{,}284&49{,}900&48{,}822{,}291&179{,}887{,}680
\end{array}
\]
All exceptional-tuple digests, reroute maxima, and seven geometry-lane
maxima are pinned in the supplemental certificate.  The largest geometry
premiums are respectively
\[
\begin{aligned}
&39155883271390100828809596211624350117737071281,\\
&39668547314355452559959739197487065940915191238,\\
&40287409045971746658167168137122038667039311900,\\
&40891353398864895345494242236029970259845626460,
\end{aligned}
\]
each below $P_{K'}$.  Applying the same floor identity as in the preceding
proposition gives $G_{K'}>0$ row by row.

For the wall, write
\[
B_3=B_2\cup A,\qquad B_4=B_2\cup C,\qquad
B_5=B_2\cup A\cup D,
\]
where $|B_2|=24$, $|A|=2$, $|C|=|D|=3$, and $A\cap C=\varnothing$.
The pairwise atlas gives two six-dimensional vanishing spaces inside the
same eight-dimensional $H_3$, each on a $29$-point union.  Their
intersection has dimension at least four on
$B_2\cup A\cup C\cup D$, whose size is $32-|C\cap D|\in\{29,30,31,32\}$.
Exact replay of all four resulting caps gives the same premium and negative
margin displayed in the proposition.  The branch-free baseline is larger
still.  This proves only a method wall.  The next proposition supersedes its
frontier conclusion by a different adjacent-support coupling.
\end{proof}

\begin{proposition}[Adjacent-flat circuit coupling]
\label{prop:kb-rank11-adjacent-flat-circuit-coupling}
Let $M$ be a finite matroid on $N$ labelled elements.  Fix $r\ge1$ and
$B\ge r$, with $N\ge B$.  Suppose every rank-$r$ flat has at most $B$
elements and every rank-$(r+1)$ flat has at most $B+1$ elements.  If $C_j$
is the number of $j$-element circuit supports, then
\[
 (r+2)C_{r+2}\le (B-r)\binom{N}{r+1}-(N-B)C_{r+1}.
 \tag{AFC}\label{eq:kb-rank11-adjacent-flat-circuit-coupling}
\]
\end{proposition}

\begin{proof}
For a rank-$r$ flat $F$, put $b=|F|$ and let $t_F$ count its independent
$r$-set bases.  Every $(r+1)$-circuit contributes $r+1$ such bases, so
\[
 \sum_F t_F(b-r)\ge (r+1)C_{r+1}.                 \tag{1}
\]
For each basis $A$ of $F$ and $x\notin F$, the independent set
$Q=A\cup\{x\}$ has at most $B-b$ exact $(r+2)$-circuit completions.  Relative
to the unrestricted cap $B-r$, this witness saves at least $b-r$.  Each
independent $(r+1)$-set has only $r+1$ rank-$r$ subsets.  Hence, writing
$c(Q)$ for its exact completion count,
\[
 \sum_F t_F(N-b)(b-r)
 \le (r+1)\sum_Q\bigl((B-r)-c(Q)\bigr).           \tag{2}
\]
Since $b\le B$, (1)--(2) imply
\[
 \sum_Q\bigl((B-r)-c(Q)\bigr)\ge (N-B)C_{r+1}.
\]
There are at most $\binom{N}{r+1}$ independent sets $Q$, while every
$(r+2)$-circuit is exposed by its $r+2$ such subsets.  Substitution gives
\eqref{eq:kb-rank11-adjacent-flat-circuit-coupling}.
\end{proof}

\begin{proposition}[Single-completion carrier]
\label{prop:kb-rank11-single-completion-carrier}
Let the ten-dimensional correction space $V$ have empty common zero set.
Fix $2\le c\le9$.  Suppose an independent $(c-1)$-set $A$ attains exactly
$M_c>0$ support-$c$ circuit completions, and let $X_A$ be that completion
set.  Then
\[
 D_A=A\cup X_A,\qquad |D_A|=M_c+c-1,
 \qquad W_A=\{f\in V:f|_A=0\},\qquad \dim W_A=11-c,
\]
and every polynomial in $W_A$ vanishes on $D_A$.  Thus $(D_A,W_A)$ is a
fixed union and vanishing space for every fixed-union collision theorem;
no second positive support stratum is required.
\end{proposition}

\begin{proof}
Independence of the evaluation functionals on $A$ gives
$\dim W_A=10-(c-1)=11-c$.  A completion $x$ makes $A\cup\{x\}$ dependent,
so evaluation at $x$ lies in the span of the evaluations at $A$.
Consequently every polynomial vanishing on $A$ also vanishes on $x$.
Completion points are distinct domain points outside $A$, whence
$|D_A|=(c-1)+M_c$.
\end{proof}

\begin{proposition}[Base-field-normalized fixed-union adjacent-support census]
\label{prop:kb-rank11-fixed-union-adjacent-support-coupling}
Let the ten-dimensional correction space $V$ have empty common zero set.
Let a fixed $g$-dimensional subspace vanish on a fixed $u$-point set $D$.
Fix $2\le d\le g-1$, and put
\[
 R=K'-u-g\ge0,\qquad N=m'-u\ge R+d-1.
\]
If $C_{e,i}$ counts support-$e$ circuits meeting $D$ in exactly $i$ points,
then, for $0\le i\le d-2$,
\[
 (d+1-i)C_{d+1,i}+(N-R-d+1+i)C_{d,i}
 \le \binom{u}{i}R\binom{N}{d-i},                 \tag{FAS1}
\]
and
\[
 C_{d,i}\le
 \left\lfloor\frac{\binom{u}{i}R\binom{N}{d-1-i}}{d-i}\right\rfloor.
 \tag{FAS2}
\]
For nonnegative integer weights $w_d,w_{d+1}$, put
\[
\begin{aligned}
 L_i&=\left\lfloor
       \frac{\binom{u}{i}R\binom{N}{d-1-i}}{d-i}\right\rfloor,\\
 A_i&=\binom{u}{i}R\binom{N}{d-i},\\
 \lambda_i&=(d+1-i)w_d-(N-R-d+1+i)w_{d+1},\\
 J_i&=\left\lfloor
 \frac{w_{d+1}A_i+\max(\lambda_i,0)L_i}{d+1-i}
 \right\rfloor.
\end{aligned}
\]
Then the complete weighted census obeys
\begin{equation}
\tag{FAS3}
\begin{aligned}
 w_dC_d+w_{d+1}C_{d+1}\le{}&\sum_{i=0}^{d-2}J_i
 +w_d\left(\binom{u}{d-1}R+\binom{u}{d}\right)\\
 &+w_{d+1}\left(
 \left\lfloor\frac{\binom{u}{d-1}RN}{2}\right\rfloor
 +\binom{u}{d}R+\binom{u}{d+1}\right).
\end{aligned}
\end{equation}
Envelopes for support-disjoint adjacent pairs may be added.
\end{proposition}

\begin{proof}
Fix an independent $i$-set $S\subset D$, contract $S$, and restrict to the
$N$ points outside $D$.  A rank-$a$ flat with basis $A$ leaves at least
$g-a$ vanishing dimensions, so the common-root bound gives flat size at most
$R+a$.  Take $a=d-1-i$ and apply
Proposition~\ref{prop:kb-rank11-adjacent-flat-circuit-coupling} with
$B=R+a$; summing over $S$ proves (FAS1).  Exposing a support-$d$ circuit by
one of its $d-i$ outside points leaves exactly the residual completion budget
$R$, proving (FAS2).  Linear optimization has only the two endpoints from
(FAS2): after multiplying (FAS1) by $w_{d+1}$, the coefficient of
$C_{d,i}$ has numerator $\lambda_i$, giving $J_i$.  The remaining strata
follow by direct exposure:
\[
\begin{array}{ll}
C_{d,d-1}\le\binom{u}{d-1}R,&C_{d,d}\le\binom{u}{d},\\
C_{d+1,d-1}\le\left\lfloor\binom{u}{d-1}RN/2\right\rfloor,
&C_{d+1,d}\le\binom{u}{d}R,\quad C_{d+1,d+1}\le\binom{u}{d+1}.
\end{array}
\]
This proves (FAS3).  Finally, support-disjoint pairs contain disjoint census
terms.
\end{proof}

\begin{proposition}[$K'=83$ adjacent-support carrier payment]
\label{prop:kb-rank11-k83-adjacent-support-payment}
In the fixed component target, every rank-nine alternative is empty at
$K'=83$.  The exact global completion premium is
\[
 P_{83}=41364793335621487128860475977676014245181683050,
\]
below the safe premium ceiling by
\[
 M_{83}=20915524776266057709711793515157915895070>0.
\]
After unchanged kernel capacity is added, demand exceeds complete capacity
by
\[
 G_{83}=104570295123758938048546273477992878508577141441163332394>0.
\]
Thus the closed rank-nine prefix is $10\le K'\le83$.
\end{proposition}

\begin{proof}
The exact support-two/support-three partition is one ordinary lane and the
72 positive offsets $M_3-M_2=1,\ldots,72$.  Every offset retains every exact
defect pair before the support-four/support-five carrier alternatives are
expanded.  Positive ordinary pairs use the exhaustive $T_{23}/A_{23}$
positions, and a lone positive stratum uses its single completion carrier.
The latter carrier is supplied by
Proposition~\ref{prop:kb-rank11-single-completion-carrier}.

Apply Proposition~\ref{prop:kb-rank11-fixed-union-adjacent-support-coupling}
to every fixed union.  The former ordinary wall
$(M_2,M_3,M_4,M_5)=(18,18,36,36)$ is paid in both
$T_{23}=(39,7)$ and $A_{23}=(38,8)$ by the support-disjoint pairs $4/5$ and
$6/7$.  A primary router and an independent implementation agree on all 73
lane maxima and coverage counts.  Across 146 implementation/lane jobs, the
global maximum is the raw-safe offset-two branch
\[
 s_2=56,\quad s_3=54,\quad s_4=41,\quad s_5=33,
\]
with premium $P_{83}$.  The supplemental manifest pins all source and raw
capture hashes.  Exact substitution in the component floor identity gives
$M_{83}$ and $G_{83}$.
\end{proof}

\begin{proposition}[$K'=84$ adjacent-support carrier payment]
\label{prop:kb-rank11-k84-adjacent-support-payment}
In the fixed component target, every rank-nine alternative is empty at
$K'=84$.  The exact global completion premium is
\[
 P_{84}=41388798786059119503097492734939028640066114130,
\]
below the safe premium ceiling by
\[
 M_{84}=44581160171407926086602515730765812413619>0.
\]
After unchanged kernel capacity is added, demand exceeds complete capacity
by
\[
 G_{84}=222890179708699305421112332292209416493280153112001413582>0.
\]
Thus the closed rank-nine prefix is $10\le K'\le84$.
\end{proposition}

\begin{proof}
Use the same proved adjacent-support census as in
Proposition~\ref{prop:kb-rank11-k83-adjacent-support-payment}.  At $K'=84$
the exact support-two/support-three partition is one ordinary lane and the
73 offsets $M_3-M_2=1,\ldots,73$.  The primary and independent routers agree
on every lane maximum and coverage key, hence on all 148 implementation/lane
jobs.  Their global maximum is the raw-safe ordinary single-carrier branch
\[
 s_2=74,\quad s_3=55,\quad s_4=45,\quad s_5=37,
\]
with premium $P_{84}$.  The supplemental manifest pins both router sources,
both lane-wave captures, the compact merger, and the component-payment
output. Exact substitution in the
component floor identity gives $M_{84}$ and $G_{84}$.
\end{proof}

\begin{proposition}[$K'=85$ best-single adjacent-support payment]
\label{prop:kb-rank11-k85-best-single-adjacent-payment}
In the fixed component target, every rank-nine alternative is empty at
$K'=85$.  The exact global completion premium is
\[
 P_{85}=41412868016209776721228891386909879523306833354,
\]
below the safe premium ceiling by
\[
 M_{85}=1793645398692419426975603430807602228515>0.
\]
After unchanged kernel capacity is added, demand exceeds complete capacity
by
\[
 G_{85}=8967598503742781003071510733325918643075973211834024001>0.
\]
Thus the closed rank-nine prefix is $10\le K'\le85$.
\end{proposition}

\begin{proof}
The exact support-two/support-three partition is one ordinary lane and the
74 positive offsets $M_3-M_2=1,\ldots,74$.  Paired ordinary traversals agree
on all $492960$ source units and give premium strictly below $P_{85}$.

For the offsets, primary and independently written raw traversals agree on
all $16028400$ source units and $112198800$ raw rows per implementation.
Exactly $15696867$ units are raw-safe.  The remaining $331533$ units occur
in offsets $1$ through $41$; offsets $42$ through $74$ are entirely
raw-safe.  The largest raw-safe value is $P_{85}$, at the offset-eleven
branch
\[
 s_2=56,\quad s_3=45,\quad s_4=58,\quad s_5=37.
\]

Expand each residual unit through the exhaustive full-completion pairwise
carrier atlas and the proved fixed-union charges.  Every adjacent-support
cap supplied by
Proposition~\ref{prop:kb-rank11-fixed-union-adjacent-support-coupling}
is an individually valid upper bound for the same profile.  Hence the
minimum of the uncharged bound and the single-edge bounds is valid; no
simultaneous composition of overlapping edges is used.  Paired exhaustive
traversals agree on all 41 residual offsets, all $331533$ unsafe units, and
$49090656$ deduplicated profiles per implementation.  Every best-single
value is at most $P_{85}$.

The supplemental manifest pins the three paired waves, all mathematical
sources, and the component-payment output.  Exact substitution in the
component floor identity gives $M_{85}$ and $G_{85}$.
\end{proof}

\begin{proposition}[Canonical kernel-basis capacity]
\label{prop:kb-rank11-kernel-basis-capacity}
In the setup of Proposition~\ref{prop:kb-rank11-component-incidence}, the
rank-deficient component lane cannot carry its dominant-lane incidence floor
for any residual dimension
\[
 10\le K'\le4598.
\]
More precisely, for $1\le d\le9$, put $r=10-d$ and let $M_d(K')$
be the support-local record cap after cancelling $r$ coordinates; for
$d=9$, take $M_9=61871313426630599$.  The complete kernel-lane
capacity is at most
\begin{equation}
 \operatorname{Cap}(K')=
 \sum_{d=1}^{9}\binom{n'}{10-d}M_d(K')
                    \binom{K'-10}{d+1}.
 \tag{CK}\label{eq:kb-rank11-kernel-basis-capacity}
\end{equation}
For every $10\le K'\le4598$, this is strictly below
\[
 \frac{495405467}{10^9}\,274980728111260126\binom{m'}{11}.
\]
At $K'=4598$, the integer demand exceeds the capacity by
\[
219272330501201744129177266158988707697316238048878827197685.
\]
At $K'=4599$ the comparison reverses, so no larger interval is claimed.
\end{proposition}

\begin{proof}
Fix a rank-deficient eleven-subset $T$, let
$r=\operatorname{rank}(\operatorname{ev}_T)$, and choose one canonical
rank basis $B\subset T$, $|B|=r$.  Rank zero is absent because the
ten-dimensional residual correction space has empty global common zero
set.  Put $H=\ker(\operatorname{ev}_B)$, so $\dim H=d=10-r$.

Solving the coordinate equations on $B$ determines the quotient
component modulo $H$.  If two selected slopes use the same $B$, every
remaining compatibility condition is affine-linear in the slope and has
two roots, hence is an identity.  Thus, after one common affine codeword-pair
translation, all records assigned to $B$ have explanations in $H$.
Every word of $H$ and both translated received columns vanish on $B$,
so exact locator cancellation preserves slopes, supports, and pair
noncontainment.  Support-local transversality bounds the resulting family by
$M_d(K')$; at $d=9$, use the uniform post-near
margin/interleaving cap printed above.

Since $T$ and $B$ have the same evaluation rank, the $d+1$
coordinates of $T\setminus B$ are common zeros of $H$.  A
$d$-dimensional subspace of degree-$<K'$ polynomials has at most
$K'-d$ common zeros.  The $r=10-d$ basis coordinates are already among
them, leaving at most $K'-10$ choices.  Hence one canonical basis carries
at most
\[
 M_d(K')\binom{K'-10}{d+1}
\]
incidences.  There are at most $\binom{n'}r$ bases of rank $r$, and
canonical assignment counts each tuple once, proving
\eqref{eq:kb-rank11-kernel-basis-capacity}.

For $d\le8$, the exact cap used in the replay is
\[
M_d(K')=\left\lfloor\max\left\{
 \frac{(1048576+J)_{\underline{d+1}}}
      {(67472+J)(67473)_{\overline{d-1}}},
 \frac{(1048576+d)_{\underline{d+1}}}
      {(67473)_{\overline d}}
\right\}\right\rfloor,
\qquad J=K'-(10-d).
\]
Exact integer evaluation of this finite sum on all $4589$ rows gives the
displayed strict interval and endpoint gap.  At $K'=4599$, capacity
exceeds the integer demand by
$95457494746881463288875361950515757435711627164872173764503$,
which records the honest method wall.
\end{proof}

\begin{corollary}[All-bases kernel-capacity compression]
\label{cor:kb-rank11-kernel-multibasis-capacity}
In the setup of Proposition~\ref{prop:kb-rank11-component-incidence}, the
rank-deficient component lane cannot carry its dominant-lane incidence floor
for any residual dimension
\[
 10\le K'\le11641.
\]
More precisely, its capacity is at most
\[
 \operatorname{Cap}_{\rm all}(K')=
 \sum_{d=1}^{9}\left\lfloor
 \frac{\binom{n'}{10-d}M_d(K')\binom{K'-10}{d+1}}{d+2}
 \right\rfloor .
\]
At $K'=11641$, the integer demand exceeds this capacity by
\[
17769453550459149385453824948016076737082337523706893862084.
\]
At $K'=11642$, the capacity exceeds demand by
\[
187031323586740190878769118921060658362307444191332937452616,
\]
so no larger interval is claimed.
\end{corollary}

\begin{proof}
The evaluation matroid on a rank-$(10-d)$ eleven-set $T$ is loopless:
a loop would be a coordinate at which the complete residual correction
space vanishes, contrary to its empty global common zero set. Fix a basis
$B$ of $T$. For every $e\in T\setminus B$, the fundamental circuit
$C(e,B)$ contains some $b_e\in B$, and
$(B\setminus\{b_e\})\cup\{e\}$ is a basis. These $d+1$ bases are pairwise
distinct and differ from $B$. Hence $T$ has at least $d+2$ rank bases.

Decorate each $(\gamma,T)$ incidence by every one of its bases. For a
fixed basis, the common quotient, cancellation, support-local record cap,
and $\binom{K'-10}{d+1}$ extension bound in the proof of
Proposition~\ref{prop:kb-rank11-kernel-basis-capacity} apply unchanged;
none used the canonical tie-break. Thus the numerator of each summand in
the displayed capacity bounds the decorated
incidences, while every undecorated incidence is represented at least
$d+2$ times. Division and the integer floor prove the capacity formula.

Exact integer replay of all $11632$ rows gives the displayed interval and
endpoint gap. The next-row calculation gives the stated reversal. The
factor $d+2$ is sharp from looplessness alone: take $9-d$ coloops and one
parallel class of size $d+2$.
\end{proof}

\begin{corollary}[Record-support/ambient kernel-capacity hybrid]
\label{cor:kb-rank11-kernel-hybrid-capacity}
For one residual record with exact support $S$, $|S|=m'$, the number of
rank-$(10-d)$ eleven-subsets of $S$ is at most
\[
 P_d(K')=\left\lfloor
 \frac{\binom{m'}{10-d}\binom{K'-10}{d+1}}{d+2}
 \right\rfloor .
\]
Combining this with the ambient all-bases capacity $A_d(K')$ from
Corollary~\ref{cor:kb-rank11-kernel-multibasis-capacity}, the kernel lane
cannot carry its dominant-lane incidence floor for
\[
 10\le K'\le11772.
\]
At $K'=11772$, the integer demand exceeds the hybrid capacity by
\[
76504076505592948633027913576880724493595282142849410185084.
\]
At $K'=11773$, the hybrid capacity exceeds demand by
\[
139343682529231472322825521514042608524569163680782450618944,
\]
so no larger interval is claimed.
\end{corollary}

\begin{proof}
For one record, decorate each rank-$(10-d)$ eleven-subset of $S$ by all
its bases. A fixed basis $B\subset S$ has at most
$\binom{K'-10}{d+1}$ same-rank extensions, and every eleven-set has at
least $d+2$ bases. There are at most $\binom{m'}{10-d}$ candidate bases,
which proves the displayed per-record bound without using $M_d(K')$.

Let $R_{\rm act}\ge N_{\min}=274980728111260126$ be the actual residual
record count and let $I_d$ be the rank-$(10-d)$ incidence count. The two
bounds give
\[
 I_d\le\min\{A_d(K'),R_{\rm act}P_d(K')\}.
\]
After division by $R_{\rm act}$, each term
$\min\{A_d/R_{\rm act},P_d\}$ is nonincreasing in $R_{\rm act}$. Hence it
suffices to compare
\[
 \sum_{d=1}^9\min\{A_d(K'),N_{\min}P_d(K')\}
\]
with the integer demand at $N_{\min}$. Exact replay of all $11763$ rows
gives the stated interval and endpoint gap. At both boundary rows the
ambient branch is selected for $d=1,2$ and the record-support branch for
$d=3,\ldots,9$. The next-row calculation gives the displayed reversal.
\end{proof}

\begin{proposition}[Shared kernel nine-shadow resource]
\label{prop:kb-rank11-kernel-nine-shadow}
For one residual record with exact support $S$, let $I_d(S)$ count the
rank-$(10-d)$ eleven-subsets of $S$. Then
\[
 \sum_{d=1}^{9}
 \frac{\binom{d+2}{2}}{\binom{K'-d-9}{2}}I_d(S)
 \le \binom{m'}9,
\]
where a zero extension coefficient forces the corresponding incidence
count to vanish.
\end{proposition}

\begin{proof}
Fix a counted corank-$d$ eleven-set $T$ and put $r=10-d$. Its evaluation
matroid $M$ is loopless. The dual $M^*$ has rank $c=d+1$ and no coloops,
and a nine-subset spans $M$ exactly when its complementary pair is
independent in $M^*$.

Ignore the loops of $M^*$ and partition its remaining elements into
parallel classes. If there are at least $c+1$ classes, representatives of
any $c+1$ classes give at least $\binom{c+1}{2}$ independent pairs. If
there are exactly $c$ classes, each has size at least two, since a singleton
class would be a coloop. The cross-class pairs then number at least
$4\binom c2\ge\binom{c+1}{2}$. Thus $T$ has at least
$\binom{d+2}{2}$ spanning nine-subsets.

Fix such a rank-$r$ nine-subset $U\subset S$. Its kernel has dimension
$d$, so generalized MDS bounds its common-zero set, and hence the closure
of $U$, by $K'-d$ coordinates. Apart from the nine coordinates of $U$,
there are at most $K'-d-9$ choices. Thus $U$ has at most
$\binom{K'-d-9}{2}$ same-rank eleven-set extensions. Double counting
spanning pairs $(T,U)$ in each rank and summing over the rank partition of
the $\binom{m'}9$ nine-subsets proves the proposition.
\end{proof}

\begin{corollary}[One-shadow kernel-capacity cut]
\label{cor:kb-rank11-kernel-nine-shadow-capacity}
Intersecting the inequality in
Proposition~\ref{prop:kb-rank11-kernel-nine-shadow} with the individual
ambient/record caps excludes the dominant kernel lane for
\[
 10\le K'\le15445.
\]
At $K'=15445$, the integer demand exceeds capacity by
\[
178044655461817065880792270525721984196903835342334290540589.
\]
At $K'=15446$, capacity exceeds demand by
\[
124087038578417364551353992932097013573495323735890481286577,
\]
so no larger interval is claimed by this one-resource method.
\end{corollary}

\begin{proof}
After division by the actual record count, corank $d$ has upper bound
$u_d=\min\{A_d/N_{\min},P_d\}$. Its weight in that inequality is
$w_d=\binom{d+2}{2}/\binom{K'-d-9}{2}$, which increases with $d$.
The exact capacity is therefore the fractional-knapsack optimum obtained
by filling coranks in order. Exact rational replay checks all $15436$
rows. At both boundary rows corank one is full, corank two is partial, and
all higher coranks vanish. The displayed next-row reversal is exact, not a
rounding artifact.
\end{proof}

\begin{proposition}[Full nine-shadow containment resource]
\label{prop:kb-rank11-kernel-nine-shadow-containment}
With $I_d(S)$ as above, put
\[
 E_0=\binom{m'-9}{2},\qquad E_1=\binom{K'-10}{2}.
\]
Then
\[
 \left(52+3\frac{E_0}{E_1}\right)I_1(S)
 +55\sum_{d=2}^{9}I_d(S)
 \le E_0\binom{m'}9.
\]
When $E_1=0$, the corank-one term is absent because $I_1(S)=0$.
\end{proposition}

\begin{proof}
Let $J_1$ count rank-nine nine-subsets of $S$. Every kernel eleven-set
contains exactly $\binom{11}{9}=55$ nine-subsets. A rank-nine nine-subset
can only extend inside its rank-nine closure, giving at most $E_1$
extensions; every lower-rank nine-subset has at most the unrestricted
support-pair count $E_0$. Hence, with $I=\sum_d I_d$,
\[
 55I\le E_1J_1+E_0\left(\binom{m'}9-J_1\right).
\]
The corank-one case of Proposition~\ref{prop:kb-rank11-kernel-nine-shadow}
gives $3I_1\le E_1J_1$. Since $E_0>E_1$, substituting this lower bound
for $J_1$ and rearranging gives the proposition.
\end{proof}

\begin{corollary}[Two-resource nine-shadow kernel-capacity cut]
\label{cor:kb-rank11-kernel-nine-shadow-containment-capacity}
The individual ambient/record caps together with
Propositions~\ref{prop:kb-rank11-kernel-nine-shadow} and
\ref{prop:kb-rank11-kernel-nine-shadow-containment} exclude the dominant
kernel lane for
\[
 10\le K'\le15670.
\]
At $K'=15670$, the integer demand exceeds capacity by
\[
60244744187647715538325354175068999745872308513185869854532.
\]
At $K'=15671$, capacity exceeds demand by
\[
291105561463347587484268984669020036510369238771859813045635,
\]
so no larger interval is claimed.
\end{corollary}

\begin{proof}
The primary exact optimizer enumerates every cap breakpoint and every
intersection of the rank-preserving and full-containment resource bounds.
It checks all $15661$ rows. At both boundary rows only coranks one and two
are positive, both resource inequalities bind, and both individual caps
are slack.

For an independent certificate, write $w_1,w_2$ for the first two weights
in Proposition~\ref{prop:kb-rank11-kernel-nine-shadow} and
$v_1=52+3E_0/E_1$. Put
\[
 D=v_1w_2-55w_1,\qquad
 \lambda=\frac{v_1-55}{D},\qquad
 \mu=\frac{w_2-w_1}{D}.
\]
Exact replay verifies $\lambda,\mu\ge0$,
$\lambda w_1+\mu v_1=\lambda w_2+55\mu=1$, and
$\lambda w_d+55\mu\ge1$ for $d\ge3$. The resulting dual bound is exact at
both boundary rows and gives the same signs and capacities as the primal
replay.
\end{proof}

\begin{proposition}[Rank-eight nine-shadow extension deficit]
\label{prop:kb-rank11-kernel-rank8-nine-shadow-deficit}
Let $U$ be a rank-eight nine-subset of one exact support $S$.  At least
\[
 L_2=\binom{67474}{2}=2276336601
\]
unordered pairs in $S\setminus U$ raise its rank to ten.  Consequently,
with
\[
 E_2=\binom{K'-11}{2},
\]
the full-containment resource of
Proposition~\ref{prop:kb-rank11-kernel-nine-shadow-containment} sharpens to
\[
 \left(52+3\frac{E_0}{E_1}\right)I_1(S)
 +\left(55+6\frac{L_2}{E_2}\right)I_2(S)
 +55\sum_{d=3}^{9}I_d(S)
 \le E_0\binom{m'}9.
 \tag{R8S}
\]
A zero extension coefficient again forces the corresponding incidence
count to vanish.
\end{proposition}

\begin{proof}
Put $C=\operatorname{cl}_S(U)$, $c=|C|$, and $X=S\setminus C$.
The generalized-MDS closure bound used in
Proposition~\ref{prop:kb-rank11-kernel-nine-shadow} gives $c\le K'-2$,
so $q=|X|\ge67474$.  Contracting by $U$ leaves a rank-two matroid on $X$.
Every parallel class $P\subset X$ has
\[
 |P|\le K'-1-c,
\]
because $C\cup P$ is contained in a rank-nine closure, whose generalized-MDS
cap is $K'-1$.  Hence every $x\in X$ has at least
\[
 q-(K'-1-c)=m'-K'+1=67473
\]
partners outside its parallel class.  Dividing the ordered count by two
gives at least $q\cdot67473/2\ge\binom{67474}{2}$ rank-raising pairs.

Let $J_2$ count rank-eight nine-subsets of $S$.  Such a subset has at most
$E_0-L_2$ kernel extensions.  Refining the first inequality in the proof of
Proposition~\ref{prop:kb-rank11-kernel-nine-shadow-containment} therefore
gives
\[
 55\sum_d I_d
 \le E_0\binom{m'}9-(E_0-E_1)J_1-L_2J_2.
\]
The shared-shadow resource gives $3I_1\le E_1J_1$ and
$6I_2\le E_2J_2$.  Substitution and rearrangement give (R8S).
\end{proof}

\begin{corollary}[Rank-eight-deficit nine-shadow capacity cut]
\label{cor:kb-rank11-kernel-rank8-nine-shadow-capacity}
The individual ambient/record caps, the shared rank-preserving resource,
and (R8S) exclude the dominant kernel lane for
\[
 10\le K'\le17608.
\]
At $K'=17608$, integer demand exceeds floored capacity by
\[
126547040539829546354916747965612889135249249684319416999204.
\]
At $K'=17609$, capacity exceeds demand by
\[
165662859003771823867021831078593815988062146919602894849014,
\]
so no larger interval is claimed by this method.
\end{corollary}

\begin{proof}
The primary replay enumerates all nonnegative vertices of the two-resource
dual and checks all $17599$ closed rows exactly.  At the endpoint and wall,
coranks one and three attain their individual caps, coranks two and four are
resource-tight, coranks five through nine vanish, and both shared resources
bind.  An independently written replay reconstructs the primal optimizer
from the printed fifteen-interval active-set ledger and verifies feasibility,
complementary slackness, and strong duality on every row.  It obtains the
same rational optima and the displayed opposite boundary signs.
\end{proof}

\begin{proposition}[Two-step nine-shadow hierarchy]
\label{prop:kb-rank11-kernel-two-step-nine-shadow}
For one exact support $S$, let $I_d(S)$ count rank-$(10-d)$
eleven-subsets.  For every $3\le d\le9$, put
\[
 s_d=\binom{d+2}{2},\qquad
 L_d=\binom{67472+d}{2},\qquad
 E_d=\binom{K'-d-9}{2},\qquad
 Q_d=\binom{11-d}{2}.
\]
Then
\[
 \frac{s_dL_d}{E_d}I_d(S)\le Q_d I_{d-2}(S).
 \tag{TS}
\]
A zero source-extension coefficient forces the corresponding source
incidence count to vanish.
\end{proposition}

\begin{proof}
Let $J_d(S)$ count rank-$(10-d)$ nine-subsets.  The shared-shadow theorem,
Proposition~\ref{prop:kb-rank11-kernel-nine-shadow}, gives
\[
 s_d I_d(S)\le E_dJ_d(S).
 \tag{1}
\]
Fix a nine-subset $U$ counted by $J_d(S)$ and put
$C=\operatorname{cl}_S(U)$, $c=|C|$, and $X=S\setminus C$.
The generalized-MDS closure cap gives $c\le K'-d$, hence
$q=|X|\ge67472+d$.  After contracting $U$, every parallel class in $X$
has size at most $K'-d+1-c$: adjoining such a class to $C$ gives a
rank-$(11-d)$ closure.  Every point of $X$ therefore has at least
$67471+d$ partners outside its parallel class.  At least
\[
 \frac{q(67471+d)}2\ge\binom{67472+d}{2}=L_d
\]
unordered pairs raise the rank of $U$ by two.

Now fix a resulting rank-$(12-d)$ eleven-subset $T$ and write
$D=T\setminus U$.  The dual-rank identity
\[
 r_T(T\setminus D)=r_T(T)-|D|+r_{T^*}(D)
\]
shows that $U$ has rank $10-d$ exactly when the complementary pair $D$
consists of two coloops of the evaluation matroid on $T$.  This matroid is
loopless, has rank $12-d$, and has eleven elements, so it has at most
$11-d$ coloops: $12-d$ coloops would make every remaining element a loop.
Thus $T$ contains at most $Q_d$ source shadows, and double counting gives
\[
 L_dJ_d(S)\le Q_dI_{d-2}(S).
 \tag{2}
\]
Combining (1) and (2) proves (TS).
\end{proof}

\begin{corollary}[Two-step nine-shadow kernel-capacity cut]
\label{cor:kb-rank11-kernel-two-step-nine-shadow-capacity}
The individual ambient/record caps, the full-containment resource, and the
seven inequalities \textup{(TS)} exclude the dominant kernel lane for
\[
 10\le K'\le18101.
\]
At $K'=18101$, integer demand exceeds floored capacity by
\[
33462159928103132226516704640419847248244116666500998762314.
\]
At $K'=18102$, capacity exceeds demand by
\[
275016496133605602641019628236447268989861205055439981187167,
\]
so no larger interval is claimed by this method.
\end{corollary}

\begin{proof}
On every newly checked row $17609\le K'\le18102$, the exact optimizer has
one active-set pattern: the corank-one cap, full containment, and all seven
two-step inequalities bind; every corank is positive; the rank-preserving
nine-shadow resource and all other individual caps are slack.  The primary
replay constructs the odd and even primal chains and solves the nine exact
dual equations by rational Gaussian elimination.  An independently written
replay instead derives the hierarchy multipliers by backward odd/even
recurrences.  Both verify feasibility and strong duality on all $494$ rows
and obtain the displayed opposite boundary signs.
\end{proof}

\begin{proposition}[Multistep kernel-shadow hierarchy]
\label{prop:kb-rank11-kernel-multistep-shadow}
For one exact support $S$, let $I_d(S)$ count rank-$(10-d)$
eleven-subsets.  For every $2\le t<d\le9$, put
\[
 s_{d,t}=\binom{d+2}{t},\quad
 L_{d,t}=\binom{67472+d}{t},\quad
 E_{d,t}=\binom{K'-d-11+t}{t},\quad
 Q_{d,t}=\binom{9-d+t}{t}.
\]
Then
\[
 \frac{s_{d,t}L_{d,t}}{E_{d,t}}I_d(S)
 \le Q_{d,t}I_{d-t}(S).
 \tag{MS}
\]
There are 28 such inequalities; the seven rows with $t=2$ are
Proposition~\ref{prop:kb-rank11-kernel-two-step-nine-shadow}.
\end{proposition}

\begin{proof}
Let $N$ be the dual of the evaluation matroid on a source eleven-set.
It is coloopless of rank $d+1$, because the evaluation matroid is loopless.
Write $f_j$ for the number of independent $j$-sets of $N$.  Every
independent $j$-set has at least $d+2-j$ one-point extensions.  Otherwise
at most $d+1-j$ elements lie outside its closure.  Rank forces equality,
and deleting any one of those elements leaves rank at most $d$, making
each such element a coloop, a contradiction.  Double counting extensions
therefore gives
\[
 (j+1)f_{j+1}\ge(d+2-j)f_j.
\]
Iteration from $f_0=1$ yields $f_t\ge\binom{d+2}{t}$.  Taking complements
gives at least $s_{d,t}$ spanning $(11-t)$-shadows per source.  A fixed
shadow has at most $E_{d,t}$ same-rank source extensions by the
generalized-MDS closure cap.

Starting from such a shadow, after $j$ independent support points have been
adjoined its closure has size at most $K'-d+j$.  Thus at least
$67472+d-j$ support points raise rank at the next step.  Dividing the
resulting ordered count by $t!$ gives at least $L_{d,t}$ rank-raising
$t$-sets.

Conversely, fix a resulting rank-$(10-d+t)$ target eleven-set.  By the
dual-rank identity, a complementary $t$-set produces a rank-$(10-d)$
shadow exactly when all its elements are coloops of the target matroid.
A loopless rank-$(10-d+t)$ matroid on eleven elements has at most
$9-d+t$ coloops, or all remaining elements would be loops.  Hence the
target contains at most $Q_{d,t}$ source shadows.  The two double counts
give \textup{(MS)}.
\end{proof}

\begin{corollary}[All-step kernel-shadow capacity cut]
\label{cor:kb-rank11-kernel-multistep-shadow-capacity}
The individual ambient/record caps, full containment, and the 28
inequalities \textup{(MS)} exclude the dominant kernel lane for
\[
 10\le K'\le18158.
\]
At $K'=18158$, integer demand exceeds floored capacity by
\[
289110608820324799941118306538399899258195112067661304310498.
\]
At $K'=18159$, capacity exceeds demand by
\[
20286290696334777989469267474876769475675508046109372076445,
\]
so no larger interval is claimed by this method.
\end{corollary}

\begin{proof}
On all 58 replayed rows $18102\le K'\le18159$, the exact certificate uses
the hierarchy tree
\[
(t,d)=(2,3),(2,4),(2,6),(2,8),(3,5),(2,7),(2,9).
\]
The corank-one cap and full-containment resource bind, all nine coranks are
positive, and every other cap and the rank-preserving shadow resource are
slack.  Seventeen hierarchy rows bind; the ten tight rows outside the tree
are checked as exact cycle identities.  The primary replay solves the nine
dual equations by rational Gaussian elimination.  An independently written
replay reconstructs both primal factors and dual multipliers by forward and
backward tree recurrences.  Both check all 28 inequalities, feasibility,
and strong duality on every row and obtain the displayed opposite boundary
signs.
\end{proof}

\begin{proposition}[Corank-one projective-normal pair cap]
\label{prop:kb-rank11-kernel-corank1-projective-pair}
After canonical rank-nine basis cancellation, the complete shortened
corank-one chart has
\[
 (n,K,m,s)=(1048577,1,67473,1).
\]
Every selected record owns at least $2(m-1)=134944$ ordered coordinate
pairs whose agreement normals are linearly independent.  Consequently the
number of records in the chart is at most
\[
 M_1\le\left\lfloor\frac{n(n-1)}{2(m-1)}\right\rfloor=8147918.
 \tag{PP}
\]
This improves the support-local transversality cap $16295594$ by $8147676$.
\end{proposition}

\begin{proof}
Represent the affine explanation family by parameter points in $\F^2$.
Agreement at a coordinate is an affine line with a normal in $\F^2$.
The common-zero bound gives $z\le K-s=0$, so all $m$ normals incident with
one selected point are nonzero.  Same-support pair noncontainment forces
them to span $\F^2$.

Let their nonempty projective-class sizes be $a_1,\ldots,a_c$.  Then
$c\ge2$, $\sum_i a_i=m$, and the number of ordered independent pairs is
\[
 m^2-\sum_i a_i^2.
\]
Convexity gives $\sum_i a_i^2\le(m-1)^2+1$, with equality for class sizes
$m-1,1$.  Thus every selected point owns at least $2(m-1)$ ordered
independent coordinate pairs.  Two independent affine lines meet in at
most one parameter point, so a coordinate pair is owned by at most one
record.  Double counting the $n(n-1)$ ordered coordinate pairs proves
\textup{(PP)}; the integer division has remainder $29760$.
\end{proof}

\begin{corollary}[Projective-pair kernel-capacity cut]
\label{cor:kb-rank11-kernel-projective-pair-capacity}
The projective cap \textup{(PP)}, the other individual ambient/record caps,
and the 28 inequalities \textup{(MS)} exclude the dominant kernel lane for
\[
 10\le K'\le377673.
\]
At $K'=377673$, integer demand exceeds floored capacity by
\[
608290099077401798561583762592584078050381528604243813748500153228.
\]
At $K'=377674$, capacity exceeds demand by
\[
1089804128361045148874283346879615159892995682385275039289561845323,
\]
so no larger interval is claimed by this resource system.
\end{corollary}

\begin{proof}
The predecessor handles $K'\le18158$.  On the remaining exact replay the
corank-one and corank-two individual caps bind, both shared nine-shadow
resources are strict, and every corank is positive.  The seven hierarchy
tree edges are
\[
 (t,d)=(2,3),(3,4),(2,5),(2,6),(2,7),(2,8),(2,9).
\]
They determine all nonroot counts.  Twenty-two of the 28 hierarchy rows
are equalities.  Backward tree recurrence gives positive edge prices and
positive prices for both root caps, with dual objective equal to the primal
sum.

The source-pinned exact replay checks all $359516$ rows
$18159\le K'\le377674$ in 64 bounded workers.  The packet verifiers
independently reconstruct the projective partition bound and the replay
start, endpoint, and adjacent wall from exact tree-path products.  These
give the displayed opposite integer signs.
\end{proof}

\begin{proposition}[Corank-two projective-normal basis cap]
\label{prop:kb-rank11-kernel-corank2-projective-basis}
After canonical rank-eight basis cancellation, the complete shortened
corank-two chart has
\[
 (n,K,m,s)=(1048578,2,67474,2).
\]
Every selected record owns at least
$3(m-1)(m-2)=13657614768$ ordered coordinate triples whose agreement
normals are linearly independent.  Consequently the number of records in
the chart is at most
\[
 M_2\le\left\lfloor
 \frac{n(n-1)(n-2)}{3(m-1)(m-2)}
 \right\rfloor=84416263.
 \tag{PB}\label{eq:kb-rank11-kernel-corank2-projective-basis}
\]
This improves the support-local transversality cap $253241283$ by
$168825020$.
\end{proposition}

\begin{proof}
Represent the affine explanation family by parameter points in $\F^3$.
The common-zero bound gives $z\le K-s=0$.  For any one-dimensional normal
span, the support-local MDS step leaves at least
\[
 w+s-1=67472+2-1=67473=m-1
\]
incident normals outside it.  Thus no two of the $m$ incident nonzero
normals are projectively equal.  Same-support pair noncontainment gives full
incident rank, so the resulting $m$ points in $\operatorname{PG}(2)$ are
not all collinear.

Choose a projective line containing $q\ge2$ of the points and put
$r=m-q\ge1$.  A collinear triple not contained in that line has zero or one
point on it.  There are at most $\binom r3$ triples of the first type, while
each pair off the line determines at most one triple of the second type.
Hence the number of collinear triples is at most
\[
 \binom q3+\binom r3+\binom r2
 =\binom q3+\binom{r+1}3
 \le\binom{m-1}3.
\]
The last inequality follows from
\[
 \binom{m-1}3-\binom q3-\binom{r+1}3
 =(r-1)\left(\binom q2-1\right)
  +(q-2)\binom{r-1}2\ge0.
\]
Therefore the number of ordered independent triples is at least
\[
 m(m-1)(m-2)-6\binom{m-1}3=3(m-1)(m-2).
\]
Three independent affine hyperplanes in $\F^3$ meet in at most one
parameter point.  Double counting the $n(n-1)(n-2)$ ordered coordinate
triples proves \textup{(PB)}; the integer division has remainder
$2935655472$.
\end{proof}

\begin{proposition}[Projective-paving caps and the shortening gap]
\label{prop:kb-rank11-kernel-projective-paving-scope}
In the completely shortened corank-$d$ chart, $1\le d\le9$, the incident
normal matroid has rank $d+1$ and every set of at most $d$ normals is
independent.  It therefore has at least $\binom{m-1}{d}$ unordered bases,
and projective-basis double counting gives the record cap
\[
 P_d=\left\lfloor
 \frac{(1048576+d)_{d+1}}
 {(d+1)(67472+d-1)_d}
 \right\rfloor.
\]
For $d=1,\ldots,9$ these caps are
\[
\begin{split}
&8147918,\ 84416263,\ 983902549,\ 12232092309,\\
&158406193634,\ 2109949210211,\ 28689347099870,\\
&396280526311830,\ 5542092977392141.
\end{split}
\tag{PC}
\]
They are complete-chart caps, not uniform shortened-chart caps.

More precisely, put $R=1048576$, $w=67472$, and, after rank-$(10-d)$
shortening, let
\[
 t=K'-10-z,\qquad 0\le t\le K'-10,
\]
where $z$ is the number of global zero normals.  For $d\ge2$ the
support-local endpoint envelope is
\[
 M_d(K')=\left\lfloor\max\{P_d,F_d(1),F_d(K'-10)\}\right\rfloor,
 \tag{IG}
\]
where, before flooring,
\[
 P_d=\frac{(R+d)_{d+1}}{(d+1)(w+d-1)_d},\qquad
 F_d(t)=\frac{(R+d+t)_{d+1}}
 {(w+d+t)(w+1)^{\overline{d-1}}}.
\]
At $K'=377674$ this gives $M_2=253238254$ and
$M_3=3935391907$.  Corank one is exceptional: its projective-pair argument
proves the uniform cap $M_1=8147918$.  This endpoint envelope is a valid
generic bound, not an obstruction to a stronger multiplicity-aware count;
the corank-two theorem below supplies such a count.
\end{proposition}

\begin{proof}
The matroid assertion follows by deletion--contraction.  If an element is
a coloop, adjoining it to every $(d-1)$-subset of the deletion already gives
$\binom{m-1}{d}$ bases.  Otherwise deletion and contraction preserve the
required independence hypotheses, and Pascal's identity closes the
induction.  Ordering the bases and double counting their affine-hyperplane
intersections proves \textup{(PC)}.

For $t\ge1$, the pointwise support-local transversality count before endpoint
maximization gives $F_d(t)$.  The sign of
$F_d(t+1)/F_d(t)-1$ is the sign of
\[
 dt+(d+1)(w+d)-R.
\]
This expression increases with $t$, so $F_d$ has one turn and its maximum
on the integer interval is at $1$ or $K'-10$.  The complete chart $t=0$
is controlled separately by $P_d$, proving \textup{(IG)}.  The same
one-turn argument for the rank-two projective-pair count is maximized at
$t=0$, proving the uniform corank-one assertion.  Exact integer evaluation
at $K'=377674$ gives the displayed audit values.
\end{proof}

\begin{proposition}[Rank-three bounded-parallel basis floor]
\label{prop:kb-rank11-rank3-bounded-parallel}
Let $M$ be a loopless rank-three matroid on $m\ge3$ elements, and suppose
every parallel class has size at most $a$.  If $b(M)$ is its number of
unordered bases, then
\[
 2b(M)\ge(m-1)(m-1-a).
 \tag{BP}\label{eq:kb-rank11-rank3-bounded-parallel}
\]
If $a\mid(m-1)$ and $m-1\ge2a$, equality is attained by one coloop over a
rank-two matroid whose parallel classes all have size $a$.
\end{proposition}

\begin{proof}
Induct on $m$.  If $e$ is a coloop, the deletion is a loopless rank-two
matroid on $q=m-1$ elements.  Writing its parallel-class sizes as $q_i$,
\[
 2b(M)=q^2-\sum_iq_i^2\ge q^2-aq=(m-1)(m-1-a).
\]

Otherwise choose $e$ in a smallest parallel class $P$, of size $c$.
Deletion preserves rank three, so induction gives
\[
 2b(M\setminus e)\ge(m-2)(m-2-a).
\]
In $M/e$, the other elements of $P$ are loops.  Every nonloop parallel
class is a union of original parallel classes and hence has size at least
$c$.  There are at least two such classes and their total size is $m-c$.
Merging classes decreases the number of independent pairs, so
\[
 b(M/e)\ge c(m-2c)\ge m-2.
\]
Indeed the simplification of $M$ has at least three points, hence $m\ge3c$,
and the residual at $m=3c$ is $(c-1)(c-2)\ge0$.  Deletion--contraction now
gives
\[
 2b(M)\ge(m-2)(m-2-a)+2(m-2)
 =(m-1)(m-1-a)+(a-1),
\]
which proves \textup{(BP)}.  Equal rank-two classes in the stated coloop
construction give equality.
\end{proof}

\begin{proposition}[Uniform corank-two projective-normal basis cap]
\label{prop:kb-rank11-kernel-corank2-uniform-projective-basis}
After canonical rank-eight basis cancellation and deletion of all global
zero normals, put $t=K'-10-z$.  Every official corank-two chart has
\[
 (n,K,m,s)=(1048578+t,2+t,67474+t,2),
 \qquad 0\le t\le1048566,
\]
and at most
\[
 M_2\le84416263
 \tag{UPB}\label{eq:kb-rank11-kernel-corank2-uniform-projective-basis}
\]
records.  The complete chart $t=0$ attains the largest value of the proved
envelope.
\end{proposition}

\begin{proof}
Write $R=1048576$ and $w=67472$.  Support-local transversality leaves at
least $w+1$ incident normals outside every one-dimensional normal span.
Thus the loopless rank-three incident normal matroid has maximum parallel
class size
\[
 a\le m-(w+1)=t+1.
\]
By Proposition~\ref{prop:kb-rank11-rank3-bounded-parallel}, every record
owns at least
\[
 6b(M)\ge3w(w+t+1)
\]
ordered independent coordinate triples.  Double counting gives
\[
 H(t)=\frac{(R+t)(R+t+1)(R+t+2)}{3w(w+t+1)}
\]
as a record envelope.  The sign of $H(t+1)/H(t)-1$ is the sign of
\[
 2t+3w+3-R.
\]
Hence $H$ has one turn and its maximum on the official integer interval is
at an endpoint.  Exact division gives
\[
 \lfloor H(0)\rfloor=84416263,
 \qquad
 \lfloor H(1048566)\rfloor=40828171,
\]
with positive next-integer gaps.  This proves \textup{(UPB)} uniformly.
\end{proof}

\begin{corollary}[Projective-basis kernel-capacity cut]
\label{cor:kb-rank11-kernel-projective-basis-capacity}
The projective caps for coranks one and two, the other individual
ambient/record caps, and the 28 inequalities \textup{(MS)} exclude the
dominant kernel lane for
\[
 10\le K'\le568338.
\]
At $K'=568338$, integer demand exceeds floored capacity by
\[
38432453444617070485037263551626410396462586389410416394578520596038.
\]
At $K'=568339$, capacity exceeds demand by
\[
36180877960369511460476382880286784896208001102094988739728829832800,
\]
so no larger interval is claimed by this resource system.
\end{corollary}

\begin{proof}
The predecessor handles $K'\le377673$.  On the remaining exact replay the
corank-one and corank-two individual caps bind, both shared nine-shadow
resources are strict, every nonroot individual cap is strict, and every
corank is positive.  The seven hierarchy tree edges are
\[
 (t,d)=(2,3),(2,4),(2,6),(2,8),(3,5),(2,7),(2,9).
\]
They give the two rooted components
\[
 1\longrightarrow3,qquad
 2\longrightarrow4\longrightarrow6\longrightarrow8,qquad
 2\longrightarrow5\longrightarrow7\longrightarrow9.
\]
Seventeen of the 28 hierarchy rows are equalities.  Backward tree recurrence
gives positive edge prices and positive prices for both root caps, with dual
objective equal to the primal sum.

The source-pinned exact replay checks all $190666$ rows
$377674\le K'\le568339$ in 64 bounded workers.  The independent verifier
enumerates every projective line split in the basis bound and reconstructs
the replay start, endpoint, and adjacent wall from exact tree-path products.
These give the displayed opposite integer signs.
\end{proof}

\begin{proposition}[Corank-three projective-frame basis cap]
\label{prop:kb-rank11-kernel-corank3-projective-frame}
After canonical rank-seven basis cancellation, the complete shortened
corank-three chart has
\[
 (n,K,m,s)=(1048579,3,67475,3).
\]
Every selected record owns at least
$4(m-1)(m-2)(m-3)=1228711865141376$ ordered coordinate quadruples
whose agreement normals are linearly independent. Consequently the number
of records in the chart is at most
\[
 M_3\le\left\lfloor
 \frac{n(n-1)(n-2)(n-3)}{4(m-1)(m-2)(m-3)}
 \right\rfloor=983902549.
 \tag{PF}\label{eq:kb-rank11-kernel-corank3-projective-frame}
\]
This improves the support-local transversality cap $3935435218$ by
$2951532669$.
\end{proposition}

\begin{proof}
Represent the affine explanation family by parameter points in $\F^4$.
The common-zero bound gives $z\le K-s=0$. For every one-dimensional normal
span, the support-local MDS step leaves at least
\[
 w+s-1=67472+3-1=67474=m-1
\]
incident normals outside it, so the incident projective points are distinct.
For every two-dimensional normal span, the same step leaves at least
\[
 w+s-2=67473=m-2
\]
normals outside it. Thus no three incident points are collinear.
Same-support pair noncontainment gives full incident rank, so these points
span $\operatorname{PG}(3)$.

Choose three incident points and let $H$ be their projective plane. Put
$q=|S\cap H|\ge3$ and $r=m-q\ge1$. Partition coplanar unordered
quadruples by the number of points in $H$. Four points in $H$ contribute
$\binom q4$, while three in $H$ and one outside contribute none. For a
fixed outside pair, its line meets $H$ at a point $x\notin S$, since
$x\in S$ would give three collinear selected points. Secants of $S\cap H$
through $x$ use disjoint point pairs, giving at most
$\lfloor q/2\rfloor\binom r2$ quadruples with two points in $H$. Every
outside triple spans a plane meeting $H$ in a line with at most two selected
points, giving at most $2\binom r3$ quadruples with one point in $H$.
All-outside quadruples contribute at most $\binom r4$. Hence the dependent
quadruple count is at most
\[
 B(q,r)=\binom q4+\left\lfloor\frac q2\right\rfloor\binom r2
        +2\binom r3+\binom r4.
\]
Vandermonde expansion gives
\[
 \binom{m-1}4-B(q,r)
 =(q-3)\binom r3
 +\left(\binom{q-1}2-\left\lfloor\frac q2\right\rfloor\right)\binom r2
 +(r-1)\binom{q-1}3\ge0.
\]
Therefore the number of ordered independent quadruples is at least
\[
 (m)_4-24\binom{m-1}4=4(m-1)(m-2)(m-3).
\]
Four independent affine hyperplanes in $\F^4$ meet in at most one
parameter point. Double counting the $(n)_4$ ordered coordinate quadruples
proves \textup{(PF)}; the integer division has remainder
$1056607358217600$.
\end{proof}

\begin{proposition}[Rank-four bounded-point/line basis floor]
\label{prop:kb-rank11-rank4-bounded-point-line}
Let $M$ be a loopless rank-four matroid on $m=a+r$ elements, where
$a\ge1$ and $r\ge3$. Suppose every parallel class has size at most $a$ and
every rank-two flat has size at most $a+1$. Put
\[
 h_a(r)=\min\left\{\left\lfloor\frac{a+1}{2}\right\rfloor,
                    \left\lfloor\frac{a+r}{4}\right\rfloor\right\},
\]
\[
 C_a(r)=(a+r-1)(r-1)(r-2),\qquad
 L_a(r)=3(a+r-h_a(r)-1)(r-2),
\]
and define $Q_a(3)=6$ and, for $r\ge4$,
\[
 Q_a(r)=\min\{C_a(r),Q_a(r-1)+L_a(r)\}.
 \tag{BPL}\label{eq:kb-rank11-rank4-bounded-point-line}
\]
Then the number $b(M)$ of unordered bases satisfies
\[
 6b(M)\ge Q_a(r).
\]
\end{proposition}

\begin{proof}
Induct on $r$. At $r=3$, rank four supplies one basis. If $M$ has a coloop
$e$, then $M\setminus e$ has rank three on $q=m-1$ elements. Sequentially
choosing an ordered independent triple gives at least $q$, $q-a=r-1$, and
$q-(a+1)=r-2$ choices, respectively. Hence $6b(M)\ge C_a(r)$.

Otherwise choose $e$ in a smallest parallel class $P$, of size $c$.
The simplification has at least four points, so $4c\le m$. A rank-two flat
through $P$ and any other parallel class contains at least $2c$ elements,
so $2c\le a+1$. Thus $c\le h_a(r)$. Deletion preserves rank four and the
two ceilings. After contracting $e$ and deleting the loops $P\setminus e$,
the resulting rank-three matroid has $a+r-c$ elements and parallel-class
ceiling $a+1-c$. Proposition~\ref{prop:kb-rank11-rank3-bounded-parallel}
therefore gives
\[
 6b(M/e)\ge3(a+r-c-1)(r-2)\ge L_a(r).
\]
Deletion--contraction gives the second branch of \textup{(BPL)}.

For exact evaluation, unfold the recurrence. Besides the no-reset path, a
last coloop reset at $j\ge4$ contributes
\[
 C_a(j)+\sum_{x=j+1}^{r}L_a(x).
\]
The difference between the candidates at $j+1$ and $j$ is
$(j-1)(3h_a(j+1)-a-2)$. Since $h_a$ is nondecreasing, the reset candidates
have one turn; the no-reset path is no larger than the $j=4$ candidate by
$6(h_a(4)-1)\ge0$. Hence only the no-reset path and the reset adjacent to
that sign change can minimize. Splitting
$h_a(x)=\min\{\lfloor(a+1)/2\rfloor,\lfloor(a+x)/4\rfloor\}$ at its
constant branch and by $a+x$ modulo four gives an exact integer evaluator.
\end{proof}

\begin{corollary}[Uniform corank-three projective-frame cap]
\label{cor:kb-rank11-kernel-corank3-uniform-projective-frame}
Every official corank-three shortened chart satisfies
\[
 M_3\le983902549.
 \tag{UPF}\label{eq:kb-rank11-kernel-corank3-uniform-projective-frame}
\]
\end{corollary}

\begin{proof}
After canonical rank-seven basis cancellation, delete all global zero
normals and put $t=K'-10-z$, $a=t+1$. The remaining chart has
\[
 (n,K,m,s)=(1048579+t,3+t,67475+t,3),\qquad r=m-a=67474.
\]
Support-local transversality leaves at least $67474$ incident normals
outside every one-dimensional span and at least $67473$ outside every
two-dimensional span. Thus the incident loopless rank-four matroid has
parallel-class ceiling $a$ and rank-two-flat ceiling $a+1$.
Proposition~\ref{prop:kb-rank11-rank4-bounded-point-line} gives at least
$4Q_{t+1}(67474)$ ordered independent coordinate quadruples per record, so
\[
 M_3(t)\le\left\lfloor
 \frac{(1048576+t)(1048577+t)(1048578+t)(1048579+t)}
      {4Q_{t+1}(67474)}
 \right\rfloor.
\]
The one-turn exact evaluator checks every $0\le t\le1048566$. Its maximum
is $983902549$ at $t=0$; the values at $t=1$ and $t=1048566$ are
$983891721$ and $951742008$. The primary verifier exhausts all $1048567$
rows with exact integers, while the independent verifier unfolds the
recurrence directly at five separated rows. A source-pinned 512 MB,
60-second-limited worker independently reports the same all-row maximum.
\end{proof}

\begin{corollary}[Projective-frame kernel-capacity cut]
\label{cor:kb-rank11-kernel-projective-frame-capacity}
The projective caps for coranks one, two, and three, the other individual
ambient/record caps, and the 28 inequalities \textup{(MS)} exclude the
dominant kernel lane for
\[
 10\le K'\le796598.
\]
At $K'=796598$, integer demand exceeds floored capacity by
\[
1063274038253455766288412818872693782800681544679740581002823089126086.
\]
At $K'=796599$, capacity exceeds demand by
\[
670721678337441589385303494237372283642375643589068751593971045368244,
\]
so no larger interval is claimed by this resource system.
\end{corollary}

\begin{proof}
The predecessor handles $K'\le568338$. On the remaining exact replay the
corank-one, corank-two, and corank-three individual caps bind, both shared
nine-shadow resources are strict, every nonroot individual cap is strict,
and every corank is positive. The six hierarchy forest edges are
\[
 (t,d)=(2,4),(2,5),(3,6),(4,7),(5,8),(6,9).
\]
They give the components
\[
 1,\qquad 2\longrightarrow4,\qquad
 3\longrightarrow5,6,7,8,9.
\]
Twelve of the 28 hierarchy rows are equalities. Backward forest recurrence
gives positive edge prices and positive prices for all three root caps,
with dual objective equal to the primal sum.

The source-pinned exact replay checks all $228261$ rows
$568339\le K'\le796599$ in 64 bounded workers. The independent verifier
enumerates every projective plane split in the frame bound and reconstructs
the replay start, endpoint, and adjacent wall from direct forest products.
These give the displayed opposite integer signs.
\end{proof}

\begin{proposition}[Shortening-weighted record/extension cap]
\label{prop:kb-rank11-kernel-shortening-weighted-extension}
Put $R=1048576$, $w=67472$, $S=K'-10$, and let $P_d$ be the printed
complete-chart cap in corank $d$. For one fixed basis of a decorated
corank-$d$ kernel incidence, define
\[
 F_d(t)=
 \frac{(R+d+t)^{\underline{d+1}}}
 {(w+d+t)(w+1)^{\overline{d-1}}}.
\]
On $796599\le K'\le1048576$, its record/extension contribution is at most
\[
 U_d(K')=
 \begin{cases}
 P_d\binom{S}{d+1},&1\le d\le3,\\[2mm]
 F_d(1)\binom{S-1}{d+1},&4\le d\le9.
 \end{cases}
 \tag{SWE}\label{eq:kb-rank11-kernel-shortening-weighted-extension}
\]
After summing over bases, every undecorated incidence is counted at least
$d+2$ times.
\end{proposition}

\begin{proof}
Let $z$ be the number of common zero normals outside the fixed basis and
put $t=S-z$. Every extension of the basis uses $d+1$ of those zero normals,
so one record has at most $\binom{S-t}{d+1}$ extensions. The same $z$ is
deleted in the support-local chart. Thus the complete chart contributes at
most $P_d\binom{S}{d+1}$, while a noncomplete chart contributes at most
$F_d(t)\binom{S-t}{d+1}$.

For $q=d+1$ and every nonzero term with $t\ge1$, the successive ratio of
the unfloored noncomplete bounds is
\[
 \frac{(R+d+t+1)(w+d+t)(S-t-q)}
 {(R+t)(w+d+t+1)(S-t)}
 <\left(1+\frac q{R+t}\right)
  \left(1-\frac q{S-t}\right)<1.
\]
The final inequality uses $R+t>S-t$. Hence the noncomplete maximum is at
$t=1$. For $d\le3$, the already proved uniform caps $P_d$ and the larger
complete extension count give the first branch of \textup{(SWE)}. For
$4\le d\le9$, exact cross multiplication at $K'=796599$ shows that the
$t=1$ branch exceeds the complete branch; their ratio increases with $S$.
The six comparisons are replayed with exact rationals. The all-bases
decoration theorem supplies the final divisor $d+2$.
\end{proof}

\begin{corollary}[Terminal shortening-weighted kernel cut]
\label{cor:kb-rank11-kernel-shortening-weighted-capacity}
The fixed-kernel branch is impossible on every official residual dimension
\[
 10\le K'\le1048576.
\]
\end{corollary}

\begin{proof}
The projective-frame capacity cut handles $K'\le796598$. On the remaining
interval, Proposition~\ref{prop:kb-rank11-kernel-shortening-weighted-extension}
gives the rational capacity
\[
 C(K')=\sum_{d=1}^{9}
 \frac{\binom{R+K'}{10-d}U_d(K')}{d+2}.
\]
Compare it with the exact unrounded demand
\[
 D(K')=\frac{495405467}{10^9}\,
 274980728111260126\binom{w+K'}{11},
\]
and put $G=D-C$. This is a polynomial of degree at most eleven. Exact
forward differencing at $K_0=796599$ gives
\[
 \Delta^jG(K_0)>0\qquad(0\le j\le11).
\]
For every integer $s\ge0$, Newton interpolation therefore gives
\[
 G(K_0+s)=\sum_{j=0}^{11}\binom{s}{j}\Delta^jG(K_0)>0.
\]
This covers the complete upper interval. Floors only decrease capacity.
The primary verifier reconstructs and hashes all twelve exact Newton
coefficients. The independent verifier instead expands $G(K_0+s)$ in the
ordinary power basis; all twelve coefficients are positive and have
SHA-256 digest
\texttt{85f67eade03100e7fc6f3ef0443a3a787f1133427a0f9d4030d638eca9080082}.
\end{proof}

\begin{proposition}[Rank-eight minimal-shortening exclusion]
\label{prop:kb-rank11-rank8-minimal-shortening-exclusion}
At residual dimension $K'=10$, the rank-eight affine-owner alternative of
Proposition~\ref{prop:kb-rank11-component-ninesubset-target} is empty.
Every nine-coordinate restriction of the correction space has rank nine.
This conclusion is sharp for the dimension-equality argument: the adjacent
row $K'=11$ admits the rank-eight chart in
Proposition~\ref{prop:kb-rank11-rank8-fixed-chart-local-cap-fence}.
\end{proposition}

\begin{proof}
The dense-root saturation theorem places the ten-dimensional correction
space $V'$ inside the residual degree-below-$K'$ Reed--Solomon space.  At
$K'=10$ both spaces have dimension ten, so
\[
 V'=\{f\in\F[X]:\deg f<10\}.
\]
For any nine distinct coordinates $b_1,\ldots,b_9$, the evaluation minor
formed by $1,X,\ldots,X^8$ has Vandermonde determinant
\[
 \prod_{1\le i<j\le9}(b_j-b_i)\ne0.
\]
Thus evaluation on the nine-set has rank nine, contradicting the rank-eight
target.  At $K'=11$, the space
$\langle1,X,\ldots,X^7,L_B,XL_B\rangle$ from the next proposition has
rank eight on $B$, so no adjacent-row propagation is claimed.
\end{proof}

\begin{proposition}[Rank-eight codimension-one circuit-shadow census]
\label{prop:kb-rank11-rank8-circuit-shadow-census}
At $K'=11$, fix a nine-set $B$ in the rank-eight affine-owner target.
There is one circuit $C_B\subset B$, independent of every extension pair,
with $2\le c:=|C_B|\le9$.  Every eleven-set
$T=B\cup\{x,y\}$ has exactly
\[
 \binom{11-c}{2}
\]
rank-eight nine-shadows, $55-\binom{11-c}{2}$ rank-nine nine-shadows,
and $c$ rank-ten bases.  Moreover
\[
 L_{C_B}\{f:\deg f<11-c\}\subset V'.
\]
Every $c\in\{2,\ldots,9\}$ occurs at the correction-space and chart-rank
level.  Thus no circuit-size branch is removed by local linear algebra.
\end{proposition}

\begin{proof}
Put $P=\{f:\deg f<11\}$.  It has dimension eleven, while $V'$ is a
ten-dimensional hyperplane.  Since evaluation of $V'$ on $B$ has rank
eight,
\[
 \ker(\operatorname{ev}_B|_{V'})
 =L_B\{a+bX:a,b\in\F\};
\]
both sides have dimension two.  A functional $\lambda$ cutting out $V'$
therefore annihilates the complete kernel of
$\operatorname{ev}_B:P\to\F^B$ and factors uniquely as
\[
 \lambda(f)=\sum_{b\in B}\mu_b f(b).
\]
Let $C_B$ be the support of $(\mu_b)_{b\in B}$.  If $|C_B|=1$, that
coordinate is an evaluation loop of $V'$.  For every one of the at least
$2578110$ distinct fixed-chart slopes, the normalized rich equation there
reduces to $E_b(\gamma)=0$.  Since $E_b(Z)$ is affine, two slopes make it
identically zero, producing a coordinate in the complete residual common
support, contrary to maximal shortening.  Hence $2\le |C_B|\le9$.

Evaluation of $P$ on the eleven distinct points of $T$ is an isomorphism.
The unique relation on the rank-ten image of $V'$ is consequently
$(\mu_b)_{b\in B}$, extended by zero at $x,y$.  Its support is the fixed
circuit $C_B$.  A nine-subset of $T$ has rank eight exactly when its omitted
pair misses this circuit; otherwise its nine columns are independent.
This gives the two shadow counts.  A ten-subset is a basis exactly when its
omitted point lies in $C_B$, giving $c$ bases.  The factorization of
$\lambda$ also places every polynomial vanishing on $C_B$ in $V'$, proving
the locator-ideal inclusion.

For sharpness, choose any $c$ points of $B$, nonzero weights on them, and
take $V'$ to be the kernel of the corresponding weighted evaluation
functional.  Its restriction to $B$ has rank eight, and no coordinate is a
loop because evaluation functionals at at most eleven distinct points are
Vandermonde-independent.  This realizes every $c=2,\ldots,9$ at the stated
linear-algebra level.  The eight-petal space in
Proposition~\ref{prop:kb-rank11-rank8-fixed-chart-local-cap-fence} is the
$c=9$ endpoint, with one rank-eight and 54 rank-nine shadows.
\end{proof}

\begin{proposition}[Rank-eight fixed-chart local-payment fence]
\label{prop:kb-rank11-rank8-fixed-chart-local-cap-fence}
At residual dimension $K'=11$, the local rank-eight output of
Proposition~\ref{prop:kb-rank11-component-ninesubset-target}, even with the
marked-extension unit retained, admits
\[
 4070352>2578110
\]
distinct exact-support records on one fixed chart.  Each record has
$\binom{67473}{2}$ rank-ten affine-owner component extensions, so the
marked weight is
\[
 9265216597693056>5869376383979174,
 \tag{R8F}\label{eq:kb-rank11-rank8-fixed-chart-local-cap-fence}
\]
where the right side is the exact weighted-selector demand at $K'=11$.
The errors have affine rank at most two, every support is pair-noncontained,
and the construction lifts reversibly to the official row.
\end{proposition}

\begin{proof}
On the residual domain, choose a nine-set $B$, put
\[
 u_0(X)=\prod_{b\in B}(X-b),\qquad u_1(X)=Xu_0(X),
\]
and take the ten-space
\[
 V'=\langle1,X,\ldots,X^7,u_0,u_1\rangle.
\]
Its evaluation rank on $B$ is eight, with kernel
$U=\langle u_0,u_1\rangle$.  For distinct $x,y\notin B$, the determinant
of kernel evaluation on $\{x,y\}$ is
\[
 u_0(x)u_0(y)(y-x)\ne0,
\]
so $B\cup\{x,y\}$ has evaluation rank ten.

Here $(n',m')=(1048587,67483)$.  Partition the complement of $B$ into
eight petals $P_e$ of size $67473$ and a remainder $R$ of size
\[
 1048587-9-8\cdot67473=508794.
\]
Choose distinct $t_e$ and owner pairs
$(a_e,b_e)=(t_eu_0,1)$.  Define the received pair by
\[
 (r_0,r_1)=
 \begin{cases}
 (0,1),&B,\\
 (a_e,1),&P_e,\\
 (s_x,0),&x\in R.
 \end{cases}
\]
Choose the $s_x$ successively so that all
\[
 \gamma_{e,x}=s_x-a_e(x)
\]
are distinct and avoid the eighteen dense slopes.  At the $j$-th
remainder coordinate, at most $64j+8\cdot18$ values are forbidden.  The
maximum
\[
 64(508794-1)+144=32562896
\]
is below the deployed base prime $2130706433$, so the greedy choice exists.

At slope $\gamma_{e,x}$ use explanation
$h_{e,x}=a_e+\gamma_{e,x}$.  Its exact support is
\[
 B\mathbin{\dot\cup}P_e\mathbin{\dot\cup}\{x\},
\]
of size $67483$.  Equality at another remainder coordinate would repeat a
slope, while equality on another petal would force
$(t_e-t_f)u_0=0$ off $B$.  For every two-set in $P_e$, adjoining it to $B$
gives a rank-ten component owned by $(a_e,1)$.  If any polynomial pair
contained the complete support, its second component would equal one on
$|B|+|P_e|=67482>10$ coordinates, hence would be the constant one, but the
received second column is zero at $x$.  Thus the support is pair-noncontained.
Finally
\[
 e_{e,x}=r_0-a_e+\gamma_{e,x}(r_1-1),
\]
so all error differences lie in $\langle u_0,r_1-1\rangle$.

There are $8\cdot508794=4070352$ records and their marked weight is the
left side of \textup{(R8F)}.  Direct evaluation of the weighted selector at
$K'=11$ gives the right side.  To lift, choose a common core $J$ of size
$K-11$, multiply the residual received values, owners, explanations, and
correction space by its locator, and set both received columns to zero on
$J$.  Evaluation ranks are unchanged off $J$, every support gains $J$, and
the same root-bound argument preserves pair noncontainment.

This witness has only four million records and its normalized deviations
do not span the complete ten-space.  It therefore fences a payment from the
printed chart-local output; it is not an unsafe line and does not realize
the dense-anchor ancestors.
\end{proof}

\begin{proposition}[Rank-eight owner-pair capacity]
\label{prop:kb-rank11-rank8-owner-pair-capacity}
In the setup of Proposition~\ref{prop:kb-rank11-component-incidence}, the
rank-eight affine-owner fixed target is impossible for every
\[
 37996\le K'\le1048576.
\]
For a fixed rank-eight nine-subset, its marked component load is at most
\[
 W_B\le981105\binom{n'-9}{2}.
 \tag{R8W}
\]
At $K'=37996$, the weighted selector demand is
$579191514708840299$, while the right side of (R8W) is
$579155144020629315$, a gap of $36370688210984$.
At $K'=37995$, the cap still exceeds demand by $18174297527234$.
\end{proposition}

\begin{proof}
Fix a rank-eight nine-subset $B$ and put
$U=\ker(\operatorname{ev}_B)$, so $\dim U=2$. Owners agreeing with the
received pair on $B$ form one affine translate of $U^2$.

Every marked extension has the form $T=B\cup\{x,y\}$ with
$\operatorname{rank}(\operatorname{ev}_T)=10$. Therefore evaluation on
$\{x,y\}$ is an isomorphism from $U$ to $\mathbb F^2$. The two equations
for the first received column determine the first $U$-coordinate of the
owner uniquely, and the two equations for the second column determine the
second $U$-coordinate uniquely. Thus one unordered coordinate pair
$\{x,y\}$ determines at most one owner point $p$.

Let $q_p$ count the full-rank coordinate pairs outside $B$ that determine
$p$, and let $t_p$ count the selected records owned by $p$. Coordinate-pair
uniqueness and fixed-owner exception disjointness give
\[
 \sum_p q_p\le\binom{n'-9}{2},
 \qquad
 t_p\le n'-m'+1=981105.
\]
Consequently
\[
 W_B\le\sum_p t_p q_p
     \le981105\binom{n'-9}{2},
\]
in the selector's marked $(\text{record},T)$ unit.

For uniformity, if $L(K')$ denotes the unrounded selector demand and
$U(K')$ the displayed cap, then
\[
 \frac{L(K')}{U(K')}
 =\text{constant}\cdot
   \frac{\binom{m'}{11}}{\binom{n'}{11}},
\]
using
$\binom{n'}9\binom{n'-9}2=55\binom{n'}{11}$.
Each factor $(m'-i)/(n'-i)$, $0\le i\le10$, strictly increases with $K'$
because $n'-m'=981104$. The exact positive gap at $K'=37996$ therefore
persists through the deployed endpoint. The failed preceding row records
the honest method wall.
\end{proof}

\begin{corollary}[Rank-eight bridge to the dense-owner terminal]
\label{cor:kb-rank11-rank8-dense-owner-bridge}
For every $22526\le K'\le37995$, a surviving rank-eight fixed chart
contains one owner pair that owns at least $200632$ selected records and
whose complete pair core has deficiency at most four.
\end{corollary}

\begin{proof}
With the notation in the preceding proof, let $q_p$ count coordinate pairs
determining owner $p$ and $t_p$ its selected records. Then
\[
 W_B\le\sum_p t_p q_p,\qquad
 \sum_p q_p\le\binom{n'-9}{2}.
\]
At $K'=22526$, exact replay gives
\[
 W_B-200631\binom{n'-9}{2}=11714977255865.
\]
Thus some $t_p\ge200632$. If that owner's pair-core deficiency were at
least five, fixed-owner exception disjointness would give
\[
 t_p\le1+\left\lfloor\frac{981104}{5}\right\rfloor
     =196221,
\]
a contradiction. Hence its deficiency is at most four.

The weighted-average ratio is again a positive constant times
$\binom{m'}{11}/\binom{n'}{11}$, so the conclusion persists through
$K'=37995$. At $K'=22525$, the comparison misses by
$1170919108090$, and no smaller interval is claimed.
\end{proof}

\begin{remark}[Typed scope of the rank-eleven packet]
\label{rem:kb-rank11-dense-locator-scope}
The density in Proposition~\ref{prop:kb-rank11-component-incidence} is a
record/coordinate-subset incidence density, not a record or component count.
The corollary is recordwise, and
Proposition~\ref{prop:kb-rank11-split-pencil-cell} is local to one fixed
cell, and Proposition~\ref{prop:kb-rank11-split-pencil-paircore} remains
local to one lifted owner plane.  Proposition~\ref{prop:kb-rank11-component-ninesubset-target}
selects one fixed populated chart, not the complete component lane.  The
deduplicated local conclusions alone do not pay that chart by
Proposition~\ref{prop:kb-rank11-fixed-chart-local-cap-fence}.  That fence
does not carry the marked weight in
Proposition~\ref{prop:kb-rank11-rank9-weighted-exclusion}, which eliminates
the rank-nine alternative from $K'=20618$ onward.
Proposition~\ref{prop:kb-rank11-rank9-residual-petal-capacity} extends that
elimination down to $K'=15635$ without transporting the original-row core
floor. Proposition~\ref{prop:kb-rank11-rank9-exact-petal-partition} solves
the integer petal partition and extends it further to $K'=15529$.
Proposition~\ref{prop:kb-rank11-weighted-selected-support-split-pencil}
adds a base-field-normalized selected-support census which avoids the
one-chart per-coordinate relaxation.  Its application in
Proposition~\ref{prop:kb-rank11-rank9-minimal-split-pencil-payment} closes
the isolated rank-nine row $K'=10$ by retaining the full component density.
Proposition~\ref{prop:kb-rank11-weighted-selected-support-core-offset}
prices the extra common-core coordinate at $K'=11$, and
Proposition~\ref{prop:kb-rank11-k11-circuit-split-pencil-payment} couples all
rank-nine shadows of each eleven-set.  Large circuits consume at least 45
shadows, while sparse representations of the global hyperplane functional
coalesce by Vandermonde independence.  This closes $K'=11$ despite the
surviving fixed-chart witness.
Proposition~\ref{prop:kb-rank11-kernel-basis-capacity}
first eliminates kernel-lane domination through $K'=4598$. Retaining every
basis decoration in Corollary~\ref{cor:kb-rank11-kernel-multibasis-capacity}
extends that interval through $K'=11641$. The record-support/ambient hybrid
of Corollary~\ref{cor:kb-rank11-kernel-hybrid-capacity} extends it through
$K'=11772$. Proposition~\ref{prop:kb-rank11-kernel-nine-shadow} couples the
corank strata and Corollary~\ref{cor:kb-rank11-kernel-nine-shadow-capacity}
extends the exact interval through $K'=15445$. Counting all contained
nine-subsets in
Proposition~\ref{prop:kb-rank11-kernel-nine-shadow-containment} supplies a
second resource; Corollary~\ref{cor:kb-rank11-kernel-nine-shadow-containment-capacity}
extends the interval through $K'=15670$, but not above the printed $15671$
method wall.
Proposition~\ref{prop:kb-rank11-kernel-rank8-nine-shadow-deficit} prices
the rank-eight shadow extension deficit, and
Corollary~\ref{cor:kb-rank11-kernel-rank8-nine-shadow-capacity} extends the
exact interval through $K'=17608$, but not above the printed $17609$ wall.
Proposition~\ref{prop:kb-rank11-kernel-two-step-nine-shadow} couples all
remaining coranks two ranks at a time, and
Corollary~\ref{cor:kb-rank11-kernel-two-step-nine-shadow-capacity} extends
the exact interval through $K'=18101$, but not above the printed $18102$
wall.
Proposition~\ref{prop:kb-rank11-kernel-multistep-shadow} supplies all 28
step sizes, and
Corollary~\ref{cor:kb-rank11-kernel-multistep-shadow-capacity} extends the
exact interval through $K'=18158$, but not above the printed $18159$ wall.
Proposition~\ref{prop:kb-rank11-kernel-corank1-projective-pair} sharpens the
active corank-one cap, and
Corollary~\ref{cor:kb-rank11-kernel-projective-pair-capacity} extends the
exact interval through $K'=377673$, but not above the printed $377674$ wall.
Proposition~\ref{prop:kb-rank11-kernel-corank2-projective-basis} sharpens the
complete-chart corank-two cap.  Proposition~\ref{prop:kb-rank11-kernel-projective-paving-scope}
records the weaker support-local shortening envelope, while
Propositions~\ref{prop:kb-rank11-rank3-bounded-parallel}
and~\ref{prop:kb-rank11-kernel-corank2-uniform-projective-basis}
prove the stronger cap uniformly.  Therefore
Corollary~\ref{cor:kb-rank11-kernel-projective-basis-capacity} extends the
exact interval through $K'=568338$, but not above the printed $568339$ wall.
Proposition~\ref{prop:kb-rank11-kernel-corank3-projective-frame} sharpens the
complete-chart corank-three cap.  Proposition~\ref{prop:kb-rank11-rank4-bounded-point-line}
and Corollary~\ref{cor:kb-rank11-kernel-corank3-uniform-projective-frame}
prove that cap uniformly.  Therefore
Corollary~\ref{cor:kb-rank11-kernel-projective-frame-capacity} extends the
exact interval through $K'=796598$, but not above the printed $796599$ wall.
Proposition~\ref{prop:kb-rank11-kernel-shortening-weighted-extension}
couples each chart's record cap to its extension count instead of maximizing
them separately. Corollary~\ref{cor:kb-rank11-kernel-shortening-weighted-capacity}
then removes the fixed-kernel branch through the deployed endpoint.
Proposition~\ref{prop:kb-rank11-rank8-owner-pair-capacity} eliminates the
rank-eight target from $K'=37996$ onward, but not below its printed wall.
Corollary~\ref{cor:kb-rank11-rank8-dense-owner-bridge} routes the upper
surviving rank-eight interval to the guarded dense-owner terminal, but does
not assign chronology or coalesce the multiple owners allowed by the
explicit fence. Proposition~\ref{prop:kb-rank11-rank8-minimal-shortening-exclusion}
removes rank eight from the shortest row.  Together with the minimal
split-pencil and circuit-shadow payments, there is no component target at
$K'=10$ or $K'=11$.  Proposition~\ref{prop:kb-rank11-quotient-line-sparse-circuit}
controls all sparse circuits selected by the codimension-two quotient line,
and Proposition~\ref{prop:kb-rank11-k12-quotient-line-payment} retains the
absolute kernel term while paying the complete $K'=12$ component
population.  Proposition~\ref{prop:kb-rank11-codimension-three-sparse-circuit}
controls the quotient-plane sparse circuits, and
Proposition~\ref{prop:kb-rank11-k13-sparse-circuit-payment} retains both
kernel coranks while paying the complete $K'=13$ component population.
Proposition~\ref{prop:kb-rank11-sparse-circuit-completion-ladder} extends
the completion dichotomy to arbitrary quotient dimension, and
Proposition~\ref{prop:kb-rank11-joint-sparse-shadow-ledger} prevents its
sparse incidences from spending the rank-nine-shadow capacity twice.
Proposition~\ref{prop:kb-rank11-k14-k21-joint-sparse-shadow-payment}
retains every kernel corank and core offset and pays the next eight rows.
Proposition~\ref{prop:kb-rank11-k22-near-saturation-payment} passes the old
$K'=22$ wall, Proposition~\ref{prop:kb-rank11-completion-defect-hierarchy}
pays $K'=23$, and
Proposition~\ref{prop:kb-rank11-k24-k40-full-deficit-payment} pays rank nine
through $K'=40$, and
Proposition~\ref{prop:kb-rank11-k41-sharp-isolated-payment} pays $K'=41$.
Proposition~\ref{prop:kb-rank11-cross-support-defect-carrier} couples the
sparse support strata, and
Proposition~\ref{prop:kb-rank11-k42-cross-support-payment} pays $K'=42$.
Proposition~\ref{prop:kb-rank11-descending-support-completion-ladder} makes
the lower-support completion alternatives exhaustive, and
Proposition~\ref{prop:kb-rank11-k43-descending-support-payment} pays
$K'=43$.  Proposition~\ref{prop:kb-rank11-completion-branch-lattice}
permits selective refinement of the heaviest leaves, and
Proposition~\ref{prop:kb-rank11-k44-branch-lattice-payment} pays $K'=44$.
Proposition~\ref{prop:kb-rank11-support45-joint-zero-carrier} couples the
terminal support-four and support-five carriers,
Proposition~\ref{prop:kb-rank11-support4-external-charge} converts that
coupling into a support-four incidence cap, and
Proposition~\ref{prop:kb-rank11-k45-full-completion-product-payment} pays
$K'=45$ after the exhaustive support-$2$ through support-$9$ completion
product.  Proposition~\ref{prop:kb-rank11-support45-deep-defect-partition}
refines the former simultaneous fallback to every exact support-four and
support-five defect pair, and
Proposition~\ref{prop:kb-rank11-k46-k53-deep-joint-payment} pays
$46\le K'\le53$ after exact Pareto compression of the other supports.
Proposition~\ref{prop:kb-rank11-small-support-self-collision-charge}
uses the positive intersection of any two same-support deletion-vanishing
spaces for supports two through five, and
Proposition~\ref{prop:kb-rank11-k54-k59-small-support-collision-payment}
pays $54\le K'\le59$ after exact defects and grouped Pareto compression.
Proposition~\ref{prop:kb-rank11-cross-support-collision-charge} couples each
nonempty source carrier of support two through five to every target support
whose two support sizes sum to at most eleven, and
Proposition~\ref{prop:kb-rank11-k60-k70-cross-support-collision-payment}
pays $60\le K'\le70$ after retaining every inherited cap and higher-support
branch.  Proposition~\ref{prop:kb-rank11-multicarrier-collision-charge}
prices target circuits relative to any fixed union carrying a positive
vanishing space,
Proposition~\ref{prop:kb-rank11-k71-carrier-position-trichotomy} exhausts
the relative positions of the attaining support-two, support-three, and
support-four carriers, and
Proposition~\ref{prop:kb-rank11-k71-carrier-trichotomy-payment} pays
$K'=71$.  Proposition~\ref{prop:kb-rank11-full-completion-pairwise-carrier-atlas}
extends the carrier classification to every full-completion overlap,
Proposition~\ref{prop:kb-rank11-flat-coupled-support45-charge} gives the
base-field-normalized joint support-four/support-five census, and
Proposition~\ref{prop:kb-rank11-k72-k73-full-carrier-atlas-payment} pays
$K'=72,73$, and
Proposition~\ref{prop:kb-rank11-k74-k78-full-carrier-atlas-payment} pays
$K'=74,\ldots,78$, and
Proposition~\ref{prop:kb-rank11-k79-k82-full-carrier-atlas-payment} pays
$K'=79,\ldots,82$ and isolates the first pairwise-atlas wall at $K'=83$.
Proposition~\ref{prop:kb-rank11-adjacent-flat-circuit-coupling} and
Proposition~\ref{prop:kb-rank11-fixed-union-adjacent-support-coupling}
convert that method wall into support-disjoint adjacent-pair payments, and
Proposition~\ref{prop:kb-rank11-k83-adjacent-support-payment} pays $K'=83$,
while Proposition~\ref{prop:kb-rank11-k84-adjacent-support-payment} pays
$K'=84$ by an independent exact-row continuation, and
Proposition~\ref{prop:kb-rank11-k85-best-single-adjacent-payment} pays
$K'=85$ after the exact raw-threshold split and residual best-single census.
Thus the remaining intervals are rank eight only on
$22\le K'\le85$, rank nine or rank eight on $86\le K'\le15528$, rank
eight only on $15529\le K'\le22525$,
dense-owner chronology only on
$22526\le K'\le37995$, and no rank-eleven component target on
$37996\le K'\le1048576$.
Proposition~\ref{prop:kb-rank11-rank8-fixed-chart-local-cap-fence} shows
that neither the distinct nor the weighted fixed-chart output alone can
pay the lower rank-eight interval.  Any continuation must retain the
normalized full-span/anchor ancestry or couple charts and owners through a
chronology-correct global incidence argument.
At the first remaining rank-nine row $K'=86$, no closure is asserted here.
The $K'=83$ pairwise wall remains a valid limitation of that restricted
method, but its frontier conclusion is superseded by the adjacent-support
payments.  No extrapolation from any exact row is used.
Proposition~\ref{prop:kb-rank11-rank8-circuit-shadow-census} gives the exact
first-row split: the full-circuit case has 54 rank-nine shadows, while every
proper circuit gives additional rank-eight shadows and a locator ideal of
dimension at least three.  The circuit-shadow payment retains precisely
those neighboring marks; it does not assert a smaller fixed-chart cap.
The fence does not instantiate the fixed
dense-anchor ancestors or an unsafe line.  The packet does not recursively
cover the remaining records,
construct the chronology owner required by
Remark~\ref{rem:mca-dense-core-owner-acceptance-contract}, pay the remaining
rank-eight or dense-owner chronology branch, move an
active-v4 atom, or close
KoalaBear.
\end{remark}

\section{The new spread residual}

\begin{theorem}[spread residual after core-relative incidence]\label{thm:residual}
Assume a deployed support-wise MCA-bad line contains a first-match primitive spread core and has more bad slopes than its row budget.  Then all of the following must hold.
\begin{enumerate}[label=\textup{(R\arabic*)}]
\item At most $58$ bad slopes use the core interpolant itself.
\item The remaining corrections are not contained in one projective correction ray.
\item If their linear correction span $W$ has dimension at most $50$ on KoalaBear or $15$ on Mersenne-31, then some $s+1$ coordinate equations have a positive-dimensional common component.
\item If $W$ does not absorb the high core, then its dimension and support excess lie beyond every transition in \cref{cor:transitions}.
\item Every low-dimensional surviving component is therefore an evaluation rank-flat or an exact polynomial clone component in the sense of \cref{lem:nonproper}; it is not an unstructured failure of Bezout counting.
\end{enumerate}
\end{theorem}

\begin{proof}
Items (R1)--(R4) are the contrapositives of \cref{thm:core-compatible,thm:ray,thm:proper,thm:clone-tolerant}.  Item (R5) is \cref{lem:nonproper}.
\end{proof}

The theorem changes the spread-abundance target.  It is no longer necessary to control arbitrary extensions of a core.  The only remaining spread mechanism is
\[
 \boxed{
 \begin{array}{c}
 \text{many projectively independent correction directions, together with}\\
 \text{a high-dimensional correction span or a positive-dimensional}\\
 \text{rank-flat/clone atlas surviving all first-match routes.}
 \end{array}
 }
\]

\begin{problem}[next spread theorem]\label{prob:next}
\emph{Intermediate status.}
This is the local correction-component form of the final
\cref{prob:mca-spread-routing}.  It is retained to state the geometric
alternatives exposed by \cref{thm:residual}, but is not a separate terminal
input.

Prove that every positive-dimensional correction component from \cref{thm:residual} either
\begin{enumerate}[label=\textup{(\alph*)}]
\item yields a rational-owner, quotient, extension, field-drop, or quantified near-sunflower atom declared in \cref{thm:partial-relative,thm:owner-localization,cor:near-sunflower};
\item has correction directions contained in a bounded union of projective rays, so \cref{thm:ray} pays it; or
\item admits a proper quotient correction space to which \cref{thm:proper} or \cref{thm:clone-tolerant} applies.
\end{enumerate}
For Mersenne-31 it is enough that the total surviving spread contribution be at most $16{,}777{,}215$; for KoalaBear the available budget is much larger.
\end{problem}


\part{Completion certificates and remaining row-sharp inputs}
\label{part:completion}

\section{Integrated theorem status}
\label{sec:status-and-obligations}
\begin{theorem}[Audited theorem status]\label{thm:audited-status}
All statements explicitly presented as theorems, propositions, lemmas, or corollaries in Parts~I--VI are unconditional under their displayed hypotheses.  In particular, the paper proves the complete support and profile compiler, the minimum-distance flag sharpening, the primitive-Q/Sidon implication and its shallow unconditional range, saturation, one-pencil moving-root BC, Q $\Rightarrow$ SP, the order-$32$ rank dictionary, dyadic descent, rational-atom extraction and coherence, owner localization, the seven-point component-free owner-pencil bound, the $58$-slope core-compatible bound, the complete correction-ray bound, and the displayed proper and clone-tolerant correction-space payments.

No theorem in this paper proves either deployed adjacent MCA-safe row or
either deployed adjacent list-safe row.  The terminal open inputs used by the
completion certificates are exactly
\cref{prob:mca-spread-routing,prob:large-owner,prob:mca-exception-routing,prob:list-completion}.
The earlier \cref{prob:spread-abundance,prob:next} are intermediate
formulations subsumed by \cref{prob:mca-spread-routing}, not additional
terminal cells.
\end{theorem}

\begin{remark}[Status of the v3 saturated-BC problem]
\label{rem:saturated-bc-status}
The v3 problem labelled \texttt{prob:saturated-bc} is superseded in v4, not
silently discharged.  Its primitive one-parameter MCA clause is proved by
\cref{cor:bc-one-pencil}.  Its higher-dimensional MCA clause remains open and
is represented jointly by
\cref{prob:mca-spread-routing,prob:large-owner,prob:mca-exception-routing};
in particular, v4 still requires exhaustive same-owner routing from every
surviving balanced-core component into that endgame.  Its separate
arbitrary-word list-interior clause is contained in
\cref{prob:list-completion}.  Thus no archived problem is being banked as a
live theorem, and no higher-dimensional balanced-core payment is claimed
merely from the one-pencil moving-root bound.
\end{remark}

\section{The final MCA certificate}
\label{sec:mca-final-certificate}
The order-$32$ method replaces the old sum over many independently budgeted MCA cells by a maximum-type endgame.

\begin{problem}[Complete spread-component routing]\label{prob:mca-spread-routing}
This is the final v4 successor to the higher-dimensional MCA clause of the
v3 \texttt{prob:saturated-bc}.  Together with
\cref{prob:large-owner,prob:mca-exception-routing}, it must provide the
exhaustive same-owner routing that the old balanced-core cell left open.  Its
row-sharp conclusion subsumes the intermediate
\cref{prob:spread-abundance,prob:next} for purposes of the final certificate.

Prove that every positive-dimensional correction component surviving \cref{thm:residual} is routed, with the same received-line owner retained, into one of the following disjoint outcomes:
\begin{enumerate}[label=\textup{(\alph*)}]
\item a rational-owner, quotient, extension, field-drop, or quantified near-sunflower atom;
\item a bounded union of complete projective correction rays paid by \cref{thm:ray};
\item a proper quotient correction space paid by \cref{thm:proper} or \cref{thm:clone-tolerant};
\item an explicitly named residual component with a literal row-sharp slope bound.
\end{enumerate}
The total spread contribution must fit the row budget.
\end{problem}

\begin{problem}[Complete MCA exception routing]\label{prob:mca-exception-routing}
Verify the exception-routing input of \cref{def:exception-routing}: denominator-root, extension, coordinate-clone, imperfect-quotient, and quantified near-sunflower failures must either be absorbed into the spread or large-owner bounds with a common owner retained, or occupy a set of at most $31$ slopes.
\end{problem}

\begin{theorem}[Exact MCA completion certificate]\label{thm:mca-completion-certificate}
Fix either deployed MCA row.  If \cref{prob:mca-spread-routing,prob:large-owner,prob:mca-exception-routing} are proved with the literal row constants in Parts~V--VI, then
\[
 B_{\rm MCA}(m)\le B^*.
\]
Together with the unsafe witness at agreement $m-1$ from CAP25 v13.2, the exact adjacent MCA threshold is $m$.
\end{theorem}

\begin{proof}
The hypotheses supply the spread-abundance and exception-routing inputs of \cref{thm:conditional-final}; \cref{prob:large-owner} supplies its atom input.  Apply that theorem and then the imported unsafe witness.
\end{proof}

\section{The final list certificate}
\label{sec:list-final-certificate}
For a list row retain the disjoint summed ledger
\begin{equation}
 U_{\rm list}
 =U_{\rm paid}+U_Q+U_{\rm list-int}+U_{\rm ext}+U_{\rm new}.
 \label{eq:list-final-ledger}
\end{equation}
Here $U_{\rm paid}$ is the exact first-match sum of all tangent, common-support, quotient, planted, field-drop, fixed-union, affine-span, common-zero, rank-flat, codimension-one, support-flag, minimum-distance, puncturing, and other certified cells; $U_Q$ is the row-sharp pruned boundary-prefix maximum; $U_{\rm list-int}$ is the high-nullity arbitrary-word interior codeword-ray term; $U_{\rm ext}$ is the exact extension projection payment; and $U_{\rm new}$ is zero unless an explicitly named residual cell survives.

\begin{problem}[Row-sharp list completion]\label{prob:list-completion}
This is the final v4 successor to the arbitrary-word list-interior clause of
the v3 \texttt{prob:saturated-bc}; it also includes the remaining prefix and
extension terms required by the v4 list ledger.

For each deployed list row, prove exact integer bounds for the remaining high-nullity prefix fiber, arbitrary-word interior, and extension projection terms so that
\[
 U_{\rm list}\le B^*.
\]
The prefix estimate must use the actual residual mass and the exact budget left after the certified compiler cells are charged.
\end{problem}

\begin{theorem}[Exact summed completion certificate]\label{thm:exact-completion-certificate}
Fix one deployed list row.  If an exhaustive first-match compiler proves, with no unresolved cell,
\[
 U_{\rm list}\le B^*,
\]
then the first safe integer agreement is the adjacent row $a_0+1$.
\end{theorem}

\begin{proof}
CAP25 v13.2 proves the unsafe inequality at $a_0$.  Exhaustive first-match coverage and the displayed upper bound give safety at $a_0+1$.  Monotonicity identifies the first safe integer agreement.
\end{proof}

\section{Conclusion}\label{sec:conclusion}
This paper gives one unified Grande Finale manuscript.  It contains the exact threshold refinements, profile and Fourier compiler, primitive Q, support and minimum-distance recursions, saturation and moving-root geometry, the order-$32$ rank pivot, rational-atom extraction and coherence, owner localization, and the current spread-core incidence theory.

The remaining MCA problem has been compressed to three row-sharp inputs: route the positive-dimensional correction components isolated by \cref{thm:residual}, prove the large-owner image bound of \cref{prob:large-owner}, and complete the finite exception routing of \cref{prob:mca-exception-routing}.  The binding numerical atom target is the Mersenne-31 full-owner factor $9.57219783037$.  The remaining list problem is the exact high-nullity prefix/interior/extension sum of \cref{prob:list-completion}.  Until those inputs are proved, no adjacent safe row is claimed.

\begin{thebibliography}{99}
\bibitem[Cho26]{CAP25v132}
P. Chojecki,
\newblock Universal field-size caps and a two-sided sandwich for mutual correlated agreement on smooth Reed--Solomon domains,
\newblock Version 13.2, manuscript, July 2026.

\bibitem[BS94]{BalogSzemeredi}
A. Balog and E. Szemer\'edi,
\newblock A statistical theorem of set addition,
\newblock \emph{Combinatorica} 14 (1994), 263--268.

\bibitem[BGM23]{BGM23}
J. Brakensiek, S. Gopi, and V. Makam,
\newblock Generic Reed--Solomon codes achieve list-decoding capacity,
\newblock \emph{Proceedings of STOC 2023}; arXiv:2206.05256.

\bibitem[BCHKS25]{BSCHKS25}
E. Ben-Sasson, D. Carmon, U. Hab\"ock, S. Kopparty, and S. Saraf,
\newblock On proximity gaps for Reed--Solomon codes,
\newblock ECCC TR25-169, 2025.

\bibitem[GG25]{GG26}
R. Goyal and V. Guruswami,
\newblock Optimal proximity gaps for subspace-design codes and (random) Reed--Solomon codes,
\newblock ECCC TR25-166, 2025.

\bibitem[GGSW26]{GGSW26}
R. Goyal, V. Guruswami, Y. Sun, and M. Wootters,
\newblock Locality of curve-decoding and improved proximity gaps,
\newblock arXiv:2607.08516, 2026.

\bibitem[GMR+22]{GreenMatolcsiRuzsaShakanZhelezov}
B. Green, D. Matolcsi, I. Z. Ruzsa, G. Shakan, and D. Zhelezov,
\newblock A weighted Pr\'ekopa--Leindler inequality and sumsets with quasicubes,
\newblock in \emph{Analysis at Large}, Springer, 2022, 125--129.

\bibitem[LW10]{LiWanSieve}
J. Li and D. Wan,
\newblock A new sieve for distinct coordinate counting,
\newblock \emph{Sci. China Math.} 53 (2010), 2351--2362.

\bibitem[Tao10]{TaoEntropy}
T. Tao,
\newblock Sumset and inverse sumset theorems for Shannon entropy,
\newblock \emph{Combin. Probab. Comput.} 19 (2010), 603--639.

\bibitem[TV06]{TaoVu}
T. Tao and V. Vu,
\newblock \emph{Additive Combinatorics},
\newblock Cambridge University Press, 2006.

\bibitem[Gow98]{GowersSzemeredi}
W. T. Gowers,
\newblock A new proof of Szemer\'edi's theorem for arithmetic progressions of length four,
\newblock \emph{Geom. Funct. Anal.} 8 (1998), 529--551.

\bibitem[GMT23]{GreenMannersTao}
B. Green, F. Manners, and T. Tao,
\newblock Sumsets and entropy revisited,
\newblock arXiv:2306.13403, 2023.

\bibitem[Goh24]{GohEnergy}
M. K. Goh,
\newblock On an entropic analogue of additive energy,
\newblock arXiv:2406.18798, 2024.

\bibitem[MRSZ22]{MatolcsiRuzsaShakanZhelezov}
D. Matolcsi, I. Z. Ruzsa, G. Shakan, and D. Zhelezov,
\newblock An analytic approach to cardinalities of sumsets,
\newblock \emph{Combinatorica} 42 (2022), 203--236.

\bibitem[Wei48]{Weil}
A. Weil,
\newblock On some exponential sums,
\newblock \emph{Proc. Natl. Acad. Sci. USA} 34 (1948), 204--207.
\end{thebibliography}
\end{document}
